% LuaLaTeX / A4実寸 / 10pt / 数式左寄せ
\documentclass[a4paper,10pt,oneside,openany,fleqn]{ltjsbook}
\usepackage[margin=22mm,headheight=16pt,headsep=8mm,footskip=11mm]{geometry}
\usepackage{luatexja-fontspec}
\setmainjfont[BoldFont=HaranoAjiGothic-Bold]{HaranoAjiMincho-Regular}
\setsansjfont[BoldFont=HaranoAjiGothic-Bold]{HaranoAjiGothic-Medium}
\usepackage{amsmath,amssymb,amsthm}
\everymath{\displaystyle}
\usepackage{graphicx,xcolor,tikz,pgfplots}
\usetikzlibrary{calc,arrows.meta,patterns}
\pgfplotsset{compat=1.18}
\usepackage{booktabs,longtable,tabularx,array,enumitem}
\usepackage[most]{tcolorbox}
\usepackage{needspace,fancyhdr}
\definecolor{blueblack}{HTML}{18374B}
\definecolor{tealblack}{HTML}{176C73}
\definecolor{softblue}{HTML}{F0F5F7}
\definecolor{softteal}{HTML}{EDF6F3}
\usepackage[unicode,colorlinks=true,linkcolor=blueblack,urlcolor=tealblack,bookmarksnumbered=true]{hyperref}
\hypersetup{pdftitle={積分の体系的解説},pdfsubject={微分とのつながりから，計算の選択と発展へ}}
\setcounter{tocdepth}{1}
\setcounter{secnumdepth}{2}
\renewcommand{\arraystretch}{1.35}
\setlength{\parskip}{0.18\baselineskip}
\setlength{\parindent}{2em}
\setlength{\mathindent}{2em}
\setlength{\emergencystretch}{2em}
\allowdisplaybreaks[2]
\raggedbottom
\newcommand{\dd}{\,\mathrm{d}}
\newcommand{\e}{\mathrm{e}}
\newcommand{\R}{\mathbb{R}}
\newcommand{\chapterref}[1]{第\ref{#1}章}
\newcommand{\sectionref}[1]{第\ref{#1}節}
\newcounter{example}[chapter]
\renewcommand{\theexample}{\thechapter.\arabic{example}}
\newenvironment{example}[1]{\refstepcounter{example}\begin{tcolorbox}[enhanced,breakable,colback=softblue,colframe=blueblack,boxrule=0.5pt,arc=0mm,left=2.5mm,right=2.5mm,top=1.5mm,bottom=1.5mm,title={例題\theexample\quad #1},fonttitle=\bfseries]}{\end{tcolorbox}}
\newcounter{result}[chapter]
\renewcommand{\theresult}{\thechapter.\arabic{result}}
\newenvironment{theorem}[1]{\refstepcounter{result}\begin{tcolorbox}[enhanced,breakable,colback=softteal,colframe=tealblack,boxrule=0.4pt,arc=0mm,title={定理\theresult\quad #1},fonttitle=\bfseries]}{\end{tcolorbox}}
\newenvironment{definition}[1]{\begin{tcolorbox}[enhanced,breakable,colback=white,colframe=tealblack,boxrule=0.4pt,arc=0mm,title={定義\quad #1},fonttitle=\bfseries]}{\end{tcolorbox}}
\newenvironment{note}{\par\smallskip\noindent\textbf{補足}\quad}{\par\smallskip}
\pagestyle{fancy}
\fancyhf{}
\fancyhead[L]{\small\nouppercase{\leftmark}}
\fancyhead[R]{\small 積分の体系的解説}
\fancyfoot[C]{\thepage}
\renewcommand{\headrulewidth}{0.3pt}
\renewcommand{\chaptermark}[1]{\markboth{第\thechapter 章\quad #1}{}}
\fancypagestyle{plain}{\fancyhf{}\fancyfoot[C]{\thepage}\renewcommand{\headrulewidth}{0pt}}
\fancypagestyle{plainhead}{\fancyhf{}\fancyfoot[C]{\thepage}\renewcommand{\headrulewidth}{0pt}}
\begin{document}
\begin{titlepage}
\thispagestyle{empty}
\vspace*{15mm}
{\sffamily\color{tealblack}\large CALCULUS\par}
\vspace{7mm}
{\sffamily\bfseries\color{blueblack}\Huge 積分の体系的解説\par}
\vspace{6mm}
{\Large 微分とのつながりから，計算の選択と発展へ\par}
\vspace{18mm}
\begin{center}
\begin{tikzpicture}[x=2.9cm,y=1.9cm,>=Stealth]
\fill[softteal] (0,0) -- plot[domain=0:1.8,samples=80] (\x,{0.6*\x*\x}) -- (1.8,0) -- cycle;
\foreach \i in {0,...,8}{\pgfmathsetmacro{\xx}{\i*0.2}\pgfmathsetmacro{\yy}{0.6*(\xx+0.2)*(\xx+0.2)}\draw[tealblack!35] (\xx,0) rectangle ({\xx+0.2},\yy);}
\draw[->,blueblack] (-0.1,0)--(2.45,0) node[right]{$x$};
\draw[->,blueblack] (0,-0.1)--(0,3.05) node[above]{$y$};
\draw[very thick,tealblack,domain=0:2.12,samples=100] plot (\x,{0.6*\x*\x});
\draw[thick,blueblack,domain=0.8:2.3] plot (\x,{1.8*\x-1.35});
\fill[blueblack] (1.5,1.35) circle(1.2pt);
\node[anchor=west,blueblack] at (2.15,2.15) {変化を調べる};
\node[tealblack] at (1.25,0.35) {細かく分けて足す};
\end{tikzpicture}
\end{center}
\vfill
{\large 高校の微積分から，その先の考え方へ\par}
\vspace{4mm}
{第1版\quad 2026年10月21日\par}
\end{titlepage}
\frontmatter
\chapter*{はじめに}
\addcontentsline{toc}{chapter}{はじめに}
\markboth{はじめに}{}
曲線で囲まれた図形の面積を求めたいとき，長方形や三角形の面積の公式だけでは困ることがある。
図形を細かく分ければ近い値は求められそうだが，どこまで分ければ正確になるのだろうか。
一方で，関数のグラフを描くときには，どこで増え，どこで減るかを知りたくなる。
この二つの問いを結び付けるのが，微分と積分である。

本書では，積分を中心に，高校で学ぶ極限・微分・積分の考え方を説明する。
面積を細かく分けるところから出発し，極限と微分を準備して，微分と積分がどのようにつながるのかを確かめる。
その後で，積分の計算と応用を考えていこう。

積分の問題を解く際には，公式を知っていても，何を使えばよいか分からないことがある。
例えば，平方根が含まれているというだけでは置換は一つに決まらない。
平方根の中の式を微分すると外側にある因子と一致するなら，その中身を新しい変数にするのが自然である。
円の方程式の形が見えるなら，三角関数で表す方法もある。
同じ「置換積分」であっても，簡単にしたい箇所が違うのである。
本書では，このように式を見て方針を決める理由を，実際の計算とともに示す。

また，一つの積分をいくつかの方法で解くと，異なる公式の関係が見えることがある。
三角関数の置換，円の幾何，双曲線関数は，別々に覚えるだけでは見えにくい関係をもっている。
基本の計算を身に付けた後では，その背景にも進めるようにした。
高度な結果を使わずに解く方法がある場合は，基本の道具でどこまで進めるかも確かめる。

\section*{対象読者と読む順序}
二次関数，整式の計算，三角関数の基本，指数・対数，数列の和を学んだ読者を想定する。
微積分を初めて学ぶ場合は，\chapterref{章:導入}から\chapterref{章:積分の応用}までを順に読み，章末問題で計算を確かめてほしい。
\chapterref{章:極限と微分}は後の説明に必要な微分をまとめているため，他の章より多くの内容を含む。
すでに微分を学んでいる場合は，\chapterref{章:定積分}から読み始め，必要になった箇所だけ\chapterref{章:極限と微分}へ戻ってよい。

\begin{center}
\begin{tabularx}{\linewidth}{@{}p{35mm}X@{}}
\toprule
読みたい内容 & 主な箇所\\
\midrule
積分の意味を知る & 第1〜3章。近似の和が定積分になり，原始関数で計算できる理由を扱う。\\
計算方法を選ぶ & 第4〜7章。式の整理，置換，部分積分，対称性を比較する。\\
積分を問題に使う & 第8〜9章。図形・運動・数列・不等式を通して，何を足すかを考える。\\
手法の背景へ進む & 第10〜11章。置換の幾何，逆関数，近似，広義積分を扱う。\\
自分で方針を決める & \chapterref{章:総合演習}。問題を解いてからヒントや巻末の解答を参照する。\\
\bottomrule
\end{tabularx}
\end{center}

\chapterref{章:積分と数列}の一般のベータ関数，第10〜11章，および\chapterref{章:積分の応用}の重心とパップス＝ギュルダンの定理は，高校の基本事項から先へ進むための発展内容である。
基本の学習に際して，これらの定理をすべて覚える必要はない。
発展内容では，高校の計算から何が広がるかと，新たに何を仮定するかを区別して読むとよい。
本書は高校数学全分野の教科書ではない。複素数平面，確率，整数などは，積分との関係が必要な範囲で扱う。

\section*{文字と説明について}
特に断らなければ，変数は実数を表す。
$\log x$は自然対数，$\e$はその底とし，三角関数の角は弧度法で測る。
$\mathbb{N}$を用いる場合は正の整数の集合とする。
不定積分に現れる$C$は積分定数であり，式ごとに値が変わってよい。
分母が零になる点などを除いて区間が分かれる場合には，それぞれの区間で定数を選ぶ。

定積分の$\dd x$は，どの変数について積分しているかを示す。
置換の途中で$\dd u=g'(x)\dd x$と書くときには，置換積分の公式を簡潔に表している。
この記号だけを根拠に，どのような操作でも分数のように行えると考えないようにしたい。

重要な公式は，その公式を使う理由と導出を本文に示した。
一方，連続関数の積分可能性など，実数の性質にまで戻る証明は，採用する事実として明示している。
発想を得るための近似と，結果が正しいことを確かめる議論も分けてある。
例題と章末問題は本書の説明のために設定したものであり，特定の入試問題の転載ではない。
章末問題と総合演習には，巻末にすべて解答を付けた。

\tableofcontents
\mainmatter
\chapter{曲線の面積と変化をどう求めるか}\label{章:導入}
\section{直線でできることを曲線へ広げる}
長方形の面積は横の長さと縦の長さの積で求められる。直線の傾きも，横の変化量に対する縦の変化量の比で求められる。ところが，境界が曲線になったり，傾きが場所ごとに変わったりすると，同じ計算をそのまま使うことはできない。この本では，直線や長方形について知っている計算から出発し，曲線の変化と面積を調べる方法を作る。

まず，二つの問いを並べてみよう。第一は，放物線 $y=x^2$ 上の点 $(1,1)$ で，曲線はどの程度の傾きで進んでいるのかという問いである。第二は，この放物線と $x$ 軸，直線 $x=1$ に囲まれた部分の面積はいくらかという問いである。第一の問いでは二点を結ぶ直線を考え，その二点を近づける。第二の問いでは細い長方形を並べ，その幅を小さくする。どちらも，計算できるものから始めて，細かくした先にある値を求めようとしている。

ここで「限りなく近づける」という表現は，まだ計算の見通しを説明するために使っている。何をどのように近づけるのかを明らかにするのが\chapterref{章:極限と微分}であり，長方形の和を積分として定義するのが\chapterref{章:定積分}である。この章では具体的な図と計算を通して，後の定義が何のために必要なのかを確かめる。

\section{二点間の変化と一点での変化}\label{節:導入変化}
関数 $y=f(x)$ において，$x$ が $a$ から $a+h$ まで変わったとする。$h\ne0$ のとき，横の変化量は $h$，縦の変化量は $f(a+h)-f(a)$ だから，二点を結ぶ直線の傾きは
\[
 \frac{f(a+h)-f(a)}{h}
\]
である。これをその区間の\textbf{平均変化率}という。「平均」と呼ぶのは，途中の変化のしかたを区別せず，始点と終点の変化だけで一つの数を作っているからである。

$f(x)=x^2$，$a=1$ では，
\[
 \frac{(1+h)^2-1}{h}
 =\frac{2h+h^2}{h}=2+h \qquad(h\ne0)
\]
となる。$h=1$ なら傾きは $3$，$h=0.1$ なら $2.1$，$h=-0.1$ なら $1.9$ である。$h$ を正の側からも負の側からも $0$ に近づけると，傾きは $2$ に近づく。この値を，放物線上の点 $(1,1)$ における接線の傾きとして採用したい。

\begin{figure}[htbp]
\centering
\begin{tikzpicture}[x=2.6cm,y=1.5cm,>=Stealth]
 \draw[->] (-0.12,0)--(2.1,0) node[right]{$x$};
 \draw[->] (0,-0.25)--(0,3.75) node[above]{$y$};
 \draw[very thick,domain=0:1.9,samples=80] plot (\x,{\x*\x}) node[right]{$y=x^2$};
 \draw[blue,thick,domain=0.45:1.75] plot (\x,{2*\x-1});
 \draw[red,thick,domain=0.55:1.8] plot (\x,{2.6*\x-1.6});
 \fill (1,1) circle(1.6pt) node[below right]{$P(1,1)$};
 \fill (1.6,2.56) circle(1.6pt) node[right]{$Q$};
 \draw[dashed] (1,0) node[below]{$1$}--(1,1);
 \draw[dashed] (1.6,0) node[below]{$1+h$}--(1.6,2.56);
 \node[blue] at (1.82,2.08) {接線};
 \node[red] at (0.65,1.6) {割線};
\end{tikzpicture}
\caption{二点を結ぶ割線の傾きを，点 $Q$ を点 $P$ に近づけて調べる。}
\label{図:導入接線}
\end{figure}

図\ref{図:導入接線}のように，曲線上の二点を結ぶ直線を\textbf{割線}と呼ぶ。割線は二点で決まるのに対し，接線の傾きは一つの点の近くの関数の振る舞いで決まる。図だけを見て「曲線に一度だけ触れる直線」と定義すると，接線が別の場所で再び曲線と交わる場合を扱えない。微分では，割線の傾きの極限を使って接線を定める。

\begin{example}{平均の速さから瞬間の速さへ}
一直線上を動く物体の時刻 $t$ 秒における位置が $s(t)=t^2$ メートルで表されている。時刻 $2$ 秒から $2+h$ 秒までの平均の速度を求め，時刻 $2$ 秒における速度を予想せよ。
\end{example}
位置の変化を経過時間で割ればよいので，平均の速度は
\[
 \frac{s(2+h)-s(2)}{h}
 =\frac{(2+h)^2-4}{h}=4+h
\]
メートル毎秒である。たとえば $h=0.1$ なら $4.1$，$h=0.01$ なら $4.01$ となる。時刻 $2$ 秒より前の区間を調べるために $h<0$ としても，同じ式が使える。どちらから近づけても $4$ に近づくので，瞬間の速度は $4$ メートル毎秒と予想できる。

この計算では $h$ を $0$ に代入していない。元の式では $h=0$ にすると分母も分子も $0$ となり，割り算ができない。$h\ne0$ のまま式を整理し，$h$ が $0$ に近づくときの値を調べたのである。\textbf{変化量を $0$ にすることと，$0$ に近づけることは別の操作}である。

\section{長方形で面積を挟む}\label{節:導入長方形}
次に，$y=x^2$ のグラフの下で，$0\le x\le1$ の範囲にある部分の面積を考える。この部分の面積を $S$ と書く。横の区間を $n$ 等分すると，一つの小区間の幅は $\frac1n$ である。$x^2$ はこの範囲で増加するから，小区間の左端の高さで作った長方形は曲線の下に収まり，右端の高さで作った長方形は曲線の上まで覆う。

左端の高さを使った和を $L_n$，右端の高さを使った和を $U_n$ とすると，
\begin{align*}
 L_n&=\frac1n\sum_{k=0}^{n-1}\left(\frac{k}{n}\right)^2,\\
 U_n&=\frac1n\sum_{k=1}^{n}\left(\frac{k}{n}\right)^2
\end{align*}
である。添字が異なるのは，左端では $0$ から $\frac{n-1}{n}$ まで，右端では $\frac1n$ から $1$ までの高さを使うためである。図\ref{図:導入長方形}では $n=4$ としている。

\begin{figure}[htbp]
\centering
\begin{tikzpicture}[x=3.7cm,y=3.1cm,>=Stealth]
 \begin{scope}
  \foreach \k in {0,1,2,3}{
   \pgfmathsetmacro{\xl}{\k/4}
   \pgfmathsetmacro{\xr}{(\k+1)/4}
   \pgfmathsetmacro{\hh}{(\k/4)^2}
   \filldraw[fill=blue!16,draw=blue!60] (\xl,0) rectangle (\xr,\hh);
  }
  \draw[->] (-0.08,0)--(1.13,0) node[right]{$x$};
  \draw[->] (0,-0.05)--(0,1.17) node[above]{$y$};
  \draw[thick,domain=0:1,samples=60] plot (\x,{\x*\x});
  \node[below] at (1,0) {$1$};
  \node at (0.5,-0.23) {左端の高さ：$L_4$};
 \end{scope}
 \begin{scope}[xshift=5.4cm]
  \foreach \k in {0,1,2,3}{
   \pgfmathsetmacro{\xl}{\k/4}
   \pgfmathsetmacro{\xr}{(\k+1)/4}
   \pgfmathsetmacro{\hh}{((\k+1)/4)^2}
   \filldraw[fill=red!13,draw=red!60] (\xl,0) rectangle (\xr,\hh);
  }
  \draw[->] (-0.08,0)--(1.13,0) node[right]{$x$};
  \draw[->] (0,-0.05)--(0,1.17) node[above]{$y$};
  \draw[thick,domain=0:1,samples=60] plot (\x,{\x*\x});
  \node[below] at (1,0) {$1$};
  \node at (0.5,-0.23) {右端の高さ：$U_4$};
 \end{scope}
\end{tikzpicture}
\caption{下からの長方形和と上からの長方形和。分割を細かくすると両者の差が小さくなる。}
\label{図:導入長方形}
\end{figure}

正方形の和の公式
\[
 \sum_{k=1}^{n}k^2=\frac{n(n+1)(2n+1)}6
\]
を使うと，
\begin{align*}
 U_n&=\frac{n(n+1)(2n+1)}{6n^3}
     =\frac13+\frac{1}{2n}+\frac{1}{6n^2},\\
 L_n&=\frac{(n-1)n(2n-1)}{6n^3}
     =\frac13-\frac{1}{2n}+\frac{1}{6n^2}
\end{align*}
となる。面積は両者の間にあるから，
\[
 L_n\le S\le U_n
\]
である。そして $n$ を大きくすれば，$L_n$ も $U_n$ も $\frac13$ に近づく。したがって $S=\frac13$ と考えられる。この議論を，曲線に囲まれた面積の定義まで含めて明確にするのが定積分である。

二つの和の差は，展開せずに求めることもできる。同じ高さの項が打ち消し合うので，
\[
 U_n-L_n
 =\frac1n\left(1^2-0^2\right)=\frac1n
\]
となる。つまり，求めたい値を幅 $\frac1n$ の範囲に閉じ込められる。このように\textbf{値そのものをただ予想するだけでなく，どれだけの幅で挟めたかを調べる}と，近似の確かさが分かる。

\begin{example}{一回の近似から誤差の範囲を読む}
上の放物線の面積を $10$ 等分した長方形で挟め。さらに，右端の和 $U_{10}$ を面積の近似値としたとき，どの程度大きく見積もる可能性があるか説明せよ。
\end{example}
$n=10$ を代入すると，
\[
 L_{10}=\frac13-\frac1{20}+\frac1{600}
       =\frac{57}{200}=0.285,
 \qquad
 U_{10}=\frac{77}{200}=0.385
\]
である。したがって $0.285\le S\le0.385$ と分かる。$U_{10}$ は常に面積以上であり，
\[
 0\le U_{10}-S\le U_{10}-L_{10}=0.1
\]
である。この情報だけからでも，大きく見積もる誤差は $0.1$ 以下と保証できる。実際の面積 $\frac13$ を知った後なら，誤差は $\frac{31}{600}$ と計算できるが，未知の値を求める途中でも使えるのは前者の評価である。

\section{面積の増え方に現れる高さ}\label{節:導入基本定理}
ここまでの二つの問いは，一見すると別の話である。接線では縦の変化を横の変化で割り，面積では縦の長さと横の長さを掛けて足している。両者をつなぐために，右端を固定せず動かしてみる。

$y=x^2$ の下で，$0$ から $x$ までの面積を $A(x)$ と書く。ただし，この段階では $x\ge0$ とする。右端を $x$ から $x+h$ へ少し動かすと，面積の増加 $A(x+h)-A(x)$ は，幅 $h$ の細い帯の面積である。$h>0$ なら，その帯の高さは $x^2$ 以上 $(x+h)^2$ 以下だから，
\[
 x^2h\le A(x+h)-A(x)\le(x+h)^2h.
\]
正の数 $h$ で割ると，
\[
 x^2\le\frac{A(x+h)-A(x)}h\le(x+h)^2.
\]
$h$ を $0$ に近づけると左右はともに $x^2$ へ近づく。面積の平均変化率も，その間にあるので $x^2$ へ近づくはずである。つまり，\textbf{面積の瞬間の増加率は，その場所の曲線の高さになる}。

この結び付きによって，長方形の和を毎回計算せずに面積を求める道が開ける。もし $A(x)=\frac{x^3}{3}$ なら，
\begin{align*}
 \frac{A(x+h)-A(x)}h
 &=\frac{(x+h)^3-x^3}{3h}\\
 &=x^2+xh+\frac{h^2}{3}
\end{align*}
であり，確かに $h$ を $0$ へ近づけると $x^2$ になる。さらに $A(0)=0$ も満たす。この計算は面積の候補が条件に合うことを示しているが，その条件で面積が一つに決まることはまだ証明していない。後で平均値の定理を学ぶと，「微分して同じ関数になる二つの関数の差は定数である」と示せる。それを使って，この候補が面積そのものであると結論できる。

ここでは，候補を見つける計算と，候補が正しいと確かめる論理を分けておこう。積分の計算で「微分の逆を探す」という方針を使えるのは，この両方がそろうからである。

\section{符号と単位を読み落とさない}
グラフの高さが負になる場合，幅と高さの積も負になる。幾何学的な面積は負にはならないが，積分で足し合わせる量には符号が付く。この違いは，運動を考えると自然に理解できる。

右向きを正とする直線上で，物体が毎秒 $2$ メートルの速度で $3$ 秒進み，その後，毎秒 $1$ メートルの速度で逆向きに $2$ 秒進んだとする。位置の変化は
\[
 2\times3+(-1)\times2=4
\]
メートルである。一方，実際に動いた道のりは
\[
 2\times3+1\times2=8
\]
メートルとなる。速度のグラフで符号を付けて面積を足すと位置の変化が得られ，速さのグラフの面積を足すと道のりが得られる。同じグラフを使っていても，何を求めているかによって計算が変わる。

単位も計算の意味を教えてくれる。位置を時間で割れば速度になり，速度に時間を掛ければ位置の変化になる。面積を横の長さで割れば縦の長さになり，縦の長さに横の長さを掛ければ面積になる。微分や積分の式が複雑になっても，単位を確かめる習慣は役に立つ。

これから，まず変化率を求める方法を整え，続いて積分を定義する。その後の各章では，式をどのように変えれば既知の微分の逆を見つけられるか，あるいは不定積分を求めずに定積分の値を決められるかを調べる。最初からすべての公式を覚える必要はない。今扱っている式で，足しているのか，割っているのか，どの変数を変えているのかを確かめながら進めよう。

\section*{章末問題}
\begin{enumerate}[label=\textbf{\thechapter.\arabic*},leftmargin=*]
 \item\label{問:導入速度} $s(t)=t^2$ とする。時刻 $3$ から $3+h$ までの平均の速度を求めよ。$h=0.2,-0.2$ の場合の値を求め，時刻 $3$ の瞬間の速度を予想せよ。
 \item\label{問:導入直線面積} $y=2x+1$ と $x$ 軸，直線 $x=0,x=2$ に囲まれた台形の面積を求めよ。また，区間 $[0,2]$ を $n$ 等分して左端・右端の高さから作る長方形和 $L_n,U_n$ を求め，それぞれが同じ面積に近づくことを確かめよ。
 \item\label{問:導入中点} $y=x^2$ の下の $0\le x\le1$ の面積を，$4$ 等分した各小区間の中央の高さで近似せよ。得られた値を $L_4,U_4$ と比較せよ。
 \item\label{問:導入三次割線} $y=x^3$ のグラフ上の二点 $(1,1)$，$(1+h,(1+h)^3)$ を結ぶ直線の傾きを求めよ。$h\ne0$ の条件を使う箇所を明記し，$h$ を $0$ に近づけたときの傾きを予想せよ。
 \item\label{問:導入拡大面積} $y=x^2$ の下の $0\le x\le2$ の面積を $S$ とする。$4$ 等分した左端と右端の長方形和を求めて $S$ を挟め。一般の $n$ 等分では，$0\le x\le1$ の場合の長方形和の何倍になるか。そこから $S$ を予想せよ。
 \item\label{問:導入符号} 直線上を動く物体の速度が，時刻 $0$ から $4$ までは毎秒 $3$ メートル，時刻 $4$ から $7$ までは毎秒 $-2$ メートルである。位置の変化と道のりをそれぞれ求めよ。二つの答えが異なる理由を，速度の符号を使って説明せよ。
\end{enumerate}

\smallskip\noindent 解答は\pageref{解答:第1章}ページ。

\chapter{極限と微分}\label{章:極限と微分}
\section{数列の極限を言葉と式で捉える}\label{節:数列極限}
前章では，長方形の幅を小さくするときに $n$ を大きくし，割線を接線へ近づけるときに $h$ を $0$ に近づけた。どちらも，変化する数がある値に近づく様子を調べている。まず，番号 $n=1,2,3,\ldots$ に従って並んだ数列 $a_1,a_2,a_3,\ldots$ から考えよう。

\begin{definition}{数列の収束}
$n$ を十分大きくすれば，$a_n$ と一定の実数 $\alpha$ との差を，どのように小さく指定した正の数よりも小さくでき，それ以降も小さいままであるとき，数列 $\{a_n\}$ は $\alpha$ に\textbf{収束する}という。
\[
 \lim_{n\to\infty}a_n=\alpha
 \qquad\text{または}\qquad a_n\to\alpha\quad(n\to\infty)
\]
と書き，$\alpha$ を極限値という。
\end{definition}

「それ以降も」という部分は大切である。たとえば $a_n=(-1)^n$ は $1$ を何度も取るが，その間に $-1$ も取るので，$1$ に収束しない。極限は，一部の項だけが近いという意味ではない。また，$a_n=\frac1n$ はどの項も $0$ ではないが，$0$ に収束する。極限値を実際に取る必要もない。

定義を数式で書くと，「任意の $\varepsilon>0$ に対して，ある自然数 $N$ が存在し，$n\ge N$ なら $|a_n-\alpha|<\varepsilon$ となる」ということである。たとえば $a_n=\frac1n$ なら，$N>\frac1\varepsilon$ となる自然数を選べばよい。すると $n\ge N$ に対して
\[
 \left|\frac1n-0\right|=\frac1n\le\frac1N<\varepsilon
\]
が成り立つ。極限の記号は，このような誤差の条件を短く表している。

数列が限りなく大きくなる場合には $a_n\to+\infty$，限りなく小さくなる場合には $a_n\to-\infty$ と書く。ただし，$\infty$ は実数の極限値ではない。このような発散と，$(-1)^n$ のような振動による発散は区別する。

収束する数列については，和，差，積の極限をそれぞれの極限から計算できる。商では分母の極限が $0$ でないことが必要である。具体的には $a_n\to\alpha$，$b_n\to\beta$ のとき，
\[
 a_n+b_n\to\alpha+\beta,\qquad
 a_nb_n\to\alpha\beta,\qquad
 \frac{a_n}{b_n}\to\frac\alpha\beta\quad(\beta\ne0)
\]
となる。商の式は十分大きな $n$ で分母が $0$ でない範囲で使う。これらの極限法則と，次の挟み撃ちの原理は，極限の定義から証明できる基本性質として用いる。

\begin{theorem}{挟み撃ちの原理}
十分大きな $n$ について $a_n\le b_n\le c_n$ が成り立ち，$a_n\to\alpha$，$c_n\to\alpha$ なら，$b_n\to\alpha$ である。
\end{theorem}

実際，$a_n,c_n$ がともに $\alpha$ との差を $\varepsilon$ より小さくすれば，その間の $b_n$ も $\alpha-\varepsilon<b_n<\alpha+\varepsilon$ に入る。前章の長方形による面積の計算は，この原理を使う議論だった。

\begin{example}{極限を求めやすい形に直す}
次の極限を求めよ。
\[
 \text{(1)}\quad\lim_{n\to\infty}\frac{3n^2-n+2}{2n^2+5},
 \qquad
 \text{(2)}\quad\lim_{n\to\infty}\bigl(\sqrt{n^2+3n}-n\bigr).
\]
\end{example}
(1) は分子も分母も大きくなるため，そのままでは商の極限法則を適用できない。最も高い次数 $n^2$ で分子と分母を割ると，
\[
 \frac{3n^2-n+2}{2n^2+5}
 =\cfrac{3-\cfrac1n+\cfrac2{n^2}}{2+\cfrac5{n^2}}
 \longrightarrow\frac32.
\]
分母の極限が $2\ne0$ となったので，ここで商の法則を使える。

(2) は大きい数同士の差である。$\sqrt{n^2+3n}$ と $n$ がともに大きくなることだけでは，差の極限は決まらない。共役な式を掛けて差を積に変えると，
\begin{align*}
 \sqrt{n^2+3n}-n
 &=\frac{3n}{\sqrt{n^2+3n}+n}\\
 &=\cfrac3{\sqrt{1+\cfrac3n}+1}\longrightarrow\frac32.
\end{align*}
ここでは平方根の連続性も使った。大きい数同士の差が，有限の値へ近づく場合がある。

\subsection{単調で有界な数列}
数列が\textbf{上に有界}であるとは，すべての $n$ について $a_n\le M$ を満たす一定の実数 $M$ が存在することをいう。下に有界であることも，同様に一定の下限によって定義する。増加を続けていても，上に有界なら無限に大きくなることはできない。

\begin{theorem}{単調有界数列の収束：採用する事実}
単調増加する数列が上に有界なら，有限の極限値へ収束する。単調減少する数列が下に有界なら，同じく有限の極限値へ収束する。
\end{theorem}

この定理の完全な証明には，実数の連続性に関する性質が必要になる。本書では実数の基礎理論までさかのぼらず，この定理を採用する。ただし，適用するときに単調性と有界性の両方を確かめることは省略しない。収束すると予想した式に極限値を代入するだけでは，収束の証明にはならないからである。

\section{等比数列と無限級数}\label{節:等比級数}
等比数列 $1,r,r^2,\ldots$ の有限和を $S_n=1+r+\cdots+r^{n-1}$ とおく。$rS_n=r+r^2+\cdots+r^n$ を引けば，中間の項が消えるため，
\[
 (1-r)S_n=1-r^n,
 \qquad
 S_n=\frac{1-r^n}{1-r}\quad(r\ne1)
\]
を得る。まず有限の項を足した式を確定させ，その後で $n$ を大きくするのが順序である。

$|r|<1$ なら $r^n\to0$ なので，
\[
 \sum_{k=0}^{\infty}r^k
 :=\lim_{n\to\infty}\sum_{k=0}^{n-1}r^k
 =\frac1{1-r}
\]
となる。無限個の項を一度に足しているのではなく，\textbf{部分和の数列の極限}を定義している。一般に $\sum_{k=1}^{\infty}a_k$ が収束するとは，部分和 $\sum_{k=1}^{n}a_k$ が有限の極限値をもつことをいう。

なぜ $|r|<1$ で $r^n\to0$ になるのかも確認できる。$0<q=|r|<1$ とし，$\frac1q=1+c$，$c>0$ とおく。二項定理から $(1+c)^n\ge1+nc$ なので，
\[
 0\le |r^n|=q^n\le\frac1{1+nc}\longrightarrow0.
\]
$r=0$ の場合は直接明らかである。

$r=1$ なら部分和は $n$ となり，$r=-1$ なら部分和は $1,0,1,0,\ldots$ と振動する。$|r|>1$ では各項 $r^n$ さえ $0$ に近づかないため，級数は収束しない。級数が収束するなら，隣り合う部分和の差である一般項は必ず $0$ に近づく。この条件は必要条件であり，一般の級数では十分条件ではない。たとえば調和級数の発散は\chapterref{章:積分と数列}で積分との比較から調べる。

\begin{example}{無限和の値と打切り誤差}
$\sum_{k=0}^{\infty}\left(-\frac13\right)^k$ の和を求めよ。また，$k=0$ から $n-1$ までの和で近似するとき，誤差の絶対値を $10^{-4}$ 未満にする自然数 $n$ を一つ挙げよ。
\end{example}
公比の絶対値は $\frac13<1$ だから，和は
\[
 S=\cfrac{1}{1+\cfrac13}=\frac34.
\]
有限和と極限値との差は，有限和の公式から
\[
 S-S_n=\cfrac{\left(-\cfrac13\right)^n}{1+\cfrac13},
 \qquad
 |S-S_n|=\frac{3}{4\cdot3^n}.
\]
$n=9$ では $\frac3{4\cdot19683}<10^{-4}$ となるので，$9$ 項まで足せばよい。交互に符号が変わるため誤差の符号も変わるが，誤差の大きさはこの一つの式で評価できる。

\section{関数の極限と連続性}\label{節:関数極限連続}
数列では自然数 $n$ を大きくした。関数では，実数 $x$ をある点 $a$ に近づける。$x\ne a$ のまま $x$ を $a$ に十分近づければ，$f(x)$ を $L$ にどれだけでも近づけられるとき，
\[
 \lim_{x\to a}f(x)=L
\]
と書く。正確には，任意の $\varepsilon>0$ に対して，ある $\delta>0$ を選び，$0<|x-a|<\delta$ なら $|f(x)-L|<\varepsilon$ となることである。関数の定義域に属する $x$ だけを考える。

$x>a$ の側から近づける右極限を $\lim_{x\to a+}f(x)$，$x<a$ の側から近づける左極限を $\lim_{x\to a-}f(x)$ と書く。両側で近づけられる点では，左右の極限が同じ値に一致して初めて両側の極限が存在する。たとえば $\frac{|x|}{x}$ は $x\to0+$ で $1$，$x\to0-$ で $-1$ だから，$x\to0$ の極限は存在しない。

\begin{definition}{連続性}
$f$ が $x=a$ で定義され，$\lim_{x\to a}f(x)=f(a)$ となるとき，$f$ は $a$ で\textbf{連続}であるという。閉区間の端点では，区間の内側からの片側極限を使う。
\end{definition}

極限は $x=a$ 以外の近くの値から決まるのに対し，連続性はその極限と実際の値 $f(a)$ が一致するかどうかも問う。多項式，分母が $0$ でない有理関数，および三角関数・指数関数・対数関数は，それぞれの定義域で連続である。和・積・商・合成による連続性の保存も基本性質として用いる。

\begin{example}{値を補って連続にする}
$x\ne2$ に対して $f(x)=\frac{x^2-4}{x-2}$ と定める。$f(2)$ をどのように定めれば，$x=2$ でも連続になるか。
\end{example}
$x\ne2$ の範囲では約分でき，
\[
 f(x)=\frac{(x-2)(x+2)}{x-2}=x+2.
\]
したがって $\lim_{x\to2}f(x)=4$ である。$f(2)=4$ と定めれば極限と関数値が一致し，連続になる。約分後の $x+2$ は $x=2$ でも値をもつが，元の分数はそのままでは定義されない。式が等しい範囲と，関数の定義域を分けて考える必要がある。

後で微分の定理を証明するために，連続関数について次の二つの存在定理を使う。

\begin{theorem}{連続関数の存在定理：採用する事実}\label{定理:連続関数存在}
$f$ が閉区間 $[a,b]$ 上で連続なら，次が成り立つ。
\begin{enumerate}[label=(\arabic*),leftmargin=*]
 \item $f$ は区間内で最大値と最小値をとる。
 \item $f(a)$ と $f(b)$ の間にある任意の値 $k$ に対して，$f(c)=k$ となる $c\in[a,b]$ が存在する。これを\textbf{中間値の定理}という。
\end{enumerate}
\end{theorem}

これらの証明も実数の連続性に基づくため，本書では採用する事実とする。条件を外すと結論は崩れる。たとえば $f(x)=x$ は開区間 $(0,1)$ では最大値も最小値もとらない。端の値 $0,1$ に近づくことはできるが，区間の中には到達する点がない。閉区間という条件は，単なる表記上の飾りではない。

\section{微分係数と導関数}\label{節:微分定義}
\begin{definition}{微分係数と導関数}
$f$ が $a$ の近くで定義され，有限の極限
\[
 f'(a)=\lim_{h\to0}\frac{f(a+h)-f(a)}h
\]
が存在するとき，$f$ は $a$ で\textbf{微分可能}であるといい，$f'(a)$ を $a$ における微分係数という。各点の微分係数を対応させた関数 $f'(x)$ を\textbf{導関数}という。
\end{definition}

微分係数は一つの点での数，導関数は点を動かしたときの関数である。$y=f(x)$ のとき $f'(x)$ を $y'$ や $\frac{\dd y}{\dd x}$ とも書く。$\frac{\dd y}{\dd x}$ は通常の分数そのものではないが，「何を何で微分しているか」を明示する記号である。

$f(x)=x^2$ なら，
\[
 \frac{f(a+h)-f(a)}h=2a+h\longrightarrow2a
\]
なので，$f'(a)=2a$，導関数は $f'(x)=2x$ である。点 $(a,f(a))$ における接線は，傾き $f'(a)$ をもつ直線として
\[
 y-f(a)=f'(a)(x-a)
\]
で表される。

微分可能なら連続である。実際，$h\ne0$ に対して
\[
 f(a+h)-f(a)
 =h\frac{f(a+h)-f(a)}h
\]
であり，右辺は $h\to0$ のとき $0\cdot f'(a)=0$ に近づく。一方，連続だから微分可能とは限らない。$f(x)=|x|$ は $0$ で連続だが，差商 $\frac{|h|}{h}$ の右極限は $1$，左極限は $-1$ なので，$0$ で微分できない。折れ曲がった点に一つの接線の傾きを与えようとしても，左右の候補が一致しないのである。

\subsection{微分を一次式で表す}
微分可能性は，
\[
 f(a+h)=f(a)+f'(a)h+h\varepsilon(h),
 \qquad \varepsilon(h)\to0\quad(h\to0)
\]
と書き直せる。実際，$h\ne0$ に対して $\varepsilon(h)=\frac{f(a+h)-f(a)}h-f'(a)$ とおけばよい。最後の $h\varepsilon(h)$ は，$h$ で割ると $0$ に近づく誤差である。「小さい変化に対して関数を接線で近似できる」という意味が，ここに数式で表されている。この形は合成関数の微分を証明するときにも役立つ。

\section{積・商・合成関数の微分}\label{節:微分法則}
基本的な関数を一つずつ定義から微分しても，複雑な式を毎回展開するのは大変である。そこで，式の作られ方と微分の関係を先に調べる。

\begin{theorem}{和・積・商の微分}
$f,g$ が微分可能で，$c$ が定数なら，
\begin{align*}
 (f+g)'&=f'+g',& (cf)'&=cf',\\
 (fg)'&=f'g+fg',&
 \left(\frac fg\right)'&=\frac{f'g-fg'}{g^2}\quad(g\ne0)
\end{align*}
が成り立つ。
\end{theorem}

和と定数倍は差商を分ければ分かる。積の公式では，二つの因子が同時に変わることに注意する。途中の項を一つ足して引けば，
\begin{align*}
 &f(x+h)g(x+h)-f(x)g(x)\\
 &\quad=\{f(x+h)-f(x)\}g(x+h)
       +f(x)\{g(x+h)-g(x)\}.
\end{align*}
これを $h$ で割って極限をとる。微分可能な $g$ は連続だから $g(x+h)\to g(x)$ となり，積の公式が得られる。積を微分して単に $f'g'$ としてはいけない。

逆数については，
\[
 \cfrac{\cfrac1{g(x+h)}-\cfrac1{g(x)}}h
 =-\frac{g(x+h)-g(x)}{h\,g(x+h)g(x)}
 \longrightarrow-\frac{g'(x)}{g(x)^2}.
\]
$g(x)\ne0$ と連続性により，十分近くでは $g(x+h)\ne0$ となる。この逆数の公式と積の公式から，商の公式も導ける。

\begin{theorem}{合成関数の微分}\label{定理:合成微分}
$u=g(x)$ が $x$ で微分可能で，$f$ が $u=g(x)$ で微分可能なら，
\[
 \frac{\dd}{\dd x}f(g(x))=f'(g(x))g'(x)
\]
である。
\end{theorem}

外側の関数の変化率に，内側の変数の変化率を掛ける公式である。証明のため，$\Delta u=g(x+h)-g(x)$ とおく。$g$ の連続性から $\Delta u\to0$ であり，$f$ の微分可能性から
\[
 f(u+\Delta u)-f(u)=f'(u)\Delta u+\Delta u\,\eta(\Delta u),
 \qquad \eta(\Delta u)\to0
\]
と書ける。$\eta(0)=0$ と定めれば，$\Delta u=0$ の場合もこの等式は使える。$h$ で割ると
\[
 \frac{f(g(x+h))-f(g(x))}h
 =\{f'(g(x))+\eta(\Delta u)\}
   \frac{g(x+h)-g(x)}h.
\]
右辺の極限は $f'(g(x))g'(x)$ になる。途中で $\Delta u$ を分母に置いていないので，内側の関数が同じ値を取る場合も含めて証明できた。

\section{基本関数を微分する}\label{節:基本微分}
\subsection{べき関数}
正の整数 $n$ について二項定理を使うと，
\[
 \frac{(x+h)^n-x^n}{h}
 =nx^{n-1}+\binom n2x^{n-2}h+\cdots+h^{n-1}.
\]
$h\to0$ とすれば，
\[
 \frac{\dd}{\dd x}x^n=nx^{n-1}
\]
を得る。定数の微分は $0$ であり，負の整数べきは逆数の微分から同じ形になる。実数べきへの拡張は，対数関数を微分した後で説明する。

\subsection{三角関数と弧度法}
三角関数を微分するとき，角は必ず弧度法で表す。基本になる極限は
\[
 \lim_{h\to0}\frac{\sin h}{h}=1
\]
である。$0<h<\frac\pi2$ では，単位円における三角形，扇形，外側の三角形の面積比較から
\[
 \sin h<h<\tan h
\]
を得る。したがって $\cos h<\frac{\sin h}{h}<1$ であり，挟み撃ちの原理から右極限は $1$ となる。$\frac{\sin h}{h}$ は $h$ の符号を変えても同じだから，左極限も $1$ である。また，
\[
 \frac{\cos h-1}{h}
 =-\cfrac{2\sin^2\cfrac h2}{h}
 =-\sin\frac h2\,\cfrac{\sin\cfrac h2}{\cfrac h2}
 \longrightarrow0.
\]
加法定理に代入すれば，
\begin{align*}
 \frac{\sin(x+h)-\sin x}{h}
 &=\sin x\frac{\cos h-1}{h}
   +\cos x\frac{\sin h}{h}\\
 &\longrightarrow\cos x.
\end{align*}
同様に余弦を微分でき，商の公式から正接も求まる。
\[
 (\sin x)'=\cos x,\qquad
 (\cos x)'=-\sin x,\qquad
 (\tan x)'=\frac1{\cos^2x}\quad(\cos x\ne0).
\]
度数法で $x$ を表すと，弧度への換算係数が微分に現れる。上の簡潔な公式が成り立つのは，弧度法を使っているからである。

\subsection{指数関数と対数関数}
指数関数を微分するためには，底を特別な値 $\e$ に選ぶ。
\begin{theorem}{自然対数の底と基本極限：採用する事実}
数列 $\left(1+\frac1n\right)^n$ は極限 $\e$ をもち，$2<\e<3$ である。この底に対する指数関数は
\[
 \lim_{h\to0}\frac{\e^h-1}{h}=1
\]
を満たす。$\e$ を底とする対数を自然対数といい，本書では $\log x$ と書く。
\end{theorem}

この極限を実数指数の定義から構成する議論は，本書の証明範囲から外す。以後の微分公式は，この基本極限と指数法則から導く。$\e^{x+h}=\e^x\e^h$ なので，
\[
 \frac{\e^{x+h}-\e^x}{h}
 =\e^x\frac{\e^h-1}{h}\longrightarrow\e^x.
\]
よって $(\e^x)'=\e^x$ である。

対数では $x>0$ とし，$k=\log(x+h)-\log x$ とおく。対数の連続性から $h\to0$ のとき $k\to0$ である。また $x+h=x\e^k$ だから，$h=x(\e^k-1)$ となる。したがって
\[
 \frac{\log(x+h)-\log x}{h}
 =\frac1x\frac{k}{\e^k-1}\longrightarrow\frac1x.
\]
以上と合成関数の微分から，
\begin{align*}
 (\e^{u(x)})'&=\e^{u(x)}u'(x),\\
 (\log u(x))'&=\frac{u'(x)}{u(x)}\quad(u(x)>0),\\
 (a^x)'&=a^x\log a\quad(a>0)
\end{align*}
を得る。最後は $a^x=\e^{x\log a}$ を使った。さらに $u(x)\ne0$ の区間では，$u$ の符号を分けると
\[
 (\log|u(x)|)'=\frac{u'(x)}{u(x)}
\]
も成り立つ。絶対値を付けることで，負の値の区間でも実数の対数として定義できる。

$x>0$ では実数 $\alpha$ に対して $x^\alpha=\e^{\alpha\log x}$ と書けるため，
\[
 (x^\alpha)'=\e^{\alpha\log x}\frac\alpha x
 =\alpha x^{\alpha-1}
\]
である。整数べきなら，元の関数が定義される負の $x$ にも既に同じ公式が成り立つ。非整数べきの場合は，まず $x>0$ を基本の定義域として考える。

\begin{example}{式の構造に合わせて微分する}
次の関数を微分し，(1) については $x=2$ における接線も求めよ。
\[
 \text{(1)}\quad f(x)=\frac{x^2+1}{x-1},\qquad
 \text{(2)}\quad g(x)=\e^{x^2}\sin(3x),\qquad
 \text{(3)}\quad y=x^x\quad(x>0).
\]
\end{example}
(1) に商の公式を使うと，
\[
 f'(x)=\frac{2x(x-1)-(x^2+1)}{(x-1)^2}
      =\frac{x^2-2x-1}{(x-1)^2}.
\]
$f(2)=5$，$f'(2)=-1$ なので，接線は $y-5=-(x-2)$，すなわち $y=-x+7$ である。別の計算として，多項式の割り算から
\[
 f(x)=x+1+\frac2{x-1},\qquad
 f'(x)=1-\frac2{(x-1)^2}
\]
としてもよい。微分する前に式を整理すると，計算が短くなる場合がある。

(2) は二つの関数の積であり，それぞれの中に合成がある。まず外側の積の公式を使い，次に内側を微分する。
\[
 g'(x)=2x\e^{x^2}\sin(3x)+3\e^{x^2}\cos(3x).
\]
$2x$ と $3$ は，それぞれ $x^2$ と $3x$ を微分した因子である。

(3) は指数も底も変わるので，定数指数の公式をそのまま使えない。$y>0$ だから対数をとり，$\log y=x\log x$ とすると，
\[
 \frac{y'}y=\log x+1,
 \qquad y'=x^x(\log x+1).
\]
このように対数をとって積やべきを分けてから微分する方法を\textbf{対数微分法}という。対数をとれる条件 $x>0$，$y>0$ を先に確かめている点も大切である。

\section{平均値の定理が局所と区間を結ぶ}\label{節:平均値定理}
微分係数は一点での変化率である。一方，関数が区間全体で増えるかどうかを知るには，離れた二点の値を比べる必要がある。この二つをつなぐのが平均値の定理である。証明の準備として，極値で何が起こるかを確かめよう。

点 $c$ の十分近くで $f(x)\le f(c)$ なら，$f$ は $c$ で\textbf{極大値}をとるという。不等号を逆にしたものが\textbf{極小値}である。極大・極小は近くの点との比較であり，区間全体での最大・最小とは区別する。

$c$ が区間の内点で，$f$ がそこで微分可能かつ極大なら，小さい $h$ に対して $f(c+h)-f(c)\le0$ である。したがって差商は $h>0$ で $0$ 以下，$h<0$ で $0$ 以上となる。左右の極限が一致するには，$f'(c)=0$ でなければならない。極小の場合も同じである。ただし，端点では左右両側を比較できず，この結論をそのまま使えない。

\begin{theorem}{ロルの定理}\label{定理:ロル}
$a<b$ とする。$f$ が閉区間 $[a,b]$ で連続，開区間 $(a,b)$ で微分可能で，$f(a)=f(b)$ なら，
\[
 f'(c)=0
\]
を満たす $c\in(a,b)$ が少なくとも一つ存在する。
\end{theorem}

定理\ref{定理:連続関数存在}により，$f$ は最大値と最小値をとる。$f$ が定数なら，どの内点でも導関数は $0$ である。定数でなければ，最大値または最小値の少なくとも一方は共通の端点値 $f(a)=f(b)$ と異なる。その値をとる点は内点なので，先ほどの議論から，その点で $f'(c)=0$ となる。これで証明できた。

\begin{theorem}{平均値の定理}\label{定理:平均値}
$a<b$ とする。$f$ が閉区間 $[a,b]$ で連続，開区間 $(a,b)$ で微分可能なら，
\[
 f'(c)=\frac{f(b)-f(a)}{b-a}
\]
を満たす $c\in(a,b)$ が少なくとも一つ存在する。
\end{theorem}

右辺は二つの端点を結ぶ割線の傾きである。そこで，その傾きの直線を $f$ から引き，両端の高さをそろえる。
\[
 g(x)=f(x)-\frac{f(b)-f(a)}{b-a}(x-a)
\]
とおけば，$g(a)=g(b)=f(a)$ である。また，$g$ は閉区間で連続，内側で微分可能だから，ロルの定理を使える。ある内点 $c$ で
\[
 0=g'(c)=f'(c)-\frac{f(b)-f(a)}{b-a}
\]
となり，求める等式を得る。

\begin{figure}[htbp]
\centering
\begin{tikzpicture}[x=2.6cm,y=1.1cm,>=Stealth]
 \draw[->] (-0.15,0)--(2.35,0) node[right]{$x$};
 \draw[->] (0,-0.4)--(0,4.7) node[above]{$y$};
 \draw[very thick,domain=0:2.12,samples=70] plot (\x,{\x*\x});
 \draw[blue,thick] (0,0)--(2,4);
 \draw[red,thick,domain=0.4:1.8] plot (\x,{2*\x-1});
 \fill (0,0) circle(1.6pt);
 \fill (2,4) circle(1.6pt) node[above left]{$(2,4)$};
 \fill (1,1) circle(1.6pt) node[below right]{$(1,1)$};
 \draw[dashed] (1,0) node[below]{$1$}--(1,1);
 \draw[dashed] (2,0) node[below]{$2$}--(2,4);
 \node[blue] at (0.63,2.05) {割線};
 \node[red] at (1.8,2.2) {接線};
 \node at (1.35,4.55) {$y=x^2$};
\end{tikzpicture}
\caption{平均変化率に等しい微分係数をもつ点が，区間の内側に存在する。}
\label{図:平均値定理}
\end{figure}

図\ref{図:平均値定理}では，$f(x)=x^2$ の区間 $[0,2]$ における平均変化率は $2$ であり，$f'(1)=2$ となる。一般には，その点が区間の中央とは限らず，一つに決まるとも限らない。

条件の役割も確認しておこう。$f(x)=|x|$ を $[-1,1]$ で考えると，両端の値は等しいが，微分可能な点の導関数は $-1$ または $1$ である。$0$ では微分できないため，ロルの定理も平均値の定理も適用できない。導関数の連続性は定理の仮定に含まれていないが，区間内のすべての点で微分可能であることは必要である。

\subsection{導関数の符号から増減が分かる理由}
区間内の任意の二点 $x_1<x_2$ に平均値の定理を使うと，その間のある $c$ に対して
\[
 f(x_2)-f(x_1)=f'(c)(x_2-x_1)
\]
となる。$x_2-x_1>0$ だから，区間内で $f'>0$ なら $f(x_2)>f(x_1)$ であり，$f$ は狭義に増加する。$f'<0$ なら狭義に減少する。$f'\ge0$ の場合には $f(x_2)\ge f(x_1)$ までが保証され，値が一定の部分もあり得る。

同じ議論から，次の結果も得られる。
\begin{theorem}{導関数が同じ関数の差}\label{定理:導関数同じ}
一つの区間のすべての内点で $F'=G'$ が成り立ち，端点を含む場合にはそこで両関数が連続なら，その区間で $F-G$ は定数である。特に，導関数が恒等的に $0$ なら，元の関数は定数である。
\end{theorem}

実際，$H=F-G$ とおくと $H'=0$ であり，任意の二点について $H(x_2)-H(x_1)=0$ となる。後の積分では，微分して同じ関数になる候補が，定数の違いを除いて同じであることをこの定理で確かめる。「一つの区間」という条件にも意味がある。たとえば $x<0$ で $H(x)=0$，$x>0$ で $H(x)=1$ と定めれば，定義域のすべての点で $H'=0$ だが，左右の区間に同じ定数を使うことはできない。

\section{増減・凹凸からグラフを読む}\label{節:増減凹凸}
$f'(c)=0$ は極値の候補を探す条件であり，極値を保証する条件ではない。$f(x)=x^3$ は $f'(0)=0$ だが，$0$ の前後でも増加し続ける。極大・極小を判定するには，その点の前後で導関数の符号がどう変わるかを調べる。$+$ から $-$ に変われば極大，$-$ から $+$ に変われば極小である。

導関数をもう一度微分した $f''=(f')'$ を\textbf{第二次導関数}という。$f''>0$ なら接線の傾き $f'$ が増加し，グラフは\textbf{下に凸}になる。$f''<0$ なら傾きが減少し，グラフは\textbf{上に凸}になる。たとえば $y=x^2$ は下に凸で，$y=-x^2$ は上に凸である。

この呼び方を，接線との位置関係から確かめられる。$f''>0$ の区間で $a$ を固定し，$x>a$ とすると，平均値の定理により
\[
 f(x)-f(a)=f'(c)(x-a)>f'(a)(x-a)
 \qquad(a<c<x)
\]
となる。$x<a$ の場合も，$x<c<a$ で $f'(c)<f'(a)$ となることと $x-a<0$ を合わせれば，同じ結論を得る。したがって，曲線は接線の上側にある。$f''<0$ では不等号が逆になる。

グラフの凹凸が切り替わる点を\textbf{変曲点}という。$f''$ が連続な場合，凹凸が切り替わる場所は $f''=0$ となる点から探せるが，$f''=0$ だけでは十分でない。$y=x^4$ は $x=0$ で $f''=0$ になっても，その前後でともに下に凸である。

\begin{example}{極値と区間全体の最大・最小}
$f(x)=x^3-3x$ の増減と凹凸を調べ，区間 $[-2,3]$ における最大値・最小値を求めよ。
\end{example}
導関数を因数分解すると，符号を調べやすい。
\[
 f'(x)=3x^2-3=3(x-1)(x+1),\qquad f''(x)=6x.
\]
したがって $x<-1$ で増加，$-1<x<1$ で減少，$x>1$ で増加する。$f(-1)=2$ が極大値，$f(1)=-2$ が極小値である。また，$x<0$ で上に凸，$x>0$ で下に凸だから，原点 $(0,0)$ は変曲点である。
\[
\begin{array}{c|ccccccc}
 x&-2&&-1&&1&&3\\ \hline
 f'(x)&&+&0&-&0&+&\\
 f(x)&-2&\nearrow&2&\searrow&-2&\nearrow&18
\end{array}
\]
閉区間での最大・最小は，極値だけを比較しては決まらない。両端も候補になるため，$f(-2)=-2$，$f(3)=18$ を加えて比較する。最大値は $18$ で $x=3$ のとき，最小値は $-2$ で $x=-2,1$ のときにとる。

\begin{figure}[htbp]
\centering
\begin{tikzpicture}[x=2.1cm,y=0.8cm,>=Stealth]
 \draw[->] (-2.3,0)--(2.35,0) node[right]{$x$};
 \draw[->] (0,-3.4)--(0,3.5) node[above]{$y$};
 \draw[very thick,domain=-2.1:2.1,samples=120] plot (\x,{\x^3-3*\x});
 \fill (-1,2) circle(1.5pt) node[above left]{$(-1,2)$};
 \fill (1,-2) circle(1.5pt) node[below right]{$(1,-2)$};
 \fill (0,0) circle(1.5pt) node[above right]{変曲点};
 \draw[dashed] (-1,0) node[below]{$-1$}--(-1,2)--(0,2);
 \draw[dashed] (1,0) node[above]{$1$}--(1,-2)--(0,-2);
 \node at (1.1,2.85) {$y=x^3-3x$};
\end{tikzpicture}
\caption{極値は増減の切り替わり，変曲点は凹凸の切り替わりを表す。}
\label{図:三次増減}
\end{figure}

図\ref{図:三次増減}は原点付近を描いたものである。見えている山の高さ $2$ が，指定区間全体の最大値とは限らないことに注意しよう。最大・最小を求めるときは，まず定義域と対象の区間を確かめ，内点の極値候補，微分できない点，含まれる端点を比較する。連続関数でも，区間が閉じていなければ最大・最小が存在しない場合がある。

\section{方程式と不等式に微分を使う}\label{節:微分方程式不等式}
方程式の根を式で表せなくても，根が存在することや，根の個数は分かることがある。ここでは，連続性が存在を保証し，増減が個数を制限する。不等式では，左右の差を一つの関数にして，その最小値を調べる。

\begin{example}{解の存在と不等式の等号条件}
\begin{enumerate}[label=(\arabic*),leftmargin=*]
 \item 方程式 $x=\cos x$ が区間 $[0,1]$ にただ一つの解をもつことを示せ。
 \item $x>0$ のとき $\log x\le x-1$ を示し，等号成立条件を求めよ。
\end{enumerate}
\end{example}
(1) では $h(x)=x-\cos x$ とおく。$h$ は $[0,1]$ で連続で，
\[
 h(0)=-1<0,\qquad h(1)=1-\cos1>0
\]
である。中間値の定理から，$h(x)=0$ となる点が $(0,1)$ に存在する。一方，この区間では $h'(x)=1+\sin x>0$ だから，$h$ は狭義に増加する。二つの異なる点でともに $0$ をとることはできず，解はただ一つである。グラフの交わりを予想しただけではなく，存在と一意性を別々に示した。

(2) は右辺から左辺を引き，$g(x)=x-1-\log x$ とおく。
\[
 g'(x)=1-\frac1x=\frac{x-1}{x}.
\]
定義域では $x>0$ なので，符号は分子 $x-1$ で決まる。$g$ は $0<x<1$ で減少し，$x>1$ で増加する。したがって，最小値は $g(1)=0$ であり，$g(x)\ge0$ を得る。両側で狭義に増減しているため，等号は $x=1$ のときに限る。

\section{接線で値を近似する}\label{節:一次近似}
\sectionref{節:微分定義}で得た式
\[
 f(a+h)=f(a)+f'(a)h+h\varepsilon(h),\qquad \varepsilon(h)\to0
\]
を使えば，$h$ が小さいとき
\[
 f(a+h)\approx f(a)+f'(a)h
\]
と近似できる。これを\textbf{一次近似}という。右辺は $a$ における接線の値である。計算しやすい点 $a$ を近くに選び，そこからの変化を傾きで見積もる。

\begin{example}{平方根の近似と誤差の向き}
$\sqrt{4.04}$ を一次近似で求め，近似値が実際の値より大きいか小さいかを確かめよ。
\end{example}
$f(x)=\sqrt{x}$ とすると，$f(4)=2$，$f'(4)=\frac1{2\sqrt4}=\frac14$ である。$4.04=4+0.04$ なので，
\[
 \sqrt{4.04}\approx2+\frac14\cdot0.04=2.01.
\]
$f''(x)=-\cfrac1{4x^{\cfrac32}}<0$ だから，グラフは上に凸で，接線による値は実際の値より大きい。数値によっても，$2.01^2=4.0401>4.04$ と確かめられる。

誤差の大きさは，有理化によって直接評価できる。
\[
 2.01-\sqrt{4.04}
 =\frac{2.01^2-4.04}{2.01+\sqrt{4.04}}
 =\frac{0.0001}{2.01+\sqrt{4.04}}.
\]
分母は $4$ より大きいから，$0<2.01-\sqrt{4.04}<0.000025$ である。一次近似の式だけでは，指定した小数桁までの正しさが自動的に保証されるわけではない。この例では別の式で誤差を評価した。より一般的な誤差の扱いは\chapterref{章:近似と広義積分}で学ぶ。

\section{逆関数・媒介変数・陰関数を微分する}\label{節:逆媒介陰微分}
\subsection{逆関数では傾きが逆数になる}
一つの区間で連続かつ狭義単調な関数 $y=f(x)$ は，その値域で逆関数 $x=g(y)$ をもつ。このとき逆関数も連続であるという基本性質を用いる。その区間の内点 $a$ について $b=f(a)$，$f'(a)\ne0$ とし，$y=b+k$ に対応する $x$ を $a+h$ と書けば，$k\to0$ のとき $h\to0$ であり，$k=f(a+h)-f(a)$ である。したがって
\[
 \frac{g(b+k)-g(b)}k
 =\frac{h}{f(a+h)-f(a)}
 =\cfrac{1}{\cfrac{f(a+h)-f(a)}h}
 \longrightarrow\frac1{f'(a)}.
\]
よって
\[
 g'(b)=\frac1{f'(a)},\qquad
 g'(y)=\frac1{f'(g(y))}
\]
となる。分母が $0$ でないことが必要である。たとえば $f(x)=x^2$ を $x>0$ に限ると，逆関数は $g(y)=\sqrt y$ であり，$g'(y)=\frac1{2\sqrt y}$ となる。$x^2$ を実数全体で考えると，正負の二つの $x$ が同じ値に対応して逆関数が定まらないので，先にどちらの枝を使うかを決めなければならない。

\subsection{媒介変数で表した曲線}
曲線が $x=x(t)$，$y=y(t)$ と表されるとき，$t$ を\textbf{媒介変数}という。両関数が連続な導関数をもち，考えている点で $x'(t)\ne0$ なら，その近くで $x(t)$ は狭義単調であり，$t$ を $x$ の関数として表せる。合成関数と逆関数の微分から
\[
 \frac{\dd y}{\dd x}
 =\frac{y'(t)}{x'(t)}
\]
を得る。$x'(t)=0$ の点ではこの商は使えず，垂直接線などを別に調べる必要がある。さらに二階まで微分できる場合は，$x$ での微分が $\frac1{x'(t)}\frac{\dd}{\dd t}$ に対応するので，
\[
 \frac{\dd^2y}{\dd x^2}
 =\frac1{x'(t)}\frac{\dd}{\dd t}
   \left(\frac{y'(t)}{x'(t)}\right)
\]
となる。$\frac{y''(t)}{x''(t)}$ ではない。

\subsection{関係式を保ったまま微分する}
$x,y$ の関係が一つの式で与えられ，$y$ を $x$ の関数として明示していない場合がある。その式が実際に微分可能な関数 $y=y(x)$ を定めている範囲では，両辺を $x$ で微分できる。これを\textbf{陰関数の微分}という。$y$ も $x$ とともに変わるので，たとえば $(y^2)'=2yy'$，$(xy)'=y+xy'$ となる。

\begin{example}{二通りの表し方から接線を求める}
\begin{enumerate}[label=(\arabic*),leftmargin=*]
 \item $x=t^2+1$，$y=t^3-t$ で表される曲線の $t=1$ における接線を求めよ。
 \item 円 $x^2+y^2=25$ の点 $(3,4)$ における接線を求めよ。
\end{enumerate}
\end{example}
(1) では $t=1$ のとき $(x,y)=(2,0)$ である。また，
\[
 \frac{\dd x}{\dd t}=2t,\qquad
 \frac{\dd y}{\dd t}=3t^2-1,\qquad
 \left.\frac{\dd y}{\dd x}\right|_{t=1}
 =\frac{3-1}{2}=1.
\]
したがって接線は $y=x-2$ である。$t$ を消去しなくても，それぞれの変化率の比で傾きを求められた。

(2) では，点 $(3,4)$ の近くで上半円 $y=\sqrt{25-x^2}$ を考えれば，$y$ が微分可能な関数であることを確かめられる。その上で，元の関係式を微分すると
\[
 2x+2yy'=0,\qquad y'=-\frac xy\quad(y\ne0).
\]
点 $(3,4)$ での傾きは $-\frac34$ だから，接線は
\[
 y-4=-\frac34(x-3),\qquad 3x+4y=25
\]
である。円全体を一つの関数 $y=f(x)$ とみなしたわけではなく，指定点の近くの枝を使っている。$y=0$ の点では上の商は使えず，実際に接線は垂直になる。

これで，一点での変化率から，区間の増減や接線の近似までを調べる道具がそろった。次章では，細かい区間での寄与を足して定積分を作り，この微分の議論とつなげる。特に，平均値の定理から得た「導関数が同じなら差は定数」という結果が，積分の計算を支えることになる。

\section*{章末問題}
\begin{enumerate}[label=\textbf{\thechapter.\arabic*},leftmargin=*]
 \item\label{問:極限と漸化式} 次の極限を求めよ。また，$a_1=1$，$a_{n+1}=\sqrt{2+a_n}$ で定まる数列が単調増加かつ上に有界であることを示し，その極限を求めよ。
 \[
  \text{(1)}\quad\lim_{n\to\infty}\frac{2n^2+3}{5n^2-n},
  \qquad
  \text{(2)}\quad\lim_{n\to\infty}\bigl(\sqrt{n^2+2n}-n\bigr).
 \]
 \item\label{問:等比誤差} $\sum_{k=0}^{\infty}\left(-\frac12\right)^k$ の和を求めよ。$k=0$ から $n-1$ までの和で近似するとき，誤差の絶対値を $10^{-3}$ 未満にする最小の自然数 $n$ を求めよ。
 \item\label{問:連続と微分可能} $x\ne0$ で $f(x)=\frac{\sin2x}{x}$，$x=0$ で $f(0)=c$ とする。$f$ が $0$ で連続となる $c$ を求めよ。また，$g(x)=|x-1|$ は $1$ で連続だが微分可能でないことを，差商を使って説明せよ。
 \item\label{問:微分計算} 次の関数を微分せよ。(2) については $x=0$ における接線も求めよ。
 \[
  \text{(1)}\quad\log(1+x^2),\qquad
  \text{(2)}\quad(x^2+1)\e^{-x},\qquad
  \text{(3)}\quad x^{\sin x}\quad(x>0).
 \]
 \item\label{問:平均値評価} 平均値の定理を使い，任意の実数 $a,b$ に対して $|\sin b-\sin a|\le|b-a|$ が成り立つことを示せ。適用する関数と区間を明記せよ。
 \item\label{問:三次関数グラフ} $f(x)=x^3-3x^2$ の増減，極値，凹凸，変曲点を求めよ。さらに，区間 $[-1,3]$ における最大値・最小値と，それらをとる $x$ を求めよ。
 \item\label{問:指数不等式方程式} 微分を用いて，すべての実数 $x$ について $\e^x\ge1+x$ を示し，等号成立条件を求めよ。また，方程式 $\log x+x=2$ は $x>0$ でただ一つの解をもち，その解が $(1,2)$ にあることを示せ。
 \item\label{問:逆媒介陰} 次を求めよ。
 \begin{enumerate}[label=(\arabic*),leftmargin=*]
  \item $f(x)=x^3+x$ の逆関数を $g$ とするとき，$g'(0)$。
  \item $x=t^2$，$y=t^3$ で表される曲線の，$t=2$ における接線。
  \item 曲線 $x^2+xy+y^2=7$ の点 $(1,2)$ における接線。ただし，その点の近くで $y$ は $x$ の微分可能な関数として表されるものとする。
 \end{enumerate}
\end{enumerate}

\smallskip\noindent 解答は\pageref{解答:第2章}ページ。

\chapter{定積分と微積分の基本定理}\label{章:定積分}
\section{細かい区間の寄与を足す}\label{節:リーマン和}
\chapterref{章:導入}では，長方形の幅を小さくして面積を求めた。そのときは等しい幅に分け，左端や右端の高さを使った。しかし，同じ図形の面積が，幅の分け方や高さを選ぶ点によって変わっては困る。そこで，まず分け方と点の選び方を明示し，それらに依存しない極限が存在するときに定積分を定義する。

$a<b$ とし，区間 $[a,b]$ を
\[
 a=x_0<x_1<\cdots<x_n=b
\]
と分ける。この点の組を\textbf{分割} $P$ といい，第 $k$ 小区間の幅を $\Delta x_k=x_k-x_{k-1}$ と書く。さらに，各小区間 $[x_{k-1},x_k]$ の中から点 $\xi_k$ を一つずつ選ぶ。$f$ が $[a,b]$ で有界なら，
\[
 S(P,\xi)=\sum_{k=1}^{n}f(\xi_k)\Delta x_k
\]
を作れる。これを\textbf{リーマン和}という。$f(\xi_k)$ が正なら高さと幅の積であり，負なら符号を付けた寄与として足し合わせる。

\begin{figure}[htbp]
\centering
\begin{tikzpicture}[x=3cm,y=1.8cm,>=Stealth]
 \filldraw[fill=blue!12,draw=blue!55] (0,0) rectangle (0.8,1.1225);
 \filldraw[fill=blue!12,draw=blue!55] (0.8,0) rectangle (1.8,1.01);
 \filldraw[fill=blue!12,draw=blue!55] (1.8,0) rectangle (3,1.5625);
 \draw[->] (-0.1,0)--(3.35,0) node[right]{$x$};
 \draw[->] (0,-0.13)--(0,2.4) node[above]{$y$};
 \draw[very thick,domain=0:3.1,samples=80] plot (\x,{1+0.25*(\x-1)^2});
 \foreach \xx/\hh/\name in {0.3/1.1225/1,1.2/1.01/2,2.5/1.5625/3}{
  \draw[blue,dashed] (\xx,0)--(\xx,\hh);
  \fill[blue] (\xx,\hh) circle(1.7pt);
  \node[blue] at (\xx,-0.45) {$\xi_{\name}$};
 }
 \node[below left] at (0,0) {$a=x_0$};
 \node[below] at (0.8,0) {$x_1$};
 \node[below] at (1.8,0) {$x_2$};
 \node[below right] at (3,0) {$x_3=b$};
 \node at (2.7,2.28) {$y=f(x)$};
\end{tikzpicture}
\caption{幅が異なっていてもよく，各小区間のどの点の高さを使ってもよい。}
\label{図:リーマン和}
\end{figure}

図\ref{図:リーマン和}のように，分割を等間隔にする必要はなく，代表点 $\xi_k$ を端点にする必要もない。分割を細かくするとは，単に小区間の個数を増やすことではない。一つだけ幅の広い区間が残っている可能性を除くため，最大の幅
\[
 \|P\|=\max_{1\le k\le n}\Delta x_k
\]
を $0$ に近づける。

\begin{definition}{定積分}
$[a,b]$ 上の有界な関数 $f$ に対し，$\|P\|\to0$ のとき，分割と代表点の選び方によらずリーマン和が同じ有限値 $I$ に近づくなら，$f$ は $[a,b]$ で\textbf{積分可能}であるという。その値を\textbf{定積分}と呼び，
\[
 \int_a^b f(x)\dd x=I
\]
と書く。
\end{definition}

正確には，任意の $\varepsilon>0$ に対して，ある $\delta>0$ を選び，$\|P\|<\delta$ を満たすすべての分割とすべての代表点について $|S(P,\xi)-I|<\varepsilon$ とできることである。特定の分け方で一つの極限が求まるだけでは，この定義を満たすことまでは証明できない。

記号 $\int$ は和を極限まで進める操作を表し，$\dd x$ はどの変数に沿った幅を足しているかを表す。$
\int_a^b f(t)\dd t$ と文字を替えても同じ数である。$x$ や $t$ はこの積分の内部で用いる変数であり，$a,b$ が固定されていれば，定積分全体は一つの数になる。

\section{連続な関数は積分できる}\label{節:積分可能性}
定義は積分が存在する場合の意味を定めたが，存在そのものをまだ保証していない。本書で主に扱う連続関数については，次の定理を用いる。

\begin{theorem}{連続関数の積分可能性：採用する事実}\label{定理:連続積分可能}
閉区間 $[a,b]$ で連続な関数は，その区間でリーマン積分可能である。
\end{theorem}

この定理の完全な証明は，本書では扱わない。考え方は，十分短い小区間では関数の高さの変動も小さくなり，どの代表点を選んでも和の差が小さくなるというものである。閉区間上の連続性を使うと，その小ささを区間全体でそろえられる。各小区間での最大の高さと最小の高さの差が $\eta$ 以下なら，上から作る和と下から作る和の差は
\[
 \sum_{k=1}^{n}\eta\Delta x_k=\eta(b-a)
\]
以下になる。このような評価をすべての分割について保証する部分も含め，上の定理として採用する。

したがって，連続関数の定積分の値を計算するときは，存在をこの定理で確かめたうえで，計算しやすい等分割と代表点を選んでよい。右端を選ぶなら，
\[
 \int_a^b f(x)\dd x
 =\lim_{n\to\infty}\frac{b-a}{n}
  \sum_{k=1}^{n}f\left(a+\frac{k(b-a)}n\right)
\]
となる。ここでは幅が $\frac{b-a}{n}$，第 $k$ 小区間の右端が $a+\frac{k(b-a)}n$ である。

\begin{example}{リーマン和から値を求める}
区間 $[1,3]$ を $n$ 等分し，各小区間の右端を使って $\int_1^3(2x-1)\dd x$ を求めよ。
\end{example}
被積分関数は連続なので，定積分は存在する。幅は $\frac2n$，右端は $1+\frac{2k}{n}$ である。したがって，リーマン和は
\begin{align*}
 S_n
 &=\frac2n\sum_{k=1}^{n}\left\{2\left(1+\frac{2k}n\right)-1\right\}\\
 &=\frac2n\sum_{k=1}^{n}\left(1+\frac{4k}n\right)\\
 &=2+\frac8{n^2}\cdot\frac{n(n+1)}2
  =6+\frac4n.
\end{align*}
よって $n\to\infty$ として，定積分は $6$ となる。この例では和の公式が使えたが，一般の関数では和を直接計算するのは難しい。後で微分との関係を使い，別の計算法を導く。

\section{定積分の基本性質}\label{節:定積分性質}
以下では，特に断らない限り，被積分関数は扱う閉区間で連続であるとする。まず，積分の端点が同じ場合と，順序が逆の場合を
\[
 \int_a^a f(x)\dd x=0,\qquad
 \int_b^a f(x)\dd x=-\int_a^b f(x)\dd x\quad(a<b)
\]
と定める。積分の向きが反対になると符号が反対になる。この約束により，区間をつなぐ式を端点の大小にかかわらず統一して書ける。

\begin{theorem}{線形性と区間の加法性}\label{定理:積分基本性質}
$\alpha,\beta$ を実数とする。関係する区間で $f,g$ が連続なら，
\begin{align*}
 \int_a^b\{\alpha f(x)+\beta g(x)\}\dd x
 &=\alpha\int_a^b f(x)\dd x+\beta\int_a^b g(x)\dd x,\\
 \int_a^b f(x)\dd x
 &=\int_a^c f(x)\dd x+\int_c^b f(x)\dd x
\end{align*}
が成り立つ。後者は $a,c,b$ の順序によらない。
\end{theorem}

線形性は，同じ分割と代表点を使った有限和について
\[
 \sum_k\{\alpha f(\xi_k)+\beta g(\xi_k)\}\Delta x_k
 =\alpha\sum_k f(\xi_k)\Delta x_k
  +\beta\sum_k g(\xi_k)\Delta x_k
\]
と分け，極限をとれば得られる。区間の加法性は，$a<c<b$ のとき，分割に $c$ を含めて和を二つに分ければ分かる。端点の順序が異なる場合は，逆向きの積分の定義から同じ式になる。積分区間を分けても，被積分関数を勝手に変えているわけではない。

\begin{theorem}{大小関係と絶対値の評価}\label{定理:積分評価}
$a\le b$ とする。区間 $[a,b]$ で連続な関数について，次が成り立つ。
\begin{align*}
 f(x)\le g(x)\ (a\le x\le b)
 &\quad\Longrightarrow\quad
 \int_a^b f(x)\dd x\le\int_a^b g(x)\dd x,\\
 \left|\int_a^b f(x)\dd x\right|
 &\le\int_a^b|f(x)|\dd x.
\end{align*}
特に，$m\le f(x)\le M$ なら，
\[
 m(b-a)\le\int_a^b f(x)\dd x\le M(b-a).
\]
\end{theorem}

最初の性質は，各小区間の幅が正であり，リーマン和の各項について大小関係が保たれることから分かる。二つ目は $-|f|\le f\le|f|$ を積分すれば得られる。最後の評価では，定数関数の積分が高さと区間の長さの積になることを使った。これらの不等式は，端点を $a\le b$ の順に置いていることが前提である。

\begin{example}{原始関数を探さずに範囲を決める}
$\int_0^1\frac1{1+x^2}\dd x$ が $\frac12$ 以上 $1$ 以下であることを示せ。
\end{example}
$0\le x\le1$ では $1\le1+x^2\le2$ であり，すべて正なので逆数の大小は逆になる。
\[
 \frac12\le\frac1{1+x^2}\le1.
\]
区間の長さは $1$ だから，これを積分して
\[
 \frac12\le\int_0^1\frac1{1+x^2}\dd x\le1
\]
を得る。積分の値を正確に計算できなくても，関数の高さを比較すれば範囲を保証できる。この例の原始関数は\chapterref{章:置換の背景}で扱う。

\subsection{関数の平均値}
$a<b$ のとき，
\[
 f_{\mathrm{av}}=\frac1{b-a}\int_a^b f(x)\dd x
\]
を $f$ の区間 $[a,b]$ における\textbf{平均値}という。等分割のリーマン和を区間の長さで割ると $\frac1n\sum_{k=1}^{n}f(\xi_k)$ になるので，有限個の値の平均を細かくしていったものと解釈できる。微分で用いた平均変化率とは別の量である。

$f$ の最小値を $m$，最大値を $M$ とすると，積分の評価から $m\le f_{\mathrm{av}}\le M$ である。連続関数は最小値と最大値の間の値をすべてとるから，ある $c\in[a,b]$ で
\[
 \int_a^b f(x)\dd x=f(c)(b-a)
\]
が成り立つ。これを\textbf{積分の平均値の定理}という。高さを一定の $f(c)$ にした長方形と，元の関数の定積分が等しくなるという意味である。

\section{符号付きの積分と面積を区別する}\label{節:積分符号面積}
$f\ge0$ のとき，リーマン和は長方形の面積の和である。この極限として，曲線の下の面積を定める。連続関数が負になる部分では，定積分への寄与は負になるため，定積分をそのまま幾何学的な面積と呼ぶことはできない。

$f$ が正の部分の面積を $S_+$，負の部分と $x$ 軸の間の面積を $S_-$ とすれば，
\[
 \int_a^b f(x)\dd x=S_+-S_-,\qquad
 \int_a^b|f(x)|\dd x=S_++S_-
\]
である。どちらも同じグラフから求めるが，符号を保つか，面積としてすべて正にするかが異なる。面積を求めるには零点を調べ，符号が変わる場所で区間を分ける。

\begin{example}{同じグラフから二つの量を求める}
$y=x-1$ の $0\le x\le3$ における定積分と，グラフと $x$ 軸の間の面積を，三角形の面積から求めよ。
\end{example}
グラフは $x=1$ で $x$ 軸と交わる。$0\le x\le1$ の下側の三角形の面積は $\frac12$，$1\le x\le3$ の上側の三角形の面積は $2$ である。したがって，
\[
 \int_0^3(x-1)\dd x=-\frac12+2=\frac32,
 \qquad
 \int_0^3|x-1|\dd x=\frac12+2=\frac52.
\]
図\ref{図:定積分符号}の下側の部分は，左の計算では引き，右の計算では足している。

\begin{figure}[htbp]
\centering
\begin{tikzpicture}[x=2.3cm,y=1.5cm,>=Stealth]
 \fill[red!15] (0,0)--(0,-1)--(1,0)--cycle;
 \fill[blue!15] (1,0)--(3,2)--(3,0)--cycle;
 \draw[->] (-0.2,0)--(3.45,0) node[right]{$x$};
 \draw[->] (0,-1.3)--(0,2.5) node[above]{$y$};
 \draw[very thick,domain=0:3.2] plot (\x,{\x-1});
 \draw[dashed] (3,0)--(3,2)--(0,2);
 \node[below] at (1,0) {$1$};
 \node[below] at (3,0) {$3$};
 \node[left] at (0,-1) {$-1$};
 \node[left] at (0,2) {$2$};
 \node[red!70!black] at (0.3,-0.32) {$S_-$};
 \node[blue!70!black] at (2.4,0.48) {$S_+$};
 \node at (2.45,2.25) {$y=x-1$};
\end{tikzpicture}
\caption{定積分は上側の面積から下側の面積を引く。幾何学的な面積は両方を足す。}
\label{図:定積分符号}
\end{figure}

\section{積分の上端を動かして微分する}\label{節:基本定理第一}
ここまで，定積分をリーマン和の極限として定義し，和から分かる性質を調べた。次に，右端を変数にして積分そのものを関数として見る。$f$ が $[a,b]$ で連続なとき，
\[
 A(x)=\int_a^x f(t)\dd t\qquad(a\le x\le b)
\]
とおく。$x$ が動くと，上端までに足し合わせる範囲が変わる。$f\ge0$ なら面積を表すが，負の値を含む場合にも符号付きの量として同じ式を使う。

\begin{figure}[htbp]
\centering
\begin{tikzpicture}[x=2.5cm,y=1.5cm,>=Stealth]
 \fill[blue!12] (0,0)--plot[domain=0:2,samples=60] (\x,{1+0.2*\x*\x})--(2,0)--cycle;
 \fill[red!18] (2,0)--plot[domain=2:2.5,samples=35] (\x,{1+0.2*\x*\x})--(2.5,0)--cycle;
 \draw[->] (-0.15,0)--(3.15,0) node[right]{$t$};
 \draw[->] (0,-0.1)--(0,3.15) node[above]{$y$};
 \draw[very thick,domain=0:3,samples=80] plot (\x,{1+0.2*\x*\x});
 \draw[dashed] (2,0)--(2,1.8);
 \draw[dashed] (2.5,0)--(2.5,2.25);
 \node[below] at (0,0) {$a$};
 \node[below] at (2,0) {$x$};
 \node[below] at (2.5,0) {$x+h$};
 \node[blue!70!black] at (1,0.55) {$A(x)$};
 \node[red!70!black] at (2.25,0.8) {増分};
 \node at (2.7,3.05) {$y=f(t)$};
\end{tikzpicture}
\caption{上端を $h$ だけ動かしたときの積分の増分は，幅 $h$ の区間の積分になる。}
\label{図:面積関数増分}
\end{figure}

図\ref{図:面積関数増分}は $h>0$ で $f$ が正の場合を描いている。幅が小さいとき，増分は高さ $f(x)$，幅 $h$ の長方形に近い。これを符号や $h$ の向きに依存しない形で示す。

\begin{theorem}{微積分の基本定理：積分を微分する}\label{定理:基本定理第一}
$f$ が $[a,b]$ で連続なら，$A(x)=\int_a^x f(t)\dd t$ は $[a,b]$ で連続であり，内点で微分可能で，
\[
 A'(x)=f(x)\qquad(a<x<b)
\]
を満たす。端点では，区間の内側からの片側微分で同じ式が成り立つ。
\end{theorem}

まず，連続関数 $f$ は閉区間で有界だから，$|f(t)|\le M$ となる $M$ がある。積分の評価と区間の加法性により，任意の $u,v\in[a,b]$ について
\[
 |A(v)-A(u)|=\left|\int_u^v f(t)\dd t\right|
 \le M|v-u|
\]
となる。よって $A$ は連続である。

微分を調べるため，$x,x+h\in[a,b]$，$h\ne0$ とする。区間の加法性から
\begin{align*}
 \frac{A(x+h)-A(x)}h-f(x)
 &=\frac1h\int_x^{x+h}f(t)\dd t-f(x)\\
 &=\frac1h\int_x^{x+h}\{f(t)-f(x)\}\dd t.
\end{align*}
ここで，定数関数 $f(x)$ の $x$ から $x+h$ までの積分が $hf(x)$ になることを使った。

任意の $\varepsilon>0$ に対し，$f$ の $x$ での連続性から，$|t-x|<\delta$ なら $|f(t)-f(x)|<\varepsilon$ となる $\delta>0$ を選べる。$0<|h|<\delta$ のとき，$x$ と $x+h$ の間のすべての $t$ がこの条件を満たす。積分区間の長さは $|h|$ だから，向きがどちらであっても
\[
 \left|\frac1h\int_x^{x+h}\{f(t)-f(x)\}\dd t\right|
 \le\frac1{|h|}\,\varepsilon|h|=\varepsilon.
\]
したがって差商は $f(x)$ に近づき，$A'(x)=f(x)$ を得る。図から予想した「積分の増加率が高さになる」という関係が，負の高さや $h<0$ の場合も含めて証明できた。

\section{原始関数から定積分を計算する}\label{節:原始関数基本定理}
基本定理は，積分によって作った関数を微分すると元に戻ることを示した。今度は，微分すると $f$ になる関数を先に一つ見つけ，それを積分の計算に使う。

\begin{definition}{原始関数と不定積分}
一つの区間で $F'(x)=f(x)$ を満たす関数 $F$ を，$f$ の\textbf{原始関数}という。すべての原始関数を表す記法として
\[
 \int f(x)\dd x=F(x)+C
\]
と書き，\textbf{不定積分}と呼ぶ。$C$ は任意の定数であり，\textbf{積分定数}という。
\end{definition}

なぜ定数を足すだけで，すべての原始関数を表せるのだろうか。$F,G$ がともに $f$ の原始関数なら，$(F-G)'=0$ である。前章の定理\ref{定理:導関数同じ}により，一つの区間では $F-G$ が定数になるからである。逆に，$F+C$ を微分すれば $F'=f$ となる。たとえば $x^2$ の原始関数は $\frac{x^3}{3}+C$ である。

ここでも区間ごとに考える必要がある。$\frac1x$ の定義域は $x<0$ と $x>0$ に分かれるので，それぞれの区間で $\log|x|+C$ と書けるが，二つの区間で同じ定数を使う必要はない。また，基本定理から連続関数には原始関数が存在するが，その原始関数が既知の関数の式で簡単に表せるとは限らない。

\begin{theorem}{微積分の基本定理：原始関数で計算する}\label{定理:基本定理第二}
$f$ が $[a,b]$ で連続で，$F$ が $[a,b]$ で連続，$(a,b)$ で $F'=f$ を満たすなら，
\[
 \int_a^b f(x)\dd x=F(b)-F(a)
\]
が成り立つ。右辺を $[F(x)]_a^b$ とも書く。
\end{theorem}

証明には $A(x)=\int_a^x f(t)\dd t$ を使う。基本定理の前半から $A'=f$ なので，$(F-A)'=0$ である。$F-A$ は閉区間で連続だから，平均値の定理により端点を含む区間全体で一定となる。$A(a)=0$ を用いれば
\[
 F(x)-A(x)=F(a),\qquad A(x)=F(x)-F(a).
\]
ここに $x=b$ を代入すると，求める公式を得る。逆向きの積分や両端が同じ場合も，それぞれの定義から同じ計算式になる。

この証明の順序を確かめておこう。定積分は，初めから原始関数の差として定義したのではない。リーマン和の極限として定義し，その積分を微分すると $f$ に戻ることを示し，最後に平均値の定理によって任意の原始関数と結び付けた。これにより，長方形の和で表される量を，微分の逆を探す計算で求められるようになった。

\begin{example}{微分して確かめてから端点を代入する}
次の定積分を求めよ。
\[
 \text{(1)}\quad\int_0^2(3x^2-4x+1)\dd x,\qquad
 \text{(2)}\quad\int_1^{\e}\frac1x\dd x,\qquad
 \text{(3)}\quad\int_0^{\frac\pi2}\cos x\dd x.
\]
\end{example}
(1) では，微分すると $3x^2-4x+1$ になる関数を項ごとに探す。$F(x)=x^3-2x^2+x$ とすれば $F'=3x^2-4x+1$ だから，
\[
 \int_0^2(3x^2-4x+1)\dd x
 =[x^3-2x^2+x]_0^2=(8-8+2)-0=2.
\]
不定積分では $F+C$ と書くが，定積分では $(F(b)+C)-(F(a)+C)$ の定数が打ち消し合う。計算結果に積分定数を付ける必要はない。

(2) は区間全体で $x>0$ だから，$\log x$ を原始関数として使える。
\[
 \int_1^{\e}\frac1x\dd x=[\log x]_1^{\e}=1-0=1.
\]
(3) では $(\sin x)'=\cos x$ より，
\[
 \int_0^{\frac\pi2}\cos x\dd x
 =[\sin x]_0^{\frac\pi2}=1-0=1.
\]
また，最初の例題で和から求めた積分は，$\int_1^3(2x-1)\dd x=[x^2-x]_1^3=6$ と短く計算できる。二つの方法が同じ値になる理由が，基本定理で保証されている。

\section{上下端が変わる積分}\label{節:変数上下端}
$\int_a^x f(t)\dd t$ の微分が分かれば，上端や下端が別の関数になっている場合も合成関数の微分で扱える。$f$ がある区間で連続で，$u(x),v(x)$ が微分可能，その値がこの区間内にあるとする。区間内の固定した点 $c$ を選び，$A(s)=\int_c^s f(t)\dd t$ とおけば，
\[
 \int_{u(x)}^{v(x)}f(t)\dd t=A(v(x))-A(u(x)).
\]
両辺を微分し，$A'=f$ を使うと，
\[
 \boxed{\frac{\dd}{\dd x}\int_{u(x)}^{v(x)}f(t)\dd t
 =f(v(x))v'(x)-f(u(x))u'(x)}
\]
を得る。下端からの寄与に負号が付くのは，区間の加法性でその部分を引いているからである。$u(x)$ と $v(x)$ の大小が途中で入れ替わっても，向き付きの積分の約束によりこの式を使える。

\begin{example}{原始関数の式を求めずに微分する}
$G(x)=\int_x^{x^2}\e^{t^2}\dd t$ を微分せよ。
\end{example}
被積分関数 $\e^{t^2}$ はすべての実数で連続である。上端 $x^2$ の微分は $2x$，下端 $x$ の微分は $1$ だから，
\[
 G'(x)=\e^{(x^2)^2}\cdot2x-\e^{x^2}\cdot1
       =2x\e^{x^4}-\e^{x^2}.
\]
上端を代入してから内側の関数の微分を掛けるので，最初の指数は $x^4$ になる。$\e^{t^2}$ の原始関数を具体的な式で表す必要はない。

ただし，この公式では被積分関数を $f(t)$ とし，積分の外側の変数 $x$ がその中には現れないとしている。たとえば
\[
 H(x)=\int_0^x xt\dd t
      =x\int_0^x t\dd t=\frac{x^3}{2}
\]
では，$H'(x)=\frac{3x^2}{2}$ となる。上端を代入した $x^2$ だけでは足りない。積分の間は $x$ を定数として外に出せても，積分全体を $x$ で微分するときには，その因子も変わるからである。

\section{積分を含む等式から関数を決める}\label{節:積分方程式}
未知の関数が積分の中にも現れる等式を，\textbf{積分方程式}という。上端が変数なら，微分によって積分を外せる場合がある。ただし，微分は定数の情報を失うため，元の等式から得られる値の条件を残し，最後に元の式へ代入して確かめる。

\begin{example}{微分と初期値から関数を求める}
連続関数 $f$ が，すべての実数 $x$ について
\[
 f(x)=2+\int_0^x f(t)\dd t
\]
を満たすとき，$f$ を求めよ。
\end{example}
$x=0$ を代入すると，まず $f(0)=2$ と分かる。$f$ が連続なので，右辺の積分は基本定理により微分可能である。したがって左辺の $f$ も微分可能となり，両辺を微分して $f'(x)=f(x)$ を得る。

この等式だけで直ちに $f=\e^x$ としてはいけない。定数倍も同じ微分の関係を満たすからである。$\e^{-x}$ を掛けた関数を微分すると，
\[
 \frac{\dd}{\dd x}\{\e^{-x}f(x)\}
 =\e^{-x}\{f'(x)-f(x)\}=0
\]
となる。導関数が $0$ なので $\e^{-x}f(x)$ は定数であり，$x=0$ の値からその定数は $2$ である。よって候補は $f(x)=2\e^x$ となる。

最後に元の式を確認する。
\[
 2+\int_0^x2\e^t\dd t
 =2+[2\e^t]_0^x=2+2\e^x-2=2\e^x.
\]
確かに成り立つので，求める関数は $f(x)=2\e^x$ である。導いた条件が必要条件であるだけでなく，元の積分方程式を満たすことまで確かめた。

ここまでで，定積分の意味と，その計算を微分に結び付ける根拠がそろった。次章からは，原始関数を見つけるための式の整理を詳しく扱う。どの公式を使う場合も，求めた原始関数を微分して被積分関数に戻るかを確認できる。複雑な計算ほど，この戻り道が役に立つ。

\section*{章末問題}
\begin{enumerate}[label=\textbf{\thechapter.\arabic*},leftmargin=*]
 \item\label{問:定積分リーマン和} 区間 $[0,2]$ を $n$ 等分し，各小区間の右端を使ったリーマン和から $\int_0^2x^2\dd x$ を求めよ。幅と代表点を明記せよ。
 \item\label{問:定積分性質} 連続関数 $f,g$ が
 \[
  \int_0^1f(x)\dd x=2,\qquad
  \int_1^3f(x)\dd x=-1,\qquad
  \int_0^3g(x)\dd x=4
 \]
 を満たす。$\int_0^3f(x)\dd x$ と $\int_3^0\{2f(x)-3g(x)\}\dd x$ を求めよ。
 \item\label{問:定積分と面積} $\int_{-1}^2(x-1)\dd x$ と，直線 $y=x-1$，$x$ 軸，直線 $x=-1,x=2$ で囲まれる部分の面積を求めよ。符号が変わる点で区間を分けて説明せよ。
 \item\label{問:基本定理計算} 基本定理を使って次を求めよ。
 \[
  \text{(1)}\quad\int_{-1}^2(x^2+2x)\dd x,\qquad
  \text{(2)}\quad\int_0^{\frac\pi2}(\sin x+\cos x)\dd x,
 \]
 \[
  \text{(3)}\quad\int_1^{\e}\left(\frac2x+x\right)\dd x.
 \]
 \item\label{問:面積関数増減} $A(x)=\int_1^x(t-2)(t+1)\dd t$ とする。$A'(x)$ を求め，区間 $[-2,3]$ における $A$ の最大値・最小値と，それらをとる $x$ を求めよ。
 \item\label{問:変数上下端} $G(x)=\int_{\cos x}^{\sin x}\sqrt{1+t^2}\dd t$ を微分せよ。下端から生じる項の符号に注意せよ。
 \item\label{問:積分の評価} $0\le x\le1$ で $1\le\sqrt{1+x^2}\le1+\frac{x^2}{2}$ が成り立つことを示し，それを用いて
 \[
  1\le\int_0^1\sqrt{1+x^2}\dd x\le\frac76
 \]
 を証明せよ。平方根の原始関数を求める必要はない。
 \item\label{問:積分方程式基本} 次の条件を満たす連続関数を求めよ。存在しない場合は，その理由を示せ。
 \begin{enumerate}[label=(\arabic*),leftmargin=*]
  \item すべての実数 $x$ で $f(x)=3-\int_0^x f(t)\dd t$。
  \item すべての $x\in[0,1]$ で $g(x)=x+\int_0^1g(t)\dd t$。
 \end{enumerate}
\end{enumerate}

\smallskip\noindent 解答は\pageref{解答:第3章}ページ。

\chapter{基本的な積分と式の整理}\label{章:基本積分}

前章では，定積分を計算するために原始関数を探せばよいことが分かった。そこで今度は，目の前の関数に対して，その原始関数をどう見つけるかを考える。たとえば $x^2$ なら，微分すると $3x^2$ になる $x^3$ を思い出し，係数を直せばよい。ところが $(x+1)^2$ や $\frac{x^2}{x+1}$ は，見た形が違うだけで，同じ方法へ進めることがある。

この章では，まず微分公式を逆向きに読む。そのうえで，展開，割り算，平方完成，部分分数，三角恒等式を使い，知っている原始関数が見える形へ式を整える。新しい積分法を次々に覚える前に，式を変えるだけで何ができるかを確かめよう。

\section{微分公式を逆向きに読む}\label{節:基本公式の逆}

$F'(x)=f(x)$ ならば，$F$ は $f$ の原始関数である。この関係を使うとき，最初から正しい係数まで思い付く必要はない。まず形の近い候補を微分し，元の関数との差を直すのでよい。たとえば
\[
 \frac{\dd}{\dd x}x^5=5x^4
\]
だから，$x^4$ の原始関数を得るには，$x^5$ を $5$ で割る。
\[
 \int x^4\dd x=\frac{x^5}{5}+C.
\]
同じ理由で，$r\ne-1$ のとき
\[
 \int x^r\dd x=\frac{x^{r+1}}{r+1}+C
\]
となる。指数を一つ増やし，増やした指数で割るのは，微分で指数が一つ減り，その指数が係数として出ることを元に戻しているのである。

ここで $r=-1$ だけは別に考える必要がある。上の式へ $r=-1$ を代入すると分母が $0$ になり，意味を持たない。また，$x^0=1$ を微分しても $x^{-1}$ は現れない。代わりに，対数の微分
\[
 \frac{\dd}{\dd x}\log x=\frac1x\qquad(x>0)
\]
を使う。$x<0$ では $\log(-x)$ を微分すれば
\[
 \frac{\dd}{\dd x}\log(-x)
 =\frac1{-x}\,(-1)=\frac1x
\]
となるから，二つの形は $\log|x|$ とまとめられる。

\begin{theorem}{基本的な不定積分}
次の公式は，右辺が微分でき，左辺の関数が定義される区間で成り立つ。
\[
\begin{array}{c|c}
 f(x)&\int f(x)\dd x\\ \hline
 x^r\quad(r\ne-1)&\frac{x^{r+1}}{r+1}+C\\[5pt]
 \frac1x&\log|x|+C\\[5pt]
 \e^x&\e^x+C\\[3pt]
 a^x\quad(a>0,\ a\ne1)&\frac{a^x}{\log a}+C\\[5pt]
 \sin x&-\cos x+C\\[3pt]
 \cos x&\sin x+C\\[3pt]
 \frac1{\cos^2x}&\tan x+C
\end{array}
\]
実数 $r$ に対するべきの公式は，まず $x>0$ で考える。整数のべきなど，負の $x$ でも定義される場合は，その定義域内の区間にも適用できる。
\end{theorem}

この表は，右の式を微分して確かめられる。たとえば $a^x$ の微分では $\log a$ が掛かるので，積分では $\log a$ で割る。$a=1$ を除くのは，この割り算ができないためであり，$1^x=1$ の積分そのものは $x+C$ である。

不定積分の $C$ は，考えている一つの区間上での任意定数である。$\frac1x$ の定義域は $0$ の左右に分かれる。左側で加える定数と右側で加える定数が，同じである必要はない。図\ref{図:対数の区間}では，この区間の分かれ方を数直線に示した。原始関数を一つ選ぶ場合と，定義域全体のすべての原始関数を述べる場合とを区別しよう。

\begin{figure}[htbp]
 \centering
 \begin{tikzpicture}[x=1cm,y=1cm]
  \draw[-{Stealth}] (-4,0)--(4,0) node[right] {$x$};
  \draw[very thick,blue] (-3.7,0)--(-0.12,0);
  \draw[very thick,blue] (0.12,0)--(3.7,0);
  \filldraw[fill=white] (0,0) circle (2pt);
  \node[below] at (0,-0.08) {$0$};
  \node[above] at (-2,0.18) {$\log(-x)+C_-$};
  \node[above] at (2,0.18) {$\log x+C_+$};
  \node[below] at (-2,-0.35) {$x<0$};
  \node[below] at (2,-0.35) {$x>0$};
 \end{tikzpicture}
 \caption{$0$ を越えて一つの区間にはできない。定数は左右で独立に選べる。}
 \label{図:対数の区間}
\end{figure}

\begin{example}{根号と分母をべきの形へ直す}
$x>0$ において，$\int\left(2\sqrt{x}+\frac3{x^2}\right)\dd x$ を求めよ。
\end{example}

根号のままでは公式と見比べにくいので，$\sqrt{x}=x^{\frac12}$，$\frac1{x^2}=x^{-2}$ と書く。二つの項は別々に積分できるから，
\[
\begin{aligned}
 \int\left(2\sqrt{x}+\frac3{x^2}\right)\dd x
 &=2\int x^{\frac12}\dd x+3\int x^{-2}\dd x\\
 &=2\cdot\frac23x^{\frac32}
   +3\cdot\frac{x^{-1}}{-1}+C\\
 &=\frac43x^{\frac32}-\frac3x+C.
\end{aligned}
\]
分母にある式をいつも対数で積分するわけではない。$\frac1x$ は指数が $-1$ なので対数になり，$\frac1{x^2}$ は指数が $-2$ なので，べきの公式を使う。答えを微分すると
\[
 \frac43\cdot\frac32x^{\frac12}
 +3x^{-2}=2\sqrt{x}+\frac3{x^2}
\]
となり，係数と符号の両方を確認できる。

和を項ごとに積分できるのは，微分が和と定数倍を保つためである。一方，積について
\[
 \int f(x)g(x)\dd x
 =\left(\int f(x)\dd x\right)
  \left(\int g(x)\dd x\right)
\]
という規則はない。右辺を微分すると積の微分の二項が現れ，ふつうは $f(x)g(x)$ にはならない。積があるときは，まず展開できるか，恒等式で和にできるかを考える。その後に学ぶ置換積分と部分積分も，この問題への別の答えである。

\section{一次式が入っているときの係数}\label{節:一次式の積分}

$\sin(3x)$ を積分する場合，$-\cos(3x)$ を候補にすると
\[
 \frac{\dd}{\dd x}\{-\cos(3x)\}=3\sin(3x)
\]
となる。合成関数の微分によって，内側の $3x$ の微分 $3$ が余分に掛かる。したがって候補を $3$ で割ればよい。
\[
 \int\sin(3x)\dd x=-\frac13\cos(3x)+C.
\]
一般に $F'=f$，$a\ne0$ ならば，
\[
 \frac{\dd}{\dd x}\left\{\frac1aF(ax+b)\right\}
 =\frac1a f(ax+b)\,a=f(ax+b).
\]
よって
\[
 \int f(ax+b)\dd x=\frac1aF(ax+b)+C
\]
となる。$a$ で割るのは書式上の約束ではなく，微分したときに現れる $a$ を打ち消すためである。

\begin{example}{内側の微分を確かめる}
次の不定積分を求めよ。
\[
 \int\left(4\e^{2x-1}+\frac3{5x+2}-2\cos(3x)\right)\dd x.
\]
\end{example}

それぞれ，指数の $2x-1$，分母の $5x+2$，角の $3x$ の微分を確かめる。これらは $2,5,3$ であるから，
\[
\begin{aligned}
 &\int\left(4\e^{2x-1}+\frac3{5x+2}-2\cos(3x)\right)\dd x\\
 &\qquad=2\e^{2x-1}+\frac35\log|5x+2|
       -\frac23\sin(3x)+C.
\end{aligned}
\]
$5x+2$ が $0$ になる点は除き，その左右の各区間で考える。対数の項を微分すると
\[
 \frac35\cdot\frac5{5x+2}=\frac3{5x+2}
\]
である。分母の係数 $5$ を見落とすと，ここで元の式へ戻らない。

一次式の場合は内側の微分が定数なので，係数の修正だけで済んだ。内側が $x^2+1$ ならば，微分は $2x$ であり，単なる定数倍ではない。次章の置換積分では，この $2x$ が積分される式の中にあるかどうかが手がかりになる。

\section{展開と割り算で和を作る}\label{節:展開と割算}

整式の積は，展開すれば各項が $x$ のべきになる。式を短く保つことと，積分しやすくすることは同じではない。積の形では微分の逆が見えなくても，和にすると基本公式を一項ずつ使える。

\begin{example}{積を展開する理由}
$\int x(2x-1)^2\dd x$ を求めよ。
\end{example}

二乗を展開してから $x$ を掛けると，
\[
 x(2x-1)^2=x(4x^2-4x+1)=4x^3-4x^2+x.
\]
したがって
\[
 \int x(2x-1)^2\dd x
 =x^4-\frac43x^3+\frac12x^2+C.
\]
候補を微分すれば $4x^3-4x^2+x$ に戻り，これをくくり直すと元の式になる。高いべきでは展開する項数が多くなるため，別の方法が有利になる。しかしこの例のように，数項で済むならば，展開は十分に見通しのよい方法である。

分数式では，分子の次数が分母以上のとき，まず整式の割り算を考える。割り算は，分母に掛ける整式と余りとに分子を分ける操作である。そうすれば分数全体が，「整式」と「分子の次数が低い分数」の和になる。

\begin{example}{分子の次数を下げる}
$x\ne-1$ において，$\int\frac{x^3+2x^2+3}{x+1}\dd x$ を求めよ。
\end{example}

分子を分母で割ると
\[
 x^3+2x^2+3=(x+1)(x^2+x-1)+4
\]
となる。この等式は，右辺を展開して確かめられる。両辺を $x+1$ で割って
\[
 \frac{x^3+2x^2+3}{x+1}
 =x^2+x-1+\frac4{x+1}
\]
とすれば，すべて既知の形である。よって
\[
 \int\frac{x^3+2x^2+3}{x+1}\dd x
 =\frac{x^3}{3}+\frac{x^2}{2}-x+4\log|x+1|+C.
\]
分数のまま眺めていたときに難しかったのは，高い次数の分子と分母が一つにまとまっていたためである。割り算によってその大部分を整式へ移すと，最後には一次式の逆数だけが残った。

約分できる式でも，もとの定義域は確かめる。たとえば $\frac{x^2-1}{x-1}=x+1$ は $x\ne1$ で成り立つ。右辺が $x=1$ でも定義されるからといって，左辺の関数が自動的にその点で定義されたことにはならない。積分する区間にその点が含まれるなら，連続に延長した関数を扱うのか，もとの式を扱うのかを明らかにする必要がある。

\section{部分分数で分母を分ける}\label{節:部分分数}

割り算を終えても，分母が二つの一次式の積になっていることがある。ここで分数の分母を勝手に分配することはできない。たとえば $\frac1{AB}$ は $\frac1A+\frac1B$ ではない。しかし，分子に合う係数を選べば，より簡単な分数の和として表せる。この書き換えを\textbf{部分分数分解}という。

\begin{example}{異なる一次因子に分ける}
$\int\frac{5x+1}{(x-1)(x+2)}\dd x$ を求めよ。
\end{example}

一次式の逆数ならば対数で積分できる。そこで
\[
 \frac{5x+1}{(x-1)(x+2)}
 =\frac A{x-1}+\frac B{x+2}
\]
と表せないかを調べる。右辺を通分すると分子は $A(x+2)+B(x-1)$ なので，求める係数は
\[
 5x+1=A(x+2)+B(x-1)
       =(A+B)x+(2A-B)
\]
を満たせばよい。係数を比べて
\[
 A+B=5,\qquad 2A-B=1
\]
を解くと $A=2$，$B=3$ となる。したがって
\[
 \int\frac{5x+1}{(x-1)(x+2)}\dd x
 =2\log|x-1|+3\log|x+2|+C.
\]
元の関数は $x=-2,1$ で定義されないので，答えはそれらをまたがない各区間で読む。

係数を決める途中で，分母を払った恒等式に $x=1,-2$ を代入してもよい。$x=1$ なら $6=3A$，$x=-2$ なら $-9=-3B$ となり，同じ係数が得られる。これは，元の分数に分母を $0$ にする値を代入したのではない。係数を決めるための整式の恒等式を調べているのである。

分母に同じ一次因子が重なっている場合も，低いべきから分けると積分できる。たとえば $\frac1{(x+1)^2}$ 自体はべきの公式で扱えるが，分子が一次式ならば，その分子を $x+1$ と定数の和にする。

\begin{example}{重なった一次因子を持つ分母}
$\int\frac{2x+3}{(x+1)^2}\dd x$ を求めよ。
\end{example}

分母にある $x+1$ を分子にも作りたいので，$2x+3=2(x+1)+1$ と書く。すると
\[
 \frac{2x+3}{(x+1)^2}
 =\frac2{x+1}+\frac1{(x+1)^2}.
\]
第一項には対数，第二項にはべきの公式を使い，
\[
 \int\frac{2x+3}{(x+1)^2}\dd x
 =2\log|x+1|-\frac1{x+1}+C
\]
を得る。二乗の分母があるからといって，すべてが逆数の原始関数になるわけではない。分解後の各項を見て，どの公式を使うかを決めている。

一般に有理関数，すなわち整式どうしの商を積分するときは，まず割り算で分子の次数を下げ，次に分母を因数分解して分けることを考える。ただし，実数の範囲で一次式に分かれない二次式もある。その場合に役立つ形を次に調べる。

\section{平方完成で二次式の中心を見る}\label{節:平方完成と積分}

二次式 $x^2-4x+5$ は
\[
 x^2-4x+5=(x-2)^2+1
\]
と平方完成できる。$x$ そのものよりも，$x-2$ を一つのまとまりとして見ると，二次式の中にある構造が明確になる。積分でも，このまとまりとその微分が式の中にあるかを調べると，候補が見つかることがある。

\begin{example}{平方根の微分と見比べる}
$\int\frac{x-2}{\sqrt{x^2-4x+5}}\dd x$ を求めよ。
\end{example}

平方完成すると，積分される式は
\[
 \frac{x-2}{\sqrt{(x-2)^2+1}}
\]
となる。分母に平方根があるので，$\sqrt{(x-2)^2+1}$ を微分してみる。合成関数の微分から
\[
 \frac{\dd}{\dd x}\sqrt{(x-2)^2+1}
 =\frac1{2\sqrt{(x-2)^2+1}}\,2(x-2)
 =\frac{x-2}{\sqrt{(x-2)^2+1}}.
\]
これがちょうど元の式なので，
\[
 \int\frac{x-2}{\sqrt{x^2-4x+5}}\dd x
 =\sqrt{x^2-4x+5}+C.
\]
根号の中は常に $1$ 以上だから，この答えは実数全体で使える。候補が見えた後に微分して一致を確かめたので，これは推測だけで終わった解答ではない。

同じ根号でも，分子に $x-2$ がない $\int\frac1{\sqrt{(x-2)^2+1}}\dd x$ では，いまの微分結果と一致しない。その違いを無視して同じ答えにはできない。また
\[
 \frac1{x^2-4x+3}=\frac1{(x-2)^2-1}
\]
ならば，平方の差なので $(x-3)(x-1)$ に分解できる。一方，$(x-2)^2+1$ は実数の一次式の積に分解できない。平方完成はいつでも積分を終わらせる操作ではないが，その後に因数分解を使えるのか，置換を考えるべきなのかを判断する材料になる。

\section{三角恒等式で積とべきをほどく}\label{節:三角式の整理}

$\sin x$ と $\cos x$ の原始関数は分かっていても，$\sin^2x$ や $\cos^2x$ は，いまの基本公式にそのまま当てはめられない。そこで，二倍角の公式
\[
 \cos(2x)=1-2\sin^2x=2\cos^2x-1
\]
を解き直して
\[
 \sin^2x=\frac{1-\cos(2x)}2,\qquad
 \cos^2x=\frac{1+\cos(2x)}2
\]
とする。左辺は二乗だが，右辺は定数と余弦の和になった。ここまで変形すれば基本公式を使える。

\begin{example}{二乗と積を和へ直す}
次の積分を求めよ。
\[
 \text{(1)}\quad\int_0^{\frac\pi2}\sin^2x\dd x,
 \qquad
 \text{(2)}\quad\int\sin(3x)\cos x\dd x.
\]
\end{example}

(1) は二倍角の公式から
\[
\begin{aligned}
 \int_0^{\frac\pi2}\sin^2x\dd x
 &=\frac12\int_0^{\frac\pi2}\{1-\cos(2x)\}\dd x\\
 &=\left[\frac x2-\frac14\sin(2x)\right]_0^{\frac\pi2}
 =\frac\pi4.
\end{aligned}
\]
$\cos(2x)$ の積分に $\frac12$ が付くため，全体の係数が $\frac14$ になる。最初から定積分を計算しているので，原始関数に任意定数を付けても上下端で差を取ると消える。

(2) は加法定理を二つ並べて考える。
\[
 \sin(3x+x)=\sin(3x)\cos x+\cos(3x)\sin x,
\]
\[
 \sin(3x-x)=\sin(3x)\cos x-\cos(3x)\sin x.
\]
加えると右辺の第二項が打ち消し合うので，
\[
 \sin(3x)\cos x=\frac12\{\sin(4x)+\sin(2x)\}.
\]
よって
\[
 \int\sin(3x)\cos x\dd x
 =-\frac18\cos(4x)-\frac14\cos(2x)+C.
\]
積を和へ直す公式を忘れても，加法定理から必要な式を作れる。微分して $\frac12\sin(4x)+\frac12\sin(2x)$ に戻ることを確かめ，もう一度積に直せば確認も終わる。

ここまでの計算では，積分記号の中の式を同じ値の別の式に書き換えてきた。どの変形を使うかは，元の式に現れた手がかりで決まる。整式の小さな積なら展開し，分子の次数が高ければ割り算をし，一次因子の積が分母にあれば部分分数を考える。二次式は平方完成でまとまりを見やすくし，三角関数の二乗や積は恒等式で和へ直す。

ただし，式を長く展開すれば必ずよいわけではない。たとえば $x(x^2+1)^{20}$ は展開できるが，$x^2+1$ を一つの変数と考えれば計算を短くできる。その見方を正当化するのが，次章で学ぶ置換積分である。

\section*{章末問題}
\begin{enumerate}[label=\textbf{\thechapter.\arabic*},leftmargin=*]
 \item\label{問:基本べき} $x>0$ において，$\int\frac{3x^2-2x+4}{\sqrt{x}}\dd x$ を求め，微分して確かめよ。
 \item\label{問:基本割算} $\int\frac{x^3+1}{x+1}\dd x$ を求めよ。式を整理した後も，もとの関数で除かれる点を述べよ。
 \item\label{問:基本部分分数} $\int\frac{5x-1}{(x-1)(x+1)}\dd x$ を求めよ。部分分数の係数を決める等式も書け。
 \item\label{問:基本重因子} $\int\frac{3x+4}{(x+1)^2}\dd x$ を求めよ。
 \item\label{問:基本平方完成} 分母を平方完成してから，$\int\frac1{x^2+6x+8}\dd x$ を求めよ。
 \item\label{問:基本三角二乗} $\int_0^{\frac\pi2}\cos^2(2x)\dd x$ を求めよ。
 \item\label{問:基本三角積} 加法定理から必要な積を和へ直す式を導き，$\int\sin(2x)\sin(5x)\dd x$ を求めよ。
\end{enumerate}

\smallskip\noindent 解答は\pageref{解答:第4章}ページ。

\chapter{置換積分}\label{章:置換積分}

$x(x^2+1)^{20}$ を積分しようとすると，展開によって多くの項が生じる。しかし，式の中で繰り返し使われているまとまりは $x^2+1$ だけである。それを $t$ と書けば，べきの部分は $t^{20}$ になる。残った $x$ は邪魔なのだろうか。実は $x^2+1$ の微分は $2x$ なので，この $x$ も変数を換えるために必要な形になっている。

置換積分では，このように式の一部分を新しい変数にして，合成関数の微分を逆に使う。計算だけでなく，どの区間をどの区間に対応させたのかも確かめる。特に平方根や三角関数から元の変数を表すときには，一つの値に一つの値が対応する範囲を選ぶ必要がある。

\section{合成関数の微分を元へ戻す}\label{節:合成形の置換}

$F'(t)=f(t)$ であれば，合成関数の微分から
\[
 \frac{\dd}{\dd x}F(g(x))=f(g(x))g'(x)
\]
となる。この右辺の形が積分される式に現れたら，原始関数は左辺の $F(g(x))$ である。これが置換積分の根拠である。

ここで，微分可能で，その導関数も連続である関数を\textbf{連続微分可能}という。\label{用語:連続微分可能}

\begin{theorem}{合成形の置換積分}
区間 $I$ で $g$ は連続微分可能とし，$g$ の値を含む区間で $f$ は連続とする。$F'=f$ ならば
\[
 \int f(g(x))g'(x)\dd x=F(g(x))+C.
\]
さらに $a,b\in I$ に対して
\[
 \int_a^b f(g(x))g'(x)\dd x
 =\int_{g(a)}^{g(b)}f(t)\dd t.
\]
この二つの式に，$g$ が単調であるという条件は要らない。
\end{theorem}

不定積分の式は，上で確かめた合成関数の微分そのものである。定積分の式も，基本定理によって両辺がともに
\[
 F(g(b))-F(g(a))
\]
となることから分かる。$g$ が途中で増加から減少へ変わっても，この証明は変わらない。後で逆関数を使う置換と比べるため，この点を覚えておこう。

実際の計算では，$t=g(x)$ と置き，$\dd t=g'(x)\dd x$ と書く。この記号は，$g'(x)$ と $\dd x$ を含む積を，$t$ による積分へ書き換えることを表している。微小な数を自由に割り算するという説明に頼る必要はなく，根拠は上の合成関数の微分にある。

\begin{example}{まとまりとその微分を見つける}
$\int x(x^2+1)^{20}\dd x$ を求めよ。
\end{example}

$x^2+1$ を $t$ と置くと，$\dd t=2x\dd x$ だから，$x\dd x=\frac12\dd t$ とまとめて書き換えられる。したがって
\[
\begin{aligned}
 \int x(x^2+1)^{20}\dd x
 &=\frac12\int t^{20}\dd t\\
 &=\frac{t^{21}}{42}+C
 =\frac{(x^2+1)^{21}}{42}+C.
\end{aligned}
\]
新しい変数 $t$ で積分した後，不定積分の答えは元の変数 $x$ に戻す。右辺を微分すると
\[
 \frac{21}{42}(x^2+1)^{20}\,2x
 =x(x^2+1)^{20}
\]
となる。

ここでは $x=0$ でも式は成り立つ。計算の途中で $\dd x=\frac{\dd t}{2x}$ として $x$ で割る必要はなかった。$x\dd x$ 全体を $\frac12\dd t$ に換えれば，余分な除外点を作らずに済む。また $x^2+1$ は実数全体では単調でないが，原始関数の微分による確認には問題がない。

新しい変数を選ぶときは，「見た目が複雑な部分」だけでなく，「その部分を微分したものが残りにあるか」を調べる。たとえば
\[
 \int\frac{\log x}{x}\dd x
\]
では，$t=\log x$ と置けば $\dd t=\frac1x\dd x$ であり，$\int t\dd t$ に変わる。したがって，$x>0$ で答えは $\frac12(\log x)^2+C$ となる。複雑に見える対数を変数にすると，同時に分母の $x$ もなくなっている。

\section{分母の微分が分子にあるとき}\label{節:対数型の置換}

$\frac1t$ の原始関数は $\log|t|$ だから，置換積分を使うと
\[
 \int\frac{g'(x)}{g(x)}\dd x=\log|g(x)|+C
\]
となる。ただし，$g(x)\ne0$ である区間で使う。$g$ が連続ならば，その区間内で符号を変えるには $0$ を通る必要があるから，符号は一定である。正なら $\log g(x)$，負なら $\log(-g(x))$ を微分すれば，同じ商に戻る。

\begin{example}{絶対値を残すべき対数}
$\int\frac{6x}{3x^2-2}\dd x$ を求めよ。
\end{example}

分母の $3x^2-2$ を微分すると，分子の $6x$ が現れる。$t=3x^2-2$ と置けば $\dd t=6x\dd x$ なので，
\[
 \int\frac{6x}{3x^2-2}\dd x
 =\int\frac1t\dd t
 =\log|t|+C
 =\log|3x^2-2|+C.
\]
分母が $0$ になる $x=\pm\sqrt{\frac23}$ は除く。この二点の間では $3x^2-2<0$ なので，絶対値を外して $\log(3x^2-2)$ とすると実数の式ではなくなる。一方，$\log|3x^2-2|$ は三つの区間のどこでも使える形である。任意定数は各区間で独立に選んでよい。

分母が二次式だから必ず部分分数にする，という決め方はしない。この例では分子を見れば分母の微分があり，一度の置換で終わる。前章の式の整理と置換積分は競合する規則ではなく，式の特徴に応じて選ぶ道具である。

ただし，$\int\frac1{3x^2-2}\dd x$ の分子は $6x$ ではないから，同じ公式をそのまま使うことはできない。分母だけを $t$ と換えて，$\dd x$ を変えずに残すと，変数の異なる式を混ぜたことになる。置換後は，積分される式と微分の記号の両方を，新しい変数で書き直す必要がある。

\section{定積分では上下端も換える}\label{節:置換と積分区間}

不定積分では最後に $x$ へ戻したが，定積分なら新しい変数のまま計算を終えられる。$t=g(x)$ と置いたときは，$x=a$ に対応する $t=g(a)$ を下端に，$x=b$ に対応する $t=g(b)$ を上端にする。大きい数を自動的に上端にするのではなく，元の端点との対応を保つ。

たとえば $t=1-x$ で $x$ が $0$ から $1$ へ動けば，$t$ は $1$ から $0$ へ動く。さらに $\dd t=-\dd x$ だから，
\[
 \int_0^1(1-x)^2\dd x
 =-\int_1^0 t^2\dd t
 =\int_0^1t^2\dd t=\frac13.
\]
微分から来た負号と，上下端を入れ換えたときの負号は，それぞれ別の理由で現れる。途中で片方だけを落とさないよう，対応と微分を先に書いておくとよい。

\begin{example}{途中で折り返す置換}
$\int_{-1}^{2}2x\cos(x^2)\dd x$ を求めよ。
\end{example}

$t=x^2$ と置けば $\dd t=2x\dd x$ であり，端点の対応は
\[
 \begin{array}{c|cc}
 x&-1&2\\ \hline
 t&1&4
 \end{array}
\]
である。したがって合成形の置換積分から
\[
 \int_{-1}^{2}2x\cos(x^2)\dd x
 =\int_1^4\cos t\dd t
 =\sin4-\sin1.
\]
角 $1,4$ はラジアンである。$x$ が $-1$ から $2$ まで増える間に，$t=x^2$ は一度 $1$ から $0$ へ減り，その後 $4$ へ増える。端点だけを見た式が正しいのは，この折り返しが定積分の向きに反映されているためである。

図\ref{図:置換の折返し}に対応を示す。$x=0$ で積分を分けて確かめると，
\[
\begin{aligned}
 \int_{-1}^{2}2x\cos(x^2)\dd x
 &=\int_1^0\cos t\dd t+\int_0^4\cos t\dd t\\
 &=-\int_0^1\cos t\dd t+\int_0^4\cos t\dd t
 =\int_1^4\cos t\dd t.
\end{aligned}
\]
往復した範囲の寄与が符号付きで打ち消し合っている。この計算は，曲線の下の面積を重複も含めて足す操作とは違う。

\begin{figure}[htbp]
 \centering
 \begin{tikzpicture}[x=1.6cm,y=0.8cm]
  \draw[-{Stealth}] (-1.35,0)--(2.35,0) node[right] {$x$};
  \draw[-{Stealth}] (0,-0.35)--(0,4.65) node[above] {$t$};
  \draw[thick,blue,domain=-1:2,samples=65] plot (\x,{\x*\x});
  \draw[dashed] (-1,0)--(-1,1)--(0,1);
  \draw[dashed] (2,0)--(2,4)--(0,4);
  \fill (-1,1) circle (1.5pt);
  \fill (2,4) circle (1.5pt);
  \node[below] at (-1,0) {$-1$};
  \node[below right] at (0,0) {$0$};
  \node[below] at (2,0) {$2$};
  \node[left] at (0,1) {$1$};
  \node[left] at (0,4) {$4$};
  \node[right] at (1.6,1.9) {$t=x^2$};
  \draw[-{Stealth},thick] (-0.88,0.85)--(-0.55,0.41);
  \draw[-{Stealth},thick] (1.12,1.38)--(1.43,2.17);
 \end{tikzpicture}
 \caption{$x$ は増え続けるが，$t=x^2$ は減少してから増加する。}
 \label{図:置換の折返し}
\end{figure}

\section{元の変数を新しい変数で表す}\label{節:逆向きの置換}

合成形がそのまま見えない場合でも，根号などを新しい変数にしてから，元の変数を解き直すと式が簡単になることがある。このときは，$t=g(x)$ から $x$ を $t$ の関数として表せる範囲を選ぶ。

\begin{example}{根号を新しい変数にする}
$x>0$ において，$\int\frac{\sqrt{x}}{1+\sqrt{x}}\dd x$ を求めよ。
\end{example}

$t=\sqrt{x}$ と置けば，$t>0$ で $x=t^2$，$\dd x=2t\dd t$ となる。分子だけでなく分母と $\dd x$ も書き換えて
\[
 \int\frac{\sqrt{x}}{1+\sqrt{x}}\dd x
 =\int\frac{t}{1+t}\,2t\dd t
 =2\int\frac{t^2}{t+1}\dd t.
\]
根号はなくなったが分数が残った。ここで前章の割り算を使い，$t^2=(t+1)(t-1)+1$ と分ける。
\[
\begin{aligned}
 2\int\frac{t^2}{t+1}\dd t
 &=2\int\left(t-1+\frac1{t+1}\right)\dd t\\
 &=t^2-2t+2\log(t+1)+C.
\end{aligned}
\]
$t+1>0$ なので，ここでは対数の絶対値を外せる。$t=\sqrt{x}$ を戻すと，答えは
\[
 x-2\sqrt{x}+2\log(1+\sqrt{x})+C.
\]
微分して確かめると
\[
 1-\frac1{\sqrt{x}}+\frac1{(1+\sqrt{x})\sqrt{x}}
 =1-\frac1{1+\sqrt{x}}
 =\frac{\sqrt{x}}{1+\sqrt{x}}
\]
となる。置換が積分を一度で終わらせなくても，根号を有理式へ変えるという目的は達している。

この例で $x=t^2$ と書けたのは，$t=\sqrt{x}>0$ と決めていたからである。逆に $t=x^2$ から $x$ を求めるなら，$x=\sqrt{t}$ と $x=-\sqrt{t}$ の二つがある。$x>0$ の区間では前者，$x<0$ の区間では後者を使う。このように，逆関数を一つに決めるために選ぶ範囲を\textbf{枝}と呼ぶことがある。

区間上の連続関数で，逆関数を持つことと狭義単調であることは同値である。ここで狭義単調とは，異なる入力に対して出力が必ず一方向に増えるか，必ず一方向に減ることをいう。したがって，$t=g(x)$ から一つの関数 $x=h(t)$ を使って戻したいなら，そのような区間を選ばなければならない。さらに逆関数の微分を $h'(t)=\frac1{g'(h(t))}$ と書く場合には，分母が $0$ でないことも必要である。

これに対し，\sectionref{節:合成形の置換}の公式では逆関数を作っていない。$F(g(x))$ をそのまま微分したので，単調性を仮定しなかった。定積分の公式
\[
 \int_\alpha^\beta f(h(t))h'(t)\dd t
 =\int_{h(\alpha)}^{h(\beta)}f(x)\dd x
\]
自体も，同じ合成形の公式であり，$h$ の単調性を必要としない。単調性が必要になるのは，逆関数を一つ選んで式を書き直す段階である。この二つを区別すれば，置換に単調性が要る場合と要らない場合を混同せずに済む。

\section{三角関数で根号を簡単にする}\label{節:三角置換の初歩}

$a>0$ として，$\sqrt{a^2-x^2}$ には平方の差がある。三角関数の関係式 $1-\sin^2\theta=\cos^2\theta$ と形が同じだから，$x=a\sin\theta$ と置くことを考える。すると
\[
 \sqrt{a^2-x^2}
 =a\sqrt{1-\sin^2\theta}
 =a\sqrt{\cos^2\theta}
 =a|\cos\theta|.
\]
最後は $a\cos\theta$ ではなく，まず絶対値付きである。$\sqrt{z^2}=|z|$ だからである。そのうえで，たとえば $-\frac\pi2\le\theta\le\frac\pi2$ を選べば $\cos\theta\ge0$ となり，絶対値を外せる。この区間では $\sin\theta$ が一対一でもあるので，$-a\le x\le a$ の各 $x$ に対して $\theta$ が一つに決まる。

\begin{figure}[htbp]
 \centering
 \begin{tikzpicture}[x=2cm,y=1.5cm]
  \draw[-{Stealth}] (-1.95,0)--(1.95,0) node[right] {$\theta$};
  \draw[-{Stealth}] (0,-1.3)--(0,1.35) node[above] {$\sin\theta$};
  \draw[thick,blue,domain=-1.570796:1.570796,samples=70]
    plot (\x,{sin(\x r)});
  \draw[dashed] (-1.570796,0)--(-1.570796,-1)--(0,-1);
  \draw[dashed] (1.570796,0)--(1.570796,1)--(0,1);
  \node[above] at (-1.570796,0) {$-\frac\pi2$};
  \node[below] at (1.570796,0) {$\frac\pi2$};
  \node[above left] at (0,0) {$0$};
  \node[left] at (0,1) {$1$};
  \node[right] at (0,-1) {$-1$};
 \end{tikzpicture}
 \caption{正弦はこの区間で狭義単調増加し，逆向きの対応を一つに決められる。}
 \label{図:正弦の枝}
\end{figure}

\begin{example}{半円の根号を三角関数へ換える}
$a>0$ とする。$\int_0^a\sqrt{a^2-x^2}\dd x$ を求めよ。
\end{example}

ここでは $x\ge0$ だから，$x=a\sin\theta$ と置き，$0\le\theta\le\frac\pi2$ を選ぶ。$x=0$ は $\theta=0$，$x=a$ は $\theta=\frac\pi2$ に対応する。さらに
\[
 \dd x=a\cos\theta\dd\theta,\qquad
 \sqrt{a^2-x^2}=a\cos\theta
\]
なので，
\[
\begin{aligned}
 \int_0^a\sqrt{a^2-x^2}\dd x
 &=a^2\int_0^{\frac\pi2}\cos^2\theta\dd\theta\\
 &=a^2\left[\frac\theta2+\frac14\sin(2\theta)\right]_0^{\frac\pi2}
 =\frac{\pi a^2}{4}.
\end{aligned}
\]
曲線は $x^2+y^2=a^2$ の上半分であり，積分は半径 $a$ の円の四分の一の面積だから，幾何学的な結果とも一致する。積分計算そのものは，三角置換と二倍角の公式で得られている。

$0\le\theta\le\pi$ を選ぶと，$\sin\theta$ は $0$ から $1$ へ増えた後，$0$ へ戻る。この区間全体を $0\le x\le a$ と一対一に対応させることはできず，後半では $\cos\theta$ も負になる。枝を指定することは，こうした重複と符号を確定するために必要なのである。

もう一つ，$1+x^2$ には $1+\tan^2\theta=\frac1{\cos^2\theta}$ が対応する。$x=\tan\theta$ と置くと，微分も $\dd x=\frac1{\cos^2\theta}\dd\theta$ になる。この一致により，分母が消える場合がある。

\begin{example}{分母と微分の係数が打ち消し合う}
$\int_0^1\frac1{1+x^2}\dd x$ を求めよ。
\end{example}

$x=\tan\theta$ と置き，$0\le\theta\le\frac\pi4$ を選ぶ。この区間で正接は狭義単調増加し，端点は $x=0,1$ に対応する。よって
\[
\begin{aligned}
 \int_0^1\frac1{1+x^2}\dd x
 &=\int_0^{\frac\pi4}
   \frac1{1+\tan^2\theta}\,
   \frac1{\cos^2\theta}\dd\theta\\
 &=\int_0^{\frac\pi4}1\dd\theta
 =\frac\pi4.
\end{aligned}
\]
不定積分として表すなら，最後に $\tan\theta=x$ を満たす $\theta$ を $x$ の関数として書くことになる。この逆関数の名前，枝，微分は\chapterref{章:置換の背景}で詳しく扱う。ここでは定積分の端点が既知の角に対応するため，逆関数の記号を導入しなくても計算を終えられる。

\section{置換の後にも式を整える}\label{節:置換後の整理}

置換を選ぶとき，「すべてがすぐ基本公式になるか」だけを条件にすると，有用な置換を見落とす。根号をなくす，指数の中身を一つにする，分母の二次式を一次式にするなど，まず何を簡単にしたいかを決める。その後で，前章の式の整理を使えばよい。

\begin{example}{残った因子を新しい変数で表す}
$\int\frac{x^3}{1+x^2}\dd x$ を求めよ。
\end{example}

分母の $1+x^2$ を $t$ と置くと，$\dd t=2x\dd x$ である。分子には $x^3$ があるが，$x^3\dd x=x^2(x\dd x)$ と分ければ，$x\dd x=\frac12\dd t$ を使える。残る $x^2$ も $t-1$ と書けるので，
\[
\begin{aligned}
 \int\frac{x^3}{1+x^2}\dd x
 &=\frac12\int\frac{t-1}{t}\dd t\\
 &=\frac12\int\left(1-\frac1t\right)\dd t
 =\frac12t-\frac12\log t+C.
\end{aligned}
\]
$t=1+x^2>0$ を戻し，定数 $\frac12$ を任意定数へまとめると
\[
 \int\frac{x^3}{1+x^2}\dd x
 =\frac{x^2}{2}-\frac12\log(1+x^2)+C.
\]
$t=x^2+1$ は全実数上で一対一ではないが，残った $x^2$ を $t-1$ と書くのに逆関数 $x=h(t)$ は使っていない。$x$ の符号を選ぶ必要もない。

別の道もある。最初に割り算をすれば
\[
 \frac{x^3}{1+x^2}=x-\frac{x}{1+x^2}
\]
なので，第一項をそのまま積分し，第二項だけに置換を使って同じ答えを得る。どちらの場合も，式のどこを簡単にするかが分かっていれば，計算の途中で方針を見失わない。

これまでの置換では，複雑なまとまりとその微分が結び付いていた。しかし $x\e^x$ のように，二つの因子を片方の合成関数の微分としてまとめにくい積もある。この場合には，積の微分を逆に読む方法が役立つ。次章では，積分する因子と微分する因子を分けて考える。

\section*{章末問題}
\begin{enumerate}[label=\textbf{\thechapter.\arabic*},leftmargin=*]
 \item\label{問:置換べき} $\int(6x-4)(3x^2-4x+7)^5\dd x$ を求めよ。
 \item\label{問:置換二重対数} $\int\frac1{x\log x}\dd x$ を求めよ。実数の範囲で，どの区間に答えを適用できるかも述べよ。
 \item\label{問:置換非単調} $\int_{-2}^1 2x\e^{x^2}\dd x$ を求めよ。置換が単調でなくても成り立つ理由を説明せよ。
 \item\label{問:置換根号分母} $\int_0^3\frac1{1+\sqrt{1+x}}\dd x$ を求めよ。
 \item\label{問:置換円弧} $a>0$ とする。$x=a\sin\theta$ と置いて，$\int_0^{\frac a2}\sqrt{a^2-x^2}\dd x$ を求めよ。$\theta$ の区間も指定せよ。
 \item\label{問:置換正接} $x=\tan\theta$ と置いて，$\int_0^{\sqrt3}\frac1{1+x^2}\dd x$ を求めよ。
 \item\label{問:置換枝の違い} $t=x^2$ と置くと $x=-1,1$ がともに $t=1$ へ移る。このことから $\int_{-1}^1 1\dd x=0$ としてはいけない理由を述べよ。一方，$\int_{-1}^1\frac{2x}{1+x^2}\dd x=0$ は，同じ置換で正しく求められることを示せ。
\end{enumerate}

\smallskip\noindent 解答は\pageref{解答:第5章}ページ。

\chapter{部分積分}\label{章:部分積分}

前章では，合成関数の微分を逆向きに読んで，複雑なまとまりを一つの変数に換えた。たとえば $x(x^2+1)^{20}$ には，$x^2+1$ の微分に当たる $2x$ が，定数倍を除いて含まれていた。ところが $x\e^x$ では，$x$ と $\e^x$ のどちらかを新しい変数にしても，もう片方と $\dd x$ が一緒に簡単な形へ変わるわけではない。

ここで，積の微分を思い出す。$x\e^x$ を微分すると，$x\e^x$ のほかに $\e^x$ が現れる。したがって，$x\e^x$ 自体は求める原始関数ではないが，余分に現れる $\e^x$ の原始関数を引けばよさそうである。このように，積を微分したときの二つの項を使って，ある積分を別の積分へ換えるのが部分積分である。

\section{積の微分から公式を作る}\label{節:部分積分の導出}

二つの関数 $u(x),v(x)$ が微分できるとき，
\[
 \{u(x)v(x)\}'=u'(x)v(x)+u(x)v'(x)
\]
である。右辺の第二項を積分したいなら，この等式を
\[
 u(x)v'(x)=\{u(x)v(x)\}'-u'(x)v(x)
\]
と書き直す。両辺の原始関数を考えると，積 $uv$ と，もう一つの積分が現れる。

\begin{theorem}{部分積分の公式}
$u,v$ がある区間で連続微分可能ならば，その区間で
\[
 \int u(x)v'(x)\dd x
 =u(x)v(x)-\int u'(x)v(x)\dd x
\]
が成り立つ。不定積分は任意定数を含めた原始関数の族として読む。特に，同じ区間の $a,b$ に対して
\[
 \int_a^b u(x)v'(x)\dd x
 =[u(x)v(x)]_a^b-\int_a^b u'(x)v(x)\dd x
\]
が成り立つ。
\end{theorem}

定積分の公式では，積の微分を $a$ から $b$ まで積分し，基本定理を使った。$[uv]_a^b$ は $u(b)v(b)-u(a)v(a)$ である。この端点の項を省くことはできない。積の微分のうち，一方の項だけを移項したことに対応する量である。

公式の左辺では $u$ をそのまま，$v'$ を微分された形で書いている。右辺へ移ると，$u$ は $u'$ へ微分され，$v'$ は原始関数 $v$ に変わる。したがって実際の選択は，「どの因子を微分し，どの因子を積分すると，残る積 $u'v$ が扱いやすくなるか」という問題になる。

ときには $\dd v=v'(x)\dd x$，$\dd u=u'(x)\dd x$ とまとめて，
\[
 \int u\dd v=uv-\int v\dd u
\]
とも書く。これは同じ公式の短い表記であり，$u$ と $v$ を関係のない変数として自由に動かしているのではない。本章では，微分する因子と積分する因子を確かめやすいように，主に $u,v'$ を書いて計算する。

\begin{example}{一次式を微分して定数にする}
$\int x\e^x\dd x$ を求めよ。
\end{example}

$x$ は微分すると $1$ になる。一方，$\e^x$ は積分しても $\e^x$ なので，この二つを組み合わせれば残る積分が簡単になる。そこで
\[
 u=x,\qquad v'=\e^x,
 \qquad u'=1,\qquad v=\e^x
\]
と選ぶ。部分積分から
\[
\begin{aligned}
 \int x\e^x\dd x
 &=x\e^x-\int 1\cdot\e^x\dd x\\
 &=x\e^x-\e^x+C
 =(x-1)\e^x+C.
\end{aligned}
\]
積 $x\e^x$ を作って終わりにすると，その微分は $x\e^x+\e^x$ となる。余分な $\e^x$ を取り除くために $\e^x$ を引いている。答えの微分も
\[
 \frac{\dd}{\dd x}\{(x-1)\e^x\}
 =\e^x+(x-1)\e^x=x\e^x
\]
となり，この考え方をそのまま確認している。

逆に $u=\e^x$，$v'=x$ と選ぶことも，公式上は正しい。その場合には
\[
 \int x\e^x\dd x
 =\frac{x^2}{2}\e^x-\frac12\int x^2\e^x\dd x
\]
となる。ところが，もとの一次式が二次式になり，残った積分はむしろ手間が増えている。部分積分の公式を使えることと，その選択で計算が進むことは別である。右辺に何が残るかを一度書いてから，因子の選び方を判断すればよい。

\section{微分すると簡単になる因子を選ぶ}\label{節:部分積分の因子選択}

整式は微分すると次数が下がるため，微分する側の候補になる。指数関数や三角関数は，積分しても同じ種類の関数のままなので，積分する側に選びやすい。ただし，この順序は絶対の規則ではない。対数を含む積では，対数を微分すると分数になり，別の因子と約分できることがある。

また，見た目が積でなくても，$1$ を掛けた積と考えられる。$\log x$ の原始関数を直接思い出せない場合でも，$\log x\cdot1$ と見れば，積分する側に既知の関数 $1$ を置ける。

\begin{example}{一つの関数を積と見る}
$x>0$ において，$\int\log x\dd x$ を求めよ。
\end{example}

$\log x$ の微分は $\frac1x$ である。これに $1$ の原始関数 $x$ が掛かれば，積は $1$ になる。そこで
\[
 u=\log x,\qquad v'=1,
 \qquad u'=\frac1x,\qquad v=x
\]
と選ぶと，
\[
\begin{aligned}
 \int\log x\dd x
 &=x\log x-\int x\cdot\frac1x\dd x\\
 &=x\log x-x+C.
\end{aligned}
\]
もとの対数を，微分して得られる分数と，積分して得られる $x$ の組合せで消したのである。微分による確認は
\[
 \frac{\dd}{\dd x}(x\log x-x)
 =\log x+x\cdot\frac1x-1=\log x
\]
となる。$\log x$ は $x>0$ で考えているので，この計算にも同じ条件が付く。

この選び方は $x\log x$ にも使える。$\log x$ を微分し，$x$ を積分すると，
\[
\begin{aligned}
 \int x\log x\dd x
 &=\frac{x^2}{2}\log x
   -\int\frac{x^2}{2}\cdot\frac1x\dd x\\
 &=\frac{x^2}{2}\log x-\frac{x^2}{4}+C.
\end{aligned}
\]
ここでは整式 $x$ を積分する側に選んだ。それでも，対数を微分して得た $\frac1x$ と掛けることで，残る積分が単なる一次式になる。関数の名前だけで順序を決めず，$u'v$ を具体的に見ることが大切である。

一方，$\int\frac{\log x}{x}\dd x$ ならば，前章で見たとおり $t=\log x$ と置くほうが短い。$\frac1x\dd x=\dd t$ となり，ただちに $\int t\dd t$ へ変わるからである。「対数があるから部分積分」と決めるのではなく，対数の微分がすでに残りの部分にあるかを先に調べよう。

因子の選択には，もう一つ条件がある。積分する側の原始関数 $v$ を，その時点で求められなければならない。たとえば $\int x\e^{x^2}\dd x$ で，$x$ を微分しようとして $v'=\e^{x^2}$ と選ぶと，そこで原始関数が分からなくなる。しかし $t=x^2$ と置けば，$x\dd x=\frac12\dd t$ なので，答えは $\frac12\e^{x^2}+C$ である。部分積分を試す前に，置換に必要な微分の因子がそろっていないかを見ると，遠回りを避けられる。

\section{同じ選択を繰り返す}\label{節:部分積分の反復}

整式と指数関数，あるいは整式と三角関数の積では，一度の部分積分で整式の次数が一つ下がる。残る積分に同じ方針を適用すれば，やがて整式が定数になる。最初から一度で計算を終わらせようとする必要はない。何回か繰り返せば，確実に既知の積分へ到達することが見えている。

\begin{example}{二次式の次数を一つずつ下げる}
$\int x^2\cos x\dd x$ を求めよ。
\end{example}

まず $u=x^2$，$v'=\cos x$ と選ぶ。$u'=2x$，$v=\sin x$ なので，
\[
 \int x^2\cos x\dd x
 =x^2\sin x-2\int x\sin x\dd x.
\]
残った積分では，$x$ を微分し，$\sin x$ を積分する。$\sin x$ の原始関数は $-\cos x$ だから，
\[
\begin{aligned}
 \int x\sin x\dd x
 &=x(-\cos x)-\int1\cdot(-\cos x)\dd x\\
 &=-x\cos x+\sin x+C.
\end{aligned}
\]
これをもとの式へ戻して，任意定数をまとめると
\[
\begin{aligned}
 \int x^2\cos x\dd x
 &=x^2\sin x-2(-x\cos x+\sin x)+C\\
 &=x^2\sin x+2x\cos x-2\sin x+C.
\end{aligned}
\]
二回目の積分で現れた負号だけでなく，その式全体に最初の $-2$ が掛かることも確かめる。途中の積分を独立に計算したときは，戻す先の係数を残しておくと，符号の取り違えを防げる。

答えを微分すると，途中で現れた項が対になって消える。
\[
\begin{aligned}
 &\frac{\dd}{\dd x}
  \{x^2\sin x+2x\cos x-2\sin x\}\\
 &\qquad=2x\sin x+x^2\cos x
       +2\cos x-2x\sin x-2\cos x\\
 &\qquad=x^2\cos x.
\end{aligned}
\]
部分積分で補った項は，このように一つ前の微分で生まれた余分な項を打ち消している。繰り返す計算は，思い付きで項を増やす操作ではなく，積の微分による余分な項を順に取り除く操作である。

一般に，$P(x)$ が $n$ 次の整式なら，$P,P',P'',\ldots$ と微分を繰り返すと，$n$ 回目に定数になり，その次には $0$ になる。したがって $P(x)\e^{ax}$ や $P(x)\sin(ax)$，$P(x)\cos(ax)$ は，$a\ne0$ のとき，整式を微分する部分積分を繰り返して計算できる。$a$ が付く場合は，指数関数や三角関数を積分するたびに $a$ で割る。次数が下がることと，係数が変わることを別々に確認すればよい。

\section{もとの積分が戻ったら等式を解く}\label{節:部分積分の循環}

指数関数と三角関数の積では，整式のように微分を繰り返して $0$ になる因子がない。その代わり，正弦と余弦は微分によって互いに移り変わる。二回の部分積分で，もとの積分が再び現れることがある。そこで計算を失敗と決めず，もとの積分を未知数として含む等式が得られていないかを調べる。

\begin{example}{指数関数と正弦の積}
$\int\e^x\sin x\dd x$ を求めよ。
\end{example}

積分定数に惑わされずに等式を整理するため，まず
\[
 I(x)=\int_0^x\e^t\sin t\dd t,\qquad
 J(x)=\int_0^x\e^t\cos t\dd t
\]
と定める。$I,J$ は $x$ の関数であり，基本定理から $I'(x)=\e^x\sin x$ である。したがって $I(x)$ が求まれば，それに任意定数を加えて不定積分の答えにできる。

第一の積分で $\sin t$ を微分し，$\e^t$ を積分すると，
\[
\begin{aligned}
 I(x)
 &=[\e^t\sin t]_0^x-\int_0^x\e^t\cos t\dd t\\
 &=\e^x\sin x-J(x).
\end{aligned}
\]
続いて $J(x)$ でも $\e^t$ を積分する側に選ぶ。$\cos t$ の微分は $-\sin t$ だから，
\[
\begin{aligned}
 J(x)
 &=[\e^t\cos t]_0^x-\int_0^x\e^t(-\sin t)\dd t\\
 &=\e^x\cos x-1+I(x).
\end{aligned}
\]
ここで $I(x)$ が戻ってきた。第二式を第一式へ代入すると
\[
\begin{aligned}
 I(x)
 &=\e^x\sin x-\{\e^x\cos x-1+I(x)\},\\
 2I(x)&=\e^x(\sin x-\cos x)+1.
\end{aligned}
\]
よって $I(x)=\frac12\e^x(\sin x-\cos x)+\frac12$ である。最後の $\frac12$ も任意定数へまとめれば，
\[
 \boxed{\int\e^x\sin x\dd x
 =\frac12\e^x(\sin x-\cos x)+C}
\]
を得る。微分による確認は
\[
 \frac12\e^x(\sin x-\cos x)
 +\frac12\e^x(\cos x+\sin x)
 =\e^x\sin x
\]
である。

ここで役立ったのは，「同じ積分が戻る」という事実だけではない。右辺へ $-I(x)$ として戻ったので，移項すると $2I(x)$ が得られたことが重要である。もし同じ式を単に逆向きにたどり，$I(x)=I(x)$ に戻っただけなら，新しい情報はない。循環した後には，何倍の $I(x)$ が残ったのかを確かめる。

不定積分のまま同じ計算を進めることもできる。その場合は，途中の等式が任意定数の違いを許していることを忘れず，最後に一つの $C$ にまとめる。定積分として $I,J$ を定めたのは，等式の両辺を同じ関数として固定し，移項の意味をはっきりさせるためである。

\section{置換で形を変えてから部分積分する}\label{節:置換と部分積分}

積分法は，問題ごとに一つだけ選ぶものではない。置換によって合成関数をほどくと，今度は二つの関数の積が現れ，部分積分が使えることがある。前章でも，根号を変数に換えた後で割り算を使った。同じように，途中の式に合わせて次の道具を選べばよい。

\begin{example}{対数を角に持つ三角関数}
$x>0$ において，$\int\sin(\log x)\dd x$ を求めよ。
\end{example}

角の $\log x$ を一つの変数 $t$ にしたい。しかし，もとの積分には $\frac1x$ がないので，$\sin t\dd t$ と置き換えることはできない。$t=\log x$ を $x=\e^t$ と解き直すと，
\[
 \dd x=\e^t\dd t
\]
である。この微分の係数も含めて換えると，
\[
 \int\sin(\log x)\dd x
 =\int\e^t\sin t\dd t
\]
となる。これは直前の例題で求めた形だから，
\[
\begin{aligned}
 \int\sin(\log x)\dd x
 &=\frac12\e^t(\sin t-\cos t)+C\\
 &=\frac x2\{\sin(\log x)-\cos(\log x)\}+C.
\end{aligned}
\]
$\log x$ は $x>0$ で狭義単調増加し，逆向きの対応 $x=\e^t$ も一つに決まっている。置換後に $\e^t$ が現れたのは，新しい問題を余計に作ったためではなく，$\dd x$ を正しく書き換えたためである。

検算では，外にある $x$ と，内側の $\log x$ の両方を微分する。
\[
\begin{aligned}
 &\frac{\dd}{\dd x}
 \left[\frac x2\{\sin(\log x)-\cos(\log x)\}\right]\\
 &\qquad=\frac12\{\sin(\log x)-\cos(\log x)\}
   +\frac x2\cdot\frac1x
     \{\cos(\log x)+\sin(\log x)\}\\
 &\qquad=\sin(\log x).
\end{aligned}
\]
$t=\log x$ と置くだけでは終わらず，置換後の積を部分積分で求め，最後に元の変数へ戻した。この順序が，それぞれの手法の役割を分けている。

\section{根号を含む積分を同じ積分へ戻す}\label{節:平方根と部分積分}

$\sqrt{1+x^2}$ も，そのままでは基本公式の形ではない。前章の三角置換から $x=\tan\theta$ を考えることはできるが，ここでは部分積分を使う別の道を調べる。手がかりは，根号を微分すると $\frac{x}{\sqrt{1+x^2}}$ になり，さらに $x$ を掛ければ分子に $x^2$ ができることである。この $x^2$ を $(1+x^2)-1$ と分けると，もとの根号が再び現れる。

\begin{example}{根号とその逆数を関係付ける}
$\int\sqrt{1+x^2}\dd x$ を求めよ。
\end{example}

まず $u=\sqrt{1+x^2}$，$v'=1$ と選ぶ。$u'=\frac{x}{\sqrt{1+x^2}}$，$v=x$ だから，部分積分から
\[
 \int\sqrt{1+x^2}\dd x
 =x\sqrt{1+x^2}-\int\frac{x^2}{\sqrt{1+x^2}}\dd x
\]
を得る。右辺の積分を見ただけでは簡単になったか分かりにくいが，分子を分けると
\[
 \frac{x^2}{\sqrt{1+x^2}}
 =\frac{(1+x^2)-1}{\sqrt{1+x^2}}
 =\sqrt{1+x^2}-\frac1{\sqrt{1+x^2}}
\]
となる。したがって，もとの積分を移項すれば，任意定数を除いて
\[
 2\int\sqrt{1+x^2}\dd x
 =x\sqrt{1+x^2}+\int\frac1{\sqrt{1+x^2}}\dd x
\]
という関係が得られる。未知の積分がもう一度現れたが，その係数をまとめることができた。残る課題は，根号の逆数の積分である。

ここで $r(x)=\sqrt{1+x^2}$ と略記して，$x+r(x)$ を微分してみる。
\[
 \{x+r(x)\}'=1+\frac{x}{r(x)}
 =\frac{x+r(x)}{r(x)}.
\]
この形には，微分する前の $x+r(x)$ が分子に現れている。そこで対数の微分を使えば，その因子が約分されて
\[
 \frac{\dd}{\dd x}\log\{x+r(x)\}
 =\frac{1}{x+r(x)}\cdot\frac{x+r(x)}{r(x)}
 =\frac1{r(x)}
\]
となる。しかも $r(x)=\sqrt{1+x^2}>|x|$ なので，$x+r(x)>0$ はすべての実数 $x$ で成り立つ。対数の中身の符号も確定した。よって
\[
 \int\frac1{\sqrt{1+x^2}}\dd x
 =\log\left(x+\sqrt{1+x^2}\right)+C.
\]
これを先ほどの関係へ戻すと，
\[
 \boxed{
 \int\sqrt{1+x^2}\dd x
 =\frac12\left\{x\sqrt{1+x^2}
       +\log\left(x+\sqrt{1+x^2}\right)\right\}+C
 }
\]
を得る。

$x+r(x)$ は，突然覚えるべき式として出てきたのではない。$r'(x)=\frac{x}{r(x)}$ と $x'=1=\frac{r(x)}{r(x)}$ を足すと，分子がもとの和になる。そのため，対数の微分に現れる「もとの式で割る」という操作が，求めたい $\frac1{r(x)}$ を作るのである。

最後に答え全体を微分すると，
\[
\begin{aligned}
 &\frac12\left\{\sqrt{1+x^2}
       +\frac{x^2}{\sqrt{1+x^2}}
       +\frac1{\sqrt{1+x^2}}\right\}\\
 &\qquad=\frac12\{\sqrt{1+x^2}+\sqrt{1+x^2}\}
 =\sqrt{1+x^2}
\end{aligned}
\]
となる。この確認により，途中で不定積分を移項した計算も，最終的な原始関数として正しいことが分かる。後で双曲線関数を学ぶと，同じ根号を別の置換で扱える。ここで得た式は，その別解と比較する基準になる。

\section{定積分では端点の項を先に見る}\label{節:定積分の部分積分}

不定積分を求めてから上下端を代入してもよいが，最初から定積分の公式を使えば，端点で消える項がすぐに分かる。公式
\[
 \int_a^b uv'\dd x=[uv]_a^b-\int_a^b u'v\dd x
\]
では，残る積分だけでなく，端点の $u(a),u(b),v(a),v(b)$ も計算の見通しに関わる。積のどちらかが $0$ ならば，その端点の寄与はなくなる。

\begin{example}{端点で消える積を確かめる}
$\int_0^{\frac\pi2}x\sin x\dd x$ を求めよ。
\end{example}

$u=x$，$v'=\sin x$ と選ぶと，$u'=1$，$v=-\cos x$ なので，
\[
\begin{aligned}
 \int_0^{\frac\pi2}x\sin x\dd x
 &=[-x\cos x]_0^{\frac\pi2}
   +\int_0^{\frac\pi2}\cos x\dd x\\
 &=0+[\sin x]_0^{\frac\pi2}=1.
\end{aligned}
\]
上端では $\cos\frac\pi2=0$，下端では $x=0$ のため，端点の項はともに $0$ になった。二つの端点で消える理由は異なるので，一方だけを見て項全体を消しているわけではない。

また，積分される関数はこの区間で非負だから，求めた値が正であることは自然である。仮に $-1$ になったら，$\sin x$ の原始関数の負号か，部分積分の前の負号を確かめるべきだと分かる。

\begin{example}{部分積分の後に割り算を使う}
$\int_0^1x\log(1+x)\dd x$ を求めよ。
\end{example}

$\log(1+x)$ を微分すると $\frac1{1+x}$ になる。$x$ の原始関数は $\frac{x^2}{2}$ だから，残る積分は有理関数になる。そこで
\[
 u=\log(1+x),\qquad v'=x
\]
と選び，
\[
 \int_0^1x\log(1+x)\dd x
 =\left[\frac{x^2}{2}\log(1+x)\right]_0^1
  -\frac12\int_0^1\frac{x^2}{1+x}\dd x
\]
を得る。第一項は $\frac12\log2$ である。第二項では，分子の次数が分母より高いので，\chapterref{章:基本積分}で学んだ割り算を使う。
\[
 x^2=(1+x)(x-1)+1,
 \qquad
 \frac{x^2}{1+x}=x-1+\frac1{1+x}.
\]
したがって
\[
\begin{aligned}
 \int_0^1\frac{x^2}{1+x}\dd x
 &=\left[\frac{x^2}{2}-x+\log(1+x)\right]_0^1\\
 &=-\frac12+\log2.
\end{aligned}
\]
これを戻すと
\[
 \int_0^1x\log(1+x)\dd x
 =\frac12\log2-\frac12\left(-\frac12+\log2\right)
 =\frac14.
\]
部分積分の直後には対数が端点に残り，別の積分にも対数が現れたが，最後に差を取ると打ち消し合った。途中の式の見た目だけで，答えにも対数が必ず残ると決めることはできない。

ここまでで，式の整理，置換積分，部分積分という三つの道具を学んだ。整式の割り算や三角恒等式で和へ直せるなら，まずその変形を試せる。合成関数と内側の微分がそろっていれば置換を考え，微分すると簡単になる因子と，積分できる因子の積なら部分積分を考える。ある操作で何が簡単になり，何が残るかを具体的に見ることが，選択の根拠になる。

ただし，定積分の値だけが欲しいときは，原始関数を最後まで求める必要がない場合もある。区間の左右を入れ換えたときに式がどう変わるか，正負の部分がどう対応するかを使えば，値が決まることがある。次の\chapterref{章:定積分の対称性}では，被積分関数と積分区間を一緒に見て，その対応を計算へ生かす。

\section*{章末問題}
\begin{enumerate}[label=\textbf{\thechapter.\arabic*},leftmargin=*]
 \item\label{問:部分一次指数} $\int x\e^{3x}\dd x$ を求めよ。微分する因子と積分する因子を示し，答えを微分して確かめよ。
 \item\label{問:部分整式対数} $x>0$ において，$\int x^2\log x\dd x$ を求めよ。$x^2$ と $\log x$ のどちらを微分する側に選ぶか，その理由も述べよ。
 \item\label{問:部分反復} $\int x^3\sin x\dd x$ を求めよ。部分積分を繰り返すたびに，整式の次数がどう変わるかを確かめよ。
 \item\label{問:部分循環} 部分積分を二回使い，$\int\e^{2x}\cos x\dd x$ を求めよ。もとの積分が戻った後の等式も書け。
 \item\label{問:部分対数余弦} $x>0$ において，$\int\cos(\log x)\dd x$ を求めよ。置換の際には，$\dd x$ が何に変わるかも書け。
 \item\label{問:部分置換定積分} $\int_0^1x\log(1+x^2)\dd x$ を求めよ。置換の後に，どの積分を部分積分で求めればよいかを示せ。
 \item\label{問:部分根号定積分} 部分積分を使い，$\int_0^1\sqrt{4+x^2}\dd x$ を求めよ。必要になる根号の逆数の原始関数は，微分によって確かめて用いよ。
 \item\label{問:部分情報だけの計算} $f$ は $[0,1]$ で連続微分可能とし，$f(0)=1$，$f(1)=2$，$\int_0^1f(x)\dd x=\frac32$ を満たす。$f$ の具体的な式を求めずに，$\int_0^1xf'(x)\dd x$ と $\int_0^1(1-x)f'(x)\dd x$ を求めよ。
\end{enumerate}

\smallskip\noindent 解答は\pageref{解答:第6章}ページ。

\chapter{定積分の対称性}\label{章:定積分の対称性}

不定積分では，微分すると与えられた関数になる式を探した。定積分でも原始関数が見つかれば計算できるが，定積分にはもう一つ情報がある。どこからどこまで積分するかという区間である。たとえば，同じ関数でも区間が原点の左右に同じ長さだけ延びていれば，左右の値を組にして考えられる。区間の端が互いに逆数なら，逆数の置換によって両端を交換できる。

この章では，置換後の積分をすぐに計算せず，置換前の積分と並べてみる。二つを加えたとき，分母が消えることもあれば，変数を掛けた部分が定数に変わることもある。そのため，原始関数を求めるには手間がかかる式から，短い計算で定積分の値が決まる。

以下では，特に断らない限り，被積分関数は積分する閉区間で連続であるとする。途中に分母が零になる点がある式に，対称性だけを見て同じ議論を当てはめることはできない。特異点を含む積分は，後の章で定義を確かめてから扱う。

\section{左右の値を組にする}\label{節:奇偶の対称性}

\begin{definition}{偶関数と奇関数}
原点に関して対称な範囲で定義された関数 $f$ について，すべての $x$ で
\[
f(-x)=f(x)
\]
が成り立つとき，$f$ を偶関数という。また，すべての $x$ で
\[
f(-x)=-f(x)
\]
が成り立つとき，$f$ を奇関数という。
\end{definition}

偶関数のグラフは $y$ 軸について対称であり，奇関数のグラフは原点について対称である。$x^2$ と $\cos x$ は偶関数，$x^3$ と $\sin x$ は奇関数である。ただし，式の中に $\sin$ があれば奇関数になるわけではない。$\sin(x^2)$ では $x$ を $-x$ に変えても $x^2$ が変わらないので偶関数である。式全体に $-x$ を代入することが，確実な判定になる。

$a>0$ とし，$[-a,a]$ の積分を零で分ける。左半分で $x=-t$ と置くと，$x=-a,0$ に対して $t=a,0$ であり，$\dd x=-\dd t$ だから，
\[
\int_{-a}^{0}f(x)\dd x
 =\int_a^0 f(-t)(-\dd t)
 =\int_0^a f(-t)\dd t.
\]
積分変数の文字をそろえれば，
\begin{equation}\label{式:左右の積分}
\int_{-a}^a f(x)\dd x
 =\int_0^a\{f(x)+f(-x)\}\dd x
\end{equation}
となる。これは偶関数でも奇関数でもない場合にも使える式である。偶関数なら中括弧が $2f(x)$ になり，奇関数なら零になる。

\begin{theorem}{対称区間での偶関数・奇関数の積分}
$a>0$ とする。$f$ が偶関数なら
\[
\int_{-a}^a f(x)\dd x=2\int_0^a f(x)\dd x,
\]
$f$ が奇関数なら
\[
\int_{-a}^a f(x)\dd x=0
\]
である。
\end{theorem}

図\ref{図:奇偶の積分}の右図では，左右の部分の幾何学的な面積は等しい。しかし左側の定積分は負，右側は正である。奇関数の積分が零という結論は，図形の面積が零という意味ではない。面積を求めるなら絶対値を積分する必要がある。

\begin{figure}[htbp]
\centering
\begin{tikzpicture}[scale=1.05,>=Stealth]
\begin{scope}[xshift=-3.1cm]
\fill[blue!12] plot[domain=-1.35:1.35,samples=50](\x,{0.65*\x*\x})--(1.35,0)--(-1.35,0)--cycle;
\draw[->](-1.7,0)--(1.8,0)node[right]{$x$};
\draw[->](0,-0.2)--(0,1.7)node[above]{$y$};
\draw[blue!70!black,thick] plot[domain=-1.45:1.45,samples=60](\x,{0.65*\x*\x});
\draw[dashed](-1.35,0)--(-1.35,1.185);
\draw[dashed](1.35,0)--(1.35,1.185);
\node[below]at(-1.35,0){$-a$};\node[below]at(1.35,0){$a$};
\node at(0,-0.65){偶関数：左右を加える};
\end{scope}
\begin{scope}[xshift=1.8cm]
\fill[blue!12](0,0)--plot[domain=0:1.35,samples=40](\x,{0.55*\x*\x*\x})--(1.35,0)--cycle;
\fill[orange!18](0,0)--plot[domain=0:-1.35,samples=40](\x,{0.55*\x*\x*\x})--(-1.35,0)--cycle;
\draw[->](-1.7,0)--(1.8,0)node[right]{$x$};
\draw[->](0,-1.55)--(0,1.7)node[above]{$y$};
\draw[blue!70!black,thick] plot[domain=-1.4:1.4,samples=60](\x,{0.55*\x*\x*\x});
\draw[dashed](-1.35,0)--(-1.35,-1.353);
\draw[dashed](1.35,0)--(1.35,1.353);
\node[above]at(-1.35,0){$-a$};\node[below]at(1.35,0){$a$};
\node at(0,-1.95){奇関数：符号付きの量が打ち消す};
\end{scope}
\end{tikzpicture}
\caption{対称性と，定積分の符号の関係。曲線はそれぞれ $x^2$ 型，$x^3$ 型である。}\label{図:奇偶の積分}
\end{figure}

\begin{example}{混じった式から奇関数の部分を除く}
次の定積分を求めよ。
\[
I=\int_{-2}^{2}\left(x^3\e^{-x^2}+2x^2+1\right)\dd x.
\]
\end{example}

$x^3\e^{-x^2}$ の原始関数を先に探すこともできるが，区間が $[-2,2]$ なので，まず $x$ を $-x$ に変えてみる。$x^3$ の符号だけが変わり，$\e^{-x^2}$ は変わらない。したがって積 $x^3\e^{-x^2}$ は奇関数で，その積分は零である。残った $2x^2+1$ は偶関数だから，
\[
I=2\int_0^2(2x^2+1)\dd x
 =2\left[\frac{2x^3}{3}+x\right]_0^2
 =\frac{44}{3}.
\]
ここでは式全体が偶関数か奇関数かを決める必要はなかった。和を分けて，対称性が分かる部分にだけ使えばよい。

さらに，一般の $f$ に対して
\[
f_{\mathrm e}(x)=\frac{f(x)+f(-x)}{2},\qquad
f_{\mathrm o}(x)=\frac{f(x)-f(-x)}{2}
\]
と置くと，$f=f_{\mathrm e}+f_{\mathrm o}$ であり，前者は偶関数，後者は奇関数になる。よって対称区間の定積分に寄与するのは $f_{\mathrm e}$ の部分だけである。これは新しい積分公式を覚えるための分解ではなく，式\eqref{式:左右の積分}が何を選び出しているかを示す分解である。

\section{区間の両端を交換する}\label{節:区間の反転}

対称の中心は原点に限らない。区間 $[a,b]$ の両端を交換するには，$x=a+b-t$ と置けばよい。対応は
\[
x=a\longleftrightarrow t=b,\qquad
x=b\longleftrightarrow t=a
\]
であり，$\dd x=-\dd t$ だから，
\begin{equation}\label{式:区間反転}
\int_a^b f(x)\dd x=\int_a^b f(a+b-x)\dd x
\end{equation}
を得る。置換の直後は積分の上下端が逆になるが，微分の負号と合わせて向きを戻している。

この式の左右は同じ値なので，加えて二で割ると
\begin{equation}\label{式:反転平均}
\int_a^b f(x)\dd x
 =\frac12\int_a^b\{f(x)+f(a+b-x)\}\dd x
\end{equation}
となる。$f(x)+f(a+b-x)$ が定数 $C$ になれば，答えは $C(b-a)$ の半分である。元の式を個別に積分できるかどうかは，この計算では問題にならない。

今度は，関数 $g$ が区間の中央について対称，すなわち
\[
g(a+b-x)=g(x)
\]
を満たすとする。このとき $xg(x)$ の積分を $I$ とすると，反転により
\[
I=\int_a^b(a+b-x)g(x)\dd x.
\]
元の式を加えると $x+(a+b-x)=a+b$ なので，
\begin{equation}\label{式:対称な重みの一次因子}
\int_a^b xg(x)\dd x
 =\frac{a+b}{2}\int_a^b g(x)\dd x
\end{equation}
となる。中央からの距離が等しい二点では $g$ の値が等しく，その二点の $x$ 座標の平均は常に区間の中央になる，という関係である。

\begin{example}{変数が掛かった三角関数}
次の定積分を求めよ。
\[
I=\int_0^\pi x\sin^3x\dd x.
\]
\end{example}

$x$ と三角関数の積を見ると部分積分を考えたくなる。しかしこの区間では $\sin(\pi-x)=\sin x$ である。$g(x)=\sin^3x$ は区間の中央について対称なので，式\eqref{式:対称な重みの一次因子}から
\[
I=\frac{\pi}{2}\int_0^\pi\sin^3x\dd x.
\]
残りは $\sin^3x=(1-\cos^2x)\sin x$ として $u=\cos x$ と置く。
\[
\int_0^\pi\sin^3x\dd x
 =\int_{-1}^1(1-u^2)\dd u
 =\left[u-\frac{u^3}{3}\right]_{-1}^1
 =\frac43.
\]
したがって $I=\frac{2\pi}{3}$ である。反転で $x$ を取り除き，置換で三角関数を取り除くという二段階で計算できた。

式\eqref{式:対称な重みの一次因子}を使う際には，$xg(x)$ 全体が対称なのではなく，$g$ が対称であることを確認する。たとえば $\int_0^1 x\e^x\dd x$ では $\e^{1-x}=\e^x$ が成り立たないので，$x$ を区間の中央 $\frac12$ に置き換えることはできない。

また，$x^2g(x)$ まで同じように中央の二乗に置き換えられるわけでもない。$x^2+(a+b-x)^2$ は $x$ を含むからである。反転して足すことで何が消えるのかを，その都度式で確かめる必要がある。

\section{正弦と余弦を交換する}\label{節:正弦余弦の交換}

区間 $\left[0,\frac{\pi}{2}\right]$ での反転は，$x$ を $\frac{\pi}{2}-x$ に変える操作である。このとき
\[
\sin\left(\frac{\pi}{2}-x\right)=\cos x,\qquad
\cos\left(\frac{\pi}{2}-x\right)=\sin x
\]
なので，正弦と余弦が入れ替わる。分子と分母に両者が現れる場合，入れ替えた式との和を調べるとよい。

\begin{example}{分母を変えずに分子を交換する}
次の定積分を求めよ。
\[
I=\int_0^{\frac{\pi}{2}}
\frac{\sin^4x}{\sin^4x+\cos^4x}\dd x.
\]
\end{example}

$\sin x$ と $\cos x$ は同時に零にならないので，分母は区間全体で正である。反転すると分子は $\cos^4x$ に変わり，分母は変わらない。したがって
\[
I=\int_0^{\frac{\pi}{2}}
\frac{\cos^4x}{\sin^4x+\cos^4x}\dd x.
\]
二つを加えれば，分子の和が分母そのものになる。
\[
2I=\int_0^{\frac{\pi}{2}}1\dd x=\frac{\pi}{2},
\qquad I=\frac{\pi}{4}.
\]
$\sin^4x$ の次数を下げたり，分母を展開したりする前に，分子の足し合わせで分母ができることに気づくと計算を減らせる。

この関係は四乗に特有ではない。$q$ を $[0,1]$ 上の連続関数とし，$q(\sin x)+q(\cos x)$ が区間内で零にならなければ，
\[
\int_0^{\frac{\pi}{2}}
\frac{q(\sin x)}{q(\sin x)+q(\cos x)}\dd x
 =\frac{\pi}{4}
\]
である。たとえば $q(t)=t^2+1$ でも同じ値になる。重要なのは指数の値ではなく，反転で分子の二つの候補が交換されることである。

図\ref{図:相補重み}は，例題の被積分関数を $w(x)$ としたグラフである。左右を反転したグラフ $w(\frac{\pi}{2}-x)$ と，同じ $x$ における高さを足すと一になる。それぞれの積分が等しいので，全体の長方形の面積を二等分している。

\begin{figure}[htbp]
\centering
\begin{tikzpicture}[x=3.5cm,y=2.3cm,>=Stealth]
\fill[blue!10](0,0)--plot[domain=0:1.570796,samples=100](\x,{pow(sin(\x r),4)/(pow(sin(\x r),4)+pow(cos(\x r),4))})--(1.570796,0)--cycle;
\draw[->](-0.08,0)--(1.73,0)node[right]{$x$};
\draw[->](0,-0.08)--(0,1.2)node[above]{$y$};
\draw[dashed](0,1)--(1.570796,1)--(1.570796,0);
\draw[dashed](0.785398,0)--(0.785398,1);
\draw[blue!75!black,thick]plot[domain=0:1.570796,samples=100](\x,{pow(sin(\x r),4)/(pow(sin(\x r),4)+pow(cos(\x r),4))});
\draw[orange!80!black,thick,dashed]plot[domain=0:1.570796,samples=100](\x,{pow(cos(\x r),4)/(pow(sin(\x r),4)+pow(cos(\x r),4))});
\node[left]at(0,1){$1$};\node[below]at(0,0){$0$};
\node[below]at(0.785398,0){$\frac{\pi}{4}$};
\node[below]at(1.570796,0){$\frac{\pi}{2}$};
\node[blue!75!black]at(1.15,0.65){$w(x)$};
\node[orange!80!black]at(0.48,1.26){$w(\frac{\pi}{2}-x)$};
\end{tikzpicture}
\caption{反転した二つの重みは，同じ場所で足すと一になる。}\label{図:相補重み}
\end{figure}

\section{周期があるとき，どこから一周期を取るか}\label{節:周期関数の積分}

$T>0$ に対して $f(x+T)=f(x)$ がすべての $x$ で成り立つとき，$T$ を $f$ の周期という。周期という言葉だけから，$T$ は最小の正の周期だとは限らない。$\sin x$ の周期として $2\pi$ も $4\pi$ も使える。

周期関数をちょうど一周期分だけ積分した値は，始点によらない。まず $0\leq a\leq T$ の場合，
\[
\int_a^{a+T}f(x)\dd x
 =\int_a^T f(x)\dd x+\int_T^{a+T}f(x)\dd x.
\]
第二項で $x=t+T$ と置くと $f(t+T)=f(t)$ だから，
\[
\int_T^{a+T}f(x)\dd x=\int_0^a f(t)\dd t.
\]
二つを足すと
\begin{equation}\label{式:一周期の積分}
\int_a^{a+T}f(x)\dd x=\int_0^T f(x)\dd x
\end{equation}
になる。一般の $a$ も，整数個の $T$ だけ平行移動すれば始点を $[0,T)$ に移せるため，同じ結論になる。区間を長さ $T$ の部分に分ければ，正の整数 $n$ に対して
\[
\int_a^{a+nT}f(x)\dd x=n\int_0^T f(x)\dd x
\]
も得られる。

周期があるというだけでは，一周期分の積分が零とは限らない。たとえば $\sin^2x$ は周期 $\pi$ を持つが，一周期の積分は $\frac{\pi}{2}$ である。符号付きの部分が打ち消し合うか，積分区間が周期の整数倍になっているかを別々に確認する。

\begin{example}{絶対値がある周期関数}
次の定積分を求めよ。
\[
I=\int_{-\frac{\pi}{6}}^{\frac{13\pi}{6}}|\sin x|\dd x.
\]
\end{example}

$\sin(x+\pi)=-\sin x$ なので，絶対値を取ると周期は $\pi$ になる。一周期分は
\[
\int_0^\pi|\sin x|\dd x=\int_0^\pi\sin x\dd x=2.
\]
与えられた区間の長さは $\frac{7\pi}{3}$ であり，$2\pi$ と $\frac{\pi}{3}$ に分けられる。まず $-\frac{\pi}{6}$ から $\frac{11\pi}{6}$ までは二周期分なので積分は $4$ である。残りでは $x=t+2\pi$ と置き，
\[
\int_{\frac{11\pi}{6}}^{\frac{13\pi}{6}}|\sin x|\dd x
 =\int_{-\frac{\pi}{6}}^{\frac{\pi}{6}}|\sin t|\dd t.
\]
$|\sin t|$ は偶関数で，$0\leq t\leq\frac{\pi}{6}$ では $\sin t\geq0$ だから，
\[
I=4+2\int_0^{\frac{\pi}{6}}\sin t\dd t
 =4+2\left(1-\frac{\sqrt3}{2}\right)
 =6-\sqrt3.
\]
周期によって長い区間を短くし，偶奇によって符号の変わる区間を半分にした。周期の整数倍に満たない残りは，始点によって値が変わるので省いてはいけない。

\section{逆数の置換で端を交換する}\label{節:逆数対称性}

区間の端が $\frac1b$ と $b$ であるときは，$x=\frac1t$ でも両端を交換できる。ただし反転 $x=a+b-t$ と異なり，微分は
\[
\dd x=-\frac{1}{t^2}\dd t
\]
になる。関数の $x$ を逆数に変えるだけでは，正しい置換にならない。

特に積分が $\frac{\dd x}{x}$ を含むと，
\[
\frac{\dd x}{x}=-\frac{\dd t}{t}
\]
となり，同じ形に戻る。$b>1$ とすると，
\begin{equation}\label{式:逆数置換の対称性}
\int_{\frac1b}^b F(x)\frac{\dd x}{x}
 =\int_{\frac1b}^b F\left(\frac1x\right)\frac{\dd x}{x}
\end{equation}
である。したがって，$F(x)+F(\frac1x)=C$ なら
\[
2\int_{\frac1b}^b F(x)\frac{\dd x}{x}
 =C\int_{\frac1b}^b\frac{\dd x}{x}
 =2C\log b
\]
となる。

\begin{example}{逆数にすると補い合う分数}
次の定積分を求めよ。
\[
I=\int_{\frac13}^3\frac{1}{x(1+x^2)}\dd x.
\]
\end{example}

区間の両端が互いに逆数であり，分母に $x$ がある。そこで $F(x)=\frac{1}{1+x^2}$ と見て，$\frac{\dd x}{x}$ を残す。
\[
F\left(\frac1x\right)=\cfrac{1}{1+\cfrac{1}{x^2}}
 =\frac{x^2}{1+x^2},
\qquad F(x)+F\left(\frac1x\right)=1.
\]
式\eqref{式:逆数置換の対称性}を使えば
\[
2I=\int_{\frac13}^3\frac{\dd x}{x}
 =\left[\log x\right]_{\frac13}^3=2\log3,
\]
よって $I=\log3$ である。

この例を直接置換しても，
\[
I=\int_3^{\frac13}
\cfrac{1}{\cfrac1t(1+\cfrac1{t^2})}
\left(-\frac{1}{t^2}\right)\dd t
=\int_{\frac13}^3\frac{t}{1+t^2}\dd t
\]
となる。元の被積分関数との和は $\frac1x$ であり，上の計算と一致する。微分因子を忘れると，この和は現れない。

逆数の対称性を，いつもの左右の対称性として見直すこともできる。$x=\e^u$ と置くと，$\frac{\dd x}{x}=\dd u$，区間は $[-\log b,\log b]$ になる。さらに $x$ を $\frac1x$ に変える操作は，$u$ を $-u$ に変える操作になる。数直線上の普通の長さでは $\frac1b$ と $b$ は $1$ から等距離ではないが，対数で測った座標では左右対称なのである。

一般に，端が正で $ab=k$ なら，$x=\frac{k}{t}$ は $[a,b]$ の両端を交換する。この場合にも $\frac{\dd x}{x}=-\frac{\dd t}{t}$ が成り立つ。端を見て置換を考えるときは，和が一定か，積が一定かが手がかりになる。

\section{足すと一になる重み}\label{節:相補重み}

次の二つの式を比べる。
\[
\frac{1}{1+\e^x},\qquad
\frac{1}{1+\e^{-x}}=\frac{\e^x}{1+\e^x}.
\]
分子の和が分母なので，二つを足すと一になる。そこで
\[
w(x)=\frac{1}{1+\e^x}
\]
と置けば $w(x)+w(-x)=1$ である。このように，対称な二点で値の和が一になる因子を，ここでは互いに補い合う重みとして考える。

$g$ が偶関数なら，
\[
I=\int_{-a}^a g(x)w(x)\dd x
 =\int_{-a}^a g(x)w(-x)\dd x
\]
である。二式を加えると
\begin{equation}\label{式:偶関数と相補重み}
\int_{-a}^a\frac{g(x)}{1+\e^x}\dd x
 =\frac12\int_{-a}^a g(x)\dd x
 =\int_0^a g(x)\dd x
\end{equation}
となる。指数関数の因子を一度も積分せずに，式から取り除けた。

さらに，$h$ が奇関数なら $h(-x)=-h(x)$ なので，$\frac{1}{1+\e^{h(x)}}$ にも同じ議論が使える。必要なのは指数が $x$ であることではなく，反転すると指数の符号が変わることである。

\begin{example}{指数が複雑でも，反転で符号が変わる}
次の定積分を求めよ。
\[
I=\int_{-2}^2\frac{x^4+\cos x}{1+\e^{x^3+x}}\dd x.
\]
\end{example}

分子 $g(x)=x^4+\cos x$ は偶関数で，指数 $h(x)=x^3+x$ は奇関数である。$x$ を $-x$ に変えると分子はそのまま，指数だけが負になる。したがって
\[
2I=\int_{-2}^2(x^4+\cos x)\dd x,
\]
よって
\[
I=\int_0^2(x^4+\cos x)\dd x
 =\left[\frac{x^5}{5}+\sin x\right]_0^2
 =\frac{32}{5}+\sin2.
\]
分子は区間全体で正である必要はない。偶関数であることだけが，上の等式に使われている。

一方，分子が奇関数のときは，同じ結論にはならない。たとえば
\[
K=\int_{-a}^a\frac{x}{1+\e^x}\dd x
\]
を左右で組にすると，
\[
K=\int_0^a x\left(\frac{1}{1+\e^x}-\frac{1}{1+\e^{-x}}\right)\dd x
 =\int_0^a\frac{x(1-\e^x)}{1+\e^x}\dd x.
\]
$a>0$ なら被積分関数は $0<x\leq a$ で負なので，$K<0$ である。「分子の奇関数を対称区間で積分するから零」と考えると，分母が偶関数でないことを見落とす。積全体の対称性か，反転後の式との関係を確かめる必要がある。

この考えは原点以外にも移せる。区間の中央を $c=\frac{a+b}{2}$ とし，$g(a+b-x)=g(x)$，$w(x)+w(a+b-x)=1$ なら
\[
\int_a^b g(x)w(x)\dd x=\frac12\int_a^b g(x)\dd x.
\]
たとえば $w(x)=\frac{1}{1+\e^{x-c}}$ はこの条件を満たす。変数を $u=x-c$ に移せば，原点に関する場合と全く同じ証明になる。

\section{対称性で何を簡単にするか}\label{節:定積分の手法選択}

ここまでの方法は，どれも新しい原始関数を増やすものではない。定積分を，値の等しい別の積分に書き換え，二つを組にする方法である。最初に区間を見て，次にその区間を保つ変換を考えると，計算の候補を絞れる。

\begin{center}
\begin{tabularx}{0.96\linewidth}{@{}lXX@{}}
\toprule
区間や式の特徴 & 試す変換 & 比べる部分\\
\midrule
$[-a,a]$ & $x\mapsto -x$ & 偶関数と奇関数の部分\\
$[a,b]$ & $x\mapsto a+b-x$ & 和が定数になるか，一次因子が消えるか\\
$\left[0,\frac{\pi}{2}\right]$ & $x\mapsto\frac{\pi}{2}-x$ & 正弦と余弦が交換されるか\\
長い区間と周期関数 & 整数周期だけ平行移動 & 完全な周期と残りの区間\\
$\left[\frac1b,b\right]$ と $\frac{\dd x}{x}$ & $x\mapsto\frac1x$ & 逆数にした関数との和\\
\bottomrule
\end{tabularx}
\end{center}

候補が見つかっても，対称性だけで最後まで解けるとは限らない。対称性によって厄介な因子を一つ除き，残った部分に置換積分や部分積分を使えばよい。

\begin{example}{反転の後で置換と部分分数を使う}
次の定積分を求めよ。
\[
I=\int_0^\pi\frac{x\sin x}{1+\sin^2x}\dd x.
\]
\end{example}

$\frac{\sin x}{1+\sin^2x}$ は $x$ を $\pi-x$ に変えても同じである。したがって，最初に $x$ を対称性で取り除く。
\[
I=\frac{\pi}{2}J,\qquad
J=\int_0^\pi\frac{\sin x}{1+\sin^2x}\dd x.
\]
$J$ では分子の $\sin x\dd x$ が $-\dd(\cos x)$ に対応する。$u=\cos x$ と置くと，$1+\sin^2x=2-u^2$ だから
\[
J=\int_{-1}^1\frac{1}{2-u^2}\dd u.
\]
この分母は $(\sqrt2-u)(\sqrt2+u)$ に分かれる。よって
\[
\frac{1}{2-u^2}
 =\frac{1}{2\sqrt2}\left(\frac{1}{\sqrt2-u}+\frac{1}{\sqrt2+u}\right),
\]
\[
J=\left[\frac{1}{2\sqrt2}
\log\frac{\sqrt2+u}{\sqrt2-u}\right]_{-1}^1
 =\frac{1}{\sqrt2}\log\frac{\sqrt2+1}{\sqrt2-1}
 =\sqrt2\log(\sqrt2+1).
\]
ここでは $-1\leq u\leq1$ で対数の分子・分母がともに正である。最後の等式では $(\sqrt2+1)(\sqrt2-1)=1$ を使った。したがって
\[
I=\frac{\pi}{\sqrt2}\log(\sqrt2+1).
\]

もし最初に部分積分を使うなら，$\frac{\sin x}{1+\sin^2x}$ の原始関数が必要になり，さらにその対数を積分する項が現れる。反転によって $x$ を除いたため，必要な計算が一度の置換と部分分数で済んだ。同じ式でも，不定積分では「微分の逆を作れるか」を見て，定積分ではそれに加えて「区間を保つ置換で二つの式を補い合わせられるか」を見るのである。

\section*{章末問題}
\begin{enumerate}[label=\textbf{\thechapter.\arabic*},leftmargin=*]
\item\label{問:奇偶の混合}
次の定積分を求めよ。
\[
\int_{-2}^2\left(\frac{x^3}{1+x^2}+2x^2-\cos x\right)\dd x.
\]
\item\label{問:一次重み反転}
反転を利用して，$\int_0^\pi x\sin^2x\dd x$ を求めよ。
\item\label{問:三角相補}
$\int_0^{\frac{\pi}{2}}\frac{\cos^6x}{\sin^6x+\cos^6x}\dd x$ を求めよ。分母が零にならないことも確かめよ。
\item\label{問:周期絶対値}
$\int_{\frac{\pi}{3}}^{\frac{10\pi}{3}}|\cos x|\dd x$ と
$\int_{\frac{\pi}{3}}^{\frac{10\pi}{3}}\cos x\dd x$ をそれぞれ求め，値が異なる理由を述べよ。
\item\label{問:逆数四乗}
$\int_{\frac14}^{4}\frac{1}{x(1+x^4)}\dd x$ を，逆数の置換を用いて求めよ。
\item\label{問:指数相補}
$\int_{-1}^1\frac{3x^2+1}{1+\e^{x^5}}\dd x$ を求めよ。
\item\label{問:中心を移す相補}
$\int_2^6\frac{(x-4)^2+2}{1+\e^{x-4}}\dd x$ を求めよ。
\item\label{問:二次重みの限界}
$I=\int_0^\pi x^2\sin x\dd x$ とする。反転によって
$I=\int_0^\pi(\pi-x)^2\sin x\dd x$ と書けるが，この二式を加えるだけで $I$ は決まらないことを説明せよ。その後，部分積分を用いて $I$ を求めよ。
\end{enumerate}

\smallskip\noindent 解答は\pageref{解答:第7章}ページ。

\chapter{面積・体積・長さと運動}\label{章:積分の応用}

定積分は，細かく分けた量を足し合わせる極限として定義した。この考えを使うと，曲線に囲まれた面積だけでなく，立体の体積や曲線の長さも求められる。ただし，どの問題にも被積分関数が最初から与えられるわけではない。細い部分を一つ取り出し，その面積，体積，長さをどの式で近似するかを決める必要がある。

その際に確かめたいことが三つある。一つは，分割した部分が対象全体を覆っているか。もう一つは，同じ部分を重複して数えていないか。そして，面積や道のりのような非負の量に，符号付きの差をそのまま使っていないかである。この章では，具体的な図形と運動を通して，それぞれの確認がどの式に現れるかを見ていく。

\section{曲線で囲まれた面積}\label{節:曲線間の面積}

$a<b$ とし，$f,g$ は $[a,b]$ で連続で，$f(x)\geq g(x)$ を満たすとする。二曲線 $y=f(x)$，$y=g(x)$ の間を，$x$ 軸に垂直な細い帯に分ける。幅が $\Delta x$ の帯の高さは，その付近での $f(x)-g(x)$ である。したがって面積は
\begin{equation}\label{式:上下の差の面積}
S=\int_a^b\{f(x)-g(x)\}\dd x
\end{equation}
になる。上下の曲線がともに $x$ 軸の下にあっても，高さは上の値から下の値を引いた差であり，同じ式が使える。

二曲線が途中で交差するなら，上下関係が変わる点で区間を分ける。一つの式で表せば
\[
S=\int_a^b|f(x)-g(x)|\dd x
\]
であるが，実際に計算するには差の符号を調べて絶対値を外す。交点の $x$ 座標は区切りの候補であり，交点を求めただけで上下が決まるわけではない。差を因数分解したり，各区間で一つの値を代入したりして確かめる。

\begin{example}{途中で上下が入れ替わる面積}
二曲線 $y=x^3-3x$ と $y=x$ に囲まれた部分の面積を求めよ。
\end{example}

交点は $x^3-3x=x$，すなわち $x(x-2)(x+2)=0$ より，$x=-2,0,2$ である。差 $x^3-4x$ は $-2<x<0$ で正，$0<x<2$ で負なので，面積は
\[
S=\int_{-2}^0(x^3-4x)\dd x
 +\int_0^2(4x-x^3)\dd x.
\]
図\ref{図:面積の上下交換}では，左右の部分が原点について対称である。積分でも第一項に $x=-t$ を使えば第二項と同じになるから，
\[
S=2\int_0^2(4x-x^3)\dd x
 =2\left[2x^2-\frac{x^4}{4}\right]_0^2=8.
\]
これを $\int_{-2}^2(x^3-4x)\dd x$ とすれば奇関数の積分なので零になるが，その値は左右の面積を一方は正，一方は負として足したものである。

\begin{figure}[htbp]
\centering
\begin{tikzpicture}[x=1.6cm,y=0.65cm,>=Stealth]
\fill[blue!12]plot[domain=-2:0,samples=60](\x,{\x*\x*\x-3*\x})--plot[domain=0:-2,samples=30](\x,\x)--cycle;
\fill[orange!17]plot[domain=0:2,samples=30](\x,\x)--plot[domain=2:0,samples=60](\x,{\x*\x*\x-3*\x})--cycle;
\draw[->](-2.3,0)--(2.4,0)node[right]{$x$};
\draw[->](0,-2.8)--(0,3.4)node[above]{$y$};
\draw[blue!70!black,thick]plot[domain=-2.13:2.13,samples=120](\x,{\x*\x*\x-3*\x});
\draw[orange!80!black,thick](-2.2,-2.2)--(2.2,2.2);
\foreach \x in {-2,2}{\draw[dashed](\x,0)--(\x,\x);\node[below]at(\x,0){$\x$};}
\node[below right]at(0,0){$O$};
\node[blue!70!black]at(-1.25,2.8){$y=x^3-3x$};
\node[orange!80!black]at(1.6,3.0){$y=x$};
\end{tikzpicture}
\caption{同じ二曲線でも，交点の左右で上側の曲線が入れ替わる。}\label{図:面積の上下交換}
\end{figure}

\subsection{横に切るほうが簡単な場合}

細い帯は縦に取る必要はない。$c\leq y\leq d$ において，右側の境界が $x=q(y)$，左側が $x=p(y)$ と表せるなら，帯の横の長さは $q(y)-p(y)$ なので
\begin{equation}\label{式:左右の差の面積}
S=\int_c^d\{q(y)-p(y)\}\dd y
\end{equation}
である。式\eqref{式:上下の差の面積}では上下の差，この式では左右の差を使う。変数の文字を交換するだけでなく，実際にどの向きに帯を取るかが変わっている。

たとえば $x=y^2$ と $x=2-y$ に囲まれる部分では，交点の $y$ 座標は $-2,1$ である。その間では直線が右側なので，
\[
S=\int_{-2}^1(2-y-y^2)\dd y
 =\left[2y-\frac{y^2}{2}-\frac{y^3}{3}\right]_{-2}^1
 =\frac92.
\]
縦に切ると放物線の上側と下側を $y=\pm\sqrt x$ に分け，さらに直線との関係も場合分けすることになる。横に切れば，一つの区間で一つの差になる。境界をどちらの変数で表すと少ない場合分けで済むかが，積分変数を選ぶ基準になる。

\section{断面積を積み重ねて体積を求める}\label{節:断面積と体積}

立体を $x$ 軸に垂直な平面で切り，位置 $x$ における断面積を $A(x)$ とする。幅 $\Delta x$ の薄い部分を，底面積 $A(x)$，高さ $\Delta x$ の柱体で近似すれば，その体積は $A(x)\Delta x$ である。断面積が連続に変化する通常の立体では，分割を細かくした極限から
\begin{equation}\label{式:断面積の体積}
V=\int_a^b A(x)\dd x
\end{equation}
を得る。断面が途中で形を変える場合も，変わる位置で有限個の区間に分け，各区間で同じ考えを使える。

この式の $A(x)$ は長さではなく面積である。長さの単位を $\mathrm{cm}$ とすれば，$A(x)$ は $\mathrm{cm}^2$，幅 $\dd x$ を掛けた体積は $\mathrm{cm}^3$ になる。図形の式が合っているか迷ったとき，単位も確認の助けになる。

\begin{example}{正方形の断面を持つ立体}
平面上で $0\leq x\leq1$，$x^2\leq y\leq x$ となる部分を底面とする。各 $x$ において，底面内の線分を一辺とする正方形が，底面に垂直な断面になるような立体の体積を求めよ。
\end{example}

底面内の線分の長さは $x-x^2$ である。この長さをそのまま積分すると底面積になるが，今求めたいのは体積なので，正方形の断面積
\[
A(x)=(x-x^2)^2=x^2-2x^3+x^4
\]
を積分する。
\[
V=\int_0^1(x^2-2x^3+x^4)\dd x
 =\left[\frac{x^3}{3}-\frac{x^4}{2}+\frac{x^5}{5}\right]_0^1
 =\frac1{30}.
\]
同じ底面でも，断面が正三角形なら面積は $\frac{\sqrt3}{4}(x-x^2)^2$ であり，体積も上の値の $\frac{\sqrt3}{4}$ 倍になる。底面の情報だけでは体積は決まらず，断面の形が必要である。

\section{回転体では軸からの距離を測る}\label{節:回転体の体積}

$0\leq g(x)\leq f(x)$ の領域 $g(x)\leq y\leq f(x)$ を $x$ 軸のまわりに回すと，位置 $x$ の断面は外半径 $f(x)$，内半径 $g(x)$ の円環になる。よって
\begin{equation}\label{式:円環断面の体積}
V=\pi\int_a^b\{f(x)^2-g(x)^2\}\dd x.
\end{equation}
下側が軸そのもの，すなわち $g(x)=0$ なら断面は円板で，$V=\pi\int_a^b f(x)^2\dd x$ となる。

回転軸が $y=c$ なら，半径は $y$ 座標そのものではなく $|y-c|$ である。たとえば $c\leq g(x)\leq f(x)$ なら，外半径は $f(x)-c$，内半径は $g(x)-c$ になる。軸が領域より上にあるときは，下側の境界のほうが軸から遠い。上下の順序をそのまま外半径・内半径の順序にしてはいけない。

\subsection{軸を横切る断面と，回転後の重なり}

縦の線分 $g(x)\leq y\leq f(x)$ が回転軸 $y=c$ を横切ると，回転後の断面には中心まで点が入る。したがって内半径は零であり，外半径は両端の遠いほうである。まとめると，
\[
R(x)=\max\{|f(x)-c|,|g(x)-c|\},
\]
\[
r(x)=
\begin{cases}
0,&g(x)\leq c\leq f(x),\\
\min\{|f(x)-c|,|g(x)-c|\},&\text{それ以外}
\end{cases}
\]
として，断面積は $\pi\{R(x)^2-r(x)^2\}$ になる。

軸より上と下を別々に回すと，同じ円板の部分を二度作ることがある。その場合，二つの体積を加えると重複して数える。先に回転後の断面を考え，できた円板や円環の和集合の面積を求めることが必要である。

\begin{example}{軸をまたぐ領域を回す}
$0\leq x\leq2$，$x-1\leq y\leq1$ で表される領域を $x$ 軸のまわりに一回転してできる立体の体積を求めよ。
\end{example}

$0\leq x\leq1$ では線分が軸をまたぐ。上端の距離は $1$，下端の距離は $1-x\leq1$ なので，断面は半径 $1$ の円板になる。$1\leq x\leq2$ では線分全体が軸の上にあり，外半径 $1$，内半径 $x-1$ の円環になる。図\ref{図:回転断面の区別}は，この違いを断面で示している。
\[
\begin{aligned}
V&=\pi\int_0^1 1\dd x
  +\pi\int_1^2\{1-(x-1)^2\}\dd x\\
 &=\pi+\pi\left[x-\frac{(x-1)^3}{3}\right]_1^2
 =\frac{5\pi}{3}.
\end{aligned}
\]
最初から $\pi\int_0^2\{1-(x-1)^2\}\dd x$ とすると，$x<1$ でも中心の円板をくり抜いてしまう。しかし元の線分は軸を含んでおり，回転後にその穴は存在しない。

\begin{figure}[htbp]
\centering
\begin{tikzpicture}[scale=1.0,>=Stealth]
\begin{scope}[xshift=-3.6cm]
\fill[blue!12](0,-1)--(0,1)--(2,1)--cycle;
\draw[->](-0.3,0)--(2.5,0)node[right]{$x$};
\draw[->](0,-1.3)--(0,1.5)node[above]{$y$};
\draw[thick](0,1)--(2,1)--(0,-1);
\draw[blue!70!black,very thick](0.5,-0.5)--(0.5,1);
\draw[orange!85!black,very thick](1.5,0.5)--(1.5,1);
\node[below]at(1,0){$1$};\node[below]at(2,0){$2$};
\node[below]at(1,-1.7){回す前の領域};
\end{scope}
\begin{scope}[xshift=0.1cm]
\fill[blue!12](0,0)circle(1);
\draw[blue!70!black,thick](0,0)circle(1);
\draw[->](0,0)--(0.707,0.707)node[midway,left]{$1$};
\fill(0,0)circle(0.025);
\node[below]at(0,-1.45){$x=\frac12$：円板};
\end{scope}
\begin{scope}[xshift=3.2cm]
\fill[orange!17,even odd rule](0,0)circle(1)(0,0)circle(0.5);
\draw[orange!85!black,thick](0,0)circle(1);
\draw[orange!85!black,thick](0,0)circle(0.5);
\draw[->](0,0)--(0.707,0.707)node[midway,right]{$1$};
\draw[->](0,0)--(-0.5,0)node[midway,above]{$\frac12$};
\fill(0,0)circle(0.025);
\node[below]at(0,-1.45){$x=\frac32$：円環};
\end{scope}
\end{tikzpicture}
\caption{軸を含む線分は円板を作り，軸から離れた線分は円環を作る。}\label{図:回転断面の区別}
\end{figure}

一つの $x$ に対して元の領域が複数の線分に分かれる場合には，各線分が作る半径の範囲を調べる。半径の範囲が重なるなら，その部分を一度だけ数えるように合わせる。この確認は，元の平面図形では重なっていない二つの部分を回す場合にも必要になる。

\subsection{円筒の薄い層として分ける}

同じ立体でも，回転軸に垂直に切る代わりに，軸に平行な薄い帯を回す方法がある。たとえば $0\leq a\leq x\leq b$，$g(x)\leq y\leq f(x)$ を $y$ 軸のまわりに回す。幅 $\Delta x$ の帯は，半径がほぼ $x$，高さがほぼ $f(x)-g(x)$ の薄い円筒の層になる。

半径 $x$ と $x+\Delta x$ の間の円環の面積は
\[
\pi\{(x+\Delta x)^2-x^2\}
 =2\pi x\Delta x+\pi(\Delta x)^2
\]
である。分割を細かくすると二次の項による誤差の和は零に近づくので，体積は
\begin{equation}\label{式:円筒殻の体積}
V=2\pi\int_a^b x\{f(x)-g(x)\}\dd x
\end{equation}
となる。これを円筒殻による方法という。$2\pi x$ は円周の長さ，$f-g$ は層の高さである。

\begin{example}{同じ体積を二通りに表す}
$0\leq x\leq1$，$x^2\leq y\leq x$ の領域を $y$ 軸のまわりに回転させた立体の体積を求めよ。
\end{example}

縦の帯を回すと，半径 $x$，高さ $x-x^2$ の円筒殻になるから，
\[
V=2\pi\int_0^1x(x-x^2)\dd x
 =2\pi\left[\frac{x^3}{3}-\frac{x^4}{4}\right]_0^1
 =\frac{\pi}{6}.
\]
横に切る方法でも確かめられる。$0\leq y\leq1$ で，領域の左右の端は $x=y$ と $x=\sqrt y$ である。したがって円環の外半径は $\sqrt y$，内半径は $y$ となり，
\[
V=\pi\int_0^1\{(\sqrt y)^2-y^2\}\dd y
 =\pi\left[\frac{y^2}{2}-\frac{y^3}{3}\right]_0^1
 =\frac{\pi}{6}.
\]
断面を円環として積み重ねるか，円筒の層として重ねるかによって，被積分関数は異なる。しかし，それぞれが立体を重複なく分けているので同じ体積になる。

式\eqref{式:円筒殻の体積}で $x\geq0$ としたのは，異なる $x$ が異なる半径に対応するためである。領域が $y$ 軸をまたぐと，$x$ と $-x$ が同じ半径になる。たとえば $-1\leq x\leq1$，$0\leq y\leq1$ の長方形を回すと半径 $1$，高さ $1$ の円柱で，体積は $\pi$ である。ところが $2\pi\int_{-1}^1|x|\dd x=2\pi$ は同じ層を二度数えている。絶対値で半径を正にしただけでは，重複は解消しない。

\section{曲線の長さと媒介変数}\label{節:曲線の弧長}

曲線 $y=f(x)$ の長さを求めるとき，細い部分を長方形で近似するのは適切ではない。曲線上の近い二点を線分で結び，折れ線の長さで近似する。$f$ が $[a,b]$ で連続微分可能，すなわち導関数も連続であるとする。小区間の幅を $\Delta x$，両端の関数値の差を $\Delta y$ とすると，線分の長さは
\[
\sqrt{(\Delta x)^2+(\Delta y)^2}
 =\sqrt{1+\left(\frac{\Delta y}{\Delta x}\right)^2}\,\Delta x.
\]
平均値の定理により，小区間内のある $\xi$ で $\frac{\Delta y}{\Delta x}=f'(\xi)$ である。したがって折れ線の長さの和は，連続関数 $\sqrt{1+f'(x)^2}$ のリーマン和になり，分割を細かくすると
\begin{equation}\label{式:グラフの弧長}
L=\int_a^b\sqrt{1+f'(x)^2}\dd x
\end{equation}
に近づく。これが曲線の弧長である。有限個の点で折れ曲がっていても，各部分が連続微分可能なら，部分ごとの長さを加えればよい。

$\int_a^b|f'(x)|\dd x$ は縦方向の変化の総量であり，弧長ではない。水平な線分では $f'=0$ だが長さは $b-a$ である。式\eqref{式:グラフの弧長}の平方根の中の $1$ は，横方向の移動も数えている。

\subsection{媒介変数で表された曲線}

円や，縦向きの接線を持つ曲線は，一つの関数 $y=f(x)$ で扱いにくい場合がある。そこで
\[
x=x(t),\qquad y=y(t),\qquad \alpha\leq t\leq\beta
\]
と表す。$t$ を少し増やしたとき，横の変化はほぼ $x'(t)\Delta t$，縦の変化はほぼ $y'(t)\Delta t$ だから，線分の長さはほぼ
\[
\sqrt{x'(t)^2+y'(t)^2}\,\Delta t
\]
である。$x,y$ が連続微分可能なら，導関数の連続性によりこの近似の誤差は区間全体で制御でき，折れ線の極限として
\begin{equation}\label{式:媒介曲線の弧長}
L=\int_\alpha^\beta\sqrt{x'(t)^2+y'(t)^2}\dd t
\end{equation}
を得る。ここでは近似誤差の評価の詳細は省き，この弧長公式を用いる。途中で速度に当たる平方根が零になる点があっても，この式自体は使える。

ただし，この積分は $t$ に沿って進んだ長さを数える。同じ線分や曲線を二度たどれば二度数える。たとえば $x=\cos t$，$y=\sin t$ を $0\leq t\leq4\pi$ で積分すると $4\pi$ だが，点の集合としての単位円の周長は $2\pi$ である。図形そのものの長さを求めるなら，一度だけたどる区間を選ぶ。

媒介変数は面積にも使える。$x(t)$ が単調増加し，$y(t)\geq0$ で，曲線が $x$ の関数として表される場合は，置換 $\dd x=x'(t)\dd t$ により
\begin{equation}\label{式:媒介曲線の面積}
S=\int_\alpha^\beta y(t)x'(t)\dd t
\end{equation}
となる。$x(t)$ が途中で折り返すなら，区間を分けて曲線が囲む部分を確認する。この式に絶対値を付けるだけで常に図形の面積になるわけではない。帯が重なれば，同じ面積を二度数えるからである。

\begin{example}{一つの媒介表示から面積と長さを求める}
$a>0$ とする。曲線
\[
x=a(t-\sin t),\qquad y=a(1-\cos t),\qquad 0\leq t\leq2\pi
\]
の長さと，この曲線と $x$ 軸に囲まれた部分の面積を求めよ。
\end{example}

$x'(t)=a(1-\cos t)$，$y'(t)=a\sin t$ である。まず弧長の平方根を計算すると
\[
\begin{aligned}
x'(t)^2+y'(t)^2
 &=a^2\{(1-\cos t)^2+\sin^2t\}\\
 &=2a^2(1-\cos t)=4a^2\sin^2\frac t2.
\end{aligned}
\]
平方根は $2a|\sin\frac t2|$ である。この区間では $0\leq\frac t2\leq\pi$ なので正弦は非負であり，
\[
L=2a\int_0^{2\pi}\sin\frac t2\dd t
 =\left[-4a\cos\frac t2\right]_0^{2\pi}=8a.
\]
平方根を取る際に絶対値が必要かどうかを，変数の範囲で確認している。

一方，$x'$ は端点以外で正なので，$x$ は区間全体で単調増加する。また $y\geq0$ で，両端では $y=0$ である。よって面積は
\[
\begin{aligned}
S&=a^2\int_0^{2\pi}(1-\cos t)^2\dd t\\
 &=a^2\int_0^{2\pi}\left(\frac32-2\cos t+\frac12\cos2t\right)\dd t
 =3\pi a^2.
\end{aligned}
\]
弧長では二方向の変化から距離を作り，面積では高さに横の幅を掛けた。同じ媒介表示でも，求める量に応じて積分の中身が異なる。

\section{速度から位置と道のりへ}\label{節:運動と積分}

一直線上の運動では，位置を $s(t)$，速度を $v(t)=s'(t)$，加速度を $a(t)=v'(t)$ とする。速度が連続なら，微積分の基本定理により
\[
s(t_1)-s(t_0)=\int_{t_0}^{t_1}v(t)\dd t.
\]
左辺は位置の差，すなわち変位である。出発点に戻れば零になるが，途中で移動した距離が零という意味ではない。

速度が正なら正の向きに進み，負なら負の向きに進む。移動した距離を向きによらず加えるには速さ $|v(t)|$ を積分して，
\begin{equation}\label{式:道のり}
\text{道のり}=\int_{t_0}^{t_1}|v(t)|\dd t
\end{equation}
とする。速さは非負，速度には向きに対応する符号がある。この違いは，面積と符号付きの定積分の違いと同じである。

速度が零になる時刻は，進む向きが変わる候補である。ただし零になるだけでは反転したとは限らない。たとえば $v(t)=(t-1)^2$ は $t=1$ で零になるが，前後で非負である。道のりを計算するときは，零点で区間を分けたうえで，各区間の符号を確かめる。

\begin{example}{変位と道のりを分けて求める}
時刻 $0\leq t\leq4$ において，速度が $v(t)=(t-1)(t-3)$ で与えられる。初期位置は $s(0)=0$ とする。$t=4$ の位置，道のり，平均速度，平均の速さを求めよ。
\end{example}

速度を積分して初期条件を使うと，
\[
s(t)=\frac{t^3}{3}-2t^2+3t.
\]
したがって $s(4)=\frac43$ である。道のりでは速度の符号を調べる。$v$ は $0\leq t<1$ で正，$1<t<3$ で負，$3<t\leq4$ で正である。折り返す位置は
\[
s(1)=\frac43,\qquad s(3)=0
\]
なので，道のりは
\[
\left|\frac43-0\right|+\left|0-\frac43\right|
 +\left|\frac43-0\right|=4.
\]
これは $\int_0^1v\dd t-\int_1^3v\dd t+\int_3^4v\dd t$ と計算しても同じである。全時間は $4$ だから，平均速度は $\cfrac{\cfrac43}{4}=\frac13$，平均の速さは $\frac44=1$ となる。

加速度は $a(t)=2t-4$ である。加速度が正でも，いつも速さが増えるわけではない。$v\ne0$ の時刻では
\[
\frac{\dd}{\dd t}|v(t)|=
\begin{cases}
a(t),&v(t)>0,\\
-a(t),&v(t)<0
\end{cases}
\]
だから，速度と加速度が同符号なら速さが増え，異符号なら減る。この例では $1<t<2$ と $3<t<4$ で速さが増える。負の向きに進みながら速くなることもあるのである。

平面上の運動 $\bigl(x(t),y(t)\bigr)$ では，速度は二成分 $\bigl(x'(t),y'(t)\bigr)$ を持ち，速さは $\sqrt{x'(t)^2+y'(t)^2}$ となる。その積分が道のりになることは，式\eqref{式:媒介曲線の弧長}そのものである。運動では同じ場所を再び通った分も道のりに含むので，曲線を重複してたどることは誤りではなく，数えるべき内容になる。

\section{連続的に変わる量の平均}\label{節:関数の平均}

\sectionref{節:定積分性質}で導入した関数の平均値を，ここでは温度のような連続的に変わる量へ使おう。まず，平均を積分で表す理由を振り返る。数 $u_1,\ldots,u_n$ の平均は和を個数で割ったものである。関数 $f$ の区間 $[a,b]$ における平均も，等間隔に値を取り，その平均の極限として考えられる。$\Delta x=\frac{b-a}{n}$ とすれば，
\[
\frac1n\sum_{k=1}^n f(a+k\Delta x)
 =\frac{1}{b-a}\sum_{k=1}^n f(a+k\Delta x)\Delta x.
\]
右辺の和はリーマン和なので，$f$ が連続なら極限は
\begin{equation}\label{式:関数の平均値}
\overline f=\frac{1}{b-a}\int_a^b f(x)\dd x
\end{equation}
となる。$f\geq0$ の場合，底辺 $b-a$，高さ $\overline f$ の長方形は，曲線の下と同じ面積になる。$f$ に負の値があっても，平均の定義は同じ式である。

\begin{example}{変化する温度の平均}
時刻 $0\leq t\leq12$ の温度を $T(t)=20+6\sin\frac{\pi t}{12}$ とする。この時間帯の平均温度を求めよ。
\end{example}

始めと終わりの温度はともに $20$ だが，途中の温度は $20$ より高い。両端の平均だけでは時間全体の平均にならない。定義に従って
\[
\begin{aligned}
\overline T
 &=\frac1{12}\int_0^{12}\left(20+6\sin\frac{\pi t}{12}\right)\dd t\\
 &=20+\frac12\left[-\frac{12}{\pi}\cos\frac{\pi t}{12}\right]_0^{12}
 =20+\frac{12}{\pi}.
\end{aligned}
\]
全時間を均等に数えて平均を取っていることが，$\frac1{12}$ に表れている。

連続な $f$ の最小値を $m$，最大値を $M$ とすると，$m\leq f(x)\leq M$ を積分して
\[
m\leq\overline f\leq M
\]
を得る。中間値の定理から，ある $c\in[a,b]$ で $f(c)=\overline f$ となる。これを積分の平均値の定理という。平均が実際のどこかの値と一致することは分かるが，一般にその点が区間の中央とは限らず，一つとは限らない。

均等に数えるのでない平均には，重みが必要になる。非負の連続関数 $w$ に対し $\int_a^b w(x)\dd x>0$ なら，重み付き平均は
\[
\frac{\int_a^b f(x)w(x)\dd x}{\int_a^b w(x)\dd x}
\]
である。分母は重みの総量であり，$w=1$ の場合が式\eqref{式:関数の平均値}になる。次に扱う重心は，位置の座標を面積や長さで重み付けして平均したものである。

\section{発展：面積の重心と線の重心}\label{節:面積重心と線重心}

密度が一様な薄い板を考える。板を小さな部分に分け，それぞれの面積を $\Delta A_i$，位置を $(x_i,y_i)$ と近似すると，重心の座標は面積を重みにした平均
\[
\overline x\approx\frac{\sum_i x_i\Delta A_i}{\sum_i\Delta A_i},\qquad
\overline y\approx\frac{\sum_i y_i\Delta A_i}{\sum_i\Delta A_i}
\]
で表される。密度が一様なら各部分の質量が面積に比例し，共通の密度が分子と分母で消えるからである。

領域が $a\leq x\leq b$，$g(x)\leq y\leq f(x)$ で表され，面積 $A>0$ とする。縦の細い長方形の面積はほぼ $(f-g)\Delta x$ で，その中心の座標は $\bigl(x,\frac{f+g}{2}\bigr)$ である。分割を細かくすると
\begin{equation}\label{式:面積の重心}
\begin{aligned}
A&=\int_a^b(f-g)\dd x,\\
\overline x&=\frac1A\int_a^b x(f-g)\dd x,\\
\overline y&=\frac1{2A}\int_a^b(f^2-g^2)\dd x
\end{aligned}
\end{equation}
を得る。最後の式は $\frac{f+g}{2}(f-g)=\frac12(f^2-g^2)$ から出ている。$\overline y$ を計算するときに上端の高さ $f$ だけを掛けてはいけない。帯全体の重心は上下の中央にある。

曲線に沿って曲げた一様な針金では，質量が比例するのは面積ではなく長さである。長さを $L>0$ とすると，重心は
\begin{equation}\label{式:線の重心}
\overline x=\frac1L\int_C x\dd s,\qquad
\overline y=\frac1L\int_C y\dd s
\end{equation}
となる。ここで $\dd s$ は曲線の小さな長さを表す記号で，媒介表示なら
\[
\int_C y\dd s
 =\int_\alpha^\beta y(t)\sqrt{x'(t)^2+y'(t)^2}\dd t
\]
という普通の定積分で計算する。一様な針金では，一度だけたどる区間を選ぶ。

半径 $a$ の上半円板と，その円弧だけの針金を比較する。どちらも左右対称なので $\overline x=0$ である。半円板では $A=\frac{\pi a^2}{2}$，$g=0$，$f(x)=\sqrt{a^2-x^2}$ だから，
\[
\int_A y\dd A
 =\frac12\int_{-a}^a(a^2-x^2)\dd x
 =\frac{2a^3}{3}.
\]
左辺は面積で重み付けした $y$ 座標の和を表す記号であり，この場合の計算方法が右辺である。したがって，半円板の重心の高さは
\begin{equation}\label{式:半円板の重心}
\overline y_{\mathrm{面}}=\frac{4a}{3\pi}.
\end{equation}
円弧のほうは $x=a\cos t$，$y=a\sin t$，$0\leq t\leq\pi$ と表すと，$\dd s=a\dd t$，$L=\pi a$ だから，
\begin{equation}\label{式:半円弧の重心}
\overline y_{\mathrm{線}}
 =\frac{1}{\pi a}\int_0^\pi(a\sin t)a\dd t
 =\frac{2a}{\pi}.
\end{equation}
円弧だけのほうが重心は高い。円板には直径の近くにも面積があるが，円弧だけの針金にはその部分がないためである。

\section{発展：パップス＝ギュルダンの二つの定理}\label{節:パップスの定理}

回転体の体積を円筒殻の式で求めると，距離に面積を掛けた積分が現れる。一方，重心も座標に面積を掛けた積分で求めた。この共通する積分を使うと，回転体の体積を面積と重心で表せる。

まず回転軸を $x$ 軸とし，領域がその上側にある場合を考える。高さ $y$ の横断面の長さを $h(y)$ とする。横の細い帯を回すと半径 $y$，軸に沿う長さ $h(y)$ の円筒殻になるので，
\[
V=2\pi\int y h(y)\dd y.
\]
他方，面積 $A$ と重心の高さ $\overline y$ は
\[
A=\int h(y)\dd y,\qquad
A\overline y=\int y h(y)\dd y
\]
を満たす。よって $V=2\pi\overline y A$ である。積分区間は領域に含まれる高さ全体であり，断面が複数の区間に分かれるなら，$h(y)$ はそれらの長さの和とする。

\begin{theorem}{パップス＝ギュルダンの体積に関する定理}
一様な平面図形の面積を $A$，その面積重心と，図形と同じ平面内にある回転軸との距離を $R$ とする。図形が軸の片側にあり，回転によって内部を重複して数えない場合，できる回転体の体積は
\[
V=2\pi R A
\]
である。すなわち，面積に，重心が一回転するときに進む円周の長さを掛ければよい。
\end{theorem}

ここでは境界が区分的に滑らかで，通常の断面積による積分ができる図形を考えている。軸が境界に触れても内部を横切らなければ，上の帯による証明を使える。軸が内部を横切る場合には，軸の両側が同じ場所を回転して重なり得る。重心が軸上にあるから体積が零になる，という使い方はできない。

\subsection{回転面の面積と線重心}

回転軸と同じ平面内にある曲線を考える。この曲線を回したときにできる面の面積では，面積重心ではなく線重心が現れる。その理由を，回転面の小さな部分から考える。軸の片側にあり，軸からの距離が $r_1,r_2$ の二点を結ぶ短い線分を回すと，一般には円すい台の側面になる。軸に平行な線分は円筒面，垂直な線分は円環または円板を作る。いずれの場合も，線分の長さを $\Delta s$ とすると面積は
\[
\pi(r_1+r_2)\Delta s
\]
である。円すい台の場合は，大きい円すいと切り取った円すいの側面積の差を，相似比で整理するとこの式を得る。円筒面と円環の場合も，それぞれの面積公式から確かめられる。線分を十分短くすれば，$r_1,r_2$ はほぼ同じ $r$ なので，これは $2\pi r\Delta s$ になる。

回転軸と同じ平面内の曲線が区分的に連続微分可能で，回転によって同じ面を重複して作らないとき，折れ線近似の極限として回転面の面積は
\begin{equation}\label{式:回転面の面積}
S=2\pi\int_C r\dd s
\end{equation}
になる。ここでは，この極限の正当化の詳細は省く。たとえば $y=f(x)\geq0$ を $x$ 軸のまわりに回すなら
\[
S=2\pi\int_a^b f(x)\sqrt{1+f'(x)^2}\dd x.
\]
ここで得られるのは，指定した曲線が作る面の面積である。立体の両端を塞ぐ円板などは，その曲線に含めていなければ別に数える。

\begin{theorem}{パップス＝ギュルダンの面積に関する定理}
一様な平面曲線の長さを $L$，その線重心と，曲線と同じ平面内にある回転軸との距離を $R$ とする。曲線が軸の片側にあり，一回転によって同じ面を重複して作らない場合，曲線が作る回転面の面積は
\[
S=2\pi R L
\]
である。
\end{theorem}

実際，軸を $x$ 軸，曲線を $y\geq0$ に取れば $r=y$ であり，線重心の定義から $\int_C y\dd s=L\overline y=LR$ となるので，式\eqref{式:回転面の面積}から従う。回転軸の位置や方向が違っていても，軸に合わせて座標を取り直せば同じ議論になる。

二つの定理で掛ける円周は，図形の端や幾何学的な中心の円周ではなく，それぞれの重心が進む円周である。また，体積の定理には面積 $A$ と面積重心，面の定理には長さ $L$ と線重心を対応させる。

\begin{example}{半円板を外側の軸のまわりに回す}
$a>0$，$b>0$ とする。上半円板 $x^2+y^2\leq a^2$，$y\geq0$ を，直線 $y=-b$ のまわりに一回転してできる立体の体積と全表面積を求めよ。
\end{example}

図\ref{図:半円の二つの重心}のように，回転軸は図形の外にあり，重心の高さに $b$ を加えたものが回転半径になる。半円板の面積重心を使うと，
\[
\begin{aligned}
V&=2\pi\left(b+\frac{4a}{3\pi}\right)\frac{\pi a^2}{2}\\
 &=\pi^2a^2b+\frac{4\pi a^3}{3}.
\end{aligned}
\]
断面で確かめるなら，外半径 $b+\sqrt{a^2-x^2}$，内半径 $b$ の円環である。したがって
\[
\begin{aligned}
V&=\pi\int_{-a}^a
\left\{\left(b+\sqrt{a^2-x^2}\right)^2-b^2\right\}\dd x\\
 &=\pi\int_{-a}^a(a^2-x^2)\dd x
  +2\pi b\int_{-a}^a\sqrt{a^2-x^2}\dd x,
\end{aligned}
\]
となり，第一項が $\frac{4\pi a^3}{3}$，第二項が $\pi^2a^2b$ で一致する。

\begin{figure}[htbp]
\centering
\begin{tikzpicture}[scale=1.25,>=Stealth]
\fill[blue!10](-2,0)arc[start angle=180,end angle=0,radius=2]--cycle;
\draw[blue!75!black,thick](2,0)arc[start angle=0,end angle=180,radius=2];
\draw[thick](-2,0)--(2,0);
\draw[->](-2.5,0)--(2.6,0)node[right]{$x$};
\draw[->](0,-1.5)--(0,2.45)node[above]{$y$};
\draw[orange!80!black,thick](-2.5,-1)--(2.5,-1)node[right]{$y=-b$};
\fill[blue!75!black](0,0.8488)circle(0.045);
\fill[red!75!black](0,1.2732)circle(0.045);
\draw[blue!75!black](0,0.8488)--(2.5,0.6)node[right]{面積重心 $\frac{4a}{3\pi}$};
\draw[red!75!black](0,1.2732)--(2.5,1.6)node[right]{線重心 $\frac{2a}{\pi}$};
\draw[<->](-2.25,-1)--(-2.25,0)node[midway,left]{$b$};
\node[below]at(-2,0){$-a$};\node[below]at(2,0){$a$};
\end{tikzpicture}
\caption{同じ半円でも，板と円弧では重心が違う。回転半径は軸から測る。}\label{図:半円の二つの重心}
\end{figure}

表面積では，境界を円弧と直径に分ける。円弧が作る面は，円弧の長さ $\pi a$ と線重心の高さ $\frac{2a}{\pi}$ を使って
\[
S_{\mathrm{円弧}}
 =2\pi\left(b+\frac{2a}{\pi}\right)\pi a
 =2\pi^2ab+4\pi a^2.
\]
これだけでは全表面積にならない。直径 $-a\leq x\leq a$，$y=0$ も，回転すると半径 $b$，軸方向の長さ $2a$ の円筒面を作る。その面積は
\[
S_{\mathrm{直径}}=2\pi b\cdot2a=4\pi ab.
\]
両者は端の円周で接するが，面積のある部分では重ならない。よって全表面積は
\[
S=2\pi^2ab+4\pi a^2+4\pi ab.
\]
$x=\pm a$ では外半径と内半径がともに $b$ となり，厚みが零なので，さらに端の円板を加える必要はない。どの境界がどの面を作るかを確かめた結果として，全表面積が決まる。

面積，体積，長さ，道のり，平均はいずれも積分で表せるが，同じ被積分関数を使うわけではない。細い帯なら高さと幅，薄い立体なら断面積と厚さ，短い曲線なら二方向の変化から作った距離が必要になる。公式を選ぶ前に，細かく分けた一つ分が何を表すかを言えるようにすると，符号や重複の誤りに気づきやすくなる。

\section*{章末問題}
\begin{enumerate}[label=\textbf{\thechapter.\arabic*},leftmargin=*]
\item\label{問:面積と円環}
$y=x^2$ と $y=3x$ に囲まれた部分の面積を求めよ。また，この部分を $x$ 軸のまわりに回転させた立体の体積を求めよ。
\item\label{問:横帯の面積}
$x=y^2$ と $x=3-2y$ に囲まれた部分の面積を，$y$ を積分変数として求めよ。
\item\label{問:外半径の交代}
$0\leq x\leq2$，$-1\leq y\leq x$ の領域を $x$ 軸のまわりに回した立体の体積を求めよ。外半径が変わる位置を明示せよ。
\item\label{問:円筒と円板}
$0\leq x\leq1$，$x^2\leq y\leq1$ の領域を $y$ 軸のまわりに回した立体の体積を，円筒殻と円板の二通りで求めよ。
\item\label{問:弧長の平方完成}
$y=\frac13(x^2+2)^{\frac32}$，$0\leq x\leq1$ の曲線の長さを求めよ。
\item\label{問:道のりと平均速度}
$0\leq t\leq3$ における速度が $v(t)=3(t-1)(t-2)$，初期位置が $s(0)=0$ の運動について，終点の位置，道のり，平均速度，平均の速さを求めよ。
\item\label{問:平均値を取る点}
$f(x)=1+x^2$ の $[0,2]$ における平均値を求め，$f(c)$ がその平均値と等しくなる $c\in[0,2]$ を求めよ。
\item\label{問:三角形と両パップス}
頂点 $(0,0),(2,0),(2,1)$ の三角形を考える。面積重心の $y$ 座標を積分で求め，パップス＝ギュルダンの定理により，三角形を $x$ 軸のまわりに回した立体の体積を求めよ。また，境界の各線分の線重心を使って，できる立体の全表面積を求めよ。回転軸にある線分の寄与も説明せよ。
\end{enumerate}

\smallskip\noindent 解答は\pageref{解答:第8章}ページ。

\chapter{積分と数列・不等式}\label{章:積分と数列}

積分を使うと，項数の増えていく和を一つの面積として見ることができる。また，積分の中身を大小比較すれば，値を正確に求められなくても積分そのものの範囲が分かる。この二つを組み合わせると，数列の極限や，複雑な積の極限まで調べられる。

この章では，まず和の中に区分求積の形を見つける。その後，積分を評価する不等式と，指数を一つずつ変えた積分を結び付ける漸化式を考える。最後に現れるベータ積分は，それまで別々に扱った多項式の積分と三角関数の積分を，一つの式の中に並べるためのものである。

\section{和の中から区分求積の形を取り出す}\label{節:和の区分求積}

連続関数$f$について，区間$[a,b]$を$n$等分し，右端を使って長方形を作ると，その面積の和は
\[
 \frac{b-a}{n}\sum_{k=1}^{n}
 f\left(a+\frac{k(b-a)}{n}\right)
\]
である。$n$を大きくすれば，この和は$\int_a^b f(x)\dd x$へ近づく。とくに$[0,1]$では
\begin{equation}\label{式:単位区間の区分求積}
 \lim_{n\to\infty}\frac1n\sum_{k=1}^{n}f\left(\frac{k}{n}\right)
 =\int_0^1 f(x)\dd x
\end{equation}
となる。

和の極限でこの公式を使うときには，$\frac{k}{n}$が横の位置，$\frac1n$が長方形の幅であることを確認する。分母に$n$が現れるというだけでは足りない。和全体を，幅と高さの積の和へ書き直す必要がある。

\begin{example}{分母をそろえて長方形の幅を見つける}
次の極限を求めよ。
\[
 \lim_{n\to\infty}\sum_{k=1}^{n}\frac{n}{n^2+k^2}.
\]
\end{example}

各項の分母から$n^2$をくくると
\[
 \frac{n}{n^2+k^2}
 =\frac1n\cdot\cfrac{1}{1+\left(\cfrac{k}{n}\right)^2}
\]
となる。したがって，これは$f(x)=\frac1{1+x^2}$の右端を使う区分求積である。この関数は$[0,1]$で連続なので
\[
 \lim_{n\to\infty}\sum_{k=1}^{n}\frac{n}{n^2+k^2}
 =\int_0^1\frac{1}{1+x^2}\dd x.
\]
$x=\tan t$とおけば，$0\leqq t\leqq\frac{\pi}{4}$で$\dd x=(1+\tan^2t)\dd t$だから，右辺は$\frac{\pi}{4}$である。和の一項は$0$へ近づくが，項数も増えるため，和全体は$0$にはならない。

式\eqref{式:単位区間の区分求積}で，右端の代わりに左端や中点を使っても，同じ積分へ収束する。ただし，和に含まれる点の範囲が変わるなら，積分区間も確かめなければならない。例えば
\[
 \frac1n\sum_{k=1}^{2n}f\left(\frac{k}{n}\right)
\]
では，幅は$\frac1n$のままで，点は$0$から$2$まで並ぶ。これは$[0,2]$の区分求積である。項数が$2n$だからといって，幅を勝手に$\frac1{2n}$へ変えてはいけない。

\begin{example}{和の範囲から積分区間を決める}
次の極限を求めよ。
\[
 \lim_{n\to\infty}\sum_{k=n+1}^{2n}\frac1k.
\]
\end{example}

一項を$\frac1k=\frac1n\cfrac1{\cfrac{k}{n}}$と書く。$k=n+1,\ldots,2n$に対応する点は
$1+\frac1n,\ldots,2$であり，$[1,2]$を$n$等分した小区間の右端である。よって
\[
 \lim_{n\to\infty}\sum_{k=n+1}^{2n}\frac1k
 =\int_1^2\frac1x\dd x=\log2.
\]
ここでは$\frac1x$が$x=0$で定義されないことは問題にならない。実際に使った積分区間は$[1,2]$であり，その上で連続だからである。

有限個の端の項を付けたり取り除いたりしても，その差が$0$へ近づくなら極限は変わらない。例えば，連続関数$f$について
\[
 \frac1n\sum_{k=0}^{n}f\left(\frac{k}{n}\right)
 -\frac1n\sum_{k=1}^{n}f\left(\frac{k}{n}\right)
 =\frac{f(0)}n\longrightarrow0
\]
である。しかし，取り除いた項数が$n$とともに増える場合や，端の関数値が大きくなる場合には，この確認を省けない。区分求積の公式を使う範囲を先に決め，残りの差を別に評価するのが確実である。

\section{単調な関数を使って和をはさむ}\label{節:和と積分の比較}

区分求積では，極限を取って和と積分を結び付けた。単調性が分かる場合には，有限の$n$でも大小関係が得られる。$f$が$[1,\infty)$で正かつ単調減少なら，$k\leqq x\leqq k+1$において
\[
 f(k+1)\leqq f(x)\leqq f(k)
\]
である。長さ$1$の区間で積分すると
\[
 f(k+1)\leqq\int_k^{k+1}f(x)\dd x\leqq f(k).
\]
この不等式を足すと
\begin{equation}\label{式:減少関数の和比較}
 \int_1^{n+1}f(x)\dd x
 \leqq\sum_{k=1}^{n}f(k)
 \leqq f(1)+\int_1^{n}f(x)\dd x
\end{equation}
を得る。右側では，$k=2,\ldots,n$について$f(k)\leqq\int_{k-1}^{k}f(x)\dd x$を足し，最初の$f(1)$だけを別に残している。

図\ref{図:和と積分の比較}では，各長方形の高さを小区間の左端の値にしている。減少関数では長方形が曲線の上に出るので，長方形の面積の和が積分の上からの評価になる。高さを右端の値に変えれば，大小関係が逆になる。

\begin{figure}[htbp]
 \centering
 \begin{tikzpicture}[x=1.15cm,y=2.9cm,>=Stealth]
  \foreach \k in {1,2,3,4} {
   \pgfmathsetmacro{\h}{1/\k}
   \fill[blue!12] (\k,0) rectangle ({\k+1},\h);
   \draw[blue!65!black] (\k,0) rectangle ({\k+1},\h);
  }
  \draw[->] (0.55,0)--(5.45,0) node[right]{$x$};
  \draw[->] (0.7,0)--(0.7,1.25) node[above]{$y$};
  \draw[thick,domain=1:5,samples=90] plot (\x,{1/\x});
  \node[anchor=west] at (2,0.86) {$y=\frac1x$};
  \foreach \k in {1,2,3,4,5}
   \node[below] at (\k,0) {$\k$};
 \end{tikzpicture}
 \caption{減少関数の左端の値で作った長方形は，曲線の上に出る。}
 \label{図:和と積分の比較}
\end{figure}

\begin{example}{調和数の増え方を調べる}
$H_n=\sum_{k=1}^{n}\frac1k$とおく。$H_n\to\infty$であること，および
\[
 \lim_{n\to\infty}\frac{H_n}{\log n}=1
\]
を示せ。後の比は$n\geqq2$について考える。
\end{example}

式\eqref{式:減少関数の和比較}を$f(x)=\frac1x$に用いると
\[
 \log(n+1)\leqq H_n\leqq1+\log n
\]
となる。下側の$\log(n+1)$が限りなく大きくなるため，$H_n$も限りなく大きくなる。次に，正の$\log n$で割れば
\[
 \frac{\log(n+1)}{\log n}
 \leqq\frac{H_n}{\log n}
 \leqq1+\frac1{\log n}.
\]
左辺は
\[
 \frac{\log(n+1)}{\log n}
 =1+\cfrac{\log\left(1+\cfrac1n\right)}{\log n}
 \longrightarrow1
\]
であり，右辺も$1$へ近づく。はさみうちの原理から求める極限が得られる。

ここから$H_n-\log n\to0$と結論してはいけない。比が$1$へ近づくということは，差が元の大きさに比べて小さいことを表すだけである。実際，上の評価から分かるのは，差が$0$以上$1$以下の範囲にあることまでである。

比較の向きを図から毎回考える必要はない。増加関数なら$f(k)\leqq f(x)\leqq f(k+1)$，減少関数ならその逆，と小区間ごとの不等式を書けばよい。添字のずれを避けるためにも，最初から大きな和の不等式を暗記するより，小区間から組み立てる方が扱いやすい。

\section{積分の大小とコーシー・シュワルツの不等式}\label{節:積分の不等式}

$f(x)\leqq g(x)$が区間全体で成り立てば，$\int_a^b f(x)\dd x\leqq\int_a^b g(x)\dd x$である。この基本的な性質を，積や平方の形に使うと，関数同士の関係を表す不等式が得られる。

\begin{theorem}{積分に対するコーシー・シュワルツの不等式}
$a<b$とし，$f,g$は$[a,b]$で実数値を取る連続関数とする。このとき
\begin{equation}\label{式:積分のコーシーシュワルツ}
 \left(\int_a^b f(x)g(x)\dd x\right)^2
 \leqq
 \left(\int_a^b f(x)^2\dd x\right)
 \left(\int_a^b g(x)^2\dd x\right)
\end{equation}
が成り立つ。$g$が恒等的に$0$でなければ，等号が成り立つための必要十分条件は，ある定数$c$について区間全体で$f(x)=cg(x)$となることである。$g$が恒等的に$0$なら，どの$f$についても等号が成り立つ。
\end{theorem}

平方は負にならないので，任意の実数$t$について
\[
 0\leqq\int_a^b\{f(x)-tg(x)\}^2\dd x
\]
である。式を短くするため
\[
 A=\int_a^b f(x)^2\dd x,\qquad
 B=\int_a^b f(x)g(x)\dd x,\qquad
 C=\int_a^b g(x)^2\dd x
\]
とおくと，右辺は$A-2Bt+Ct^2$になる。

$C>0$のとき，$t=\frac BC$を代入すれば
\[
 0\leqq A-\frac{B^2}{C}
\]
だから，$B^2\leqq AC$が得られる。等号が成り立つ場合には
\[
 \int_a^b\left\{f(x)-\frac BCg(x)\right\}^2\dd x=0.
\]
被積分関数は連続かつ非負である。もし一点で正なら，連続性によってその近くの小区間でも正となり，積分が正になってしまう。したがって区間全体で$f(x)=\frac BCg(x)$である。逆に$f=cg$なら直接代入して等号を確認できる。

$C=0$の場合には，同じ連続性の議論によって$g$が区間全体で$0$となる。このとき$B=0$でもあるので，不等式の両辺はともに$0$である。$C$で割る計算の前にこの場合を分けたのは，$0$による除算を避けるためである。

この不等式を使うときには，評価したい積分の中身を$f(x)g(x)$という積に見る。片方を$1$としてよいことも多い。例えば$g(x)=1$とすれば
\begin{equation}\label{式:平均と二乗平均}
 \left(\int_a^b f(x)\dd x\right)^2
 \leqq(b-a)\int_a^b f(x)^2\dd x.
\end{equation}
区間の長さで割ると，関数の平均値の平方は，平方の平均値を超えないことが分かる。等号は$f$が定数関数のときに限る。

\begin{example}{正の関数とその逆数を組み合わせる}
$h$は$[a,b]$で連続で，すべての$x$について$h(x)>0$を満たすとする。
\[
 \left(\int_a^b h(x)\dd x\right)
 \left(\int_a^b\frac1{h(x)}\dd x\right)
 \geqq(b-a)^2
\]
を示し，等号条件を求めよ。
\end{example}

二つの積分をコーシー・シュワルツの不等式の右辺にそろえたい。そこで$f=\sqrt h$，$g=\frac1{\sqrt h}$とおけば，$f^2=h$，$g^2=\frac1h$となる。$h$は閉区間で連続かつ正なので，この二つの関数も連続である。さらに$fg=1$だから
\[
 (b-a)^2=\left(\int_a^b1\dd x\right)^2
 \leqq\left(\int_a^b h\dd x\right)
       \left(\int_a^b\frac1h\dd x\right).
\]
等号条件$f=cg$は$\sqrt h=\frac c{\sqrt h}$，すなわち$h=c$を意味する。したがって，等号は$h$が正の定数関数のとき，またそのときに限って成り立つ。

例えば$h(x)=1+x$，$0\leqq x\leqq1$なら，$h$は定数ではないから
\[
 \frac32\log2>1
\]
となる。ここでは対数の値を小数で計算せずに，$\log2>\frac23$という厳密な下からの評価が得られた。

\section{指数を変えた積分を結ぶ}\label{節:積分漸化式}

積分に自然数$n$が含まれているとき，毎回展開して計算するより，$n$を一つ小さくした積分と結び付ける方が簡単な場合がある。部分積分で$x^n$を微分すれば$n x^{n-1}$が現れるので，指数の異なる積分の間に関係を作れる。

\begin{example}{部分積分から漸化式を作る}
非負整数$n$に対して
\[
 A_n=\int_0^1 x^n\e^{-x}\dd x
\]
とおく。$A_n$の漸化式を求め，$A_2$を計算せよ。また$A_n\to0$を示せ。
\end{example}

$n\geqq1$のとき，$x^n$を微分し，$\e^{-x}$を積分する部分積分を行うと
\begin{align*}
 A_n
 &=\left[-x^n\e^{-x}\right]_0^1
   +n\int_0^1x^{n-1}\e^{-x}\dd x\\
 &=-\e^{-1}+nA_{n-1}.
\end{align*}
$x=0$の端点の項が$0$になるのは$n\geqq1$だからである。最初の値は別に
\[
 A_0=\int_0^1\e^{-x}\dd x=1-\e^{-1}
\]
と求める。よって
\[
 A_1=1-2\e^{-1},\qquad A_2=2-5\e^{-1}.
\]
一方，極限を調べるには，漸化式より積分の評価が短い。$[0,1]$では$0<\e^{-x}\leqq1$なので
\[
 0<A_n\leqq\int_0^1x^n\dd x=\frac1{n+1}\longrightarrow0.
\]
漸化式は値を順に求めるのに役立ち，積分表示は符号や大きさを調べるのに役立つ。同じ数列でも，問いに応じて使いやすい表し方が違う。

漸化式を一般項まで解くこともできる。$A_n=nA_{n-1}-\e^{-1}$を$n!$で割ると
\[
 \frac{A_n}{n!}=\frac{A_{n-1}}{(n-1)!}-\frac{\e^{-1}}{n!}.
\]
$1$から$n$まで足せば中間の項が消え，$0!=1$として
\[
 A_n=n!\left(1-\e^{-1}\sum_{k=0}^{n}\frac1{k!}\right)
\]
を得る。ただし，この形だけを見て$n\to\infty$を取ろうとすると，括弧の中がどれほど速く$0$に近づくかも調べる必要がある。先ほどの積分による評価は，その難しさを避けて直接$A_n\to0$を示している。

\section{三角関数のべきとウォリスの公式}\label{節:ウォリス公式}

次に
\[
 I_n=\int_0^{\frac{\pi}{2}}\sin^n x\dd x
 \qquad(n=0,1,2,\ldots)
\]
を考える。最初の二つは
\[
 I_0=\frac{\pi}{2},\qquad I_1=1
\]
である。$\sin^n x$を$\sin^{n-1}x\sin x$と分け，後ろの$\sin x$を積分すると，微分された$\sin^{n-1}x$から$\cos x$が現れ，$\cos^2x$を使う形になる。

$n\geqq2$として部分積分を行うと
\begin{align*}
 I_n
 &=\left[-\sin^{n-1}x\cos x\right]_0^{\frac{\pi}{2}}
 +(n-1)\int_0^{\frac{\pi}{2}}\sin^{n-2}x\cos^2x\dd x\\
 &=(n-1)\int_0^{\frac{\pi}{2}}\sin^{n-2}x(1-\sin^2x)\dd x\\
 &=(n-1)(I_{n-2}-I_n).
\end{align*}
両端の項は$0$である。最後の式の$I_n$を左辺に移すと
\begin{equation}\label{式:ウォリス漸化式}
 I_n=\frac{n-1}{n}I_{n-2}\qquad(n\geqq2)
\end{equation}
を得る。指数が$2$ずつ下がるので，偶数の場合は$I_0$へ，奇数の場合は$I_1$へつながる。

\begin{theorem}{ウォリスの積分公式}
$m\geqq1$に対して
\begin{align}
 I_{2m}
 &=\frac{1\cdot3\cdot5\cdots(2m-1)}{2\cdot4\cdot6\cdots2m}
   \frac{\pi}{2},\label{式:ウォリス偶数}\\
 I_{2m+1}
 &=\frac{2\cdot4\cdot6\cdots2m}{3\cdot5\cdot7\cdots(2m+1)}.
 \label{式:ウォリス奇数}
\end{align}
\end{theorem}

例えば$I_6=\frac56I_4=\frac56\frac34I_2=\frac56\frac34\frac12I_0$であり，一般の場合もまったく同じように式\eqref{式:ウォリス漸化式}を繰り返せばよい。偶数の積と奇数の積を短く書くために
\[
 (2m)!!=2\cdot4\cdots2m,\qquad
 (2m-1)!!=1\cdot3\cdots(2m-1)
\]
と定義することがある。これは\textbf{二重階乗}と呼ばれ，階乗を二回行うという意味ではない。この記号なら
\[
 I_{2m}=\frac{(2m-1)!!}{(2m)!!}\frac{\pi}{2},\qquad
 I_{2m+1}=\frac{(2m)!!}{(2m+1)!!}
\]
となる。

\begin{example}{偶数と奇数を分けて計算する}
次の二つの積分を求めよ。
\[
 \int_0^{\frac{\pi}{2}}\sin^6x\dd x,\qquad
 \int_0^{\frac{\pi}{2}}\sin^7x\dd x.
\]
\end{example}

偶数の場合は$I_0=\frac{\pi}{2}$，奇数の場合は$I_1=1$へ戻るので
\begin{align*}
 I_6&=\frac56\frac34\frac12\frac{\pi}{2}
     =\frac{5\pi}{32},\\
 I_7&=\frac67\frac45\frac23\cdot1
     =\frac{16}{35}.
\end{align*}
偶数と奇数で最後の値が違うことが，答えに$\pi$が現れるかどうかの違いになっている。積分区間が$\left[0,\frac{\pi}{2}\right]$であることも大切である。任意の区間では，部分積分の端点の項が同じように消えるとは限らない。

\subsection{隣り合う積分の比をはさむ}

公式は積分の値を積で表したものである。今度は，積分の大小を使ってその積の極限を求める。$0\leqq\sin x\leqq1$だから，指数を大きくすると値は小さくなる。したがって
\[
 I_{n+2}\leqq I_{n+1}\leqq I_n.
\]
$I_n>0$で割り，式\eqref{式:ウォリス漸化式}を使うと
\[
 \frac{n+1}{n+2}
 =\frac{I_{n+2}}{I_n}
 \leqq\frac{I_{n+1}}{I_n}\leqq1.
\]
両端が$1$へ近づくため
\begin{equation}\label{式:ウォリス隣接比}
 \lim_{n\to\infty}\frac{I_{n+1}}{I_n}=1
\end{equation}
である。各$I_n$の極限を先に求めなくても，比の極限は分かる。

ここで，次の有限の積を作る。
\[
 P_m=\prod_{k=1}^{m}\frac{(2k)^2}{(2k-1)(2k+1)}.
\]
$\prod$は，指定された範囲の項をすべて掛ける記号である。最初の数項を書けば
\[
 P_m=\frac{2\cdot2}{1\cdot3}
      \frac{4\cdot4}{3\cdot5}\cdots
      \frac{2m\cdot2m}{(2m-1)(2m+1)}.
\]
偶数と奇数の積をそれぞれまとめ，式\eqref{式:ウォリス偶数}と式\eqref{式:ウォリス奇数}を使うと
\begin{align*}
 \frac{I_{2m+1}}{I_{2m}}
 &=\frac{2}{\pi}
   \frac{\{(2m)!!\}^2}{(2m-1)!!(2m+1)!!}\\
 &=\frac{2}{\pi}P_m.
\end{align*}
左辺は$1$へ近づくので，次の極限が得られる。
\begin{equation}\label{式:ウォリスの積}
 \boxed{\displaystyle
 \lim_{m\to\infty}\prod_{k=1}^{m}
 \frac{4k^2}{4k^2-1}=\frac{\pi}{2}}
\end{equation}
これが\textbf{ウォリスの積}である。無限個の数を一度に掛けたと考えるのではなく，有限の積$P_m$の数列の極限として定義している。

証明の途中の不等式を残しておけば，近似の誤差も評価できる。$n=2m$の場合には
\[
 \frac{2m+1}{2m+2}
 \leqq\frac{I_{2m+1}}{I_{2m}}\leqq1
\]
だから
\[
 \frac{2m+1}{2m+2}\frac{\pi}{2}
 \leqq P_m\leqq\frac{\pi}{2}.
\]
従って，正の数を掛けて整理すれば
\[
 2P_m\leqq\pi\leqq\frac{2m+2}{2m+1}\,2P_m.
\]
公式が極限を与えるだけでなく，有限の計算の結果がどれほど確かな範囲を与えるかも分かる。

\subsection{中央の二項係数との関係}

偶数公式の分数には，階乗を使った別の形がある。
\[
 (2m)!!=2^m m!,\qquad
 (2m-1)!!=\frac{(2m)!}{2^m m!}
\]
だから
\begin{equation}\label{式:ウォリスと二項係数}
 I_{2m}=\frac{\pi}{2}\frac{(2m)!}{2^{2m}(m!)^2}
 =\frac{\pi}{2}\frac{\binom{2m}{m}}{4^m}.
\end{equation}
ここで$\binom{2m}{m}=\frac{(2m)!}{(m!)^2}$は二項係数である。

さらに偶数公式と奇数公式を掛けると，多くの因子が消えて
\[
 I_{2m}I_{2m+1}=\frac{\pi}{2(2m+1)}
\]
となる。これと$\frac{I_{2m+1}}{I_{2m}}\to1$を組み合わせれば
\[
 mI_{2m}^2
 =\frac{m\pi}{2(2m+1)}\frac{I_{2m}}{I_{2m+1}}
 \longrightarrow\frac{\pi}{4}.
\]
$I_{2m}>0$なので平方根を取ることができ，
\[
 \sqrt m\,I_{2m}\longrightarrow\frac{\sqrt\pi}{2}.
\]
式\eqref{式:ウォリスと二項係数}に戻すと
\begin{equation}\label{式:中央二項係数の極限}
 \lim_{m\to\infty}\frac{\sqrt m}{4^m}\binom{2m}{m}
 =\frac1{\sqrt\pi}
\end{equation}
を得る。階乗を含む数列の大きさが，三角関数の積分を経由して分かったことになる。

\section{多項式の積分からベータ積分へ}\label{節:ベータ積分}

\subsection{非負整数の指数}

非負整数$m,n$に対して
\[
 J_{m,n}=\int_0^1x^m(1-x)^n\dd x
\]
を考える。$m,n$が小さければ展開で求められる。しかし指数が大きくなると，展開した項数も多くなる。ここでは$x^m$と$(1-x)^n$をそのまま残し，部分積分で指数を一つずつ変える。

$n\geqq1$として，$x^m$を積分し，$(1-x)^n$を微分すると
\begin{align*}
 J_{m,n}
 &=\left[\frac{x^{m+1}}{m+1}(1-x)^n\right]_0^1
   +\frac{n}{m+1}\int_0^1x^{m+1}(1-x)^{n-1}\dd x\\
 &=\frac{n}{m+1}J_{m+1,n-1}.
\end{align*}
両端で積が$0$になるので，端点の項は消える。この式を$n$回繰り返すと
\[
 J_{m,n}
 =\frac{n!}{(m+1)(m+2)\cdots(m+n)}J_{m+n,0}.
\]
最後の積分は
\[
 J_{m+n,0}=\int_0^1x^{m+n}\dd x=\frac1{m+n+1}
\]
だから
\begin{equation}\label{式:整数ベータ積分}
 \boxed{\displaystyle
 \int_0^1x^m(1-x)^n\dd x=\frac{m!n!}{(m+n+1)!}}
\end{equation}
となる。$n=0$の場合も，右辺を直接整理すれば同じ式で表される。

\begin{example}{展開せずに多項式の積分を求める}
次の積分を求めよ。
\[
 \int_0^1x^3(1-x)^2\dd x.
\]
\end{example}

式\eqref{式:整数ベータ積分}で$m=3,n=2$とすると
\[
 \int_0^1x^3(1-x)^2\dd x
 =\frac{3!2!}{6!}=\frac1{60}.
\]
展開で検算すると
\[
 \int_0^1(x^3-2x^4+x^5)\dd x
 =\frac14-\frac25+\frac16=\frac1{60}
\]
となり，一致する。公式は展開を不要にするが，指数がどこに入るかを確かめるため，小さい場合を計算してみることにも意味がある。

\subsection{端点で大きくなる関数の積分}

指数を正の実数へ広げるため，$p,q>0$として
\[
 x^{p-1}(1-x)^{q-1}
\]
を考えたい。例えば$p=\frac12$なら$x^{p-1}=\frac1{\sqrt x}$となり，$x=0$で値が定まらない。閉区間で連続な関数の積分という，これまでの定義をそのまま使うことはできない。

そこで，問題のある端点をいったん避けて積分し，その端点へ近づける。$0<c<1$を一つ固定して
\begin{equation}\label{式:両端広義積分の定義}
 \int_0^1 f(x)\dd x
 =\lim_{\varepsilon\downarrow0}\int_{\varepsilon}^{c}f(x)\dd x
  +\lim_{\delta\downarrow0}\int_c^{1-\delta}f(x)\dd x
\end{equation}
と定める。ただし，右辺の\emph{二つの極限がそれぞれ有限の値を持つ場合}に限る。このような積分を\textbf{広義積分}といい，有限の値を持つことを収束するという。$c$を別の内点に変えても，その二点間の通常の積分が一方から他方へ移るだけなので，和の値は変わらない。

両端へ同じ速さで近づける$\int_{\varepsilon}^{1-\varepsilon}$の極限だけでは，この定義を満たしたことにはならない。一方の端で正に発散し，他方の端で負に発散するものが打ち消し合うおそれがあるためである。この章のベータ積分は非負関数だが，後にも使えるよう，両端を独立に調べる定義にしておく。

まず，$p>0$なら
\[
 \int_{\varepsilon}^{c}x^{p-1}\dd x
 =\frac{c^p-\varepsilon^p}{p}
 \longrightarrow\frac{c^p}{p}
\]
である。$p$が$1$より小さくても，$\varepsilon^p\to0$となるので収束する。関数の高さが無限に大きくなることと，面積が無限に大きくなることは同じではない。

ベータ積分の左半分では，$x\in\left[0,\frac12\right]$において$(1-x)^{q-1}$が連続で，ある正の定数$M$以下である。従って
\[
 0\leqq x^{p-1}(1-x)^{q-1}\leqq Mx^{p-1}.
\]
下端$\varepsilon$を小さくすると積分値は増加するが，上の評価により$\frac{M}{p}\left(\frac12\right)^p$以下にとどまる。単調で上に有界な数列の極限と同じ考え方で，この極限は有限に存在する。右半分も$u=1-x$と置けば$q>0$によって同様に収束する。したがって$p,q>0$なら，両端の広義積分がともに定まる。

\begin{definition}{ベータ関数}
$p,q>0$に対して
\begin{equation}\label{式:ベータ関数の定義}
 B(p,q)=\int_0^1x^{p-1}(1-x)^{q-1}\dd x
\end{equation}
と定める。端点で関数が定まらない場合には，式\eqref{式:両端広義積分の定義}による広義積分とする。
\end{definition}

整数の場合との対応は，指数に$1$のずれがある。式\eqref{式:整数ベータ積分}は
\[
 B(m+1,n+1)=\frac{m!n!}{(m+n+1)!}
\]
と書ける。$B(m,n)$の指数は$m-1,n-1$であり，$m,n$そのものではない。記号の引数と被積分関数の指数を取り違えないようにする。

\subsection{対称性と漸化式}

$u=1-x$と置くと二つの指数が入れ替わるため
\begin{equation}\label{式:ベータ対称性}
 B(p,q)=B(q,p)
\end{equation}
である。広義積分の場合も，端点を避けた区間で置換し，各端点の極限を取れば同じ結論になる。

また$x+(1-x)=1$を掛ければ
\begin{equation}\label{式:ベータ和の関係}
 B(p+1,q)+B(p,q+1)=B(p,q).
\end{equation}
もう一つの関係を得るために$x^p(1-x)^q$を微分する。
\[
 \frac{\dd}{\dd x}\{x^p(1-x)^q\}
 =p x^{p-1}(1-x)^q-qx^p(1-x)^{q-1}.
\]
これを$\varepsilon$から$1-\delta$まで積分し，両端をそれぞれ極限へ近づける。$p,q>0$なので$x^p(1-x)^q$は両端で$0$へ近づき，右辺の二つの積分はすでに収束すると分かっている。よって
\[
 pB(p,q+1)=qB(p+1,q).
\]
式\eqref{式:ベータ和の関係}と連立すると
\begin{equation}\label{式:ベータ漸化式}
 B(p+1,q)=\frac{p}{p+q}B(p,q),\qquad
 B(p,q+1)=\frac{q}{p+q}B(p,q)
\end{equation}
を得る。この漸化式は整数以外にも成り立つ。整数のときには，$B(1,1)=1$から出発して階乗の公式を再び導くこともできる。

\subsection{三角関数の積分とのつながり}

$x=\sin^2t$と置くと，$0<t<\frac{\pi}{2}$で$\dd x=2\sin t\cos t\dd t$である。$\sin t,\cos t$が正なので，実数の指数もそのまま扱える。式\eqref{式:ベータ関数の定義}は
\begin{equation}\label{式:ベータ三角表示}
 B(p,q)=2\int_0^{\frac{\pi}{2}}
 (\sin t)^{2p-1}(\cos t)^{2q-1}\dd t
\end{equation}
となる。厳密には端点を避けて置換した後で極限を取る。左辺の両端の積分が収束しているので，対応する右辺の極限も有限である。

$p=q=\frac12$なら両方の指数が$0$になり，
\[
 B\left(\frac12,\frac12\right)
 =2\int_0^{\frac{\pi}{2}}1\dd t=\pi.
\]
つまり
\[
 \int_0^1\frac1{\sqrt{x(1-x)}}\dd x=\pi
\]
である。両端で高さが限りなく大きくなるにもかかわらず，積分は有限になる具体例である。

一方，$p=\frac{n+1}{2}$，$q=\frac12$とすれば
\[
 B\left(\frac{n+1}{2},\frac12\right)=2I_n.
\]
ベータ関数の漸化式で$p$を$1$増やすことは，$\sin$の指数$n$を$2$増やすことに対応する。実際
\[
 B\left(\frac{n+3}{2},\frac12\right)
 =\frac{n+1}{n+2}
  B\left(\frac{n+1}{2},\frac12\right)
\]
は$I_{n+2}=\frac{n+1}{n+2}I_n$そのものである。整数の多項式積分とウォリスの積分は，ベータ積分の異なる場合として結び付いた。

\section*{章末問題}
\begin{enumerate}[label=\textbf{\thechapter.\arabic*},leftmargin=*]
 \item\label{問:区分求積分子が添字}
 次の極限を求めよ。
 \[
  \lim_{n\to\infty}\sum_{k=1}^{n}\frac{k}{n^2+k^2}.
 \]
 \item\label{問:区分求積根号}
 次の極限を求めよ。
 \[
  \lim_{n\to\infty}\cfrac1{n^{\cfrac32}}
      \sum_{k=1}^{n}\sqrt{n+k}.
 \]
 \item\label{問:逆平方和の残り}
 自然数$N>n\geqq1$に対して
 \[
  \frac1{n+1}-\frac1{N+1}
  \leqq\sum_{k=n+1}^{N}\frac1{k^2}
  \leqq\frac1n-\frac1N
 \]
 を示せ。これを使って$\sum_{k=1}^{\infty}\frac1{k^2}$が収束し，その第$n$項より後の和$R_n$が
 $\frac1{n+1}\leqq R_n\leqq\frac1n$を満たすことを示せ。
 \item\label{問:積分制約と二乗}
 $f$は$[0,1]$で連続で，$\int_0^1xf(x)\dd x=1$を満たすとする。
 $\int_0^1f(x)^2\dd x$の最小値を求め，それを与える関数を求めよ。
 \item\label{問:対数の積分漸化式}
 非負整数$n$に対して$K_n=\int_1^{\e}(\log x)^n\dd x$とおく。
 $K_0$と，$n\geqq1$における漸化式を求め，$K_3$を計算せよ。
 \item\label{問:ウォリス積の誤差}
 $P_m=\prod_{k=1}^{m}\frac{4k^2}{4k^2-1}$とする。
 \[
 0\leqq\frac{\pi}{2}-P_m\leqq\frac{\pi}{4m+4}
 \]
 を示せ。また$\lim_{m\to\infty}\sqrt m\int_0^{\frac{\pi}{2}}\cos^{2m}x\dd x$を求めよ。
 \item\label{問:整数ベータ二項係数}
 非負整数$m,n$に対して
 \[
 \int_0^1\binom{m+n}{m}x^m(1-x)^n\dd x=\frac1{m+n+1}
 \]
 を示せ。次に$\int_0^1x^2(1-x)^4\dd x$を求めよ。
 \item\label{問:半整数ベータ}
 広義積分$\int_0^1\sqrt{\frac{x}{1-x}}\dd x$の収束を両端で確かめ，ベータ関数の漸化式を使って値を求めよ。さらに$x=\sin^2t$による置換でも検算せよ。
\end{enumerate}

\smallskip\noindent 解答は\pageref{解答:第9章}ページ。

\chapter{置換の背景にある関数と図形}\label{章:置換の背景}

\chapterref{章:置換積分}では，$x=\sin\theta$や$x=\tan\theta$と置いて，根号や分母を簡単にした。しかし，不定積分の答えを$x$へ戻そうとすると，「正弦が$x$になる角」を関数として書く必要がある。また，$\sin x$と$\cos x$が分母で混ざった式では，一方を新しい変数にするだけでは，他方が根号として残ることがある。

この章では，まず逆三角関数を定義して，角へ戻る操作を明確にする。続いて，三角関数を同時に有理式へ変える置換を考える。その式は円上の点を直線の傾きで表す式でもあり，整数の直角三角形とも結び付く。最後に，平方の差が一定になる双曲線関数を導入し，同じ根号積分を三つの方法で比べる。使う前提は，これまでの三角関数・指数関数の性質と微分，置換積分，部分積分である。新しい関数の微分や積分公式は，それらから導く。

\section{角を一つに決める逆三角関数}\label{節:逆三角関数}

$\sin\theta=\frac12$を満たす角は，$\frac\pi6$，$\frac{5\pi}6$など無数にある。このままでは，$x$から角を返す関数にはならない。\sectionref{節:逆向きの置換}で扱った枝の選択を，ここでも行う。

\begin{definition}{逆三角関数とその値の範囲}
次の区間に制限した三角関数の逆関数を，それぞれ$\arcsin$，$\arccos$，$\arctan$と書く。
\[
\begin{array}{c|c|c}
\text{逆関数}&\text{入力の範囲}&\text{返す角の範囲}\\ \hline
\arcsin x&[-1,1]&\left[-\frac\pi2,\frac\pi2\right]\\[3pt]
\arccos x&[-1,1]&[0,\pi]\\[3pt]
\arctan x&\R&\left(-\frac\pi2,\frac\pi2\right)
\end{array}
\]
例えば$\theta=\arcsin x$とは，$\sin\theta=x$かつ$-\frac\pi2\leqq\theta\leqq\frac\pi2$であることを表す。
\end{definition}

正弦と正接は指定した区間で狭義単調増加し，余弦は狭義単調減少するので，角は一つに決まる。正接の範囲の端を含めないのは，そこで$\cos\theta=0$になり，正接が定義されないためである。$\sin^{-1}x$などの表記もあるが，逆数との混同を避け，本書では$\arcsin x$を使う。

定義から$\sin(\arcsin x)=x$は$-1\leqq x\leqq1$で成り立つ。一方，$\arcsin(\sin\theta)=\theta$は，$\theta$が選んだ範囲にあるときに限る。左辺が返す値は，必ずその範囲内だからである。

\begin{example}{逆向きに戻すと角が変わる場合}
次の値を求めよ。
\[
\arcsin\left(\sin\frac{5\pi}6\right),\qquad
\arccos\left(\cos\frac{7\pi}6\right),\qquad
\arctan\left(\tan\frac{3\pi}4\right).
\]
\end{example}

第一の正弦は$\frac12$であり，$\arcsin$の値の範囲にある角は$\frac\pi6$である。第二の余弦は$-\frac{\sqrt3}2$であり，$[0,\pi]$の中では$\frac{5\pi}6$になる。第三の正接は$-1$であり，$\left(-\frac\pi2,\frac\pi2\right)$の中では$-\frac\pi4$になる。従って答えは順に
\[
\frac\pi6,\qquad\frac{5\pi}6,\qquad-\frac\pi4
\]
である。逆関数の記号を打ち消す前に，戻った角が指定された範囲にあるかを確かめる。

\subsection{微分の分母と符号を決める}

$y=\arcsin x$，$-1<x<1$とおく。$\sin y=x$だから，逆関数の微分の公式，または両辺の微分によって
\[
\cos y\,y'=1
\]
となる。ここで$-\frac\pi2<y<\frac\pi2$なので$\cos y>0$である。従って$\cos y=\sqrt{1-x^2}$と符号まで決まり，
\[
(\arcsin x)'=\frac1{\sqrt{1-x^2}}\qquad(-1<x<1)
\]
を得る。端点では分母が$0$になり，この有限な微分係数の公式は使えない。

$y=\arccos x$の場合には，$0<y<\pi$で$\sin y>0$であり，$-\sin y\,y'=1$となる。そのため
\[
(\arccos x)'=-\frac1{\sqrt{1-x^2}}\qquad(-1<x<1)
\]
である。負号は，余弦が選んだ区間で減少することとも一致する。$y=\arctan x$では$\tan y=x$を微分して
\[
(1+\tan^2y)y'=1,\qquad
(\arctan x)'=\frac1{1+x^2}\qquad(x\in\R)
\]
となる。これらは積分公式を仮定して得たものではなく，もとの三角関数の微分から得た公式である。

合成関数の微分を使えば，$a>0$に対して
\begin{equation}\label{式:逆三角関数の積分}
\begin{aligned}
\int\frac1{a^2+x^2}\dd x
 &=\frac1a\arctan\frac xa+C,\\
\int\frac1{\sqrt{a^2-x^2}}\dd x
 &=\arcsin\frac xa+C\qquad(|x|<a)
\end{aligned}
\end{equation}
を得る。第一の式では，$x$を$a$で割ることと，外側の係数$\frac1a$の両方が必要である。第二の式では，根号から出る$a$と合成関数の微分の$\frac1a$が打ち消し合う。

\begin{example}{平方完成と逆三角関数をつなぐ}
次の不定積分を求めよ。
\[
\text{(1)}\quad\int\frac1{x^2-4x+8}\dd x,
\qquad
\text{(2)}\quad\int\frac1{\sqrt{9-4x^2}}\dd x.
\]
\end{example}

(1) の分母は$(x-2)^2+4$である。中心を$2$へ移し，長さの尺度を$2$でそろえるため，$u=\frac{x-2}{2}$と置く。$\dd x=2\dd u$だから
\[
\int\frac1{x^2-4x+8}\dd x
=\frac12\int\frac1{1+u^2}\dd u
=\frac12\arctan\frac{x-2}{2}+C.
\]
(2) では$u=\frac{2x}{3}$と置くと，$\dd x=\frac32\dd u$，$\sqrt{9-4x^2}=3\sqrt{1-u^2}$である。よって
\[
\int\frac1{\sqrt{9-4x^2}}\dd x
=\frac12\int\frac1{\sqrt{1-u^2}}\dd u
=\frac12\arcsin\frac{2x}{3}+C.
\]
もとの関数が実数で定義される区間は$-\frac32<x<\frac32$であり，この範囲で答えを用いる。

\section{逆数の三角関数を使う理由}\label{節:相反三角関数}

三角置換では$\frac1{\cos\theta}$や$\frac1{\sin\theta}$が繰り返し現れる。これらを一つの関数として書くと，平方や微分の関係を読み取りやすい。
\begin{definition}{正割・余割・余接}
分母が$0$でないところで
\[
\sec x=\frac1{\cos x},\qquad
\csc x=\frac1{\sin x},\qquad
\cot x=\frac{\cos x}{\sin x}
\]
と定義する。順に正割，余割，余接という。ここでの逆数と，前節の逆関数は異なる操作である。
\end{definition}

$\sin^2x+\cos^2x=1$を$\cos^2x$，$\sin^2x$で割ると，それぞれ
\[
1+\tan^2x=\sec^2x,\qquad
1+\cot^2x=\csc^2x
\]
となる。微分も，すでに知っている商の微分から得られる。例えば
\[
(\sec x)'=\frac{\sin x}{\cos^2x}=\sec x\tan x.
\]
同様に計算して
\begin{equation}\label{式:相反三角関数の微分}
(\csc x)'=-\csc x\cot x,\qquad
(\cot x)'=-\csc^2x
\end{equation}
を得る。この記法により，$x=\tan\theta$と置いたときの$\dd x$は$\sec^2\theta\dd\theta$，$\sqrt{1+x^2}$は$|\sec\theta|$と短く書ける。ただし，記号を短くしても，絶対値を外すための符号の確認は残る。

\begin{example}{正割と余割を対数へ結び付ける}
$\int\sec x\dd x$と$\int\csc x\dd x$を求めよ。
\end{example}

$\sec x$の微分は$\sec x\tan x$，$\tan x$の微分は$\sec^2x$である。この二つを足すと
\[
(\sec x+\tan x)'=\sec x(\sec x+\tan x)
\]
となり，元の$\sec x$が因子に現れる。また
\[
(\sec x+\tan x)(\sec x-\tan x)=1
\]
だから，$\sec x+\tan x$は定義されるところで$0$にならない。従って
\[
\int\sec x\dd x
=\int\frac{(\sec x+\tan x)'}{\sec x+\tan x}\dd x
=\log|\sec x+\tan x|+C.
\]
これは，分母の微分を分子に作る対数型の積分である。

余割では引き算を考えると
\[
(\csc x-\cot x)'=-\csc x\cot x+\csc^2x
=\csc x(\csc x-\cot x).
\]
さらに$(\csc x-\cot x)(\csc x+\cot x)=1$なので，
\[
\int\csc x\dd x=\log|\csc x-\cot x|+C
\]
となる。どちらも，もとの三角関数が定義される一つの区間で使う。特異点をまたぐ不定積分の定数を，一つにそろえる必要はない。

\section{半角を使って三角関数を有理式にする}\label{節:半角置換}

$\int R(\sin x,\cos x)\dd x$のように，正弦と余弦を足す・掛ける・割ることで作った式を考える。このような式を，二変数の\textbf{有理式}$R$で表すと書く。$t=\sin x$と置くだけでは，$\cos x$を$\pm\sqrt{1-t^2}$と表すことになり，根号が残る。両方を一度に有理式へ変えるには，半分の角の正接を使う。

$t=\tan\frac x2$と置く。二倍角の公式と$1+\tan^2\frac x2=\sec^2\frac x2$から
\begin{equation}\label{式:半角置換}
\begin{aligned}
\sin x&=\cfrac{2\sin\cfrac x2\cos\cfrac x2}
 {\cos^2\cfrac x2+\sin^2\cfrac x2}
 =\frac{2t}{1+t^2},\\
\cos x&=\cfrac{\cos^2\cfrac x2-\sin^2\cfrac x2}
 {\cos^2\cfrac x2+\sin^2\cfrac x2}
 =\frac{1-t^2}{1+t^2},\\
\frac{\dd t}{\dd x}&=\frac12\sec^2\frac x2
 =\frac{1+t^2}{2},\qquad
\dd x=\frac{2}{1+t^2}\dd t
\end{aligned}
\end{equation}
を得る。最初の二つでは分子・分母を$\cos^2\frac x2$で割っている。この操作が可能な，$x\ne(2k+1)\pi$の区間を選ぶ。ここで$k$は整数である。

この置換を\textbf{半角置換}，または\textbf{ワイエルシュトラス置換}という。例えば$-\pi<x<\pi$では$t$は実数全体を狭義単調増加して動き，逆向きには$x=2\arctan t$と書ける。一般には
\[
(2k-1)\pi<x<(2k+1)\pi,\qquad
x=2\arctan t+2k\pi
\]
である。正弦・余弦・$\dd x$のすべてが$t$の有理式になるので，置換後は割り算や部分分数などの代数的な整理を使える。

\begin{example}{分母に混ざった正弦を取り除く}
$\int_0^{\frac\pi2}\frac1{2+\sin x}\dd x$を求めよ。
\end{example}

分母$2+\sin x$の微分は$\cos x$であり，分子には現れていない。そこで半角置換を使う。$t=\tan\frac x2$と置くと，端点は$t=0,1$へ移り，
\[
\begin{aligned}
\int_0^{\frac\pi2}\frac1{2+\sin x}\dd x
&=\int_0^1\cfrac1{2+\cfrac{2t}{1+t^2}}
 \frac2{1+t^2}\dd t\\
&=\int_0^1\frac1{t^2+t+1}\dd t.
\end{aligned}
\]
ここで分母を平方完成すると$t^2+t+1=\left(t+\frac12\right)^2+\frac34$である。よって
\[
\begin{aligned}
\int_0^1\frac1{t^2+t+1}\dd t
&=\frac2{\sqrt3}\left[\arctan\frac{2t+1}{\sqrt3}\right]_0^1\\
&=\frac2{\sqrt3}\left(\frac\pi3-\frac\pi6\right)
=\frac\pi{3\sqrt3}.
\end{aligned}
\]

もし積分区間が$[0,2\pi]$なら，$x=\pi$で$t$が定義されなくなる。端点だけを$t=0,0$へ換えて答えを$0$としてはいけない。元の被積分関数は$\pi$で連続でも，置換に使った座標はそこで途切れる。この場合には$\pi$の前後に分け，$t\to+\infty$と$t\to-\infty$を別々に扱う。章末問題で確かめる。

\section{同じ有理式が円から現れる}\label{節:円の有理パラメータ}

前節の二式は
\[
X=\frac{1-t^2}{1+t^2},\qquad
Y=\frac{2t}{1+t^2}
\]
と書くと，$X^2+Y^2=1$を満たす。正弦と余弦を使ったのだから当然ともいえるが，この式は三角関数なしでも導ける。

図\ref{図:円の有理パラメータ}のように，単位円上の点$A=(-1,0)$を固定し，傾き$t$の直線$Y=t(X+1)$を引く。円との交点を求めるため代入すると
\[
\begin{aligned}
X^2+t^2(X+1)^2&=1,\\
(X+1)\{(1+t^2)X+(t^2-1)\}&=0.
\end{aligned}
\]
一つの解$X=-1$は，固定した点$A$である。もう一つの交点$P$は，上の$X,Y$で表される。$t$を動かして点を表すので，$t$を\textbf{媒介変数}，この表示を円の\textbf{有理パラメータ表示}という。

\begin{figure}[htbp]
\centering
\begin{tikzpicture}[x=2.25cm,y=2.25cm,>=Stealth]
 \draw[->] (-1.4,0)--(1.45,0) node[right]{$X$};
 \draw[->] (0,-1.2)--(0,1.35) node[above]{$Y$};
 \draw[thick,tealblack] (0,0) circle(1);
 \draw[thick,blueblack] (-1.23,-0.115)--(1.15,1.075);
 \draw[dashed] (0.6,0)--(0.6,0.8)--(0,0.8);
 \draw[blueblack] (0,0)--(0.6,0.8);
 \draw[->] (0.26,0) arc[start angle=0,end angle=53.1301,radius=0.26];
 \node at (0.35,0.17) {$x$};
 \draw[->] (-0.65,0) arc[start angle=0,end angle=26.5651,radius=0.35];
 \node at (-0.42,0.15) {$\frac x2$};
 \fill (-1,0) circle(1.2pt) node[below left]{$A=(-1,0)$};
 \fill (0.6,0.8) circle(1.2pt);
 \draw[thin,gray] (0.63,0.79)--(1.04,0.56);
 \node[anchor=west] at (1.04,0.52) {$P=(\cos x,\sin x)$};
 \node[above right] at (0.97,0.985) {$Y=t(X+1)$};
 \node[below left] at (0,0) {$O$};
\end{tikzpicture}
\caption{円上の固定点から引いた直線の傾きで，もう一つの交点を表す。図は$t=\frac12$の場合。}
\label{図:円の有理パラメータ}
\end{figure}

$P=(X,Y)\ne A$からは$t=\frac{Y}{X+1}$と戻せるので，有限な実数$t$で，円上の$A$以外のすべての点を一度ずつ表せる。$P=(\cos x,\sin x)$なら
\[
t=\frac{\sin x}{1+\cos x}=\tan\frac x2
\]
である。最後の等式は二倍角の公式から確かめられる。図の上半円では，二等辺三角形$AOP$の角を見ても，直線$AP$が正の$X$方向となす角は$\frac x2$になる。半角が現れる理由を，円の上の角としても読めた。表せなかった点$A$は，ちょうど$x=(2k+1)\pi$に対応する。

\subsection{有理数の点から整数の直角三角形へ}

$t=\frac nm$を整数$m,n$で表すと，
\[
X=\frac{m^2-n^2}{m^2+n^2},\qquad
Y=\frac{2mn}{m^2+n^2}.
\]
従って
\begin{equation}\label{式:ピタゴラス数}
(m^2-n^2)^2+(2mn)^2=(m^2+n^2)^2
\end{equation}
である。$m>n>0$なら三つの整数は正であり，直角三角形の三辺になる。$a^2+b^2=c^2$を満たす正の整数の組を\textbf{ピタゴラス数}という。

逆に整数の直角三角形$(a,b,c)$があれば，点$\left(\frac ac,\frac bc\right)$は単位円上にある。その傾きは
\[
t=\cfrac{\cfrac bc}{1+\cfrac ac}=\frac{b}{c+a}
\]
と有理数になる。これを$\frac nm$と表せば，上の公式が元の三辺と同じ比を与える。従って，この表示は偶然いくつかの例を作るだけでなく，すべての整数の直角三角形の辺の比を表す。得られた三整数に共通因数があれば，それで割ればよい。

三辺の最大公約数が$1$のピタゴラス数を，\textbf{原始ピタゴラス数}という。整数$m>n>0$が互いに素で，一方が偶数，他方が奇数ならば，式\eqref{式:ピタゴラス数}は原始ピタゴラス数を与える。逆に，すべての原始ピタゴラス数は，直角を挟む二辺の交換を除いてこの形で表せる。この主張を，整数の素因数分解の一意性を使って確かめよう。

まず上の条件で作った三整数に，共通の素因数$p$があると仮定する。$m^2-n^2$と$m^2+n^2$はともに奇数なので，$p$は奇素数である。この二数の和と差から，$p$は$2m^2$と$2n^2$を割る。従って$p$は$m,n$の両方を割り，互いに素という条件に反する。共通の素因数がないので，三整数の最大公約数は$1$である。

逆に原始ピタゴラス数$(a,b,c)$を考える。二辺$a,b$がともに偶数なら原始ではなく，ともに奇数なら$a^2+b^2$を$4$で割った余りが$2$となり，平方数$c^2$にはなれない。従って一方だけが偶数であり，必要なら辺を交換して$b$を偶数，$a$を奇数とする。このとき$c$も奇数である。また$a,c$に共通の素因数があれば，$b^2=c^2-a^2$，従って$b$もその素因数で割れ，原始という条件に反する。よって$a,c$は互いに素である。

ここで
\[
A=\frac{c+a}{2},\qquad B=\frac{c-a}{2}
\]
と置くと，$A>B>0$の整数であり，
\[
AB=\frac{c^2-a^2}{4}=\left(\frac b2\right)^2.
\]
$A,B$の共通因数は$A+B=c$と$A-B=a$も割るので，$A,B$は互いに素である。互いに素な二数の積が平方数なら，素因数は片方にしか現れず，その指数はすべて偶数になる。従って$A=m^2$，$B=n^2$と書ける。和と差を戻せば
\[
a=m^2-n^2,\qquad b=2mn,\qquad c=m^2+n^2.
\]
$A,B$が互いに素だから$m,n$も互いに素であり，$c$が奇数だから$m,n$の偶奇は異なる。これで逆向きも示せた。円の表示が辺の比を与えることに加えて，約分の要らない整数を作る条件まで分かった。

\begin{example}{傾きから三辺を作り，三辺から傾きを戻す}
$t=\frac23$からピタゴラス数を作れ。また三辺が$8,15,17$の三角形に対応する傾きを求めよ。
\end{example}

$m=3,n=2$なので
\[
m^2-n^2=5,\qquad2mn=12,\qquad m^2+n^2=13.
\]
確かに$5^2+12^2=13^2$である。後半では$P=\left(\frac{15}{17},\frac8{17}\right)$と選ぶと
\[
t=\frac8{17+15}=\frac14.
\]
$m=4,n=1$を公式へ戻せば$(15,8,17)$となる。二つの直角を挟む辺のどちらを$X$座標にするかで，対応する傾きは変わる。

\section{平方の差を作る双曲線関数}\label{節:双曲線関数}

円では$X^2+Y^2=1$を使った。これに対して，$\sqrt{1+x^2}$を簡単にしたい場合には，$Y^2-X^2=1$という平方の差を使えればよい。指数関数の和と差が，その形を作る。

\begin{definition}{双曲線関数}
実数$u$に対して
\[
\sinh u=\frac{\e^u-\e^{-u}}2,\qquad
\cosh u=\frac{\e^u+\e^{-u}}2,\qquad
\tanh u=\frac{\sinh u}{\cosh u}
\]
と定義する。順に双曲線正弦，双曲線余弦，双曲線正接という。
\end{definition}

$\cosh u>0$なので$\tanh u$はすべての実数で定義される。和と差の積から
\[
\begin{aligned}
\cosh^2u-\sinh^2u
&=\frac{(\e^u+\e^{-u})^2-(\e^u-\e^{-u})^2}{4}\\
&=\e^u\e^{-u}=1
\end{aligned}
\]
となる。従って$(X,Y)=(\sinh u,\cosh u)$は，双曲線$Y^2-X^2=1$の上側にある。図\ref{図:双曲線の媒介表示}のように，$u$が正なら$X>0$，負なら$X<0$である。この$u$は，円のときのような通常の中心角として定義したものではない。

\begin{figure}[htbp]
\centering
\begin{tikzpicture}[x=1.5cm,y=1.5cm,>=Stealth]
 \draw[->] (-2.4,0)--(2.55,0) node[right]{$X$};
 \draw[->] (0,-0.15)--(0,2.8) node[above]{$Y$};
 \draw[dashed,gray] (-2.2,2.2)--(0,0)--(2.2,2.2);
 \draw[very thick,tealblack,domain=-2.2:2.2,samples=80]
  plot(\x,{sqrt(1+\x*\x)});
 \draw[dashed,blueblack] (1.175201,0)--(1.175201,1.543081)--(0,1.543081);
 \fill[blueblack] (1.175201,1.543081) circle(1.3pt);
 \node[anchor=west] at (0.25,2.55) {$Y^2-X^2=1$};
 \node[anchor=east] at (-0.3,1.65) {$u<0$};
 \node[anchor=west] at (0.35,1.95) {$u>0$};
 \node[anchor=west] at (1.4,0.58) {$(\sinh u,\cosh u)$};
 \draw[->,blueblack] (1.75,0.82)--(1.2,1.48);
 \node[left] at (0,1) {$1$};
 \node[below left] at (0,0) {$0$};
\end{tikzpicture}
\caption{指数関数の和と差によって，双曲線の上側を一度ずつ表す。点は$u=1$の位置。}
\label{図:双曲線の媒介表示}
\end{figure}

定義を微分すると
\[
(\sinh u)'=\frac{\e^u+\e^{-u}}2=\cosh u,
\qquad
(\cosh u)'=\frac{\e^u-\e^{-u}}2=\sinh u.
\]
余弦の場合と違って，$\cosh$を微分しても負号は付かない。商の微分から
\begin{equation}\label{式:双曲線正接の微分}
(\tanh u)'=\frac{\cosh^2u-\sinh^2u}{\cosh^2u}
=\frac1{\cosh^2u}=1-\tanh^2u
\end{equation}
である。加えて指数関数の定義式を展開すると
\begin{equation}\label{式:双曲線の二倍公式}
\sinh(2u)=2\sinh u\cosh u,\qquad
\cosh(2u)=2\cosh^2u-1
\end{equation}
となる。例えば後者の右辺は
\[
\frac{(\e^u+\e^{-u})^2}{2}-1
=\frac{\e^{2u}+\e^{-2u}}2
\]
であり，左辺の定義に一致する。

\subsection{双曲線正弦から元へ戻す}

$\cosh u>0$なので，$\sinh u$は実数全体で狭義単調増加する。また$u\to\infty$で$\sinh u\to\infty$，$u\to-\infty$で$\sinh u\to-\infty$だから，どの実数$x$も$x=\sinh u$と一意に書ける。この逆関数を$\operatorname{arsinh}x$と書く。

$x=\sinh u$なら，$\cosh u=\sqrt{1+x^2}$である。定義式を足すと$\sinh u+\cosh u=\e^u$だから
\begin{equation}\label{式:逆双曲線正弦}
u=\operatorname{arsinh}x=\log\bigl(x+\sqrt{1+x^2}\bigr).
\end{equation}
対数の中は，$\sqrt{1+x^2}>|x|$によって常に正である。逆関数の微分から
\[
\frac{\dd u}{\dd x}=\frac1{\cosh u}
=\frac1{\sqrt{1+x^2}}
\]
となる。ここでも積分の値から逆関数を定義したのではなく，指数関数による定義から順に導いた。

\begin{example}{根号の分母を双曲線正弦で取り除く}
$\int_0^3\frac1{\sqrt{4+x^2}}\dd x$を求めよ。
\end{example}

$x=2\sinh u$と置けば
\[
\dd x=2\cosh u\dd u,\qquad
\sqrt{4+x^2}=2\cosh u.
\]
分母と微分の係数が打ち消し合う。端点$x=0$は$u=0$へ，$x=3$は
\[
u=\operatorname{arsinh}\frac32
=\log\frac{3+\sqrt{13}}2
\]
へ移る。従って積分値は$\log\frac{3+\sqrt{13}}2$である。三角置換$x=2\tan\theta$では$\sec\theta$の積分が残るが，この置換では$1$の積分になる。

$\cosh u$は偶関数なので，実数全体では逆関数を持たない。しかし$u\geqq0$に制限すれば，$1$から無限大へ狭義単調増加する。$x=\cosh u$を使う置換では，このように枝を選ぶ。$\sinh$と$\cosh$で，逆向きに戻すときの条件は同じではない。

\begin{example}{平方の差を双曲線正接で取り除く}
$\int_0^{\frac12}\frac1{1-x^2}\dd x$を求めよ。
\end{example}

$\tanh u=\frac{\e^{2u}-1}{\e^{2u}+1}$は実数全体で狭義単調増加し，値の範囲は$(-1,1)$である。$x=\tanh u$と置けば，式\eqref{式:双曲線正接の微分}によって$\dd x=(1-x^2)\dd u$となる。元へ戻す式は
\[
x(\e^{2u}+1)=\e^{2u}-1,\qquad
\e^{2u}=\frac{1+x}{1-x},\qquad
u=\frac12\log\frac{1+x}{1-x}
\]
である。従って端点は$0$と$\frac12\log3$へ移り，求める積分値は$\frac12\log3$となる。部分分数で
\[
\frac1{1-x^2}=\frac12\left(\frac1{1+x}+\frac1{1-x}\right)
\]
と分けても同じ対数が現れる。双曲線正接の逆関数が，その二つの対数を一つにまとめている。

\section{一つの根号積分を三つの置換で比べる}\label{節:根号積分の三方法}

ここまでの道具を，同じ積分に使ってみよう。違いは答えそのものではなく，途中のどの式が簡単になるかにある。

\begin{example}{三角・双曲線・代数による三通りの計算}
$\int\sqrt{1+x^2}\dd x$を，三角置換，双曲線関数の置換，根号を含む代数的な置換によって求めよ。
\end{example}

\subsection{三角置換では正割の三乗になる}

$x=\tan\theta$，$-\frac\pi2<\theta<\frac\pi2$と置く。この範囲では$\sec\theta>0$なので
\[
\sqrt{1+x^2}=\sec\theta,\qquad
\dd x=\sec^2\theta\dd\theta,\qquad
\int\sqrt{1+x^2}\dd x=\int\sec^3\theta\dd\theta.
\]
根号はなくなったが，三乗の積分が残った。$J=\int\sec^3\theta\dd\theta$と書き，$\sec\theta$を微分する側，$\sec^2\theta$を積分する側にして部分積分する。
\[
\begin{aligned}
J&=\sec\theta\tan\theta
-\int\sec\theta\tan^2\theta\dd\theta\\
&=\sec\theta\tan\theta
-\int\sec\theta(\sec^2\theta-1)\dd\theta\\
&=\sec\theta\tan\theta-J+\int\sec\theta\dd\theta.
\end{aligned}
\]
任意定数を最後にまとめるものとして，この等式を$J$について解く。前に求めた正割の積分から
\[
J=\frac12\sec\theta\tan\theta
+\frac12\log(\sec\theta+\tan\theta)+C.
\]
この範囲では$\sec\theta>|\tan\theta|$だから，対数の中は正である。$x$へ戻して
\begin{equation}\label{式:一と二乗の根号積分}
\int\sqrt{1+x^2}\dd x
=\frac12\left\{x\sqrt{1+x^2}
+\log\bigl(x+\sqrt{1+x^2}\bigr)\right\}+C
\end{equation}
を得る。三角置換だけで計算が終わるのではなく，その後に部分積分を使った。

\subsection{双曲線置換では二乗の公式を使う}

$x=\sinh u$と置くと，実数$x$全体に対して$u$が一つに決まり，
\[
\sqrt{1+x^2}=\cosh u,\qquad
\dd x=\cosh u\dd u
\]
となる。今回は微分も同じ$\cosh u$になるので，積分は二乗で済む。式\eqref{式:双曲線の二倍公式}より
\[
\begin{aligned}
\int\sqrt{1+x^2}\dd x
&=\int\cosh^2u\dd u\\
&=\frac12\int\{1+\cosh(2u)\}\dd u\\
&=\frac u2+\frac14\sinh(2u)+C\\
&=\frac12\{u+\sinh u\cosh u\}+C.
\end{aligned}
\]
$u=\log(x+\sqrt{1+x^2})$と戻せば，式\eqref{式:一と二乗の根号積分}と一致する。

\subsection{根号と元の変数の和を置く}

第三の方法では，関数の名前を使わずに
\[
t=x+\sqrt{1+x^2}
\]
と置く。この和を選ぶ理由は，差との積が$1$になるためである。
\[
\bigl(\sqrt{1+x^2}+x\bigr)
\bigl(\sqrt{1+x^2}-x\bigr)=1.
\]
従って$\sqrt{1+x^2}-x=\frac1t$であり，和と差を取れば
\[
x=\frac12\left(t-\frac1t\right),\qquad
\sqrt{1+x^2}=\frac12\left(t+\frac1t\right),\qquad
\dd x=\frac12\left(1+\frac1{t^2}\right)\dd t.
\]
$t>0$であり，この$x$の式は$t>0$から実数全体への狭義単調増加な対応である。根号も微分も有理式へ変わったので，
\[
\begin{aligned}
\int\sqrt{1+x^2}\dd x
&=\frac14\int\left(t+\frac2t+\frac1{t^3}\right)\dd t\\
&=\frac{t^2}{8}+\frac12\log t-\frac1{8t^2}+C.
\end{aligned}
\]
ここで
\[
x\sqrt{1+x^2}=\frac14\left(t^2-\frac1{t^2}\right)
\]
だから，再び同じ答えになる。

実は第二，第三の置換は直接つながっている。$x=\sinh u$なら
\[
t=x+\sqrt{1+x^2}=\sinh u+\cosh u=\e^u.
\]
つまり第三の方法は，双曲線の媒介変数$u$の代わりに$\e^u$を変数として，式を有理式へ直したものである。三角置換は円の関係式$1+\tan^2\theta=\sec^2\theta$を使い，双曲線置換は$\cosh^2u-\sinh^2u=1$を使う。同じ平方根に対して，どの恒等式と微分の組が短い計算を作るかが違う。

最後に答えを微分して確認しておく。$s=\sqrt{1+x^2}$と書けば$s'=\frac xs$であり，
\[
\frac{\dd}{\dd x}\log(x+s)
=\cfrac{1+\cfrac xs}{x+s}=\frac1s.
\]
従って式\eqref{式:一と二乗の根号積分}の右辺の微分は
\[
\frac12\left(s+\frac{x^2}{s}+\frac1s\right)
=\frac12(s+s)=\sqrt{1+x^2}
\]
となる。三つの計算は同じ原始関数へ到達した。

\section*{章末問題}
\begin{enumerate}[label=\textbf{\thechapter.\arabic*},leftmargin=*]
\item\label{問:逆三角関数の枝}
$\arcsin(\sin\frac{4\pi}3)$，$\arccos(\cos\frac{5\pi}4)$，$\arctan(\tan(-\frac{2\pi}3))$を求めよ。また$y=\arcsin(x^2)$を$-1<x<1$で微分せよ。
\item\label{問:逆三角関数の積分練習}
$\int\frac1{x^2+2x+5}\dd x$と$\int_0^1\frac1{\sqrt{4-x^2}}\dd x$を求めよ。
\item\label{問:余接と正接の和}
定義から$(\cot x)'=-\csc^2x$を示せ。さらに$\frac1{\sin x\cos x}=\tan x+\cot x$を使って，$\int\frac1{\sin x\cos x}\dd x$を求め，答えを用いる区間の条件を述べよ。
\item\label{問:半角置換の区間分割}
半角置換を使って$\int_0^{2\pi}\frac1{2+\sin x}\dd x$を求めよ。$x=\pi$の前後を分け，置換後の無限の端点は極限として扱うこと。
\item\label{問:円の点とピタゴラス数}
円上の点$\left(\frac{12}{13},\frac5{13}\right)$に対応する傾き$t$を求めよ。$t=\frac nm$から式\eqref{式:ピタゴラス数}で作られる三整数を求め，共通因数で割れ。また$m,n$がともに奇数なら，三整数がすべて偶数になることを示せ。
\item\label{問:双曲線二倍角と積分}
指数関数の定義から$\sinh(2u)=2\sinh u\cosh u$を示せ。さらに$\int_0^{\log2}\cosh^2u\dd u$を求めよ。
\item\label{問:双曲線余弦の置換}
$x>1$で$\int\frac1{\sqrt{x^2-1}}\dd x$を，$x=\cosh u$によって求めよ。$u$の範囲を指定し，元の変数へ戻す対数の式も導け。
\item\label{問:双曲線置換と代数置換}
$\int_0^2\sqrt{4+x^2}\dd x$を，$x=2\sinh u$と$t=x+\sqrt{4+x^2}$の二通りの置換で求めよ。また二つの変数の間に$t=2\e^u$が成り立つことを確かめよ。
\end{enumerate}

\smallskip\noindent 解答は\pageref{解答:第10章}ページ。

\chapter{近似と広義積分}\label{章:近似と広義積分}

これまでは，積分を計算できる形へ変えることを中心に考えてきた。ところが，原始関数を知っている関数だけが積分の対象になるわけではない。また，原始関数が分かっていても，その値を小数で求めるには別の計算が必要になる。この章では，積分を「正確な値を表す式」として使うだけでなく，「近似したときに残る量を表す式」として使う。

もう一つ，定積分の定義で有限にしていた区間を広げる。上端をどこまでも遠ざけたり，被積分関数が限りなく大きくなる端点へ近づけたりするとき，途中までの積分が一定の値に近づくとは限らない。そこで必要になるのは，先に式を代入することではなく，どの極限を調べているかを明確にすることである。近似の誤差を評価することと，広義積分が収束するかを判定することは，どちらも積分の大小関係によって支えられている。

\section{有限の多項式で関数を表す}\label{節:有限テイラー}

\subsection{接線の次に何を加えるか}

微分係数による近似は
\[
 f(x)\simeq f(a)+f'(a)(x-a)
\]
という形であった。右辺は，点$a$で関数の値と傾きを合わせた一次式である。さらに二階微分の値も合わせたいなら，$(x-a)^2$の項を加える。$c(x-a)^2$を二回微分すると$2c$になるので，その係数は$f''(a)$そのものではなく，$f''(a)$の半分になる。同じ考えを続けると，$n$回までの微分の値を合わせた多項式は
\[
 P_n(x)=\sum_{k=0}^{n}\frac{f^{(k)}(a)}{k!}(x-a)^k
\]
となる。ここで$f^{(0)}=f$，$0!=1$とする。

この多項式が関数にどの程度近いかは，微分の値を合わせただけでは分からない。必要なのは，差$f(x)-P_n(x)$を調べる式である。まず微積分の基本定理から
\[
 f(x)=f(a)+\int_a^x f'(t)\dd t
\]
と書ける。積分の中で$f'(t)$を$f'(a)$に置き換えれば接線の式が現れるが，その置き換えで失われる分も残しておきたい。そのために，次の部分積分を行う。
\[
\begin{aligned}
 \int_a^x (x-t)f''(t)\dd t
 &=\bigl[(x-t)f'(t)\bigr]_a^x
   +\int_a^x f'(t)\dd t\\
 &=-(x-a)f'(a)+f(x)-f(a).
\end{aligned}
\]
したがって
\[
 f(x)=f(a)+f'(a)(x-a)
      +\int_a^x(x-t)f''(t)\dd t.
\]
最後の積分は，接線近似を正確な等式にするための項である。二階微分が小さい区間ではこの積分も小さくなる，という関係がここに現れている。

\begin{theorem}{有限テイラー公式と積分剰余}
$n$を$0$以上の整数とし，$f$は$a$と$x$を含む区間で$n+1$回連続微分可能であるとする。このとき
\[
 f(x)=\sum_{k=0}^{n}\frac{f^{(k)}(a)}{k!}(x-a)^k+R_n(x),
 \qquad
 R_n(x)=\frac{1}{n!}\int_a^x(x-t)^n f^{(n+1)}(t)\dd t
\]
が成り立つ。$R_n(x)$を剰余項という。
\end{theorem}

公式の仕組みを確かめる。$n=0$は微積分の基本定理そのものである。$n\geq1$のとき，剰余項を部分積分すると
\[
\begin{aligned}
 R_n(x)
 &=\left[\frac{(x-t)^n}{n!}f^{(n)}(t)\right]_a^x
   +\frac{1}{(n-1)!}\int_a^x(x-t)^{n-1}f^{(n)}(t)\dd t\\
 &=-\frac{f^{(n)}(a)}{n!}(x-a)^n+R_{n-1}(x).
\end{aligned}
\]
つまり，一つ前の剰余項から$n$次の項を取り出した残りが$R_n$である。この関係を順に使えば，公式が得られる。$x<a$でも，向きのある定積分の部分積分として同じ計算が成り立つ。

この導出で部分積分を選んだ理由も見ておこう。$(x-t)^n$を$t$で微分すれば次数が一つ下がり，$f^{(n+1)}(t)$を積分すれば微分の回数が一つ減る。その結果，一つ前の剰余項と結び付く。形式だけを覚えるよりも，この二つが同時に一段ずつ戻ることを覚えておくとよい。

\subsection{誤差の大きさと向き}

$a$と$x$の間で$|f^{(n+1)}(t)|\leq M$なら，積分の絶対値の評価から
\[
 |R_n(x)|\leq\frac{M}{n!}\int_{\min(a,x)}^{\max(a,x)}|x-t|^n\dd t
 =\frac{M|x-a|^{n+1}}{(n+1)!}
\]
を得る。多項式をどこまで計算すれば十分かは，この上限と，求められた精度を比較して決められる。$x$が$a$に近いほど高い次数の項が小さくなるだけでなく，最後に残る項も具体的な不等式で押さえられる。

なお，$x>a$では$(x-t)^n\geq0$である。したがって$f^{(n+1)}$の符号が分かれば，剰余項の符号も分かる。誤差の絶対値だけでなく，近似値が真の値より大きいか小さいかを判断できることがある。

\begin{example}{対数を三次式で近似する}
$0\leq x\leq\frac12$の範囲で$\log(1+x)$を三次式で近似し，誤差を評価せよ。
\end{example}

中心を$a=0$に選ぶ。$\log(1+x)$は$x=0$で値が$0$になり，各階の微分の値も簡単だからである。微分すると
\[
 f'(x)=\frac{1}{1+x},\quad
 f''(x)=-\frac{1}{(1+x)^2},\quad
 f'''(x)=\frac{2}{(1+x)^3},\quad
 f^{(4)}(x)=-\frac{6}{(1+x)^4}
\]
なので
\[
 \log(1+x)=x-\frac{x^2}{2}+\frac{x^3}{3}+R_3(x).
\]
$0\leq t\leq x$では$-6\leq f^{(4)}(t)<0$である。したがって
\[
 -\frac{x^4}{4}\leq R_3(x)\leq0,
\]
すなわち
\[
 x-\frac{x^2}{2}+\frac{x^3}{3}-\frac{x^4}{4}
 \leq\log(1+x)
 \leq x-\frac{x^2}{2}+\frac{x^3}{3}
\]
となる。例えば$x=\frac15$なら
\[
 \frac{137}{750}-\frac{1}{2500}
 \leq\log\frac65\leq\frac{137}{750}.
\]
三次式の値を計算してから，残りの最大値を差し引けば，真の値を含む区間が得られる。

\begin{example}{原始関数の公式を使わずに積分を近似する}
$A=\int_0^1\e^{-x^2}\dd x$を有理数で上下から評価せよ。
\end{example}

$\e^{-x^2}$を直接何回も微分すると，多項式と指数関数の積が増えていく。ここでは，まず一変数の関数$\e^{-u}$を$u=0$で展開し，最後に$u=x^2$を代入する。$0\leq u\leq1$に対して
\[
 \e^{-u}=1-u+\frac{u^2}{2}-\frac{u^3}{6}+r_3(u).
\]
四階微分は$\e^{-u}$であり，この区間で$0<\e^{-u}\leq1$だから
\[
 0\leq r_3(u)\leq\frac{u^4}{24}.
\]
$u=x^2$を代入して積分すると
\[
\begin{aligned}
 A&=\int_0^1\left(1-x^2+\frac{x^4}{2}-\frac{x^6}{6}\right)\dd x
       +\int_0^1r_3(x^2)\dd x\\
  &=1-\frac13+\frac1{10}-\frac1{42}
       +\int_0^1r_3(x^2)\dd x.
\end{aligned}
\]
ここで
\[
 0\leq\int_0^1r_3(x^2)\dd x
 \leq\frac1{24}\int_0^1x^8\dd x=\frac1{216}
\]
なので
\[
 \boxed{\frac{26}{35}\leq A\leq\frac{26}{35}+\frac1{216}}.
\]
積分区間全体で成立する誤差評価を用意したので，その不等式を積分できた。何点かで近似がよいことを確かめるだけでは，この結論にはならない。

\subsection{有限の公式と無限級数を分ける}

有限テイラー公式は，剰余項を含めた等式である。これに対して
\[
 f(x)=\sum_{k=0}^{\infty}\frac{f^{(k)}(a)}{k!}(x-a)^k
\]
と書くには，$n\to\infty$のとき$R_n(x)\to0$になることを確かめる必要がある。有限の段階で正しい公式を何度でも書けることと，その極限で剰余項が消えることは別の主張である。

この違いは，等比数列の和にすでに現れている。$x\ne1$に対して
\[
 \frac1{1-x}=1+x+\cdots+x^n+\frac{x^{n+1}}{1-x}
\]
は有限の恒等式である。$|x|<1$なら最後の項は$0$に近づくので，無限級数の和も$\frac1{1-x}$になる。しかし$x=2$なら，最後の項は$-2^{n+1}$であり，小さくならない。有限の恒等式は正しいままであっても，剰余項を捨てることはできない。

テイラー公式で$|f^{(n+1)}(t)|\leq M_n$と評価したときも，同じことを調べる。$|x-a|^{n+1}$だけを見るのではなく，
\[
 \frac{M_n|x-a|^{n+1}}{(n+1)!}
\]
全体が$0$へ近づくかを調べなければならない。微分を繰り返すにつれて$M_n$が大きくなる場合もあるからである。

例えば$0\leq x\leq1$での$\log(1+x)$については，$(n+1)$階微分の絶対値が$n!$以下なので
\[
 |R_n(x)|\leq\frac{x^{n+1}}{n+1}\leq\frac1{n+1}.
\]
よって，この区間では剰余項が$0$に近づき，$x=1$も含めてテイラー級数が元の関数を表すことを証明できる。

\begin{note}
何回でも微分できる関数が，必ずテイラー級数と一致するわけでもない。例えば
\[
 h(x)=
 \begin{cases}
  \e^{-\frac1{x^2}}&(x\ne0),\\
  0&(x=0)
 \end{cases}
\]
は，$x=0$でのすべての階の微分係数が$0$になる。それでも$x\ne0$では$h(x)>0$である。そのため，$0$を中心に作ったテイラー多項式はすべて$0$であり，元の関数と一致しない。

この性質の理由は，指数関数の減少が任意のべきより速いことにある。$z>0$では，有限テイラー公式の剰余項の正値性から$\e^z\geq\frac{z^N}{N!}$である。$N$を十分大きく選べば，任意の固定した$m$について
\[
 \lim_{x\to0}|x|^{-m}\e^{-\frac1{x^2}}=0
\]
が従う。$x\ne0$で$h$を繰り返し微分すると，$x$の負の整数べきの有限和に$\e^{-\frac1{x^2}}$を掛けた形になる。この極限を，各階の微分係数の定義に順に使うと，$x=0$でどの階の微分係数も$0$になることが分かる。無限に滑らかであることと，テイラー級数で表せることは区別されるのである。
\end{note}

\section{台形公式とシンプソン公式}\label{節:数値積分}

\subsection{曲線を直線に置き換える}

定積分を近似する最初の方法は，区間を細かくして長方形の面積を足すことであった。両端の関数値を使えば，長方形の代わりに台形を使える。区間$[a,b]$で，両端の点を通る直線は
\[
 L(x)=\frac{b-x}{b-a}f(a)+\frac{x-a}{b-a}f(b)
\]
である。この直線の下の符号付き面積は
\[
 T=\int_a^bL(x)\dd x=\frac{b-a}{2}\{f(a)+f(b)\}
\]
となる。これが一つの区間に対する台形公式である。

図\ref{図:台形近似}では，曲線が下に凸なので弦が曲線より上にある。この場合，台形公式の値は積分の値より大きくなる。ただし，関数によって凸凹が変わる場合もある。そこで，図だけでなく二階微分による誤差の式を求めよう。以下の台形公式の誤差評価では，$f$ は $[a,b]$ で二回連続微分可能であるとする。

\begin{figure}[htbp]
\centering
\begin{tikzpicture}[x=1.7cm,y=1.4cm]
 \fill[blue!10] (0.5,0)--(2.5,0)--(2.5,1.9625)--(0.5,0.4625)--cycle;
 \draw[-{Latex}] (-0.1,0)--(2.9,0) node[right] {$x$};
 \draw[-{Latex}] (0,-0.1)--(0,2.35) node[above] {$y$};
 \draw[dashed] (0.5,0)--(0.5,0.4625);
 \draw[dashed] (2.5,0)--(2.5,1.9625);
 \draw[thick,orange!80!black] (0.5,0.4625)--(2.5,1.9625);
 \draw[thick,blue!70!black,domain=0.15:2.65,samples=80] plot (\x,{0.4+0.25*\x*\x});
 \fill (0.5,0.4625) circle (1.5pt);
 \fill (2.5,1.9625) circle (1.5pt);
 \node[below] at (0.5,0) {$a$};
 \node[below] at (2.5,0) {$b$};
 \node[above left] at (1.55,1.25) {$L(x)$};
 \node[right] at (2.05,1.18) {$f(x)$};
\end{tikzpicture}
\caption{端点を結ぶ直線の積分が台形公式になる。下に凸の曲線では，弦と曲線の間の面積だけ過大になる。}
\label{図:台形近似}
\end{figure}

$w(x)=(x-a)(b-x)$と置く。端点で$w$が$0$となり，$w''(x)=-2$となることを利用して，二回部分積分する。
\[
\begin{aligned}
 \int_a^b w(x)f''(x)\dd x
 &=-\int_a^bw'(x)f'(x)\dd x\\
 &=-[w'(x)f(x)]_a^b+\int_a^bw''(x)f(x)\dd x\\
 &=(b-a)\{f(a)+f(b)\}-2\int_a^bf(x)\dd x.
\end{aligned}
\]
したがって，誤差を「積分の値から近似値を引いたもの」と決めれば
\[
 \boxed{\int_a^bf(x)\dd x-T
 =-\frac12\int_a^b(x-a)(b-x)f''(x)\dd x}.
\]
$f''\geq0$なら右辺は$0$以下になるので，図で見た大小関係が式からも分かる。また$|f''(x)|\leq M_2$なら
\[
 \left|\int_a^bf(x)\dd x-T\right|
 \leq\frac{M_2}{2}\int_a^b(x-a)(b-x)\dd x
 =\frac{M_2(b-a)^3}{12}.
\]

区間全体を$N$等分し，$h=\frac{b-a}{N}$，$x_j=a+jh$とすると，各小区間の台形を足した値は
\[
 T_N=h\left\{\frac{f(a)+f(b)}2+\sum_{j=1}^{N-1}f(x_j)\right\}.
\]
内部の点は隣り合う二つの台形で半分ずつ使われるため，係数が$1$になる。各小区間の誤差は$\frac{M_2h^3}{12}$以下なので，合計では
\[
 \left|\int_a^bf(x)\dd x-T_N\right|
 \leq\frac{M_2(b-a)h^2}{12}
 =\frac{M_2(b-a)^3}{12N^2}.
\]
分割数を二倍にすれば，この誤差上限は四分の一になる。

\subsection{三点を放物線で結ぶ}

次に，区間の両端に加えて中点の値も使う。中点を$m$，半分の幅を$h$とし，区間を$[m-h,m+h]$と書く。$y=x-m$によって中心を$0$へ移して，三点を通る二次式を
\[
 q(y)=A+By+Cy^2
\]
とする。$q(0)=f(m)$より$A=f(m)$である。また，$q(h)$と$q(-h)$を足せば$B$の項が消えるので
\[
 C=\frac{f(m-h)-2f(m)+f(m+h)}{2h^2}.
\]
積分でも$By$は奇関数として消える。したがって
\[
\begin{aligned}
 \int_{-h}^hq(y)\dd y
 &=2hA+\frac{2h^3}{3}C\\
 &=\frac h3\{f(m-h)+4f(m)+f(m+h)\}.
\end{aligned}
\]
これがシンプソン公式である。中点の係数$4$は，二次式の積分を計算した結果として決まる。

この公式は三次式についても正確である。中心を移した三次式のうち，三次と一次の項は奇関数なので，積分も両端と中点の重み付き和も$0$になる。残った二次と定数の項については，すでに正確である。この「三次式まで誤差がない」という性質により，誤差は四階微分で評価できる。

\subsection{シンプソン公式の誤差を導く}

単に公式の係数を覚えるだけでは，どれだけ細かく分ければよいかは決められない。ここでは，有限テイラー公式から誤差の定数も導く。座標を移して区間を$[-h,h]$とし，$f$は四回連続微分可能とする。$-h$での三次テイラー多項式を$P$とすると
\[
 f(x)=P(x)+\frac16\int_{-h}^x(x-t)^3f^{(4)}(t)\dd t.
\]
また，$f$の原始関数を$F$，$-h$での$F$の四次テイラー多項式を$Q$とする。$F'=f$だから$Q'=P$であり，$F^{(5)}=f^{(4)}$である。$F$の四次の有限テイラー公式を$x=h$で使い，
\[
 F(h)-F(-h)=\int_{-h}^h f(x)\dd x,\qquad
 Q(h)-Q(-h)=\int_{-h}^h P(x)\dd x
\]
を合わせると
\[
 \int_{-h}^hf(x)\dd x
 =\int_{-h}^hP(x)\dd x
  +\frac1{24}\int_{-h}^h(h-t)^4f^{(4)}(t)\dd t
\]
を得る。この方法なら，二重積分の積分順序を入れ替える定理を新しく使う必要はない。

一方，最初の式を$x=-h,0,h$で使い，シンプソン公式の重みを掛けて足す。$P$に対する誤差は$0$なので，残るのは
\[
 \int_{-h}^hf(x)\dd x
 -\frac h3\{f(-h)+4f(0)+f(h)\}
 =\int_{-h}^hK(t)f^{(4)}(t)\dd t
\]
である。ただし，係数をまとめた$K$は
\[
 K(t)=
 \begin{cases}
 -\frac{(h+t)^3(h-3t)}{72}&(-h\leq t\leq0),\\[4pt]
 -\frac{(h-t)^3(h+3t)}{72}&(0\leq t\leq h).
 \end{cases}
\]
例えば$t\geq0$では，中点$f(0)$の剰余項に$t$は現れないので，係数は
\[
 \frac{(h-t)^4}{24}-\frac h{18}(h-t)^3
 =-\frac{(h-t)^3(h+3t)}{72}
\]
となる。$t\leq0$では，これに中点からの項$-\frac{2h}{9}(-t)^3$を加え，整理すると上の式になる。

$K$は偶関数であり，全区間で$K\leq0$である。その積分の絶対値は
\[
\begin{aligned}
 \int_{-h}^h|K(t)|\dd t
 &=\frac1{36}\int_0^h(h-t)^3(h+3t)\dd t\\
 &=\frac{h^5}{36}\int_0^1(1-u)^3(1+3u)\dd u
 =\frac{h^5}{90}.
\end{aligned}
\]
したがって$|f^{(4)}|\leq M_4$なら，一組の区間について
\[
 \left|\int_{m-h}^{m+h}f(x)\dd x
 -\frac h3\{f(m-h)+4f(m)+f(m+h)\}\right|
 \leq\frac{M_4h^5}{90}
\]
となる。さらに$f^{(4)}\geq0$なら，積分の値から近似値を引いた誤差は$0$以下である。

\begin{theorem}{合成シンプソン公式と誤差}
$N$を正の偶数とし，$h=\frac{b-a}{N}$，$x_j=a+jh$とする。$f$が$[a,b]$で四回連続微分可能で，$|f^{(4)}(x)|\leq M_4$なら
\[
 S_N=\frac h3\left\{f(a)+f(b)
 +4\sum_{j=1}^{\frac N2}f(x_{2j-1})
 +2\sum_{j=1}^{\frac N2-1}f(x_{2j})\right\}
\]
に対して
\[
 \left|\int_a^bf(x)\dd x-S_N\right|
 \leq\frac{M_4(b-a)h^4}{180}
 =\frac{M_4(b-a)^5}{180N^4}
\]
が成り立つ。
\end{theorem}

小区間を二つずつ組にすると組数は$\frac N2$である。各組の誤差上限$\frac{M_4h^5}{90}$を足すことで，この評価が得られる。組どうしの境目は左右の組で一回ずつ使われるため係数が$2$，各組の中点は係数が$4$になる。分割数を二倍にすると誤差上限は十六分の一になるが，これは四階微分が区間全体で有界という条件の下での結論である。端点で微分が発散する関数に，そのままこの上限を使うことはできない。

\begin{example}{同じ点の値から二つの近似を作る}
$\int_0^1\frac1{1+x}\dd x$について，四等分した台形公式とシンプソン公式を計算し，誤差を評価せよ。
\end{example}

分点は$0,\frac14,\frac12,\frac34,1$で，関数値は順に$1,\frac45,\frac23,\frac47,\frac12$である。したがって
\[
 T_4=\frac14\left(\frac34+\frac45+\frac23+\frac47\right)
 =\frac{1171}{1680},
\]
\[
 S_4=\frac1{12}\left\{\frac32+4\left(\frac45+\frac47\right)+\frac43\right\}
 =\frac{1747}{2520}.
\]
二階微分は$\frac2{(1+x)^3}$，四階微分は$\frac{24}{(1+x)^5}$で，いずれも正である。そのため，どちらの近似値も真の値以上になる。上限$M_2=2$，$M_4=24$を使えば
\[
 T_4-\frac1{96}\leq\int_0^1\frac1{1+x}\dd x\leq T_4,
\]
\[
 S_4-\frac1{1920}\leq\int_0^1\frac1{1+x}\dd x\leq S_4.
\]
この積分は$\log2$に等しいが，その小数値を知らなくても近似の保証は作れる。同じ五つの関数値を使っていても，曲線をどの次数の式で置き換えるかによって誤差の評価が変わる。

\section{商の極限とロピタルの定理}\label{節:ロピタル}

\subsection{微分した商を調べてよい条件}

極限を計算していると，分子も分母も$0$へ近づく商が現れる。その場合，分子と分母がどれだけ速く$0$へ近づくかを比べる必要がある。微分係数がその速さを表すことから，分子と分母の微分を比べる方法が考えられる。ただし，元の商の形を確かめず，機械的に分子と分母を微分してはいけない。

\begin{theorem}{ロピタルの定理の一側極限版}
$x$が区間の一端$c$へ近づく一側極限を考える。$c$は有限でも$+\infty,-\infty$でもよい。端点を除くその区間で$f,g$は微分可能であり，$g'(x)\ne0$とする。さらに，次のいずれかを満たすとする。
\begin{enumerate}[label=(\roman*),leftmargin=*]
 \item $f(x)\to0$かつ$g(x)\to0$である。
 \item $|f(x)|\to\infty$かつ$|g(x)|\to\infty$である。
\end{enumerate}
このとき
\[
 \lim_{x\to c}\frac{f'(x)}{g'(x)}=L
\]
が存在すれば
\[
 \lim_{x\to c}\frac{f(x)}{g(x)}=L
\]
である。ここで$L$は実数のほか，$+\infty,-\infty$でもよい。元の商は端点の十分近くで定義されている。
\end{theorem}

両側極限を求めるときは，左右の一側極限にそれぞれ適用し，結果が一致することを確認する。また，$\frac{f'}{g'}$をさらに微分して計算したい場合は，その新しい商についても同じ条件を改めて確認する。微分後の商が$0$と$0$の商になるとは限らない。

\subsection{コーシーの平均値の定理による根拠}

通常の平均値の定理では，関数の増分を区間の幅で割った。ここでは，二つの関数の増分を比べる。$f,g$が$[u,v]$で連続，内部で微分可能で，$g'\ne0$なら，ある$\xi\in(u,v)$が存在して
\[
 \frac{f(v)-f(u)}{g(v)-g(u)}=\frac{f'(\xi)}{g'(\xi)}
\]
となる。これをコーシーの平均値の定理という。

この式もロルの定理から導ける。関数
\[
 H(t)=\{f(v)-f(u)\}\{g(t)-g(u)\}
      -\{g(v)-g(u)\}\{f(t)-f(u)\}
\]
を作ると，$H(u)=H(v)=0$である。よって$H'(\xi)=0$となる点があり，微分して整理すると求める等式になる。なお，$g(v)=g(u)$であればロルの定理からどこかで$g'=0$となってしまうので，分母$g(v)-g(u)$は$0$でない。

まず有限の端点$c$で$0$と$0$の形になる場合を考える。端点に$f(c)=g(c)=0$を補えば連続になるので，$c$と$x$の間でコーシーの平均値の定理を使える。すると
\[
 \frac{f(x)}{g(x)}=\frac{f'(\xi_x)}{g'(\xi_x)}
\]
となり，$\xi_x$は$c$と$x$の間にある。$x\to c$なら$\xi_x\to c$だから，右辺の極限が$L$であることから結論が従う。

無限大の形では，少し見方を変える。例えば$x\to+\infty$で$|g(x)|\to\infty$，$L$が実数の場合，$\frac{f'(t)}{g'(t)}$が$L$に十分近い範囲の始点$B$を固定する。$x>B$に対して
\[
 D_x=\frac{f(x)-f(B)}{g(x)-g(B)}
\]
は，その範囲のどこかでの$\frac{f'}{g'}$に等しいので，$L$に十分近い。元の商は
\[
 \frac{f(x)}{g(x)}
 =D_x\left(1-\frac{g(B)}{g(x)}\right)+\frac{f(B)}{g(x)}
\]
と書ける。$B$を固定した後で$x\to+\infty$とすれば，固定した値を含む二つの比はそれぞれ$0$へ近づく。そのため元の商も$L$へ近づく。有限の端点で$|g|\to\infty$となる場合も，$B$を端点の十分近くに固定して同じ議論を使える。

$0$と$0$の形で端点が$+\infty$の場合は，十分大きい$x<y$の間に平均値の定理を使う。増分の比が$L$に近いことを確かめた後，$x$を固定して$y\to+\infty$とすると，その比は$\frac{f(x)}{g(x)}$へ近づく。$L=+\infty$の場合は「$L$との差が小さい」の代わりに「任意の大きな定数より大きい」を用い，$L=-\infty$では不等号を逆にする。同じ議論で拡張された極限も得られる。

\begin{example}{消える次数を比較する}
$\lim_{x\to0}\cfrac{\log(1+x)-x+\cfrac{x^2}{2}}{x^3}$を求めよ。
\end{example}

分子も分母も$0$に近づき，$x\ne0$で分母の微分$3x^2$は$0$でない。分子を微分すると
\[
 \frac1{1+x}-1+x=\frac{x^2}{1+x}
\]
とまとめられるので
\[
 \cfrac{\left(\log(1+x)-x+\cfrac{x^2}{2}\right)'}{(x^3)'}
 =\frac1{3(1+x)}\longrightarrow\frac13.
\]
よって求める極限は$\frac13$である。微分後に式を整理すると$x^2$が約分できるため，三回の微分は必要ない。

同じ極限は有限テイラー公式でも説明できる。三次の項まで書けば，分子は$\frac{x^3}{3}+R_3(x)$となる。$0$の近くでは$|R_3(x)|\leq C|x|^4$と評価できるので，$x^3$で割った剰余項は$0$へ近づく。どちらの方法も，分子の最初に残る項が三次であることを調べている。

ロピタルの定理の条件は省略できない。例えば$x\to0$で$\frac{x^2+1}{x^2}$を考えると，微分した商は$1$だが，元の商は$+\infty$へ発散する。元の分子は$0$にも無限大にも近づかず，定理の仮定を満たしていない。また$x\to+\infty$で
\[
 \frac{x+\sin x}{x}\longrightarrow1
\]
であっても，微分した商$1+\cos x$には極限がない。微分した商の極限が存在しないことから，元の商の極限が存在しないとは結論できない。

\section{端点を極限にする積分}\label{節:広義積分の定義}

\subsection{有限の途中まで積分してから極限をとる}

$\int_0^1\frac1{\sqrt{x}}\dd x$では，$x=0$で被積分関数が定義されず，その近くで値が限りなく大きくなる。それでも，$\varepsilon>0$で切り取った区間では
\[
 \int_{\varepsilon}^1\frac1{\sqrt{x}}\dd x
 =2-2\sqrt{\varepsilon}
\]
と計算できる。この値が$\varepsilon\to0+$で$2$へ近づくことを利用して，積分の意味を広げる。

\begin{definition}{広義積分}
$f$が$[a,b)$の各閉部分区間で積分可能で，端点$b$に特異性をもつとき，
\[
 \int_a^bf(x)\dd x=\lim_{c\to b-}\int_a^cf(x)\dd x
\]
と定める。ただし，右辺が有限の実数として存在するとき，この広義積分は収束するという。左端に特異性がある場合も同様に，一側極限で定める。

上端が無限大の場合は
\[
 \int_a^{\infty}f(x)\dd x=\lim_{R\to\infty}\int_a^Rf(x)\dd x
\]
と定める。下端が$-\infty$の場合も同様である。有限の極限が存在しなければ発散するという。
\end{definition}

ここで，$+\infty$へ発散する場合も，振動して極限がない場合も，「収束しない」という点では同じである。一方，発散の仕方を区別したいときには「$+\infty$へ発散する」などと具体的に述べる。

\subsection{複数の特異点は独立に調べる}

両端に特異性がある場合は，内部の一点$c$を固定して
\[
 \int_a^bf(x)\dd x
 =\int_a^cf(x)\dd x+\int_c^bf(x)\dd x
\]
と定める。ただし，右辺の二つの広義積分がそれぞれ収束することが必要である。両端を同じ速さで動かして相殺することは，この定義には含まれない。分割点$c$を別の内部の点へ変えても，その間の通常の定積分が一方に加わり，他方から引かれるだけなので，値は変わらない。

区間内部に特異点$d$があれば，$d$の左右にも分ける。また，実数全体の積分は，例えば
\[
 \int_{-\infty}^{\infty}f(x)\dd x
 =\int_{-\infty}^{0}f(x)\dd x+\int_0^{\infty}f(x)\dd x
\]
として，左右を独立に調べる。複数の特異点があるときは，それぞれに別の切り取り幅や上限を用意して考えると，何を確認すべきかが明確になる。

この違いを$\frac1x$で確かめよう。$\varepsilon>0$に対して
\[
 \int_{-1}^{-\varepsilon}\frac1x\dd x=\log\varepsilon\longrightarrow-\infty,
 \qquad
 \int_{\varepsilon}^{1}\frac1x\dd x=-\log\varepsilon\longrightarrow+\infty.
\]
したがって$\int_{-1}^1\frac1x\dd x$は広義積分として収束しない。しかし，同じ$\varepsilon$を使って足せば常に$0$になる。このような対称的な切り取り方による極限を，コーシーの主値という。
\[
 \operatorname{PV}\int_{-1}^1\frac1x\dd x
 =\lim_{\varepsilon\to0+}
 \left(\int_{-1}^{-\varepsilon}\frac1x\dd x
       +\int_{\varepsilon}^{1}\frac1x\dd x\right)=0.
\]
主値が存在しても，広義積分が収束するとは限らない。実際，左を$\varepsilon$，右を$2\varepsilon$で切れば，二つの積分の和は$-\log2$になる。切り取り方を結び付けて得られた相殺だけでは，各側の収束を示したことにならない。

無限区間でも同じ注意が要る。例えば$\frac{x}{1+x^2}$は奇関数なので，任意の$R$に対して$\int_{-R}^R\frac{x}{1+x^2}\dd x=0$である。一方，右半分の積分は$\frac12\log(1+R^2)$となり発散する。奇関数であることを使う前に，広義積分の収束を確認する必要がある。

\section{べき関数との比較と条件収束}\label{節:広義積分の判定}

\subsection{基準になる二つの積分}

無限大での収束を調べる基準は
\[
 \int_1^{\infty}\frac1{x^p}\dd x
\]
である。$p\ne1$なら，まず有限の$R$まで積分して
\[
 \int_1^R x^{-p}\dd x=\frac{R^{1-p}-1}{1-p}.
\]
$p>1$なら$R^{1-p}\to0$なので収束し，値は$\frac1{p-1}$となる。$p<1$なら$R^{1-p}\to\infty$なので発散する。$p=1$では$\log R\to\infty$であり，やはり発散する。

これに対して，有限端点$0$の近くでは
\[
 \int_{\varepsilon}^{1}x^{-p}\dd x
 =\frac{1-\varepsilon^{1-p}}{1-p}\qquad(p\ne1)
\]
である。今度は$\varepsilon\to0+$なので，$p<1$のときに収束し，値は$\frac1{1-p}$になる。$p\geq1$では発散する。

\begin{center}
\begin{tabular}{lll}
\toprule
積分 & 収束する条件 & 収束時の値\\
\midrule
$\int_1^{\infty}x^{-p}\dd x$ & $p>1$ & $\frac1{p-1}$\\[4pt]
$\int_0^1x^{-p}\dd x$ & $p<1$ & $\frac1{1-p}$\\
\bottomrule
\end{tabular}
\end{center}

同じべき関数でも，不等号の向きが逆になる。遠くで減る速さを調べる場合と，端点で大きくなる速さを調べる場合では，問題が異なるからである。覚えるときは，$\frac1x$がどちらでも境目になることと，実際の原始関数の極限を対応させたい。

\subsection{比較判定法}

$0\leq f(x)\leq g(x)$が端点の十分近くで成り立つとする。このとき，$g$の広義積分が収束すれば$f$の広義積分も収束する。逆向きには，$f$の広義積分が$+\infty$へ発散すれば，$g$の広義積分も$+\infty$へ発散する。

例えば無限区間では，$\int_A^R f(x)\dd x$は$R$とともに増加する。しかも
\[
 0\leq\int_A^Rf(x)\dd x\leq\int_A^Rg(x)\dd x
 \leq\int_A^{\infty}g(x)\dd x
\]
と一定の上限をもつ。上に有界な単調増加の量は有限の極限をもつので，$f$の積分も収束する。端点付近でだけ比較できればよいのは，それ以外の閉区間の積分は有限の一定値を加えるだけだからである。

特に$f,g>0$で
\[
 \lim\frac{f(x)}{g(x)}=c,\qquad0<c<\infty
\]
となるなら，端点の十分近くで$\frac c2g(x)\leq f(x)\leq\frac{3c}{2}g(x)$と挟める。そのため二つの積分は同時に収束するか，同時に発散する。この極限比較では，$c$が正の有限値であるという条件が大切である。比の極限が$0$の場合は，$f$が$g$より小さいことしか分からず，両者の収束が一致するとは限らない。

\begin{example}{収束を先に確かめてから置換する}
$\int_0^1\frac1{\sqrt{x(1+x)}}\dd x$の収束を確かめ，値を求めよ。
\end{example}

特異性があるのは$x=0$だけである。$0<x\leq1$で
\[
 0<\frac1{\sqrt{x(1+x)}}\leq\frac1{\sqrt{x}}
\]
となり，右辺は$0$の近くで積分可能だから，求める広義積分は収束する。

計算では，まず$\varepsilon\leq x\leq1$に限って$x=t^2$と置く。平方根の$x$と$\dd x=2t\dd t$が約分されるので
\[
 \int_{\varepsilon}^1\frac1{\sqrt{x(1+x)}}\dd x
 =2\int_{\sqrt{\varepsilon}}^1\frac1{\sqrt{1+t^2}}\dd t.
\]
右辺の被積分関数は$t=0$でも連続である。$\varepsilon\to0+$として
\[
 2\int_0^1\frac1{\sqrt{1+t^2}}\dd t
 =2\left[\log\left(t+\sqrt{1+t^2}\right)\right]_0^1
 =2\log(1+\sqrt2)
\]
を得る。置換はまず通常の定積分に行い，その後で端点の極限をとった。

\subsection{符号の変化が収束を生む場合}

関数が正負に変わると，絶対値を積分した場合と，元の関数を積分した場合とで，収束の結果が異なることがある。

\begin{definition}{絶対収束と条件収束}
広義積分$\int f(x)\dd x$について，$\int|f(x)|\dd x$が収束するとき絶対収束するという。$\int f(x)\dd x$は収束するが，$\int|f(x)|\dd x$は発散するとき，条件収束するという。
\end{definition}

絶対収束なら元の積分も収束する。例えば無限区間では，遠い二点$S<R$の間の積分が
\[
 \left|\int_S^Rf(x)\dd x\right|\leq\int_S^R|f(x)|\dd x
\]
で押さえられる。右辺は$S,R$が十分大きければいくらでも小さくなるので，途中までの積分はコーシーの収束条件を満たす。これは実数の完備性による収束判定であり，「途中までの値が互いにいくらでも近くなるなら，有限の極限をもつ」という事実を使っている。

\begin{example}{振動する積分の条件収束}
$\int_1^{\infty}\frac{\sin x}{x}\dd x$は条件収束することを示せ。
\end{example}

$\frac1x$だけでは無限大で積分可能でない。しかし，$\sin x$を積分すると有界な$-\cos x$になり，$\frac1x$を微分すると$-\frac1{x^2}$になる。この二つを利用するため，部分積分を選ぶ。
\[
\begin{aligned}
 \int_1^R\frac{\sin x}{x}\dd x
 &=\left[-\frac{\cos x}{x}\right]_1^R
   -\int_1^R\frac{\cos x}{x^2}\dd x\\
 &=\cos1-\frac{\cos R}{R}-\int_1^R\frac{\cos x}{x^2}\dd x.
\end{aligned}
\]
境界項$\frac{\cos R}{R}$は$0$へ近づく。また$\frac{|\cos x|}{x^2}\leq\frac1{x^2}$なので，最後の積分は絶対収束する。よって元の積分は収束する。

絶対値を付けた積分は別に調べる。整数$k\geq1$について，$[k\pi,(k+1)\pi]$では$x\leq(k+1)\pi$だから
\[
 \int_{k\pi}^{(k+1)\pi}\frac{|\sin x|}{x}\dd x
 \geq\frac1{(k+1)\pi}\int_{k\pi}^{(k+1)\pi}|\sin x|\dd x
 =\frac2{(k+1)\pi}.
\]
この下限を足すと調和級数の定数倍になり，発散する。したがって絶対値の積分は発散し，元の積分は条件収束である。ここでは積分の値そのものを求めていない。収束を証明することと，値を計算することは別の課題である。

\section{原始関数が存在することと式で書けること}\label{節:非初等原始関数}

$\e^{-x^2}$の原始関数を，これまでの置換や部分積分で求めようとしても，なじみのある有限個の式にはたどり着かない。だからといって，原始関数が存在しないわけではない。微積分の基本定理を使えば
\[
 F(x)=\int_0^x\e^{-t^2}\dd t
\]
はすべての実数で定義され，$F'(x)=\e^{-x^2}$となる。これ自体が一つの原始関数である。

「初等関数で表せない」というときは，数，変数，代数的演算，指数関数，対数関数，三角関数などを有限回組み合わせた通常の初等関数の範囲では表せない，という意味である。$\e^{-x^2}$には初等関数で表された原始関数がないことが知られている。ただし，その不可能性の証明には，本書の範囲を超える理論が必要である。ここではこの事実を採用し，置換をいくつか試して失敗したことを証明の代わりにはしない。

積分で定義した$F$についても，多くのことが分かる。被積分関数が正なので$F$は単調増加である。被積分関数が偶関数なので$F$は奇関数である。さらに
\[
 F''(x)=-2x\e^{-x^2}
\]
から，$x>0$では上に凸，$x<0$では下に凸である。有限区間での値はテイラー公式や数値積分で評価できるし，$x\geq1$では$\e^{-x^2}\leq\e^{-x}$なので，$\int_0^{\infty}\e^{-x^2}\dd x$の収束も比較によって分かる。

積分を含む定義は，計算を途中で諦めた記号ではない。微分，対称性，大小関係，近似，極限を調べるための出発点になる。原始関数の表示が難しいときには，求めたいものが「不定積分の式」なのか，「特定の区間での値」なのか，「誤差付きの近似」なのか，「収束するという事実」なのかを分けて考える。その区別により，使う道具も変わってくる。

\section*{章末問題}
\begin{enumerate}[label=\textbf{\thechapter.\arabic*},leftmargin=*]
 \item\label{問:十一対数近似}
 $P(x)=x-\frac{x^2}{2}+\frac{x^3}{3}$とする。$0\leq x\leq\frac12$で
 $0\leq P(x)-\log(1+x)\leq\frac{x^4}{4}$を示し，$\log\frac{11}{10}$を幅$\frac1{40000}$以下の有理数の区間に挟め。
 \item\label{問:十一指数積分近似}
 $\e^{-u}$の二次テイラー多項式を使い，$\int_0^{\frac12}\e^{-x^2}\dd x$の近似値と誤差の上限を求めよ。近似値が真の値の上下どちらにあるかも述べよ。
 \item\label{問:十一数値積分}
 $\int_0^2x^4\dd x$を四等分した台形公式とシンプソン公式で近似せよ。真の値との差を計算し，二階微分，四階微分による誤差上限と比較せよ。
 \item\label{問:十一ロピタル}
 次の極限をロピタルの定理で求め，適用条件を述べよ。
 \[
 \text{(a)}\ \lim_{x\to0}\frac{\e^x-1-x}{x^2},
 \qquad
 \text{(b)}\ \lim_{x\to\infty}\frac{\log x}{x^{\alpha}}\quad(\alpha>0).
 \]
 \item\label{問:十一両端比較}
 実数$p,q$に対し，$\int_0^1\frac1{x^p(1-x)^q}\dd x$が収束するための必要十分条件を求めよ。両端を独立に調べること。
 \item\label{問:十一対数尾部}
 実数$q$に対し，$\int_{\e}^{\infty}\frac1{x(\log x)^q}\dd x$の収束条件を求め，収束する場合の値を求めよ。
 \item\label{問:十一主値区別}
 $\int_{-1}^1\frac1x\dd x$と$\int_{-1}^1\frac{1-\cos x}{x}\dd x$について，広義積分としての収束と値を調べよ。前者の対称的な主値も求め，後者との違いを説明せよ。
 \item\label{問:十一積分定義関数}
 $F(x)=\int_0^x\e^{-t^2}\dd t$とする。$F$の奇偶，増減，凹凸を調べよ。また$x\geq0$で$x-\frac{x^3}{3}\leq F(x)\leq x$を示し，$\lim_{x\to\infty}F(x)$が有限であることを証明せよ。
\end{enumerate}

\smallskip\noindent 解答は\pageref{解答:第11章}ページ。

\chapter{総合演習}\label{章:総合演習}

ここでは，問題のすぐ前に使う公式を置かず，式を見て方針を選ぶ練習をする。計算を始める前に，求めるものが原始関数，定積分の値，関数の性質，極限，あるいは誤差付きの近似のどれであるかを確かめたい。同じ積分記号でも，それによって使いやすい手法は変わる。

基礎問題は一つか二つの基本操作を確実につなぐ問題，標準問題は区間や式の見方を選ぶ問題である。発展問題では複数の章の内容を使うが，新しい一般定理を調べることは求めない。微分した式を元へ戻す，平方の積分を考える，端点の極限を独立に調べるなど，本書で扱った操作から進められる。ヒントは問題の後にまとめ，解答は巻末に置いた。最初の一手を自分で考えてから，必要な番号だけを参照してほしい。

\section{基礎問題}\label{節:総合基礎}

\begin{enumerate}[label=\textbf{\thechapter.\arabic*},leftmargin=*]
 \item\label{問:総合一割り算}
 次の定積分を求めよ。
 \[
  \int_0^1\frac{x^3+x^2+2x+1}{1+x^2}\dd x.
 \]
 \item\label{問:総合二置換}
 次の定積分を求め，置換後の区間も示せ。
 \[
  \int_0^1x\sqrt{1+x^2}\dd x.
 \]
 \item\label{問:総合三部分積分}
 次の定積分を求めよ。
 \[
  \int_1^{\e}x\log x\dd x.
 \]
 \item\label{問:総合四動く端点}
 $x>0$に対して$F(x)=\int_x^{x^2}\sqrt{1+t^2}\dd t$とする。
 $F'(x)$を求め，曲線$y=F(x)$の$x=1$における接線を求めよ。
 \item\label{問:総合五相補重み}
 $a>0$とする。次の積分を$a$で表せ。
 \[
  \int_{-a}^{a}\frac{1+x^2}{1+\e^x}\dd x.
 \]
 \item\label{問:総合六面積体積}
 曲線$y=x^2$と直線$y=x$で囲まれる領域の面積を求めよ。
 また，この領域を$x$軸のまわりに一回転してできる立体の体積を求めよ。
\end{enumerate}

\section{標準問題}\label{節:総合標準}

\begin{enumerate}[label=\textbf{\thechapter.\arabic*},leftmargin=*,start=7]
 \item\label{問:総合七区間の最適化}
 実数$a$に対して
 \[
  H(a)=\int_a^{a+1}(x^2-2x)^2\dd x
 \]
 とする。$H(a)$の最小値と，それを与える$a$をすべて求めよ。
 \item\label{問:総合八置換と部分積分}
 次の定積分を求めよ。
 \[
  \int_0^1x^3\log(1+x^2)\dd x.
 \]
 \item\label{問:総合九区分求積}
 次の極限を求めよ。
 \[
  \lim_{n\to\infty}\sum_{k=1}^{n}\frac{k^2}{n(n^2+k^2)}.
 \]
 \item\label{問:総合十軸を横切る回転}
 $0\leqq x\leqq1$，$x-1\leqq y\leqq x$で定まる領域を$x$軸のまわりに一回転してできる立体の体積を求めよ。
 各$x$での回転断面が円板になるか円環になるかを説明し，半径が切り替わる点を求めること。
 \item\label{問:総合十一漸化式}
 非負整数$n$に対して$K_n=\int_0^1\frac{x^n}{1+x}\dd x$とする。
 $K_0$と，$n\geqq1$で成り立つ漸化式を求め，$K_3$を計算せよ。
 さらに$K_n\to0$を示し，それを使って
 \[
  \lim_{n\to\infty}\sum_{k=1}^{n}\frac{(-1)^{k-1}}{k}=\log2
 \]
 を示せ。
 \item\label{問:総合十二両端判定}
 実数$\alpha$に対して，広義積分
 \[
  \int_0^{\infty}\frac{x^\alpha}{1+x^2}\dd x
 \]
 の収束条件を求めよ。また$\alpha=0$の場合の値を求めよ。
\end{enumerate}

\section{発展問題}\label{節:総合発展}

\begin{enumerate}[label=\textbf{\thechapter.\arabic*},leftmargin=*,start=13]
 \item\label{問:総合十三積分方程式}
 $[0,1]$で連続な関数$f$が，すべての$x\in[0,1]$について
 \[
  f(x)=1+\int_0^x(x-t)f(t)\dd t
 \]
 を満たすとする。$f$を求め，求めた関数が確かにこの等式を満たすことを確かめよ。
 微分方程式の一般的な解法は用いず，積の微分と基本定理から求めてよい。
 \item\label{問:総合十四二つの制約}
 $[0,1]$で連続な実数値関数$f$が
 \[
  \int_0^1f(x)\dd x=1,\qquad
  \int_0^1xf(x)\dd x=0
 \]
 を満たすとする。$\int_0^1f(x)^2\dd x$の最小値と，それを与える関数を求めよ。
 \item\label{問:総合十五媒介曲線}
 $0\leqq t\leqq\pi$で
 \[
  x=t-\frac12\sin2t,\qquad y=\sin^2t
 \]
 と表される曲線を考える。この曲線と$x$軸で囲まれる領域の面積，曲線の長さ，領域を$x$軸のまわりに一回転してできる立体の体積を求めよ。
 面積と体積の積分で，同じ部分を重複して数えていないことも確かめること。
 \item\label{問:総合十六近似と条件収束}
 次の二つの問いに，原始関数の初等的な表示を求めずに答えよ。
 \begin{enumerate}[label=(\arabic*),leftmargin=*]
  \item 有限テイラー公式を使い，$A=\int_0^1\sin(x^2)\dd x$を幅$\frac1{10000}$以下の有理数の区間に挟め。
  \item 広義積分$\int_1^{\infty}\sin(x^2)\dd x$が条件収束することを示せ。
 \end{enumerate}
 \item\label{問:総合十七ウォリス}
 非負整数$n$に対して
 \[
  C_n=\int_0^1(1-x^2)^n\dd x
 \]
 とする。$C_0$と$C_n$の漸化式を求め，$C_n$を有限の積で表せ。
 次にウォリスの積分との関係を使って，
 \[
  \lim_{n\to\infty}\sqrt n\,C_n,
  \qquad
  \lim_{n\to\infty}n^{\frac32}(C_n-C_{n+1})
 \]
 を求めよ。
 \item\label{問:総合十八相補と半角}
 $n$を正の整数とする。次の広義積分を，$n$を含む有限の積で表せ。
 \[
  J_n=\int_{-\pi}^{\pi}
  \frac{\sin^{2n}x}{(1+\e^x)\sqrt{1+\cos x}}\dd x.
 \]
 最初に両端での収束を確かめること。対称性，半角公式，置換積分，部分積分を組み合わせれば，ベータ関数やガンマ関数の公式を使わずに求められる。
\end{enumerate}

\clearpage
\section{ヒント}\label{節:総合ヒント}

以下では，計算の結論ではなく，最初の変形や見るべき条件を示す。解き直すときには，その変形を選んだ理由も言葉にしてみよう。

\begin{enumerate}[label=\textbf{\thechapter.\arabic*},leftmargin=*]
 \item 分子を$1+x^2$で割ると，余りの項が分母の微分と結び付く。
 \item 根号の中を$u=1+x^2$と置く。$x\dd x$をそのまま$\dd u$で表す。
 \item $\log x$を微分する側に選ぶと，残る積分は整式になる。
 \item $G(u)=\int_0^u\sqrt{1+t^2}\dd t$と置いて，$F(x)=G(x^2)-G(x)$と書く。
 \item $x=-t$で書き直した積分と元の積分を足す。$1+x^2$は偶関数である。
 \item 交点を求めた後で，上下関係と，回転軸からの二つの距離を調べる。
 \item 端点を微分すると$H'(a)$が得られる。値の計算には$y=a-\frac12$と置き，積分区間の中心から測ると整理しやすい。
 \item $u=1+x^2$と置けば，$x^3\dd x=\frac12(u-1)\dd u$になる。その後で部分積分する。
 \item 一項を$\frac1n$と$\frac{k}{n}$だけを使って表す。積分に直したら，分子を分母と定数の差にする。
 \item 区間$[x-1,x]$は常に$0$を含む。回転後の半径は，両端の軸からの距離の大きい方である。
 \item $K_n+K_{n-1}$では分母が消える。極限には$0\leqq K_n\leqq\int_0^1x^n\dd x$を使う。
 \item $0$付近では$x^\alpha$，無限大では$x^{\alpha-2}$と比較し，二つの条件を同時に満たす範囲を求める。
 \item 積分を$x\int_0^xf(t)\dd t-\int_0^xtf(t)\dd t$に分けて微分する。$f''=f$が得られたら，$\e^{-x}(f+f')$の微分を考える。
 \item まず二つの制約を満たす一次関数$\ell(x)$を求める。$g=f-\ell$と置いたとき，$\int_0^1\ell g\dd x$がどうなるかを調べる。
 \item $x'(t)=2\sin^2t$を使う。弧長では$\sqrt{x'(t)^2+y'(t)^2}$を整理してから積分する。
 \item (1) $\sin u$の六次テイラー多項式では，偶数次の係数が$0$になる。七階微分の符号と大きさを調べる。(2) $u=x^2$と置く。収束には部分積分，絶対値の発散には長さ$\pi$の区間ごとの下からの評価を使う。
 \item $x=\sin t$によって$C_n$を$\cos$の奇数べきの積分にする。隣り合うウォリス積分の比と積を組み合わせる。
 \item 指数関数を含む重みは，$x$を$-x$にすると相補的になる。その後，$x=2t$，さらに$u=\sin t$と置く。残る積分の$(1-u^2)$の指数は整数なので，部分積分で一つずつ下げられる。
\end{enumerate}

\smallskip\noindent 解答は\pageref{解答:第12章}ページ。

\appendix
\renewcommand{\chaptermark}[1]{\markboth{付録\thechapter\quad #1}{}}
\chapter{章末問題・総合演習の解答}\label{章:解答}
\section{第1章の解答}\label{解答:第1章}
\paragraph{1.1}
平均の速度は，位置の変化を時間の変化で割って
\[
 \frac{s(3+h)-s(3)}h
 =\frac{(3+h)^2-9}{h}
 =\frac{6h+h^2}{h}=6+h \qquad(h\ne0)
\]
である。$h=0.2$ のとき $6.2$，$h=-0.2$ のとき $5.8$ となる。$h$ を正負どちらから $0$ に近づけても $6+h$ は $6$ に近づくので，瞬間の速度は毎秒 $6$ メートルと予想できる。ここで扱ったように，零に近い負の $h$ を取ったときも分子と分母がともに負となるため，同じ平均変化率の式を使える。

\paragraph{1.2}
台形の平行な二辺の長さは $1$ と $5$，その間の距離は $2$ だから，面積は
\[
 \frac{1+5}{2}\times2=6
\]
である。$n$ 等分した幅は $\frac2n$，左端の座標は $\frac{2k}{n}$（$k=0,\ldots,n-1$）なので，
\begin{align*}
 L_n
 &=\frac2n\sum_{k=0}^{n-1}\left(1+\frac{4k}{n}\right)\\
 &=\frac2n\left\{n+\frac4n\frac{n(n-1)}2\right\}
 =6-\frac4n.
\end{align*}
右端については添字を $1$ から $n$ までとし，
\begin{align*}
 U_n
 &=\frac2n\sum_{k=1}^{n}\left(1+\frac{4k}{n}\right)\\
 &=\frac2n\left\{n+\frac4n\frac{n(n+1)}2\right\}
 =6+\frac4n.
\end{align*}
$\frac4n$ が $0$ に近づくので，どちらも $6$ に近づく。

\paragraph{1.3}
各小区間の中央は $\frac18,\frac38,\frac58,\frac78$ である。幅はすべて $\frac14$ だから，求める和は
\[
 M_4=\frac14\left\{\left(\frac18\right)^2
 +\left(\frac38\right)^2+\left(\frac58\right)^2
 +\left(\frac78\right)^2\right\}
 =\frac{1+9+25+49}{256}=\frac{21}{64}.
\]
一方，
\[
 L_4=\frac14\left(0+\frac1{16}+\frac4{16}+\frac9{16}\right)=\frac7{32},
 \qquad
 U_4=\frac14\left(\frac1{16}+\frac4{16}+\frac9{16}+1\right)=\frac{15}{32}.
\]
よって $L_4=\frac{14}{64}<M_4=\frac{21}{64}<\frac{30}{64}=U_4$ である。中央の高さは各区間の左端と右端の高さの間にあるため，この大小関係は計算前にも分かる。ただし，中央の和が真の面積より大きいか小さいかは，この大小関係だけでは決まらない。

\paragraph{1.4}
割線の傾きは
\[
 \frac{(1+h)^3-1}{h}
 =\frac{3h+3h^2+h^3}{h}
 =3+3h+h^2.
\]
最後の約分には $h\ne0$ を使った。$h$ を $0$ に近づけると $3h,h^2$ はともに $0$ に近づくので，接線の傾きは $3$ と予想できる。

\paragraph{1.5}
$4$ 等分した幅は $\frac12$ であるから，
\begin{align*}
 L_4&=\frac12\left\{0^2+\left(\frac12\right)^2+1^2+\left(\frac32\right)^2\right\}=\frac74,\\
 U_4&=\frac12\left\{\left(\frac12\right)^2+1^2+\left(\frac32\right)^2+2^2\right\}=\frac{15}{4}.
\end{align*}
したがって $\frac74\le S\le\frac{15}{4}$ である。一般の $n$ 等分では，幅は $\frac2n$，高さは $\left(\frac{2k}{n}\right)^2$ となる。幅が $2$ 倍，高さが $4$ 倍なので，長方形の面積もその和も $8$ 倍になる。$0\le x\le1$ の場合の和が $\frac13$ に近づくことから，$S=\frac83$ と予想できる。横幅を $2$ 倍にしたから面積も $2$ 倍，とはならない。高さの変化も含める必要がある。

\paragraph{1.6}
位置の変化は速度に符号を付けたまま時間を掛けて足すので，
\[
 3\times4+(-2)\times3=6
\]
メートルである。道のりは移動した長さを足すので，
\[
 3\times4+2\times3=18
\]
メートルである。後半の速度は負であり，前半に進んだ向きとは逆に $6$ メートル戻っている。位置の変化ではこの分を引くが，道のりではこの移動も正の長さとして足す。

\section{第2章の解答}\label{解答:第2章}
\paragraph{2.1}
(1) 分子・分母を $n^2$ で割ると，
\[
 \frac{2n^2+3}{5n^2-n}
 =\cfrac{2+\cfrac3{n^2}}{5-\cfrac1n}
 \longrightarrow\frac25.
\]
分母の極限は $5\ne0$ なので，商の極限法則を適用できる。

(2) 根号を含む差を有理化する。
\[
 \sqrt{n^2+2n}-n
 =\frac{2n}{\sqrt{n^2+2n}+n}
 =\cfrac2{\sqrt{1+\cfrac2n}+1}
 \longrightarrow1.
\]
二つの項がそれぞれ大きくなることだけでは差の極限は分からないが，この形なら分母が $2$ に近づくことが分かる。

漸化式では，まず $1\le a_n<2$ を帰納的に示す。$a_1=1$ はこの条件を満たす。$1\le a_n<2$ なら，
\[
 1<\sqrt3\le a_{n+1}=\sqrt{2+a_n}<2
\]
だから，次の項も条件を満たす。さらに，
\[
 a_{n+1}^2-a_n^2
 =2+a_n-a_n^2=(2-a_n)(a_n+1)>0.
\]
両項が正なので $a_{n+1}>a_n$ である。したがって数列は単調増加かつ上に有界であり，有限の極限 $L$ をもつ。

$1\le L\le2$ であり，平方根の連続性から $L=\sqrt{2+L}$ となる。これを平方すると
\[
 L^2-L-2=(L-2)(L+1)=0.
\]
範囲に合うのは $L=2$ である。漸化式に極限を代入する前に，収束することを確かめた点が大切である。

\paragraph{2.2}
公比 $r=-\frac12$ は $|r|<1$ を満たすので，無限和は
\[
 S=\cfrac1{1+\cfrac12}=\frac23.
\]
最初の $n$ 項の和を $S_n$ とすると，
\[
 S-S_n=\cfrac{\left(-\cfrac12\right)^n}{1+\cfrac12},\qquad
 |S-S_n|=\frac{2}{3\cdot2^n}.
\]
この値は $n$ とともに減少する。$n=9$ では $\frac1{768}>\frac1{1000}$，$n=10$ では $\frac1{1536}<\frac1{1000}$ となる。したがって，条件を満たす最小の自然数は $n=10$ である。

\paragraph{2.3}
$x\ne0$ で
\[
 \frac{\sin2x}{x}=2\frac{\sin2x}{2x}
\]
だから，基本極限により $\lim_{x\to0}f(x)=2$ となる。連続となるにはこの極限が $f(0)=c$ と一致する必要があり，答えは $c=2$ である。

$g$ については $g(1)=0$ であり，$g(1+h)=|h|\to0$ なので $1$ で連続である。しかし，$h\ne0$ に対して差商は
\[
 \frac{g(1+h)-g(1)}h=\frac{|h|}{h}
 =\begin{cases}1&(h>0),\\-1&(h<0)\end{cases}
\]
となる。右極限と左極限が異なるため，$g'(1)$ は存在しない。連続性は値のつながりを調べ，微分可能性は差商の極限まで調べるので，両者は同じ条件ではない。

\paragraph{2.4}
(1) 外側の対数の微分に，内側の $1+x^2$ の微分を掛ける。
\[
 \frac{\dd}{\dd x}\log(1+x^2)=\frac{2x}{1+x^2}.
\]
$1+x^2>0$ なので，すべての実数 $x$ で定義される。

(2) 積の微分と合成関数の微分から，
\begin{align*}
 \frac{\dd}{\dd x}\{(x^2+1)\e^{-x}\}
 &=2x\e^{-x}-(x^2+1)\e^{-x}\\
 &=-(x-1)^2\e^{-x}.
\end{align*}
$x=0$ で関数値は $1$，微分係数は $-1$ だから，接線は $y-1=-x$，すなわち $y=1-x$ である。

(3) $y=x^{\sin x}>0$ とおき，対数をとると $\log y=\sin x\log x$ となる。両辺を微分して
\[
 \frac{y'}y=\cos x\log x+\frac{\sin x}{x}
\]
を得るので，
\[
 y'=x^{\sin x}\left(\cos x\log x+\frac{\sin x}{x}\right)
 \qquad(x>0)
\]
である。指数が変数を含むため，定数指数のべきの公式だけでは微分できない。

\paragraph{2.5}
$a=b$ なら両辺とも $0$ である。$a<b$ とすると，$f(x)=\sin x$ は $[a,b]$ で連続，$(a,b)$ で微分可能である。平均値の定理により，ある $c\in(a,b)$ について
\[
 \sin b-\sin a=(b-a)\cos c
\]
が成り立つ。$|\cos c|\le1$ だから，
\[
 |\sin b-\sin a|=|b-a|\,|\cos c|\le|b-a|.
\]
$a>b$ の場合は $a,b$ を入れ替えて同じ議論を使えばよい。

\paragraph{2.6}
導関数は
\[
 f'(x)=3x^2-6x=3x(x-2),\qquad f''(x)=6x-6
\]
である。したがって，$x<0$ で増加，$0<x<2$ で減少，$x>2$ で増加する。$x=0$ で極大値 $0$，$x=2$ で極小値 $-4$ をとる。また，$x<1$ で上に凸，$x>1$ で下に凸なので，変曲点は $(1,f(1))=(1,-2)$ である。

閉区間 $[-1,3]$ では，極値をとる点と両端を比較する。
\[
\begin{array}{c|rrrr}
 x&-1&0&2&3\\ \hline
 f(x)&-4&0&-4&0
\end{array}
\]
したがって最大値は $0$ で，$x=0,3$ のときにとる。最小値は $-4$ で，$x=-1,2$ のときにとる。

\paragraph{2.7}
$h(x)=\e^x-1-x$ とおくと，$h'(x)=\e^x-1$ である。$x<0$ では負，$x>0$ では正なので，$h$ は $x=0$ で最小値 $h(0)=0$ をとる。よって $\e^x-1-x\ge0$，すなわち $\e^x\ge1+x$ である。$0$ の両側で狭義に増減しているから，等号は $x=0$ に限る。

方程式については $p(x)=\log x+x-2$ とおく。$p$ は $x>0$ で連続かつ微分可能で，
\[
 p'(x)=\frac1x+1>0.
\]
したがって狭義に増加し，$0$ をとる点は高々一つである。一方，
\[
 p(1)=-1<0,\qquad p(2)=\log2>0
\]
だから，中間値の定理により $(1,2)$ に少なくとも一つの解がある。以上より，$x>0$ で解はただ一つであり，その解は $(1,2)$ にある。

\paragraph{2.8}
(1) $f'(x)=3x^2+1>0$ なので $f$ は狭義に増加する。また，$f(x)$ は $x\to\pm\infty$ のときそれぞれ $\pm\infty$ に向かうため，実数全体を値域とする逆関数をもつ。$f(0)=0$，$f'(0)=1$ から，
\[
 g'(0)=\frac1{f'(0)}=1.
\]
逆関数そのものを解の公式で表す必要はない。

(2) $t=2$ で $(x,y)=(4,8)$ であり，
\[
 \frac{\dd x}{\dd t}=2t,\qquad
 \frac{\dd y}{\dd t}=3t^2,\qquad
 \left.\frac{\dd y}{\dd x}\right|_{t=2}
 =\frac{12}{4}=3.
\]
よって接線は $y-8=3(x-4)$，すなわち $y=3x-4$ である。

(3) $y$ を $x$ の関数として両辺を微分すると，
\[
 2x+y+xy'+2yy'=0.
\]
$y'$ を含む項を集めて，
\[
 (x+2y)y'=-(2x+y),\qquad
 y'=-\frac{2x+y}{x+2y}
\]
を得る。点 $(1,2)$ では分母は $5\ne0$ で，傾きは $-\frac45$ となる。したがって接線は
\[
 y-2=-\frac45(x-1),\qquad 4x+5y=14
\]
である。元の関係式に点を代入して $1+2+4=7$ となることも確認できる。

\section{第3章の解答}\label{解答:第3章}

\paragraph{3.1}
一つの小区間の幅は$\Delta x=\frac2n$，第$k$区間の右端は$x_k=\frac{2k}{n}$である。したがってリーマン和は
\[
 \sum_{k=1}^n\left(\frac{2k}{n}\right)^2\frac2n
 =\frac8{n^3}\sum_{k=1}^nk^2
 =\frac{4n(n+1)(2n+1)}{3n^3}.
\]
$n\to\infty$で右辺は$\frac83$へ近づく。$x^2$は閉区間で連続なので，この極限が定積分の値であり，$\int_0^2x^2\dd x=\frac83$となる。

\paragraph{3.2}
区間の加法性から$\int_0^3f(x)\dd x=2+(-1)=1$である。次に積分の向きを直してから線形性を使うと，
\[
 \int_3^0\{2f(x)-3g(x)\}\dd x
 =-\left\{2\int_0^3f(x)\dd x-3\int_0^3g(x)\dd x\right\}
 =-(2-12)=10.
\]
向きを逆にした負号は，被積分関数全体の積分に掛かる。

\paragraph{3.3}
$x-1$は$[-1,1]$で$0$以下，$[1,2]$で$0$以上である。下側の三角形は底辺と高さがともに$2$で面積$2$，上側の三角形は底辺と高さがともに$1$で面積$\frac12$となる。したがって定積分は$-2+\frac12=-\frac32$，幾何学的な面積は$2+\frac12=\frac52$である。積分で計算しても
\[
 \int_{-1}^{2}(x-1)\dd x
 =\left[\frac{x^2}{2}-x\right]_{-1}^{2}=-\frac32
\]
と一致する。面積を積分で求めるなら，負の区間の符号を反転して足す。

\paragraph{3.4}
いずれも区間内で連続であり，基本定理を適用できる。
\begin{align*}
 \text{(1)}\quad\int_{-1}^2(x^2+2x)\dd x
 &=\left[\frac{x^3}{3}+x^2\right]_{-1}^2
 =\frac{20}{3}-\frac23=6,\\
 \text{(2)}\quad\int_0^{\frac\pi2}(\sin x+\cos x)\dd x
 &=[-\cos x+\sin x]_0^{\frac\pi2}=1-(-1)=2,\\
 \text{(3)}\quad\int_1^{\e}\left(\frac2x+x\right)\dd x
 &=\left[2\log x+\frac{x^2}{2}\right]_1^{\e}
 =2+\frac{\e^2-1}{2}=\frac{\e^2+3}{2}.
\end{align*}
(3)の区間は正の側にあるので，対数を$\log x$と書ける。

\paragraph{3.5}
基本定理から$A'(x)=(x-2)(x+1)$である。$[-2,-1]$で増加，$[-1,2]$で減少，$[2,3]$で増加するため，閉区間での最大・最小の候補は両端と$x=-1,2$である。比較するために
\[
 A(x)=\left[\frac{t^3}{3}-\frac{t^2}{2}-2t\right]_1^x
 =\frac{x^3}{3}-\frac{x^2}{2}-2x+\frac{13}{6}
\]
と求めると，
\[
 A(-2)=\frac32,\quad A(-1)=\frac{10}{3},\quad
 A(2)=-\frac76,\quad A(3)=\frac23.
\]
よって最大値は$x=-1$で$\frac{10}{3}$，最小値は$x=2$で$-\frac76$である。導関数から増減を先に調べたので，すべての点で値を比較する必要はない。

\paragraph{3.6}
上端は$\sin x$，下端は$\cos x$である。それぞれの微分まで掛けると
\[
 G'(x)=\sqrt{1+\sin^2x}\cos x
       -\sqrt{1+\cos^2x}(-\sin x)
 =\sqrt{1+\sin^2x}\cos x+\sqrt{1+\cos^2x}\sin x.
\]
下端を微分したときの負号と，端点の公式にある負号が掛かって正になる。上下端の大小が入れ替わっても同じ式が成り立つ。

\paragraph{3.7}
$x^2\ge0$だから$\sqrt{1+x^2}\ge1$である。また両辺が非負であり，
\[
 \left(1+\frac{x^2}{2}\right)^2-(1+x^2)=\frac{x^4}{4}\ge0
\]
なので，平方の大小から$\sqrt{1+x^2}\le1+\frac{x^2}{2}$を得る。$[0,1]$上でこの不等式を積分すると
\[
 1=\int_0^1 1\dd x
 \le\int_0^1\sqrt{1+x^2}\dd x
 \le\int_0^1\left(1+\frac{x^2}{2}\right)\dd x
 =1+\frac16=\frac76.
\]
近似の式を単に代入したのではなく，区間全体の不等式を確かめてから積分している。

\paragraph{3.8}
(1) $x=0$を代入すると$f(0)=3$である。右辺は連続な$f$の積分なので微分可能であり，等式から$f$も微分可能となる。微分すれば$f'=-f$だから
\[
 (\e^xf(x))'=\e^x\{f(x)+f'(x)\}=0.
\]
従って$\e^xf(x)$は定数で，初期値から$f(x)=3\e^{-x}$となる。元の式にも
\[
 3-\int_0^x3\e^{-t}\dd t
 =3-[-3\e^{-t}]_0^x=3\e^{-x}
\]
と戻るので，これが求める関数である。

(2) 積分$K=\int_0^1g(t)\dd t$は$x$に依存しない定数である。条件は$g(x)=x+K$と書ける。しかし，この式を$0$から$1$まで積分すると
\[
 K=\int_0^1(x+K)\dd x=\frac12+K
\]
となり矛盾する。したがって条件を満たす連続関数は存在しない。微分だけなら$g'=1$を得るが，それだけでは失われた定数の条件を確かめられない。

\section{第4章の解答}\label{解答:第4章}

\paragraph{4.1}
分母の $\sqrt{x}=x^{\frac12}$ で各項を割ると，
\[
 \frac{3x^2-2x+4}{\sqrt{x}}
 =3x^{\frac32}-2x^{\frac12}+4x^{-\frac12}.
\]
各項へべきの公式を使って
\[
 \int\frac{3x^2-2x+4}{\sqrt{x}}\dd x
 =\frac65x^{\frac52}-\frac43x^{\frac32}+8\sqrt{x}+C.
\]
この右辺を微分すれば
\[
 \frac65\cdot\frac52x^{\frac32}
 -\frac43\cdot\frac32x^{\frac12}
 +8\cdot\frac12x^{-\frac12}
 =3x^{\frac32}-2x^{\frac12}+4x^{-\frac12}
\]
となる。通分すると元の式へ戻る。$x>0$ という条件のもとで，すべてのべきと分母が定義される。

\paragraph{4.2}
立方の和の因数分解を使えば
\[
 x^3+1=(x+1)(x^2-x+1)
\]
だから，$x\ne-1$ で
\[
 \frac{x^3+1}{x+1}=x^2-x+1.
\]
よって
\[
 \int\frac{x^3+1}{x+1}\dd x
 =\frac{x^3}{3}-\frac{x^2}{2}+x+C.
\]
右辺の整式は $x=-1$ でも定義されるが，もとの関数はその点では定義されない。したがって，もとの関数の原始関数として述べるときは，$x<-1$ と $x>-1$ の各区間で読む。両区間の任意定数は独立でよい。

\paragraph{4.3}
\[
 \frac{5x-1}{(x-1)(x+1)}
 =\frac A{x-1}+\frac B{x+1}
\]
と置く。通分した分子が一致する条件は
\[
 5x-1=A(x+1)+B(x-1)=(A+B)x+(A-B)
\]
である。したがって $A+B=5$，$A-B=-1$ から $A=2$，$B=3$ となる。
\[
 \int\frac{5x-1}{(x-1)(x+1)}\dd x
 =2\log|x-1|+3\log|x+1|+C.
\]
$x=-1,1$ のいずれも含まない各区間で成り立つ。微分した二項を通分すれば分子が $2(x+1)+3(x-1)=5x-1$ となる。

\paragraph{4.4}
分子を $3x+4=3(x+1)+1$ と書けば
\[
 \frac{3x+4}{(x+1)^2}
 =\frac3{x+1}+\frac1{(x+1)^2}.
\]
したがって
\[
 \int\frac{3x+4}{(x+1)^2}\dd x
 =3\log|x+1|-\frac1{x+1}+C.
\]
最後の項の微分が $\frac1{(x+1)^2}$ になることに注意する。$x=-1$ をまたがない区間で用いる。

\paragraph{4.5}
平方完成すると
\[
 x^2+6x+8=(x+3)^2-1=(x+2)(x+4).
\]
平方の差なので因数分解できる。さらに
\[
 \frac1{(x+2)(x+4)}
 =\frac12\left(\frac1{x+2}-\frac1{x+4}\right)
\]
と分けられる。右辺を通分した分子は $\frac12\{(x+4)-(x+2)\}=1$ である。よって
\[
 \int\frac1{x^2+6x+8}\dd x
 =\frac12\log|x+2|-\frac12\log|x+4|+C.
\]
除外点は $x=-4,-2$ である。対数の差を $\frac12\log\left|\frac{x+2}{x+4}\right|$ とまとめてもよい。

\paragraph{4.6}
角を $2x$ として二倍角の公式を使うと
\[
 \cos^2(2x)=\frac{1+\cos(4x)}2.
\]
したがって
\[
\begin{aligned}
 \int_0^{\frac\pi2}\cos^2(2x)\dd x
 &=\left[\frac x2+\frac18\sin(4x)\right]_0^{\frac\pi2}\\
 &=\frac\pi4+\frac18\sin(2\pi)-0=\frac\pi4.
\end{aligned}
\]
被積分関数は非負だから，得られた値も非負である。こうした符号の確認は，途中の符号間違いを見つける助けになる。

\paragraph{4.7}
加法定理から
\[
 \cos(5x-2x)=\cos(5x)\cos(2x)+\sin(5x)\sin(2x),
\]
\[
 \cos(5x+2x)=\cos(5x)\cos(2x)-\sin(5x)\sin(2x).
\]
第一式から第二式を引くと
\[
 \sin(2x)\sin(5x)=\frac12\{\cos(3x)-\cos(7x)\}
\]
を得る。それぞれの角の微分で割って
\[
 \int\sin(2x)\sin(5x)\dd x
 =\frac16\sin(3x)-\frac1{14}\sin(7x)+C.
\]
微分して $\frac12\cos(3x)-\frac12\cos(7x)$ へ戻し，導いた恒等式を用いれば検算できる。

\section{第5章の解答}\label{解答:第5章}

\paragraph{5.1}
括弧の中を $t=3x^2-4x+7$ と置くと，$\dd t=(6x-4)\dd x$ である。したがって
\[
 \int(6x-4)(3x^2-4x+7)^5\dd x
 =\int t^5\dd t
 =\frac{(3x^2-4x+7)^6}{6}+C.
\]
微分すれば，外側の六乗から $6$ が現れて分母と打ち消し合い，内側の微分 $6x-4$ が掛かって元の式に戻る。多項式なので全実数上で成り立つ。

\paragraph{5.2}
$t=\log x$ と置くと $\dd t=\frac1x\dd x$ だから，
\[
 \int\frac1{x\log x}\dd x
 =\int\frac1t\dd t
 =\log|\log x|+C.
\]
実数の対数のために $x>0$ が必要で，さらに分母の $\log x$ が $0$ となる $x=1$ を除く。よって適用する区間は $0<x<1$ と $1<x$ に分かれ，それぞれの任意定数は独立である。$0<x<1$ では $\log x<0$ なので，外側の対数の絶対値を外せない。

\paragraph{5.3}
$t=x^2$，$\dd t=2x\dd x$ と置く。$x=-2$ は $t=4$，$x=1$ は $t=1$ に対応するので，
\[
 \int_{-2}^1 2x\e^{x^2}\dd x
 =\int_4^1\e^t\dd t
 =[\e^t]_4^1=\e-\e^4.
\]
合成形の置換は，$\frac{\dd}{\dd x}\e^{x^2}=2x\e^{x^2}$ に基づいており，$x^2$ が単調である必要はない。区間を分けても
\[
 \int_4^0\e^t\dd t+\int_0^1\e^t\dd t
 =\int_4^1\e^t\dd t
\]
となる。被積分関数が負になる部分があるため，定積分が負でも矛盾しない。

\paragraph{5.4}
$t=\sqrt{1+x}$ と置けば，$x=t^2-1$，$\dd x=2t\dd t$ である。端点は $x=0,3$ から $t=1,2$ へ移る。したがって
\[
\begin{aligned}
 \int_0^3\frac1{1+\sqrt{1+x}}\dd x
 &=2\int_1^2\frac{t}{1+t}\dd t\\
 &=2\int_1^2\left(1-\frac1{1+t}\right)\dd t\\
 &=2[t-\log(1+t)]_1^2\\
 &=2-2\log\frac32.
\end{aligned}
\]
$1\le t\le2$ では $1+t>0$ なので，対数の絶対値は不要である。

\paragraph{5.5}
$x=a\sin\theta$ と置き，$0\le\theta\le\frac\pi6$ を選ぶ。この区間で $\cos\theta>0$ だから，$\sqrt{a^2-x^2}=a\cos\theta$ と書ける。また $\dd x=a\cos\theta\dd\theta$ なので，
\[
\begin{aligned}
 \int_0^{\frac a2}\sqrt{a^2-x^2}\dd x
 &=a^2\int_0^{\frac\pi6}\cos^2\theta\dd\theta\\
 &=a^2\left[\frac\theta2+\frac14\sin(2\theta)\right]_0^{\frac\pi6}\\
 &=a^2\left(\frac\pi{12}+\frac{\sqrt3}{8}\right).
\end{aligned}
\]
最後は $\sin\frac\pi3=\frac{\sqrt3}{2}$ を用いた。

\paragraph{5.6}
$0\le\theta\le\frac\pi3$ で $x=\tan\theta$ と置く。端点は $0,\sqrt3$ に対応し，この区間では $\cos\theta\ne0$ である。よって
\[
 \int_0^{\sqrt3}\frac1{1+x^2}\dd x
 =\int_0^{\frac\pi3}\frac1{1+\tan^2\theta}
   \frac1{\cos^2\theta}\dd\theta
 =\int_0^{\frac\pi3}1\dd\theta=\frac\pi3.
\]
分母と微分に現れる因子が等しいことを使っている。

\paragraph{5.7}
変数を換えるときは，端点だけでなく被積分関数と微分の記号も書き換える必要がある。$t=x^2$ では $\dd t=2x\dd x$ であり，$\dd x$ をそのまま $\dd t$ にすることはできない。また $-1\le x\le1$ 全体で $x=\sqrt{t}$ と戻すこともできない。負の $x$ は $x=-\sqrt{t}$ という別の枝に属する。したがって，$1$ の積分を同じ端点の $t$ の積分として表す根拠がない。実際には $\int_{-1}^1 1\dd x=2$ である。

一方，第二の積分には $2x\dd x$ があるから，合成形の置換をそのまま使える。
\[
 \int_{-1}^1\frac{2x}{1+x^2}\dd x
 =\int_1^1\frac1{1+t}\dd t=0.
\]
原始関数 $\log(1+x^2)$ の上下端の値がともに $\log2$ になることからも確かめられる。同じ端点へ移ることだけでなく，置換公式の形を満たしていることが結論を支えている。

\section{第6章の解答}\label{解答:第6章}

\paragraph{6.1}
$x$ は微分すると定数になり，$\e^{3x}$ は積分しても指数関数のままである。したがって
\[
 u=x,\qquad v'=\e^{3x},\qquad
 u'=1,\qquad v=\frac13\e^{3x}
\]
と選ぶ。部分積分により
\[
\begin{aligned}
 \int x\e^{3x}\dd x
 &=\frac x3\e^{3x}-\frac13\int\e^{3x}\dd x\\
 &=\frac x3\e^{3x}-\frac19\e^{3x}+C
 =\frac{3x-1}{9}\e^{3x}+C.
\end{aligned}
\]
$\e^{3x}$ を積分するたびに $3$ で割るため，最後の項の係数は $\frac19$ になる。右辺を微分すれば
\[
 \frac13\e^{3x}
 +\frac{3x-1}{9}\cdot3\e^{3x}
 =\frac{1+3x-1}{3}\e^{3x}=x\e^{3x}
\]
となる。

\paragraph{6.2}
$\log x$ を微分すれば $\frac1x$ になり，$x^2$ の原始関数 $\frac{x^3}{3}$ と掛けると整式が残る。そこで $u=\log x$，$v'=x^2$ と選ぶ。
\[
\begin{aligned}
 \int x^2\log x\dd x
 &=\frac{x^3}{3}\log x
  -\int\frac{x^3}{3}\cdot\frac1x\dd x\\
 &=\frac{x^3}{3}\log x-\frac13\int x^2\dd x\\
 &=\frac{x^3}{3}\log x-\frac{x^3}{9}+C.
\end{aligned}
\]
$x^2$ を微分する側に選ぶこともできるが，その場合は $\log x$ の原始関数を求めたうえで，残る積にも対数が現れる。上の選択では一回の部分積分で対数が積分記号の中からなくなるため，見通しがよい。答えを微分すると，$\frac{x^3}{3}\cdot\frac1x=\frac{x^2}{3}$ と $-\frac{x^3}{9}$ の微分が打ち消し合う。

\paragraph{6.3}
整式を微分する側に選ぶと，最初の部分積分で三次式が二次式になる。
\[
 \int x^3\sin x\dd x
 =-x^3\cos x+3\int x^2\cos x\dd x.
\]
続いて同じ方針で
\[
 \int x^2\cos x\dd x
 =x^2\sin x-2\int x\sin x\dd x
\]
とし，二次式を一次式にする。さらに
\[
\begin{aligned}
 \int x\sin x\dd x
 &=-x\cos x+\int\cos x\dd x\\
 &=-x\cos x+\sin x+C
\end{aligned}
\]
とすれば，整式の因子は定数になった。これらを逆の順序に戻すと
\[
\begin{aligned}
 \int x^3\sin x\dd x
 &=-x^3\cos x
   +3\{x^2\sin x-2(-x\cos x+\sin x)\}+C\\
 &=-x^3\cos x+3x^2\sin x+6x\cos x-6\sin x+C.
\end{aligned}
\]
微分すると，$x^3\sin x$ 以外の項は
\[
 -3x^2\cos x+3x^2\cos x,
 \quad 6x\sin x-6x\sin x,
 \quad 6\cos x-6\cos x
\]
の組になって消える。三次，一次と飛ばさず，三次，二次，一次，定数と順に次数が下がっている。

\paragraph{6.4}
本文と同じく，積分定数を固定するため
\[
 I(x)=\int_0^x\e^{2t}\cos t\dd t,
 \qquad
 J(x)=\int_0^x\e^{2t}\sin t\dd t
\]
と定める。どちらも指数関数を積分する側に選ぶと，
\[
\begin{aligned}
 I(x)
 &=\left[\frac12\e^{2t}\cos t\right]_0^x
   +\frac12\int_0^x\e^{2t}\sin t\dd t\\
 &=\frac12\e^{2x}\cos x-\frac12+\frac12J(x),\\
 J(x)
 &=\left[\frac12\e^{2t}\sin t\right]_0^x
   -\frac12\int_0^x\e^{2t}\cos t\dd t\\
 &=\frac12\e^{2x}\sin x-\frac12I(x).
\end{aligned}
\]
第二式を第一式へ代入して
\[
 I(x)=\frac12\e^{2x}\cos x-\frac12
       +\frac14\e^{2x}\sin x-\frac14I(x).
\]
したがって
\[
 \frac54I(x)=\frac14\e^{2x}(2\cos x+\sin x)-\frac12,
\]
すなわち $I(x)=\frac15\e^{2x}(2\cos x+\sin x)-\frac25$ である。定数をまとめれば
\[
 \boxed{\int\e^{2x}\cos x\dd x
 =\frac15\e^{2x}(2\cos x+\sin x)+C}
\]
となる。右辺を微分すると
\[
 \frac15\e^{2x}
 \{2(2\cos x+\sin x)-2\sin x+\cos x\}
 =\e^{2x}\cos x
\]
である。$\e^{2x}$ の原始関数に付く $\frac12$ を，二回とも保つことが必要である。

\paragraph{6.5}
$t=\log x$ と置く。$x>0$ だから $x=\e^t$ と戻せて，$\dd x=\e^t\dd t$ となる。したがって
\[
 \int\cos(\log x)\dd x
 =\int\e^t\cos t\dd t.
\]
右辺では $\cos t$ を微分し，$\e^t$ を積分する。本文で得た $\int\e^t\sin t\dd t$ の結果も使うと
\[
\begin{aligned}
 \int\e^t\cos t\dd t
 &=\e^t\cos t+\int\e^t\sin t\dd t\\
 &=\e^t\cos t+\frac12\e^t(\sin t-\cos t)+C\\
 &=\frac12\e^t(\sin t+\cos t)+C.
\end{aligned}
\]
よって元の変数へ戻して
\[
 \int\cos(\log x)\dd x
 =\frac x2\{\sin(\log x)+\cos(\log x)\}+C.
\]
微分すると，括弧をそのまま残す項と，括弧を微分する項が加わる。前者は $\frac12\{\sin(\log x)+\cos(\log x)\}$，後者は $\frac12\{\cos(\log x)-\sin(\log x)\}$ なので，和は $\cos(\log x)$ になる。

\paragraph{6.6}
$t=1+x^2$ と置くと，$\dd t=2x\dd x$ であり，端点は $x=0,1$ から $t=1,2$ へ移る。したがって
\[
 \int_0^1x\log(1+x^2)\dd x
 =\frac12\int_1^2\log t\dd t.
\]
ここで $\log t$ を微分し，$1$ を積分する部分積分を使う。
\[
\begin{aligned}
 \frac12\int_1^2\log t\dd t
 &=\frac12\left\{[t\log t]_1^2-\int_1^2 1\dd t\right\}\\
 &=\frac12(2\log2-1)
 =\log2-\frac12.
\end{aligned}
\]
もとの式にある $x$ は，対数を微分するために外に残したのではなく，置換の $x\dd x=\frac12\dd t$ として使った。置換後は対数一つの積分になるので，そこで部分積分を行えばよい。

\paragraph{6.7}
$I=\int_0^1\sqrt{4+x^2}\dd x$ と置く。根号を微分し，$1$ を積分すると，
\[
 I=[x\sqrt{4+x^2}]_0^1
   -\int_0^1\frac{x^2}{\sqrt{4+x^2}}\dd x.
\]
分子を $x^2=(4+x^2)-4$ と分けると
\[
 \frac{x^2}{\sqrt{4+x^2}}
 =\sqrt{4+x^2}-\frac4{\sqrt{4+x^2}}.
\]
よって，$L=\int_0^1\frac1{\sqrt{4+x^2}}\dd x$ と書けば
\[
 I=\sqrt5-I+4L,
 \qquad 2I=\sqrt5+4L.
\]
残る $L$ では，$r(x)=\sqrt{4+x^2}$ とすると $r'(x)=\frac{x}{r(x)}$ である。また $r(x)>|x|$ なので $x+r(x)>0$ が成り立ち，
\[
 \frac{\dd}{\dd x}\log\{x+r(x)\}
 =\cfrac{1+\cfrac{x}{r(x)}}{x+r(x)}
 =\frac1{r(x)}
\]
となる。したがって
\[
\begin{aligned}
 L&=[\log(x+\sqrt{4+x^2})]_0^1\\
  &=\log(1+\sqrt5)-\log2
   =\log\frac{1+\sqrt5}{2}.
\end{aligned}
\]
これを戻して
\[
 \boxed{\int_0^1\sqrt{4+x^2}\dd x
 =\frac{\sqrt5}{2}+2\log\frac{1+\sqrt5}{2}}
\]
を得る。下端 $x=0$ で対数の値は $\log2$ であり，$0$ にはならない。本文の $\sqrt{1+x^2}$ の場合との違いを確認する必要がある。

\paragraph{6.8}
$u=x$，$v'=f'(x)$ とすれば $u'=1$，$v=f(x)$ を使える。したがって
\[
\begin{aligned}
 \int_0^1xf'(x)\dd x
 &=[xf(x)]_0^1-\int_0^1f(x)\dd x\\
 &=2-\frac32=\frac12.
\end{aligned}
\]
もう一つは $u=1-x$，$v'=f'(x)$ とする。今度は $u'=-1$ だから
\[
\begin{aligned}
 \int_0^1(1-x)f'(x)\dd x
 &=[(1-x)f(x)]_0^1+\int_0^1f(x)\dd x\\
 &=(0-1)+\frac32=\frac12.
\end{aligned}
\]
端点の積と $f$ の定積分が分かっているため，$f$ の具体的な式は不要であった。二つの結果を加えると $1$ になり，
\[
 \int_0^1\{x+(1-x)\}f'(x)\dd x
 =\int_0^1f'(x)\dd x=f(1)-f(0)=1
\]
とも一致する。

\section{第7章の解答}\label{解答:第7章}

\paragraph{7.1}
$\frac{x^3}{1+x^2}$ は，分子が奇関数，分母が零にならない偶関数なので奇関数である。したがって，対称区間での積分は零になる。残る二項は偶関数だから，求める値は
\[
2\int_0^2(2x^2-\cos x)\dd x
 =2\left[\frac{2x^3}{3}-\sin x\right]_0^2
 =\frac{32}{3}-2\sin2.
\]
最初の分数を割り算して積分する必要はない。ただし，分母に零がないことは確認している。

\paragraph{7.2}
積分を $I$ とする。$x=\pi-t$ と置くと，$\sin^2(\pi-t)=\sin^2t$ より
\[
I=\int_0^\pi(\pi-x)\sin^2x\dd x.
\]
元の式と加えて，
\[
2I=\pi\int_0^\pi\sin^2x\dd x
 =\frac{\pi}{2}\int_0^\pi(1-\cos2x)\dd x
 =\frac{\pi^2}{2}.
\]
よって $I=\frac{\pi^2}{4}$ である。$\pi$ が二度現れるのは，一つが一次因子 $x$ の平均 $\frac{\pi}{2}$，もう一つが $\sin^2x$ の積分 $\frac{\pi}{2}$ に由来するからである。

\paragraph{7.3}
$\sin^2x+\cos^2x=1$ なので，$\sin x$ と $\cos x$ は同時に零にならない。それぞれの六乗は非負だから，分母は正である。積分を $I$ として $x=\frac{\pi}{2}-t$ と置けば
\[
I=\int_0^{\frac{\pi}{2}}
\frac{\sin^6x}{\sin^6x+\cos^6x}\dd x.
\]
二式を加えると
\[
2I=\int_0^{\frac{\pi}{2}}1\dd x=\frac{\pi}{2},
\]
したがって $I=\frac{\pi}{4}$ である。

\paragraph{7.4}
$|\cos x|$ は周期 $\pi$ を持つ。区間の長さは $3\pi$ なので，始点に関係なく三周期分である。一周期分は，零点 $\frac{\pi}{2}$ で符号を分けて
\[
\int_0^\pi|\cos x|\dd x
 =\int_0^{\frac{\pi}{2}}\cos x\dd x
  -\int_{\frac{\pi}{2}}^\pi\cos x\dd x
 =1+1=2.
\]
よって最初の積分は $6$ である。一方，絶対値のない積分は
\[
\left[\sin x\right]_{\frac{\pi}{3}}^{\frac{10\pi}{3}}
 =-\frac{\sqrt3}{2}-\frac{\sqrt3}{2}
 =-\sqrt3.
\]
$\cos x$ 自体の周期は $2\pi$ であり，この区間はその整数周期分ではない。また絶対値のある積分はグラフと軸の間の面積をすべて正に数えるが，絶対値のない積分は軸の下の部分を負に数える。この二つの違いを区別する。

\paragraph{7.5}
積分を $I$ とする。$x=\frac1t$ なら $\dd x=-\frac1{t^2}\dd t$ なので，
\[
I=\int_{\frac14}^4\frac{t^3}{1+t^4}\dd t.
\]
積分変数の文字を $x$ に戻し，元の式と加える。
\[
2I=\int_{\frac14}^4
\left(\frac{1}{x(1+x^4)}+\frac{x^3}{1+x^4}\right)\dd x
 =\int_{\frac14}^4\frac{1}{x}\dd x
 =2\log4.
\]
よって $I=\log4$ である。変形の途中では，分子を通分すると $1+x^4$ になり，分母の同じ因子が消える。

\paragraph{7.6}
$3x^2+1$ は偶関数，$x^5$ は奇関数である。反転した重みは
\[
\frac{1}{1+\e^{-x^5}}=\frac{\e^{x^5}}{1+\e^{x^5}}
\]
だから，元の重みと加えると一になる。求める積分は
\[
\frac12\int_{-1}^1(3x^2+1)\dd x
 =\int_0^1(3x^2+1)\dd x
 =\left[x^3+x\right]_0^1=2.
\]
指数 $x^5$ の微分 $5x^4$ が分子にないため，指数を新変数にする置換は直ちには簡単にならない。ここでは指数の奇偶を見るのが有効である。

\paragraph{7.7}
区間の中央は $4$ である。$u=x-4$ と置くと，両端 $2,6$ は $-2,2$ に移り，
\[
\int_2^6\frac{(x-4)^2+2}{1+\e^{x-4}}\dd x
 =\int_{-2}^2\frac{u^2+2}{1+\e^u}\dd u.
\]
分子は偶関数であり，指数は奇関数なので，この値は
\[
\int_0^2(u^2+2)\dd u
 =\left[\frac{u^3}{3}+2u\right]_0^2
 =\frac{20}{3}
\]
である。原点に関する対称性に移してから使うことで，分子と分母の変化を同時に確認できる。

\paragraph{7.8}
反転した式と元の式を加えると
\[
2I=\int_0^\pi\{x^2+(\pi-x)^2\}\sin x\dd x
 =2I-2\pi\int_0^\pi x\sin x\dd x
  +\pi^2\int_0^\pi\sin x\dd x.
\]
ここで $\int_0^\pi x\sin x\dd x=\pi$，$\int_0^\pi\sin x\dd x=2$ だから，右辺は $2I$ に戻る。$x^2$ の項が残るため，$I$ を決める新しい等式にはならない。

部分積分では，$x^2$ を微分し，$\sin x$ を積分する。
\[
I=\left[-x^2\cos x\right]_0^\pi
  +2\int_0^\pi x\cos x\dd x
 =\pi^2+2\int_0^\pi x\cos x\dd x.
\]
もう一度部分積分すると
\[
\int_0^\pi x\cos x\dd x
 =\left[x\sin x\right]_0^\pi-\int_0^\pi\sin x\dd x
 =-2.
\]
したがって $I=\pi^2-4$ である。反転は正しい操作でも，それだけで未知の積分が求まるとは限らない。変数の次数を下げる部分積分が必要な場面である。

\section{第8章の解答}\label{解答:第8章}

\paragraph{8.1}
まず交点と上下関係を調べる。$x^2=3x$ より交点の $x$ 座標は $0,3$ であり，$0\leq x\leq3$ では $3x-x^2=x(3-x)\geq0$ である。したがって面積は，上の曲線から下の曲線を引いて
\[
\begin{aligned}
A&=\int_0^3(3x-x^2)\dd x\\
 &=\left[\frac{3x^2}{2}-\frac{x^3}{3}\right]_0^3
 =\frac{27}{2}-9=\frac92.
\end{aligned}
\]
この区間では両曲線とも $x$ 軸の上側にある。回転後の断面は，外半径 $3x$，内半径 $x^2$ の円環なので，体積は
\[
\begin{aligned}
V&=\pi\int_0^3\{(3x)^2-(x^2)^2\}\dd x\\
 &=\pi\left[3x^3-\frac{x^5}{5}\right]_0^3
 =\pi\left(81-\frac{243}{5}\right)
 =\frac{162\pi}{5}.
\end{aligned}
\]
面積では高さの差を使い，体積では円板の面積の差を使う。帯の高さ $3x-x^2$ を一つの半径として二乗すると，別の断面を数えることになる。

\paragraph{8.2}
横の帯を使うので，交点も $y$ 座標で求める。
\[
y^2=3-2y
\quad\Longleftrightarrow\quad
(y+3)(y-1)=0.
\]
よって積分区間は $-3\leq y\leq1$ である。この間では
\[
(3-2y)-y^2=(1-y)(y+3)\geq0
\]
だから，直線 $x=3-2y$ が右側，放物線 $x=y^2$ が左側にある。求める面積は
\[
\begin{aligned}
A&=\int_{-3}^1(3-2y-y^2)\dd y\\
 &=\left[3y-y^2-\frac{y^3}{3}\right]_{-3}^1\\
 &=\left(3-1-\frac13\right)-(-9-9+9)
 =\frac{32}{3}.
\end{aligned}
\]
横の帯の幅を右端から左端を引いた差で表したため，放物線を二つの平方根に分ける必要がない。

\paragraph{8.3}
各 $x$ で線分 $-1\leq y\leq x$ は回転軸を含むので，回転後の断面に穴はない。半径は，上端までの距離 $x$ と下端までの距離 $1$ のうち大きいほうである。したがって外半径は
\[
R(x)=
\begin{cases}
1,&0\leq x\leq1,\\
x,&1\leq x\leq2.
\end{cases}
\]
外半径を決める端が交代するのは $x=1$ である。断面の円板を積み重ねると
\[
\begin{aligned}
V&=\pi\int_0^1 1\dd x+\pi\int_1^2 x^2\dd x\\
 &=\pi+\pi\left[\frac{x^3}{3}\right]_1^2
 =\pi+\frac{7\pi}{3}=\frac{10\pi}{3}.
\end{aligned}
\]
上側と下側を別々に回した円板は重なる。二つの円板の面積を加えるのでなく，大きいほうの円板を一度だけ数えている。

\paragraph{8.4}
円筒殻を使うときは，$y$ 軸に平行な縦の帯を取る。回転半径は $x$，層の高さは $1-x^2$ だから，
\[
\begin{aligned}
V&=2\pi\int_0^1x(1-x^2)\dd x\\
 &=2\pi\left[\frac{x^2}{2}-\frac{x^4}{4}\right]_0^1
 =2\pi\left(\frac12-\frac14\right)=\frac{\pi}{2}.
\end{aligned}
\]
円板を使うときは，$y$ 軸に垂直な横の帯を取る。$0\leq y\leq1$ で，元の領域の左右の端は $x=0$ と $x=\sqrt y$ である。回転後の断面は半径 $\sqrt y$ の円板なので，
\[
V=\pi\int_0^1(\sqrt y)^2\dd y
 =\pi\left[\frac{y^2}{2}\right]_0^1=\frac{\pi}{2}.
\]
二つの式では積分変数と分け方が異なるが，どちらも同じ立体を重複なく数えている。

\paragraph{8.5}
グラフの弧長では $\sqrt{1+y'^2}$ を積分する。まず合成関数の微分により
\[
y'=\frac13\cdot\frac32(x^2+2)^{\frac12}\cdot2x
 =x\sqrt{x^2+2}.
\]
このまま根号を積分しようとせず，中身を整理する。
\[
1+y'^2=1+x^2(x^2+2)=x^4+2x^2+1=(x^2+1)^2.
\]
$x^2+1>0$ なので，平方根は $|x^2+1|=x^2+1$ となる。したがって
\[
L=\int_0^1(x^2+1)\dd x
 =\left[\frac{x^3}{3}+x\right]_0^1=\frac43.
\]
弧長の根号が完全平方になっていることが，計算を簡単にする理由である。

\paragraph{8.6}
速度を展開すると $v(t)=3t^2-9t+6$ である。これを積分し，$s(0)=0$ を使えば位置は
\[
s(t)=t^3-\frac92t^2+6t
\]
となる。終点の位置は
\[
s(3)=27-\frac{81}{2}+18=\frac92.
\]
道のりを求めるには，速度の符号を区別する。$v(t)=3(t-1)(t-2)$ は，$0\leq t<1$ と $2<t\leq3$ で正，$1<t<2$ で負である。折り返す位置は
\[
s(1)=1-\frac92+6=\frac52,\qquad
s(2)=8-18+12=2.
\]
各区間では一方向に動くので，道のりはそれぞれの位置の差の絶対値の和になり，
\[
D=\left|\frac52-0\right|
 +\left|2-\frac52\right|
 +\left|\frac92-2\right|
 =\frac52+\frac12+\frac52=\frac{11}{2}.
\]
これは $\int_0^1v\dd t-\int_1^2v\dd t+\int_2^3v\dd t$ と同じである。全時間は $3$ なので，平均速度と平均の速さはそれぞれ
\[
\frac{s(3)-s(0)}{3}=\frac32,\qquad
\frac{D}{3}=\frac{11}{6}
\]
となる。戻った分は変位では差し引かれ，道のりでは加えられるため，この運動では平均の速さのほうが大きい。

\paragraph{8.7}
関数の平均値は積分を区間の長さで割って求める。
\[
\overline f=\frac12\int_0^2(1+x^2)\dd x
 =\frac12\left[x+\frac{x^3}{3}\right]_0^2
 =\frac12\left(2+\frac83\right)=\frac73.
\]
$f(c)=\overline f$ となる条件は
\[
1+c^2=\frac73
\quad\Longleftrightarrow\quad
c^2=\frac43.
\]
$c\in[0,2]$ だから負の解は除き，
\[
c=\frac{2}{\sqrt3}=\frac{2\sqrt3}{3}
\]
を得る。区間の中央 $c=1$ では $f(1)=2$ なので，平均値を取る点とは異なる。

\paragraph{8.8}
三角形は $0\leq x\leq2$，$0\leq y\leq\frac{x}{2}$ と表せる。まず面積を求めると
\[
A=\int_0^2\frac{x}{2}\dd x
 =\left[\frac{x^2}{4}\right]_0^2=1.
\]
縦の帯の中心の高さは $\frac{x}{4}$ であり，帯の面積はほぼ $\frac{x}{2}\dd x$ だから，面積重心の高さは
\[
\begin{aligned}
\overline y_{\mathrm{面}}
 &=\frac1A\int_0^2\frac{x}{4}\cdot\frac{x}{2}\dd x\\
 &=\frac18\left[\frac{x^3}{3}\right]_0^2
 =\frac13.
\end{aligned}
\]
三角形は軸の上側にあり，軸は境界の一辺に重なるが内部を横切らない。したがって体積の定理が使えて，
\[
V=2\pi\overline y_{\mathrm{面}}A
 =2\pi\cdot\frac13\cdot1=\frac{2\pi}{3}.
\]

全表面積では，境界の三本の線分を別々に考え，それぞれの線重心を使う。一様な線分の線重心は中点なので，次のようになる。

$x$ 軸上の線分は長さ $2$，線重心の高さ $0$ である。回転しても軸上の線分のままで，面積のある面を作らない。その寄与は $2\pi\cdot0\cdot2=0$ である。

縦の線分 $x=2$，$0\leq y\leq1$ は長さ $1$，線重心の高さ $\frac12$ である。回転すると半径 $1$ の円板になり，面の定理から
\[
S_{\mathrm{底面}}=2\pi\cdot\frac12\cdot1=\pi
\]
となる。端点が回転軸に触れていても，面を重複して作ることはない。

斜めの線分 $y=\frac{x}{2}$，$0\leq x\leq2$ は長さ $\sqrt5$，線重心の高さ $\frac12$ である。線重心を積分で確かめると，$\dd s=\frac{\sqrt5}{2}\dd x$ より
\[
\overline y_{\mathrm{線}}
 =\frac1{\sqrt5}\int_0^2\frac{x}{2}\cdot\frac{\sqrt5}{2}\dd x
 =\frac14\left[\frac{x^2}{2}\right]_0^2=\frac12.
\]
この線分が作る円すいの側面の面積は
\[
S_{\mathrm{側面}}=2\pi\cdot\frac12\cdot\sqrt5=\pi\sqrt5.
\]
底面と側面は円周で接するが，面積のある部分では重ならない。よって全表面積は
\[
S=\pi+\pi\sqrt5=\pi(1+\sqrt5)
\]
である。体積に使った面積重心の高さ $\frac13$ と，境界の各線分に使った線重心の高さを混同しない。

\section{第9章の解答}\label{解答:第9章}

\paragraph{9.1}
分母から$n^2$をくくり，分子の$k$を$n\cdot\frac{k}{n}$と書くと
\[
 \frac{k}{n^2+k^2}
 =\frac1n\cfrac{\cfrac{k}{n}}{1+\left(\cfrac{k}{n}\right)^2}.
\]
従って，連続関数$f(x)=\frac{x}{1+x^2}$の$[0,1]$上の区分求積であり，求める極限は
\[
 \int_0^1\frac{x}{1+x^2}\dd x
 =\left[\frac12\log(1+x^2)\right]_0^1
 =\frac12\log2.
\]

\paragraph{9.2}
$\sqrt{n+k}=\sqrt n\sqrt{1+\frac{k}{n}}$なので，もとの和は
\[
 \frac1n\sum_{k=1}^{n}\sqrt{1+\frac{k}{n}}
\]
となる。$f(x)=\sqrt{1+x}$は$[0,1]$で連続だから，極限は
\[
 \int_0^1\sqrt{1+x}\dd x
 =\left[\frac23(1+x)^{\frac32}\right]_0^1
 =\frac23(2\sqrt2-1).
\]
分子の根号から出る$\sqrt n$と，前にある$n^{\frac32}$を合わせて，区分求積の幅$\frac1n$が現れる。

\paragraph{9.3}
$f(x)=\frac1{x^2}$は$x>0$で減少するので，$k\geqq2$に対して
\[
 \int_k^{k+1}\frac1{x^2}\dd x
 \leqq\frac1{k^2}
 \leqq\int_{k-1}^{k}\frac1{x^2}\dd x.
\]
$k=n+1,\ldots,N$について足せば
\[
 \int_{n+1}^{N+1}\frac1{x^2}\dd x
 \leqq\sum_{k=n+1}^{N}\frac1{k^2}
 \leqq\int_n^N\frac1{x^2}\dd x.
\]
原始関数は$-\frac1x$だから，問題の不等式を得る。

部分和$S_N=\sum_{k=1}^{N}\frac1{k^2}$は単調増加する。一方，$n=1$の右側の不等式を使うと
\[
 S_N\leqq1+1-\frac1N<2\qquad(N\geqq2)
\]
なので，上に有界である。従って有限の極限$S$を持つ。これを無限級数の和とし，$R_n=S-S_n$とおく。最初の不等式で$n$を固定して$N\to\infty$とすると
\[
 \frac1{n+1}\leqq R_n\leqq\frac1n
\]
を得る。収束の証明では，和の正確な値を求める必要はなかった。

\paragraph{9.4}
コーシー・シュワルツの不等式を$f(x)$と$x$に使うと
\[
 1=\left(\int_0^1xf(x)\dd x\right)^2
 \leqq\left(\int_0^1f(x)^2\dd x\right)
       \left(\int_0^1x^2\dd x\right)
 =\frac13\int_0^1f(x)^2\dd x.
\]
よって積分は$3$以上である。等号条件は$f(x)=cx$であり，制約に代入すると
\[
 1=\int_0^1cx^2\dd x=\frac c3
\]
から$c=3$となる。実際$f(x)=3x$なら制約を満たし，$\int_0^1(3x)^2\dd x=3$である。従って最小値は$3$，それを与える関数は$f(x)=3x$である。

\paragraph{9.5}
まず$K_0=\int_1^{\e}1\dd x=\e-1$である。$n\geqq1$では$(\log x)^n$を微分する部分積分によって
\begin{align*}
 K_n
 &=\left[x(\log x)^n\right]_1^{\e}
   -n\int_1^{\e}x(\log x)^{n-1}\frac1x\dd x\\
 &=\e-nK_{n-1}.
\end{align*}
従って
\[
 K_1=1,\qquad K_2=\e-2,\qquad K_3=\e-3(\e-2)=6-2\e.
\]
$n=0$を漸化式に形式的に代入すると$K_0=\e$という誤りが生じる。$x=1$の端点の項が$0$になるという議論は$n\geqq1$に対して行ったものである。

\paragraph{9.6}
本文の$I_n=\int_0^{\frac{\pi}{2}}\sin^n x\dd x$について
\[
 \frac{2m+1}{2m+2}
 \leqq\frac{I_{2m+1}}{I_{2m}}\leqq1,
 \qquad
 P_m=\frac{\pi}{2}\frac{I_{2m+1}}{I_{2m}}
\]
である。従って
\[
 \frac{\pi}{2}\frac{2m+1}{2m+2}
 \leqq P_m\leqq\frac{\pi}{2}.
\]
$\frac{\pi}{2}$から各辺を引き，順序をそろえると
\[
 0\leqq\frac{\pi}{2}-P_m
 \leqq\frac{\pi}{2}\left(1-\frac{2m+1}{2m+2}\right)
 =\frac{\pi}{4m+4}.
\]
次に$u=\frac{\pi}{2}-x$とおけば
\[
 \int_0^{\frac{\pi}{2}}\cos^{2m}x\dd x=I_{2m}.
\]
本文で$I_{2m}I_{2m+1}=\frac{\pi}{2(2m+1)}$および$\frac{I_{2m+1}}{I_{2m}}\to1$を示したので，
\[
 mI_{2m}^2
 =\frac{m\pi}{2(2m+1)}\frac{I_{2m}}{I_{2m+1}}
 \longrightarrow\frac{\pi}{4}.
\]
積分値が正であることを使って平方根を取り，求める極限は$\frac{\sqrt\pi}{2}$となる。

\paragraph{9.7}
整数の場合のベータ積分を使うと
\begin{align*}
 \int_0^1\binom{m+n}{m}x^m(1-x)^n\dd x
 &=\frac{(m+n)!}{m!n!}\frac{m!n!}{(m+n+1)!}\\
 &=\frac1{m+n+1}.
\end{align*}
また$m=2,n=4$の場合は
\[
 \int_0^1x^2(1-x)^4\dd x
 =\frac{2!4!}{7!}=\frac1{105}.
\]

\paragraph{9.8}
$0<x\leqq\frac12$では$1-x\geqq\frac12$だから
\[
 0\leqq\sqrt{\frac{x}{1-x}}\leqq\sqrt2\sqrt x.
\]
右辺は$0$まで積分できる。$\frac12\leqq x<1$では
\[
 0\leqq\sqrt{\frac{x}{1-x}}\leqq\frac1{\sqrt{1-x}},
\]
右辺の積分は$[-2\sqrt{1-x}]_{\frac12}^{1}=\sqrt2$で有限である。各端点に近づけた積分値は単調増加し，これらの有限な値で上から抑えられるので，両端で収束する。

被積分関数は$x^{\frac32-1}(1-x)^{\frac12-1}$なので，積分は$B\left(\frac32,\frac12\right)$である。漸化式より
\[
 B\left(\frac32,\frac12\right)
 =\frac12 B\left(\frac12,\frac12\right)=\frac{\pi}{2}.
\]
置換$x=\sin^2t$でも，開区間では平方根が$\tan t$となり，
\[
 \int_0^1\sqrt{\frac{x}{1-x}}\dd x
 =2\int_0^{\frac{\pi}{2}}\sin^2t\dd t
 =2\cdot\frac{\pi}{4}=\frac{\pi}{2}
\]
と一致する。置換は端点を避けた区間で行い，その後で極限を取る。

\section{第10章の解答}\label{解答:第10章}

\paragraph{10.1}
$\sin\frac{4\pi}3=-\frac{\sqrt3}2$であり，$\arcsin$が返す角は$\left[-\frac\pi2,\frac\pi2\right]$にあるので
\[
\arcsin\left(\sin\frac{4\pi}3\right)=-\frac\pi3.
\]
$\cos\frac{5\pi}4=-\frac{\sqrt2}2$であり，$\arccos$が返す角は$[0,\pi]$にあるので
\[
\arccos\left(\cos\frac{5\pi}4\right)=\frac{3\pi}4.
\]
また$-\frac{2\pi}3$に$\pi$を加えると$\frac\pi3$になり，正接の値は変わらない。この角は$\left(-\frac\pi2,\frac\pi2\right)$の中にあるから
\[
\arctan\left(\tan\left(-\frac{2\pi}3\right)\right)=\frac\pi3.
\]
微分では合成関数の公式を使う。$-1<x<1$なら$0\leqq x^2<1$なので，逆正弦の微分の分母は正であり，
\[
\frac{\dd}{\dd x}\arcsin(x^2)
=\frac{2x}{\sqrt{1-(x^2)^2}}
=\frac{2x}{\sqrt{1-x^4}}
\]
となる。

\paragraph{10.2}
最初の分母を平方完成すると
\[
x^2+2x+5=(x+1)^2+4.
\]
$u=\frac{x+1}{2}$と置けば$\dd x=2\dd u$だから
\[
\int\frac1{x^2+2x+5}\dd x
=\frac12\int\frac1{1+u^2}\dd u
=\frac12\arctan\frac{x+1}{2}+C.
\]
分母は実数全体で正なので，この答えは実数全体で用いることができる。

第二の積分は，$\arcsin\frac x2$の微分が$\frac1{\sqrt{4-x^2}}$であることから
\[
\int_0^1\frac1{\sqrt{4-x^2}}\dd x
=\left[\arcsin\frac x2\right]_0^1
=\frac\pi6.
\]
三角置換$x=2\sin\theta$を使っても，$0\leqq\theta\leqq\frac\pi6$で分母と$\dd x$の係数がともに$2\cos\theta$となり，同じ値を得る。

\paragraph{10.3}
$\sin x\ne0$のところで商を微分すると
\[
(\cot x)'
=\left(\frac{\cos x}{\sin x}\right)'
=\frac{-\sin^2x-\cos^2x}{\sin^2x}
=-\frac1{\sin^2x}=-\csc^2x.
\]
さらに$\sin x\cos x\ne0$なら
\[
\tan x+\cot x
=\frac{\sin^2x+\cos^2x}{\sin x\cos x}
=\frac1{\sin x\cos x}.
\]
それぞれ分母の微分を使って積分すると
\[
\begin{aligned}
\int\frac1{\sin x\cos x}\dd x
&=\int\tan x\dd x+\int\cot x\dd x\\
&=-\log|\cos x|+\log|\sin x|+C\\
&=\log|\tan x|+C.
\end{aligned}
\]
もとの分母は$x=\frac{k\pi}{2}$で$0$になる。従って，整数$k$について$\left(\frac{k\pi}{2},\frac{(k+1)\pi}{2}\right)$のような，これらの点を含まない一つの区間で用いる。任意定数は区間ごとに独立である。

\paragraph{10.4}
$2+\sin x\geqq1$なので，元の被積分関数は$[0,2\pi]$全体で連続である。ただし$t=\tan\frac x2$は$x=\pi$で定義されない。そこで
\[
I_1=\int_0^\pi\frac1{2+\sin x}\dd x,\qquad
I_2=\int_\pi^{2\pi}\frac1{2+\sin x}\dd x
\]
に分ける。$x=\pi$から少し離れた区間で置換してから極限を取ると，本文と同じ変形によって
\[
I_1=\lim_{R\to\infty}\int_0^R\frac1{t^2+t+1}\dd t,\qquad
I_2=\lim_{S\to\infty}\int_{-S}^0\frac1{t^2+t+1}\dd t
\]
となる。$x\to\pi-0$では$t\to+\infty$，$x\to\pi+0$では$t\to-\infty$だからである。二つの極限を独立に求める。

原始関数の一つは
\[
F(t)=\frac2{\sqrt3}\arctan\frac{2t+1}{\sqrt3}.
\]
$\arctan t$の極限は$t\to+\infty$で$\frac\pi2$，$t\to-\infty$で$-\frac\pi2$だから
\[
\begin{aligned}
I_1&=\frac2{\sqrt3}\left(\frac\pi2-\frac\pi6\right)
=\frac{2\pi}{3\sqrt3},\\
I_2&=\frac2{\sqrt3}\left(\frac\pi6+\frac\pi2\right)
=\frac{4\pi}{3\sqrt3}.
\end{aligned}
\]
従って求める値は$I_1+I_2=\frac{2\pi}{\sqrt3}$である。無限の端点が現れたのは置換によるものであり，元の積分に特異点があるわけではない。

\paragraph{10.5}
固定点$(-1,0)$から与えられた点へ引く直線の傾きは
\[
t=\cfrac{\cfrac5{13}}{1+\cfrac{12}{13}}
=\frac5{25}=\frac15.
\]
$m=5,n=1$とすれば
\[
m^2-n^2=24,\qquad2mn=10,\qquad m^2+n^2=26.
\]
三整数の最大公約数$2$で割って，$(12,5,13)$を得る。実際$12^2+5^2=13^2$であり，斜辺で割ると元の円上の点へ戻る。

$m,n$がともに奇数なら，$m^2,n^2$もともに奇数である。従って差$m^2-n^2$と和$m^2+n^2$は偶数であり，$2mn$も偶数である。$t=\frac nm$を既約分数にしても，公式から得る三整数に共通因数がないとは限らない。

\paragraph{10.6}
定義を使って積を計算すると
\[
2\sinh u\cosh u
=2\frac{\e^u-\e^{-u}}2\frac{\e^u+\e^{-u}}2
=\frac{\e^{2u}-\e^{-2u}}2
=\sinh(2u).
\]
また$\cosh^2u=\frac12\{1+\cosh(2u)\}$なので
\[
\int_0^{\log2}\cosh^2u\dd u
=\left[\frac u2+\frac14\sinh(2u)\right]_0^{\log2}.
\]
$\e^{2\log2}=4$，$\e^{-2\log2}=\frac14$より
\[
\sinh(2\log2)=\cfrac{4-\cfrac14}{2}=\frac{15}{8}.
\]
従って積分値は$\frac12\log2+\frac{15}{32}$である。

\paragraph{10.7}
$u>0$で$\cosh u$は$1$から無限大へ狭義単調増加するので，$x>1$と一対一に対応する。この範囲なら$\sinh u>0$であり，
\[
x=\cosh u,\qquad
\sqrt{x^2-1}=\sqrt{\sinh^2u}=\sinh u,\qquad
\dd x=\sinh u\dd u.
\]
従って積分は$\int1\dd u=u+C$になる。元へ戻すには
\[
x+\sqrt{x^2-1}=\cosh u+\sinh u=\e^u
\]
を使えばよい。右辺は正なので
\[
\int\frac1{\sqrt{x^2-1}}\dd x
=\log\bigl(x+\sqrt{x^2-1}\bigr)+C\qquad(x>1).
\]
この問題では$u>0$の枝を選んだため，平方根を$\sinh u$と書けた。$u<0$の枝なら平方根は$-\sinh u$となり，同じ途中式をそのまま使うことはできない。

\paragraph{10.8}
まず$x=2\sinh u$と置く。$x=0$は$u=0$，$x=2$は$u=\operatorname{arsinh}1=\log(1+\sqrt2)$へ移る。後者を$U$と書けば
\[
\sinh U=1,\qquad\cosh U=\sqrt2,\qquad\sinh(2U)=2\sqrt2.
\]
$\sqrt{4+x^2}=2\cosh u$，$\dd x=2\cosh u\dd u$だから
\[
\begin{aligned}
\int_0^2\sqrt{4+x^2}\dd x
&=4\int_0^U\cosh^2u\dd u\\
&=\left[2u+\sinh(2u)\right]_0^U\\
&=2\log(1+\sqrt2)+2\sqrt2.
\end{aligned}
\]

次に$t=x+\sqrt{4+x^2}$と置く。和と差の積が$4$なので
\[
\sqrt{4+x^2}-x=\frac4t.
\]
和と差を取って
\[
x=\frac12\left(t-\frac4t\right),\qquad
\sqrt{4+x^2}=\frac12\left(t+\frac4t\right),\qquad
\dd x=\frac12\left(1+\frac4{t^2}\right)\dd t
\]
を得る。端点は$t=2$と$t=2+2\sqrt2$である。従って
\[
\begin{aligned}
\int_0^2\sqrt{4+x^2}\dd x
&=\frac14\int_2^{2+2\sqrt2}
 \left(t+\frac8t+\frac{16}{t^3}\right)\dd t\\
&=\left[\frac{t^2}{8}+2\log t-\frac2{t^2}\right]_2^{2+2\sqrt2}.
\end{aligned}
\]
べきの項を直接代入してもよいが，
\[
\frac12x\sqrt{4+x^2}=\frac{t^2}{8}-\frac2{t^2}
\]
と見れば，上端で$2\sqrt2$，下端で$0$になる。対数の差は
\[
2\log(2+2\sqrt2)-2\log2=2\log(1+\sqrt2)
\]
なので，第一の方法と一致する。さらに$x=2\sinh u$を代入すれば
\[
t=2\sinh u+2\cosh u=2\e^u
\]
となる。二つの置換は，変数をこの式で換えたものになっている。

\section{第11章の解答}\label{解答:第11章}

\paragraph{11.1}
$f(x)=\log(1+x)$の四階微分は
\[
 f^{(4)}(t)=-\frac6{(1+t)^4}
\]
である。$0\leqq t\leqq x\leqq\frac12$では$-6\leqq f^{(4)}(t)\leqq0$なので，三次の積分剰余は
\[
 -\int_0^x(x-t)^3\dd t
 \leqq R_3(x)=\frac16\int_0^x(x-t)^3f^{(4)}(t)\dd t
 \leqq0.
\]
左辺は$-\frac{x^4}{4}$である。$\log(1+x)=P(x)+R_3(x)$に戻すと
\[
 0\leqq P(x)-\log(1+x)\leqq\frac{x^4}{4}
\]
を得る。$x=\frac1{10}$では
\[
 P\left(\frac1{10}\right)
 =\frac1{10}-\frac1{200}+\frac1{3000}
 =\frac{143}{1500},\qquad
 \frac{x^4}{4}=\frac1{40000}.
\]
したがって，求める区間の一つは
\[
 \boxed{\frac{143}{1500}-\frac1{40000}
 \leqq\log\frac{11}{10}\leqq\frac{143}{1500}}
\]
である。両端は有理数であり，区間の幅はちょうど$\frac1{40000}$である。

\paragraph{11.2}
$g(u)=\e^{-u}$について$g'''(u)=-\e^{-u}$なので，$u\geqq0$では二次の積分剰余が負になる。さらに$|g'''(u)|\leqq1$より
\[
 \e^{-u}=1-u+\frac{u^2}{2}+r_2(u),
 \qquad -\frac{u^3}{6}\leqq r_2(u)\leqq0.
\]
$u=x^2$を代入して$0$から$\frac12$まで積分すると，近似値は
\[
 Q=\int_0^{\frac12}\left(1-x^2+\frac{x^4}{2}\right)\dd x
 =\frac12-\frac1{24}+\frac1{320}
 =\frac{443}{960}.
\]
真の値を$A$とすれば
\[
 0\leqq Q-A
 \leqq\int_0^{\frac12}\frac{x^6}{6}\dd x
 =\frac1{42}\left(\frac12\right)^7
 =\frac1{5376}.
\]
したがって$\frac{443}{960}$は真の値を上から近似し，誤差は$\frac1{5376}$以下である。

\paragraph{11.3}
分割幅は$h=\frac12$であり，分点$0,\frac12,1,\frac32,2$での関数値は
$0,\frac1{16},1,\frac{81}{16},16$である。これをそれぞれの公式に代入すると
\[
 \begin{aligned}
 T_4&=\frac12\left(8+\frac1{16}+1+\frac{81}{16}\right)
 =\frac{113}{16},\\
 S_4&=\frac16\left\{16+4\left(\frac1{16}+\frac{81}{16}\right)+2\right\}
 =\frac{77}{12}.
 \end{aligned}
\]
真の値は$A=\left[\frac{x^5}{5}\right]_0^2=\frac{32}{5}$なので
\[
 T_4-A=\frac{53}{80},\qquad S_4-A=\frac1{60}.
\]
いずれも過大な近似である。$[0,2]$では$f''(x)=12x^2\leqq48$，$f^{(4)}(x)=24$だから，誤差上限は
\[
 |A-T_4|\leqq\frac{48\cdot2^3}{12\cdot4^2}=2,
 \qquad
 |A-S_4|\leqq\frac{24\cdot2^5}{180\cdot4^4}=\frac1{60}.
\]
台形公式では実際の誤差は上限より小さい。シンプソン公式では四階微分が一定なので，誤差を表す積分で絶対値を評価しても余分な差が生じず，上限と実際の誤差が一致する。

\paragraph{11.4}
(a) 分子，分母とも$0$へ近づき，$x\ne0$では分母の微分$2x$は$0$でない。左右それぞれでロピタルの定理を使う準備をする。微分後の商も$0$と$0$の形で，その分母の微分は$2\ne0$だから，もう一度適用できる。
\[
 \lim_{x\to0}\frac{\e^x-1}{2x}
 =\lim_{x\to0}\frac{\e^x}{2}=\frac12.
\]
この極限が存在するので，最初の商にも定理を適用でき，答えは$\frac12$である。

(b) $x>0$で両関数は微分可能であり，$\alpha>0$より
$\log x\to\infty$，$x^\alpha\to\infty$である。また分母の微分は$\alpha x^{\alpha-1}>0$なので
\[
 \frac{(\log x)'}{(x^\alpha)'}
 =\cfrac{\cfrac1x}{\alpha x^{\alpha-1}}
 =\frac1{\alpha x^\alpha}\longrightarrow0.
\]
よって求める極限は$0$である。

\paragraph{11.5}
$f(x)=\frac1{x^p(1-x)^q}$は$0<x<1$で正である。まず左端では
\[
 \frac{f(x)}{x^{-p}}=(1-x)^{-q}\longrightarrow1
 \qquad(x\to0+).
\]
極限比較により，左端で収束する条件は$\int_0^{\frac12}x^{-p}\dd x$と同じで，$p<1$である。
右端では
\[
 \frac{f(x)}{(1-x)^{-q}}=x^{-p}\longrightarrow1
 \qquad(x\to1-).
\]
$u=1-x$と置けば基準となる積分は$\int_0^{\frac12}u^{-q}\dd u$なので，右端の条件は$q<1$である。
両方の極限が独立に有限でなければならないから，必要十分条件は
\[
 \boxed{p<1\quad\text{かつ}\quad q<1}
\]
である。どちらかが満たされなければ，その端点の積分は$+\infty$へ発散する。

\paragraph{11.6}
まず有限の$R>\e$で$u=\log x$と置く。$\dd u=\frac1x\dd x$，下端は$u=1$，上端は$u=\log R$だから
\[
 \int_{\e}^R\frac1{x(\log x)^q}\dd x
 =\int_1^{\log R}u^{-q}\dd u.
\]
$R\to\infty$なら$\log R\to\infty$である。従って，これは無限大におけるべき型積分そのもので，収束条件は$q>1$である。その場合は
\[
 \int_{\e}^{\infty}\frac1{x(\log x)^q}\dd x
 =\lim_{R\to\infty}\frac{(\log R)^{1-q}-1}{1-q}
 =\boxed{\frac1{q-1}}.
\]
$q=1$では$\log\log R$，$q<1$では$(\log R)^{1-q}$が無限大に近づくため発散する。

\paragraph{11.7}
最初の積分では，左右を別に計算すると
\[
 \int_{-1}^{-\varepsilon}\frac1x\dd x=\log\varepsilon\to-\infty,
 \qquad
 \int_{\delta}^{1}\frac1x\dd x=-\log\delta\to+\infty.
\]
従って広義積分は発散する。一方，$\delta=\varepsilon$という対称な切り取り方では二つの和が常に$0$なので，対称的な主値は$0$である。

二番目の被積分関数を$g(x)=\frac{1-\cos x}{x}$とする。$1-\cos x=2\sin^2\frac x2$と$|\sin t|\leqq|t|$から
\[
 |g(x)|\leqq\frac{|x|}{2}\qquad(x\ne0)
\]
を得る。従って，左右それぞれで絶対値の積分が収束する。また$g(x)\to0$なので，$g(0)=0$と補えば連続関数になる。この$g$は奇関数だから
\[
 \int_{-1}^1\frac{1-\cos x}{x}\dd x=0.
\]
後者では各側が有限に収束するため，奇関数の対称性を使える。前者の$0$は，発散する二つの側を特定の切り取り方で相殺して得た主値である。

\paragraph{11.8}
被積分関数は偶関数である。$t=-u$による置換で
\[
 F(-x)=\int_0^{-x}\e^{-t^2}\dd t
 =-\int_0^x\e^{-u^2}\dd u=-F(x)
\]
となるので，$F$は奇関数である。基本定理より
\[
 F'(x)=\e^{-x^2}>0,\qquad F''(x)=-2x\e^{-x^2}.
\]
従って実数全体で狭義に単調増加し，$x<0$では下に凸，$x>0$では上に凸である。原点で凹凸が変わる。

$u\geqq0$に対して$\e^{-u}\leqq1$である。また$G(u)=\e^{-u}-1+u$と置くと，$G(0)=0$，$G'(u)=1-\e^{-u}\geqq0$なので$G(u)\geqq0$である。従って
\[
 1-u\leqq\e^{-u}\leqq1.
\]
$u=t^2$を代入し，$x\geqq0$として$0$から$x$まで積分すれば
\[
 \boxed{x-\frac{x^3}{3}\leqq F(x)\leqq x}
\]
を得る。

この上限$x$だけでは無限大での有界性は示せない。そこで$t\geqq1$では$t^2\geqq t$となることを使う。$x\geqq1$なら
\[
 F(x)=\int_0^1\e^{-t^2}\dd t+\int_1^x\e^{-t^2}\dd t
 \leqq1+\int_1^x\e^{-t}\dd t
 \leqq1+\e^{-1}.
\]
$F$は増加し，この一定値を超えないので，$x\to\infty$で有限の極限をもつ。

\section{第12章の解答}\label{解答:第12章}

\paragraph{12.1}
分母の次数より分子の次数が高いので，先に多項式の割り算をする。
\[
 x^3+x^2+2x+1=(x+1)(1+x^2)+x
\]
だから
\[
 \begin{aligned}
 \int_0^1\frac{x^3+x^2+2x+1}{1+x^2}\dd x
 &=\int_0^1\left(x+1+\frac{x}{1+x^2}\right)\dd x\\
 &=\left[\frac{x^2}{2}+x+\frac12\log(1+x^2)\right]_0^1\\
 &=\boxed{\frac32+\frac12\log2}.
 \end{aligned}
\]
余りの分子$x$が分母の微分$2x$の定数倍になったため，対数の形で積分できた。

\paragraph{12.2}
$u=1+x^2$と置くと$\dd u=2x\dd x$であり，$x=0,1$はそれぞれ$u=1,2$に対応する。従って
\[
 \int_0^1x\sqrt{1+x^2}\dd x
 =\frac12\int_1^2u^{\frac12}\dd u
 =\left[\frac13u^{\frac32}\right]_1^2
 =\boxed{\frac{2\sqrt2-1}{3}}.
\]
積分変数を$u$へ変えた時点で，端点も$u$の値へ変えている。

\paragraph{12.3}
$\log x$を微分し，$x$を積分すると，残りは整式の積分になる。部分積分により
\[
 \begin{aligned}
 \int_1^{\e}x\log x\dd x
 &=\left[\frac{x^2}{2}\log x\right]_1^{\e}
   -\frac12\int_1^{\e}x\dd x\\
 &=\frac{\e^2}{2}-\frac{\e^2-1}{4}
 =\boxed{\frac{\e^2+1}{4}}.
 \end{aligned}
\]

\paragraph{12.4}
$G(u)=\int_0^u\sqrt{1+t^2}\dd t$と置く。基本定理より$G'(u)=\sqrt{1+u^2}$であり，$F(x)=G(x^2)-G(x)$だから，合成関数の微分によって
\[
 F'(x)=2x\sqrt{1+x^4}-\sqrt{1+x^2}.
\]
上端だけでなく，下端からも項が現れる。$x=1$では上下端が一致するため$F(1)=0$，また$F'(1)=\sqrt2$である。従って接線は
\[
 \boxed{y=\sqrt2(x-1)}.
\]

\paragraph{12.5}
求める積分を$I$とする。$x=-t$と置き，$1+x^2$が偶関数であることを使えば
\[
 I=\int_{-a}^{a}\frac{1+t^2}{1+\e^{-t}}\dd t.
\]
元の式と足すと，
\[
 \frac1{1+\e^x}+\frac1{1+\e^{-x}}=1
\]
によって重みが消える。従って
\[
 2I=\int_{-a}^{a}(1+x^2)\dd x
 =2a+\frac{2a^3}{3},\qquad
 \boxed{I=a+\frac{a^3}{3}}.
\]
不定積分を求めずに値が決まるのは，区間と分子が反転に対して変わらないためである。

\paragraph{12.6}
交点は$x^2=x$より$x=0,1$である。$0\leqq x\leqq1$では$x^2\leqq x$なので，面積$S$は
\[
 S=\int_0^1(x-x^2)\dd x
 =\left[\frac{x^2}{2}-\frac{x^3}{3}\right]_0^1
 =\boxed{\frac16}.
\]
回転すると，$x$を固定した断面は外半径$x$，内半径$x^2$の円環になる。従って体積$V$は
\[
 V=\pi\int_0^1\{x^2-(x^2)^2\}\dd x
 =\pi\left(\frac13-\frac15\right)
 =\boxed{\frac{2\pi}{15}}.
\]
面積では上下の高さの差を積分し，体積では二つの半径の平方の差を積分している。

\paragraph{12.7}
積分の両端を微分すると
\[
 \begin{aligned}
 H'(a)&=\{(a+1)^2-2(a+1)\}^2-(a^2-2a)^2\\
 &=(a^2-1)^2-(a^2-2a)^2\\
 &=(2a-1)(2a^2-2a-1).
 \end{aligned}
\]
従って極値の候補は$a=\frac12,\frac{1-\sqrt3}{2},\frac{1+\sqrt3}{2}$である。候補の値を比べ，実数全体での最小値を確かめるため，積分区間の中心を使って式を整理する。

$y=a-\frac12$とし，$x=a+\frac12+s$と置くと，$-\frac12\leqq s\leqq\frac12$，$x-1=y+s$である。$x^2-2x=(x-1)^2-1$なので
\[
 H(a)=\int_{-\frac12}^{\frac12}\{(y+s)^2-1\}^2\dd s.
\]
展開後の$s$の奇数次の項は積分すると消える。残る項は
\[
 \begin{aligned}
 H(a)
 &=\int_{-\frac12}^{\frac12}
   \{s^4+(6y^2-2)s^2+(y^2-1)^2\}\dd s\\
 &=\frac1{80}+\frac{6y^2-2}{12}+(y^2-1)^2\\
 &=y^4-\frac32y^2+\frac{203}{240}
 =\left(y^2-\frac34\right)^2+\frac{17}{60}.
 \end{aligned}
\]
最後の平方は$0$以上で，$y=\pm\frac{\sqrt3}{2}$のときに限り$0$になる。従って
\[
 \boxed{\min H(a)=\frac{17}{60},\qquad
 a=\frac{1-\sqrt3}{2},\ \frac{1+\sqrt3}{2}}.
\]
二つの区間が同じ最小値を与えるのは，被積分関数が$x=1$に関して対称であることにも対応する。

\paragraph{12.8}
対数の中を$u=1+x^2$と置くと，$\dd u=2x\dd x$，$x^2=u-1$だから
\[
 x^3\dd x=\frac12(u-1)\dd u.
\]
区間は$[1,2]$となり，求める積分は$\frac12\int_1^2(u-1)\log u\dd u$である。ここで$\log u$を微分する部分積分を行う。
\[
 \begin{aligned}
 \frac12\int_1^2(u-1)\log u\dd u
 &=\frac12\left[\left(\frac{u^2}{2}-u\right)\log u\right]_1^2
   -\frac12\int_1^2\left(\frac u2-1\right)\dd u\\
 &=\left[-\frac{u^2}{8}+\frac u2\right]_1^2
 =\frac12-\frac38
 =\boxed{\frac18}.
 \end{aligned}
\]
端点の項は，$u=1$では$\log1=0$，$u=2$では$\frac{u^2}{2}-u=0$となるため，両方とも消える。

\paragraph{12.9}
一項を
\[
 \frac{k^2}{n(n^2+k^2)}
 =\frac1n\cfrac{\left(\cfrac kn\right)^2}{1+\left(\cfrac kn\right)^2}
\]
と書く。$f(x)=\frac{x^2}{1+x^2}$は$[0,1]$で連続なので，この和の極限はその区間での積分である。分子を$x^2=(1+x^2)-1$と分けると
\[
 \begin{aligned}
 \int_0^1\frac{x^2}{1+x^2}\dd x
 &=\int_0^1\left(1-\frac1{1+x^2}\right)\dd x\\
 &=\left[x-\arctan x\right]_0^1
 =\boxed{1-\frac{\pi}{4}}.
 \end{aligned}
\]

\paragraph{12.10}
各$x\in[0,1]$で，縦の線分は$y=x-1$から$y=x$まであり，常に$x$軸を含む。従って回転断面の中心には穴がなく，円板になる。端点の軸からの距離はそれぞれ$1-x$，$x$なので，半径$r(x)$は
\[
 r(x)=
 \begin{cases}
  1-x&\left(0\leqq x\leqq\frac12\right),\\
  x&\left(\frac12\leqq x\leqq1\right).
 \end{cases}
\]
切り替わりは$x=1-x$，すなわち$x=\frac12$で起こる。従って
\[
 \begin{aligned}
 V&=\pi\int_0^{\frac12}(1-x)^2\dd x
   +\pi\int_{\frac12}^{1}x^2\dd x\\
 &=2\pi\int_{\frac12}^{1}x^2\dd x
 =\frac{2\pi}{3}\left(1-\frac18\right)
 =\boxed{\frac{7\pi}{12}}.
 \end{aligned}
\]
上下の部分を別々に回して体積を足すと，同じ円板内の点を重複して数える。ここでは回転後にできる断面そのものを一枚ずつ数えた。

\paragraph{12.11}
まず
\[
 K_0=\int_0^1\frac1{1+x}\dd x=\log2.
\]
$n\geqq1$では，隣り合う二つの積分を足すと分母が消える。
\[
 K_n+K_{n-1}
 =\int_0^1\frac{x^n+x^{n-1}}{1+x}\dd x
 =\int_0^1x^{n-1}\dd x=\frac1n.
\]
従って漸化式は$K_n=\frac1n-K_{n-1}$であり，順に
\[
 K_1=1-\log2,\qquad
 K_2=\log2-\frac12,\qquad
 \boxed{K_3=\frac56-\log2}
\]
となる。

また$0\leqq x\leqq1$では$0\leqq\frac{x^n}{1+x}\leqq x^n$だから
\[
 0\leqq K_n\leqq\frac1{n+1}\longrightarrow0.
\]
ここで漸化式の添字を$k$とし，両辺に$(-1)^{k-1}$を掛けて$k=1,\ldots,n$について足す。
中間の$K_k$が消え，
\[
 \sum_{k=1}^n\frac{(-1)^{k-1}}k
 =K_0+(-1)^{n-1}K_n
 =\log2+(-1)^{n-1}K_n
\]
となる。右辺の最後の項は$0$へ近づくので，求める極限は$\log2$である。同時に，部分和と$\log2$との差の絶対値が$\frac1{n+1}$以下であることも分かる。

\paragraph{12.12}
積分を$x=1$で分け，$0$と無限大を独立に調べる。$x\to0+$では
\[
 \cfrac{\cfrac{x^\alpha}{1+x^2}}{x^\alpha}
 =\frac1{1+x^2}\longrightarrow1.
\]
従って$0$付近では$x^\alpha$との極限比較が使え，収束条件は$\alpha>-1$である。
一方，$x\to\infty$では
\[
 \cfrac{\cfrac{x^\alpha}{1+x^2}}{x^{\alpha-2}}
 =\frac{x^2}{1+x^2}\longrightarrow1.
\]
無限大で$x^{\alpha-2}$が積分可能になる条件は$\alpha-2<-1$，すなわち$\alpha<1$である。
従って全体の収束条件は
\[
 \boxed{-1<\alpha<1}.
\]
$\alpha=0$では，端点を避けて計算してから極限を取り，
\[
 \int_0^{\infty}\frac1{1+x^2}\dd x
 =\lim_{R\to\infty}\arctan R-\lim_{\varepsilon\to0+}\arctan\varepsilon
 =\boxed{\frac{\pi}{2}}.
\]

\paragraph{12.13}
最初から$f$が微分可能だと仮定せず，右辺の積分から確かめる。
\[
 A(x)=\int_0^xf(t)\dd t,\qquad
 B(x)=\int_0^xtf(t)\dd t
\]
と置くと，$f$が連続であるため$A'(x)=f(x)$，$B'(x)=xf(x)$である。与式は
\[
 f(x)=1+xA(x)-B(x)
\]
と書けるので，右辺が微分可能であり，
\[
 f'(x)=A(x)+xf(x)-xf(x)=A(x)
\]
を得る。さらに$f''(x)=f(x)$である。また$x=0$を代入すると
\[
 f(0)=1,\qquad f'(0)=0.
\]

ここからは，$f''=f$を満たす関数を積の微分によって求める。
\[
 \frac{\dd}{\dd x}\{\e^{-x}(f(x)+f'(x))\}
 =\e^{-x}(f''(x)-f(x))=0
\]
だから，括弧を含む積は定数である。$x=0$での値は$1$なので
$f'(x)+f(x)=\e^x$となる。今度は$\e^xf(x)$を微分すれば
\[
 \frac{\dd}{\dd x}\{\e^xf(x)\}=\e^x(f'(x)+f(x))=\e^{2x}.
\]
$0$から$x$まで積分し，$f(0)=1$を使うと
\[
 \e^xf(x)-1=\frac{\e^{2x}-1}{2},\qquad
 \boxed{f(x)=\frac{\e^x+\e^{-x}}2}.
\]

微分して得た条件だけでは，元の積分の等式を満たすかの確認が残る。
求めた関数を$g(x)=\frac{\e^x+\e^{-x}}2$と書くと，$g''=g$，$g(0)=1$，$g'(0)=0$である。従って部分積分によって
\[
 \begin{aligned}
 \int_0^x(x-t)g(t)\dd t
 &=\int_0^x(x-t)g''(t)\dd t\\
 &=\bigl[(x-t)g'(t)\bigr]_0^x+\int_0^xg'(t)\dd t\\
 &=-xg'(0)+g(x)-g(0)=g(x)-1.
 \end{aligned}
\]
よって元の等式も成立する。また，条件を満たすどの$f$も上の計算から同じ式になるため，解は一つに定まる。

\paragraph{12.14}
二つの制約を満たす一次関数を先に探す。$\ell(x)=a+bx$を代入すると
\[
 a+\frac b2=1,\qquad \frac a2+\frac b3=0.
\]
これを解くと$a=4$，$b=-6$なので，$\ell(x)=4-6x$である。
任意の条件を満たす$f$について$g(x)=f(x)-\ell(x)$と置けば
\[
 \int_0^1g(x)\dd x=0,\qquad
 \int_0^1xg(x)\dd x=0.
\]
従って
\[
 \int_0^1\ell(x)g(x)\dd x
 =4\int_0^1g(x)\dd x-6\int_0^1xg(x)\dd x=0.
\]
この関係を平方の展開に使うと
\[
 \begin{aligned}
 \int_0^1f(x)^2\dd x
 &=\int_0^1\ell(x)^2\dd x
   +2\int_0^1\ell(x)g(x)\dd x+\int_0^1g(x)^2\dd x\\
 &=\int_0^1(16-48x+36x^2)\dd x+\int_0^1g(x)^2\dd x\\
 &=4+\int_0^1g(x)^2\dd x\geqq4.
 \end{aligned}
\]
$g$は連続なので，最後の積分が$0$になるのは$g(x)=0$が全区間で成り立つ場合に限る。
実際，ある点で$g\ne0$なら，その近くでは$g^2$が正の一定値以上になり，積分も正になる。従って
\[
 \boxed{\min\int_0^1f(x)^2\dd x=4,\qquad f(x)=4-6x}.
\]
制約から候補を見つけた後で，どの関数との差も平方の積分として表したことが，最小値の証明になっている。

\paragraph{12.15}
微分すると
\[
 x'(t)=1-\cos2t=2\sin^2t,\qquad
 y'(t)=2\sin t\cos t.
\]
$0<t<\pi$では$x'(t)>0$なので，$x$は$0$から$\pi$まで狭義に増加する。また$y\geqq0$であり，両端で$y=0$となる。このため曲線は各$x$に対して一つの高さを持ち，曲線の下の領域を$x$方向に積分して重複することはない。

面積$S$は
\[
 S=\int_0^\pi y(t)x'(t)\dd t
 =2\int_0^\pi\sin^4t\dd t.
\]
反転$t\mapsto\pi-t$で両半分の積分は等しい。ウォリスの公式より
\[
 \int_0^{\frac\pi2}\sin^4t\dd t
 =\frac34\frac12\frac\pi2=\frac{3\pi}{16}
\]
なので，
\[
 \boxed{S=4\cdot\frac{3\pi}{16}=\frac{3\pi}{4}}.
\]

曲線の長さ$L$では，根号の中を先にまとめる。
\[
 \begin{aligned}
 \sqrt{x'(t)^2+y'(t)^2}
 &=\sqrt{4\sin^4t+4\sin^2t\cos^2t}\\
 &=2|\sin t|=2\sin t\qquad(0\leqq t\leqq\pi).
 \end{aligned}
\]
従って
\[
 \boxed{L=\int_0^\pi2\sin t\dd t=4}.
\]

回転体の断面は半径$y(t)$の円板であり，$x$の増分は$x'(t)\dd t$である。従って体積は
\[
 \begin{aligned}
 V&=\pi\int_0^\pi y(t)^2x'(t)\dd t
 =2\pi\int_0^\pi\sin^6t\dd t\\
 &=4\pi\int_0^{\frac\pi2}\sin^6t\dd t
 =4\pi\cdot\frac{5\pi}{32}
 =\boxed{\frac{5\pi^2}{8}}.
 \end{aligned}
\]
ここでも$x(t)$の単調性により，同じ$x$での円板を二回数えていない。

\paragraph{12.16}
(1) $g(u)=\sin u$の$0$を中心とする六次テイラー多項式は
\[
 P_6(u)=u-\frac{u^3}{6}+\frac{u^5}{120}
\]
である。偶数次の係数は$0$だが，六次までの有限テイラー公式として剰余を評価する。
七階微分は$g^{(7)}(u)=-\cos u$なので，$0\leqq u\leqq1$では
$-1\leqq g^{(7)}(u)\leqq0$である。従って
\[
 \sin u=P_6(u)+R_6(u),\qquad
 -\frac{u^7}{5040}\leqq R_6(u)\leqq0.
\]
$u=x^2$を代入して積分した多項式の値は
\[
 Q=\int_0^1\left(x^2-\frac{x^6}{6}+\frac{x^{10}}{120}\right)\dd x
 =\frac13-\frac1{42}+\frac1{1320}
 =\frac{2867}{9240}.
\]
残りの積分は
\[
 -\frac1{5040}\int_0^1x^{14}\dd x
 \leqq A-Q\leqq0
\]
だから
\[
 \boxed{\frac{2867}{9240}-\frac1{75600}
 \leqq A\leqq\frac{2867}{9240}}.
\]
両端は有理数であり，幅$\frac1{75600}$は$\frac1{10000}$より小さい。有限の式に対する区間全体の誤差を積分したので，無限級数と積分の交換は用いていない。

(2) 有限の$R>1$に対して$u=x^2$と置けば
\[
 \int_1^R\sin(x^2)\dd x
 =\frac12\int_1^{R^2}\frac{\sin u}{\sqrt u}\dd u.
\]
右辺で部分積分を行う。$T>1$として
\[
 \begin{aligned}
 \int_1^T\frac{\sin u}{\sqrt u}\dd u
 &=\left[-\frac{\cos u}{\sqrt u}\right]_1^T
   -\frac12\int_1^T\cfrac{\cos u}{u^{\cfrac32}}\dd u\\
 &=\cos1-\frac{\cos T}{\sqrt T}
   -\frac12\int_1^T\cfrac{\cos u}{u^{\cfrac32}}\dd u.
 \end{aligned}
\]
境界項$\frac{\cos T}{\sqrt T}$は$0$へ近づく。また
$\cfrac{|\cos u|}{u^{\cfrac32}}\leqq u^{-\frac32}$であり，右辺のべき関数は無限大で積分可能だから，最後の積分は絶対収束する。従って元の広義積分は収束する。

絶対値を付けた場合も同じ置換が使え，
\[
 \int_1^R|\sin(x^2)|\dd x
 =\frac12\int_1^{R^2}\frac{|\sin u|}{\sqrt u}\dd u.
\]
整数$k\geqq1$について，$u\in[k\pi,(k+1)\pi]$なら
$\frac1{\sqrt u}\geqq\frac1{\sqrt{(k+1)\pi}}$なので
\[
 \int_{k\pi}^{(k+1)\pi}\frac{|\sin u|}{\sqrt u}\dd u
 \geqq\frac1{\sqrt{(k+1)\pi}}
       \int_{k\pi}^{(k+1)\pi}|\sin u|\dd u
 =\frac2{\sqrt{(k+1)\pi}}.
\]
この下限の和は発散する。実際，$t^{-\frac12}$の減少性から
\[
 \sum_{k=1}^N\frac1{\sqrt{k+1}}
 \geqq\int_2^{N+2}\frac1{\sqrt t}\dd t
 =2(\sqrt{N+2}-\sqrt2)\longrightarrow\infty.
\]
従って絶対値の積分は発散し，$\int_1^{\infty}\sin(x^2)\dd x$は条件収束する。

\paragraph{12.17}
$n=0$では$C_0=1$である。$n\geqq1$で，$(1-x^2)^n$を微分し，$1$を積分する部分積分を行う。
\[
 \begin{aligned}
 C_n
 &=\left[x(1-x^2)^n\right]_0^1
   +2n\int_0^1x^2(1-x^2)^{n-1}\dd x\\
 &=2n\int_0^1\{1-(1-x^2)\}(1-x^2)^{n-1}\dd x\\
 &=2n(C_{n-1}-C_n).
 \end{aligned}
\]
従って漸化式と有限積の表示は
\[
 \boxed{C_n=\frac{2n}{2n+1}C_{n-1},\qquad
 C_n=\prod_{k=1}^n\frac{2k}{2k+1}}.
\]

極限を求めるため，$x=\sin t$と置くと
\[
 C_n=\int_0^{\frac\pi2}\cos^{2n+1}t\dd t.
\]
$W_m=\int_0^{\frac\pi2}\cos^mt\dd t$と書けば$C_n=W_{2n+1}$である。
$t\mapsto\frac\pi2-t$により$W_m$は$\sin$のべきのウォリス積分と等しい。
従って偶数，奇数の積の公式から
\[
 W_{2n}W_{2n+1}=\frac{\pi}{2(2n+1)}.
\]
また$0\leqq\cos t\leqq1$なので$W_{2n+2}\leqq W_{2n+1}\leqq W_{2n}$であり，漸化式を使うと
\[
 \frac{2n+1}{2n+2}
 =\frac{W_{2n+2}}{W_{2n}}
 \leqq\frac{W_{2n+1}}{W_{2n}}\leqq1.
\]
従って隣り合う二つの比は$1$へ近づく。積と比を合わせて
\[
 nC_n^2
 =nW_{2n}W_{2n+1}\frac{W_{2n+1}}{W_{2n}}
 \longrightarrow\frac\pi4.
\]
$C_n>0$なので平方根を取り，
\[
 \boxed{\lim_{n\to\infty}\sqrt n\,C_n=\frac{\sqrt\pi}{2}}.
\]
最後に漸化式から
\[
 C_n-C_{n+1}
 =\left(1-\frac{2n+2}{2n+3}\right)C_n
 =\frac{C_n}{2n+3}.
\]
元の数列の極限をこの式へ戻せば
\[
 \boxed{\lim_{n\to\infty}n^{\frac32}(C_n-C_{n+1})
 =\lim_{n\to\infty}\frac{n}{2n+3}\sqrt n\,C_n
 =\frac{\sqrt\pi}{4}}.
\]

\paragraph{12.18}
まず両端の収束を調べる。$-\pi<x<\pi$では$\cos\frac x2>0$だから，半角公式により
\[
 \begin{aligned}
 h_n(x):=\frac{\sin^{2n}x}{\sqrt{1+\cos x}}
 &=\frac{2^{2n}}{\sqrt2}
   \sin^{2n}\frac x2\,\cos^{2n-1}\frac x2.
 \end{aligned}
\]
$n\geqq1$より$2n-1\geqq1$である。$x\to\pi-$，$x\to-\pi+$のそれぞれで$\cos\frac x2\to0$となり，$\sin^{2n}\frac x2$は有界なので，$h_n(x)\to0$となる。
さらに$\frac1{1+\e^x}$は両端を含めて連続である。従って元の被積分関数は，各端点で値$0$を補えば連続になる。左右の広義積分はそれぞれ有限に収束し，ここからは対称性や置換を使ってよい。

$h_n$は偶関数である。$w(x)=\frac1{1+\e^x}$と置けば$w(x)+w(-x)=1$だから，$x=-t$による反転で得た式を元の積分と足して
\[
 \begin{aligned}
 2J_n
 &=\int_{-\pi}^{\pi}h_n(x)\{w(x)+w(-x)\}\dd x\\
 &=\int_{-\pi}^{\pi}h_n(x)\dd x
 =2\int_0^\pi h_n(x)\dd x.
 \end{aligned}
\]
従って指数関数を含む重みは消え，$J_n=\int_0^\pi h_n(x)\dd x$となる。

先ほどの半角公式を使い，$t=\frac x2$と置くと
\[
 J_n=2^{2n}\sqrt2\int_0^{\frac\pi2}
          \sin^{2n}t\cos^{2n-1}t\dd t.
\]
次に$u=\sin t$と置く。$\dd u=\cos t\dd t$なので，$\cos$の残りの偶数べきは$(1-u^2)^{n-1}$になり，
\[
 J_n=2^{2n}\sqrt2\int_0^1u^{2n}(1-u^2)^{n-1}\dd u.
\]
最後の積分は多項式の積分である。しかし，展開すると項数が$n$個になるため，ここでは$(1-u^2)$の指数を部分積分で一つずつ下げる。

非負整数$r,s$に対して
\[
 L_{r,s}=\int_0^1u^{2r}(1-u^2)^s\dd u
\]
と置く。$s\geqq1$なら
\[
 \begin{aligned}
 L_{r,s}
 &=\left[\frac{u^{2r+1}}{2r+1}(1-u^2)^s\right]_0^1
   +\frac{2s}{2r+1}\int_0^1u^{2r+2}(1-u^2)^{s-1}\dd u\\
 &=\frac{2s}{2r+1}L_{r+1,s-1}.
 \end{aligned}
\]
端点の項は，$u=0$では$u^{2r+1}=0$，$u=1$では$(1-u^2)^s=0$なので消える。$s=0$の初期値は別に
\[
 L_{r,0}=\int_0^1u^{2r}\dd u=\frac1{2r+1}
\]
である。漸化式を$s$回使い，最後にこの初期値へ戻すと
\[
 L_{r,s}
 =\frac{2^s s!}{(2r+1)(2r+3)\cdots(2r+2s+1)}.
\]
$r=n$，$s=n-1$を代入し，前の係数$2^{2n}\sqrt2$を掛けることで
\[
 \boxed{J_n=
 \frac{2^{3n-1}\sqrt2\,(n-1)!}
 {(2n+1)(2n+3)\cdots(4n-1)}}.
\]
分母には$n$個の奇数因子がある。例えば$n=1$では分母は$3$だけになり，$J_1=\frac{4\sqrt2}{3}$となる。
端点の収束を確かめた後，対称性で指数関数を除き，半角公式と置換で整数べきの積分へ直した。最後の有限積は，その多項式の積分を部分積分で順に結んだ結果である。

\chapter{公式と解法の参照表}\label{章:参照表}
この付録は，本文で導いた結果を探すためのものである。公式の条件や選び方に迷ったときは，対応する章へ戻ってほしい。特に，積分区間内に分母が零になる点がある場合には，原始関数への代入だけでは積分を計算できない。

\section{微分と積分を結ぶ公式}
$f$が$[a,b]$で連続，$F'=f$とすると，
\[
 \int_a^b f(x)\dd x=F(b)-F(a),\qquad
 \frac{\dd}{\dd x}\int_a^x f(t)\dd t=f(x).
\]
変数の上端・下端が微分可能で，必要な範囲で$f$が連続なら，
\[
 \frac{\dd}{\dd x}\int_{u(x)}^{v(x)}f(t)\dd t
 =f(v(x))v'(x)-f(u(x))u'(x).
\]
これらは\chapterref{章:定積分}の基本定理と合成関数の微分法から得られる。被積分関数に積分変数$t$とは別に$x$も入るときには，上の最後の式だけでは不十分である。

\section{基本的な不定積分}
以下は，各被積分関数が定義される一つの区間上で用いる。$C$はその区間での任意定数である。$a>0$，$b\ne0$とする。
\begin{longtable}{@{}ll@{}}
\toprule
被積分関数 & 原始関数の一つ\\
\midrule
\endhead
$x^p\ (p\ne-1)$ & $\frac{x^{p+1}}{p+1}$\\[4pt]
$\frac1x$ & $\log|x|$\\[4pt]
$\e^{bx}$ & $\frac1b\e^{bx}$\\[4pt]
$\sin bx$ & $-\frac1b\cos bx$\\[4pt]
$\cos bx$ & $\frac1b\sin bx$\\[4pt]
$\frac1{\cos^2x}$ & $\tan x$\\[4pt]
$\frac{f'(x)}{f(x)}\quad(f(x)\ne0)$ & $\log|f(x)|$\\[4pt]
$\frac1{a^2+x^2}$ & $\frac1a\arctan\frac xa$\\[4pt]
$\frac1{\sqrt{a^2-x^2}}\quad(|x|<a)$ & $\arcsin\frac xa$\\[4pt]
\bottomrule
\end{longtable}
一般の実数$p$では，まず$x>0$でべき関数を考える。逆三角関数の記号と範囲は\chapterref{章:置換の背景}で定義した。高校の基本計算では，対応する三角置換を使って定積分を求めることもできる。

\section{置換と部分積分}
$f$が連続，$g$が連続微分可能であるとき，
\[
 \int_\alpha^\beta f(g(t))g'(t)\dd t
 =\int_{g(\alpha)}^{g(\beta)}f(x)\dd x.
\]
この形の置換公式に$g$の単調性は不要である。ただし，もとの$x$から$t$へ一意に戻す方法として$x=g(t)$を使う場合は，区間と枝を選ぶ必要がある。\chapterref{章:置換積分}を参照。

$u,v$が連続微分可能なら，
\[
 \int_a^b u(x)v'(x)\dd x
 =[u(x)v(x)]_a^b-\int_a^b u'(x)v(x)\dd x.
\]
右側に残る積分が簡単になるかを考えて，微分する因子と積分する因子を決める。二回の部分積分で元の積分が戻った場合には，得られた等式を元の積分について解く。\chapterref{章:部分積分}を参照。

\section{定積分に固有の判断}
\begin{longtable}{@{}p{39mm}p{85mm}p{20mm}@{}}
\toprule
式に見えるもの & まず確かめること & 本文\\
\midrule
\endhead
対称な区間$[-a,a]$ & $f(-x)$を計算し，偶関数・奇関数に分けられるかを調べる。 & \chapterref{章:定積分の対称性}\\
区間$[a,b]$と両端の対称性 & $x\mapsto a+b-x$で変形し，元の被積分関数と足すと簡単になるかを見る。 & \chapterref{章:定積分の対称性}\\
$\sin x$と$\cos x$が対称 & $x\mapsto\frac\pi2-x$で分子と分母がどう変わるかを調べる。 & \chapterref{章:定積分の対称性}\\
分母$1+\e^x$と偶関数の因子 & $\frac1{1+\e^x}+\frac1{1+\e^{-x}}=1$を使い，左右を合わせる。 & \chapterref{章:定積分の対称性}\\
$k$と$n$を含む和 & 幅$\frac1n$と高さ$f(\frac kn)$を取り出す。区間も確認する。 & \chapterref{章:積分と数列}\\
自然数$n$を含む積分 & 部分積分などで隣の添字と結ぶか，積分のまま上下を評価する。 & \chapterref{章:積分と数列}\\
原始関数が見つからない & 値の表示・近似・大小比較・収束判定のどれが必要かを分ける。 & 第9・11章\\
端点や内部で関数が発散 & 各特異点の片側を独立に切り取り，通常の積分を計算してから極限を取る。 & \chapterref{章:近似と広義積分}\\
\bottomrule
\end{longtable}

\section{積分で量を表す}
$a<b$とする。必要な連続性・滑らかさを仮定し，曲線や立体の同じ部分を重複して数えない。
\begin{align*}
 \text{二曲線間の面積}\quad&A=\int_a^b|f(x)-g(x)|\dd x,\\
 \text{断面積からの体積}\quad&V=\int_a^b A(x)\dd x,\\
 \text{円環断面の体積}\quad&V=\pi\int_a^b\{R(x)^2-r(x)^2\}\dd x,\\
 \text{グラフの長さ}\quad&L=\int_a^b\sqrt{1+f'(x)^2}\dd x,\\
 \text{媒介曲線の長さ}\quad&L=\int_\alpha^\beta\sqrt{x'(t)^2+y'(t)^2}\dd t,\\
 \text{位置の変化}\quad&s(b)-s(a)=\int_a^b v(t)\dd t,\\
 \text{道のり}\quad&D=\int_a^b|v(t)|\dd t,\\
 \text{関数の平均値}\quad&\overline f=\frac1{b-a}\int_a^b f(x)\dd x.
\end{align*}
円環の式の$R,r$は回転軸からの距離であり，$0\le r\le R$である。回転させる領域が軸をまたぐとき，符号付きの高さをそのまま外半径・内半径に使わない。\chapterref{章:積分の応用}を参照。

\section{累乗の積分}
非負整数$m,n$に対して，
\[
 \int_0^1x^m(1-x)^n\dd x=\frac{m!n!}{(m+n+1)!}.
\]
$W_n=\int_0^{\frac\pi2}\sin^n x\dd x$とおくと，
\[
 W_0=\frac\pi2,\quad W_1=1,\quad
 W_n=\frac{n-1}{n}W_{n-2}\quad(n\ge2).
\]
どちらも\chapterref{章:積分と数列}で部分積分から導いた。添字を下げる前に，境界項と開始できる$n$を確認する。

\backmatter
\chapter*{用語索引}
\addcontentsline{toc}{chapter}{用語索引}
\markboth{用語索引}{}
主な用語と，定義・説明を置いた箇所を示す。頁番号と用語名から本文へ移動できる。
\begingroup
\renewcommand{\arraystretch}{1.2}
\begin{longtable}{@{}p{90mm}p{35mm}r@{}}
\toprule
用語 & 主な参照先 & 頁\\
\midrule
\endfirsthead
\toprule
用語 & 主な参照先 & 頁\\
\midrule
\endhead
\hyperref[節:一次近似]{一次近似} & \sectionref{節:一次近似} & \pageref{節:一次近似}\\
\hyperref[節:運動と積分]{位置の変化と道のり} & \sectionref{節:運動と積分} & \pageref{節:運動と積分}\\
\hyperref[節:逆媒介陰微分]{陰関数の微分} & \sectionref{節:逆媒介陰微分} & \pageref{節:逆媒介陰微分}\\
\hyperref[節:ウォリス公式]{ウォリスの公式} & \sectionref{節:ウォリス公式} & \pageref{節:ウォリス公式}\\
\hyperref[式:円筒殻の体積]{円筒殻} & 式\eqref{式:円筒殻の体積} & \pageref{式:円筒殻の体積}\\
\hyperref[節:回転体の体積]{回転体の体積} & \sectionref{節:回転体の体積} & \pageref{節:回転体の体積}\\
\hyperref[式:回転面の面積]{回転面の面積} & 式\eqref{式:回転面の面積} & \pageref{式:回転面の面積}\\
\hyperref[節:等比級数]{級数と部分和} & \sectionref{節:等比級数} & \pageref{節:等比級数}\\
\hyperref[節:曲線の弧長]{曲線の長さ} & \sectionref{節:曲線の弧長} & \pageref{節:曲線の弧長}\\
\hyperref[節:奇偶の対称性]{奇関数・偶関数} & \sectionref{節:奇偶の対称性} & \pageref{節:奇偶の対称性}\\
\hyperref[節:逆媒介陰微分]{逆関数の微分} & \sectionref{節:逆媒介陰微分} & \pageref{節:逆媒介陰微分}\\
\hyperref[節:逆三角関数]{逆三角関数} & \sectionref{節:逆三角関数} & \pageref{節:逆三角関数}\\
\hyperref[節:和の区分求積]{区分求積法} & \sectionref{節:和の区分求積} & \pageref{節:和の区分求積}\\
\hyperref[節:原始関数基本定理]{原始関数} & \sectionref{節:原始関数基本定理} & \pageref{節:原始関数基本定理}\\
\hyperref[節:円の有理パラメータ]{原始ピタゴラス数} & \sectionref{節:円の有理パラメータ} & \pageref{節:円の有理パラメータ}\\
\hyperref[節:広義積分の定義]{広義積分} & \sectionref{節:広義積分の定義} & \pageref{節:広義積分の定義}\\
\hyperref[節:積分の不等式]{コーシー・シュワルツの不等式} & \sectionref{節:積分の不等式} & \pageref{節:積分の不等式}\\
\hyperref[節:微分法則]{合成関数の微分} & \sectionref{節:微分法則} & \pageref{節:微分法則}\\
\hyperref[節:周期関数の積分]{周期関数の積分} & \sectionref{節:周期関数の積分} & \pageref{節:周期関数の積分}\\
\hyperref[節:数列極限]{収束と発散} & \sectionref{節:数列極限} & \pageref{節:数列極限}\\
\hyperref[節:数値積分]{シンプソン公式} & \sectionref{節:数値積分} & \pageref{節:数値積分}\\
\hyperref[節:面積重心と線重心]{重心（面積・線）} & \sectionref{節:面積重心と線重心} & \pageref{節:面積重心と線重心}\\
\hyperref[節:部分積分の循環]{循環する部分積分} & \sectionref{節:部分積分の循環} & \pageref{節:部分積分の循環}\\
\hyperref[節:相反三角関数]{正割・余割・余接} & \sectionref{節:相反三角関数} & \pageref{節:相反三角関数}\\
\hyperref[節:定積分性質]{積分の平均値の定理} & \sectionref{節:定積分性質} & \pageref{節:定積分性質}\\
\hyperref[節:原始関数基本定理]{積分定数} & \sectionref{節:原始関数基本定理} & \pageref{節:原始関数基本定理}\\
\hyperref[節:積分方程式]{積分方程式} & \sectionref{節:積分方程式} & \pageref{節:積分方程式}\\
\hyperref[節:積分漸化式]{積分漸化式} & \sectionref{節:積分漸化式} & \pageref{節:積分漸化式}\\
\hyperref[節:双曲線関数]{双曲線関数} & \sectionref{節:双曲線関数} & \pageref{節:双曲線関数}\\
\hyperref[節:相補重み]{相補的な重み} & \sectionref{節:相補重み} & \pageref{節:相補重み}\\
\hyperref[節:数値積分]{台形公式} & \sectionref{節:数値積分} & \pageref{節:数値積分}\\
\hyperref[式:中央二項係数の極限]{中央二項係数} & 式\eqref{式:中央二項係数の極限} & \pageref{式:中央二項係数の極限}\\
\hyperref[章:置換積分]{置換積分} & \chapterref{章:置換積分} & \pageref{章:置換積分}\\
\hyperref[章:定積分の対称性]{定積分の対称性} & \chapterref{章:定積分の対称性} & \pageref{章:定積分の対称性}\\
\hyperref[節:有限テイラー]{テイラー公式と剰余項} & \sectionref{節:有限テイラー} & \pageref{節:有限テイラー}\\
\hyperref[節:等比級数]{等比級数} & \sectionref{節:等比級数} & \pageref{節:等比級数}\\
\hyperref[節:微分定義]{導関数} & \sectionref{節:微分定義} & \pageref{節:微分定義}\\
\hyperref[節:半角置換]{半角置換} & \sectionref{節:半角置換} & \pageref{節:半角置換}\\
\hyperref[節:曲線の弧長]{媒介曲線の面積・長さ} & \sectionref{節:曲線の弧長} & \pageref{節:曲線の弧長}\\
\hyperref[節:パップスの定理]{パップス＝ギュルダンの定理} & \sectionref{節:パップスの定理} & \pageref{節:パップスの定理}\\
\hyperref[節:非初等原始関数]{非初等原始関数} & \sectionref{節:非初等原始関数} & \pageref{節:非初等原始関数}\\
\hyperref[節:基本定理第一]{微積分の基本定理} & \sectionref{節:基本定理第一} & \pageref{節:基本定理第一}\\
\hyperref[節:微分定義]{微分係数} & \sectionref{節:微分定義} & \pageref{節:微分定義}\\
\hyperref[節:原始関数基本定理]{不定積分} & \sectionref{節:原始関数基本定理} & \pageref{節:原始関数基本定理}\\
\hyperref[節:部分分数]{部分分数} & \sectionref{節:部分分数} & \pageref{節:部分分数}\\
\hyperref[章:部分積分]{部分積分} & \chapterref{章:部分積分} & \pageref{章:部分積分}\\
\hyperref[節:部分積分の反復]{部分積分の反復} & \sectionref{節:部分積分の反復} & \pageref{節:部分積分の反復}\\
\hyperref[節:平均値定理]{平均値の定理（微分）} & \sectionref{節:平均値定理} & \pageref{節:平均値定理}\\
\hyperref[節:関数の平均]{平均値（関数）} & \sectionref{節:関数の平均} & \pageref{節:関数の平均}\\
\hyperref[節:変数上下端]{変数上下端} & \sectionref{節:変数上下端} & \pageref{節:変数上下端}\\
\hyperref[節:広義積分の判定]{べき関数による収束判定} & \sectionref{節:広義積分の判定} & \pageref{節:広義積分の判定}\\
\hyperref[節:ベータ積分]{ベータ関数} & \sectionref{節:ベータ積分} & \pageref{節:ベータ積分}\\
\hyperref[節:曲線間の面積]{面積（曲線で囲まれた領域）} & \sectionref{節:曲線間の面積} & \pageref{節:曲線間の面積}\\
\hyperref[節:リーマン和]{リーマン和} & \sectionref{節:リーマン和} & \pageref{節:リーマン和}\\
\hyperref[節:関数極限連続]{連続性} & \sectionref{節:関数極限連続} & \pageref{節:関数極限連続}\\
\hyperref[用語:連続微分可能]{連続微分可能} & \sectionref{節:合成形の置換} & \pageref{用語:連続微分可能}\\
\hyperref[節:ロピタル]{ロピタルの定理} & \sectionref{節:ロピタル} & \pageref{節:ロピタル}\\
\hyperref[節:半角置換]{ワイエルシュトラス置換} & \sectionref{節:半角置換} & \pageref{節:半角置換}\\
\hyperref[節:和と積分の比較]{和と積分の比較} & \sectionref{節:和と積分の比較} & \pageref{節:和と積分の比較}\\
\bottomrule
\end{longtable}
\endgroup

\chapter*{参考資料}
\markboth{参考資料}{}
\addcontentsline{toc}{chapter}{参考資料}
本書の説明と問題は，この構成に合わせて作成した。以下は，扱う範囲や標準的な定理の条件，公式の表記を確認した資料である。発展内容がすべて高校の学習範囲であることを意味するものではない。
\begin{enumerate}[leftmargin=*]
\item 文部科学省『高等学校学習指導要領（平成30年告示）解説 数学編・理数編』，2018年。数学II・IIIの極限・微分・積分の範囲の確認。掲載ページ：\url{https://www.mext.go.jp/a_menu/shotou/new-cs/1407074.htm}。
\item Gilbert Strang and Edwin Herman, \textit{Calculus Volume 1}, OpenStax, 2016. 第5章第3節「The Fundamental Theorem of Calculus」。基本定理の条件と面積関数。\url{https://openstax.org/books/calculus-volume-1/pages/5-3-the-fundamental-theorem-of-calculus}。
\item Gilbert Strang and Edwin Herman, \textit{Calculus Volume 2}, OpenStax, 2016. 第3章第3節「Trigonometric Substitution」。三角置換の範囲と符号。\url{https://openstax.org/books/calculus-volume-2/pages/3-3-trigonometric-substitution}。
\item NIST, \textit{Digital Library of Mathematical Functions}, \S5.12, ``Beta Function''. ベータ積分と三角関数の積分による表示。\url{https://dlmf.nist.gov/5.12}。
\end{enumerate}
オンライン資料の確認日：2026年10月19日。本書で証明を省略して採用する事実は，該当箇所でその旨を示した。

\end{document}
